% Options for packages loaded elsewhere
\PassOptionsToPackage{unicode}{hyperref}
\PassOptionsToPackage{hyphens}{url}
\PassOptionsToPackage{dvipsnames,svgnames,x11names}{xcolor}
%
\documentclass[
  letterpaper,
  DIV=11,
  numbers=noendperiod]{scrartcl}

\usepackage{amsmath,amssymb}
\usepackage{lmodern}
\usepackage{iftex}
\ifPDFTeX
  \usepackage[T1]{fontenc}
  \usepackage[utf8]{inputenc}
  \usepackage{textcomp} % provide euro and other symbols
\else % if luatex or xetex
  \usepackage{unicode-math}
  \defaultfontfeatures{Scale=MatchLowercase}
  \defaultfontfeatures[\rmfamily]{Ligatures=TeX,Scale=1}
\fi
% Use upquote if available, for straight quotes in verbatim environments
\IfFileExists{upquote.sty}{\usepackage{upquote}}{}
\IfFileExists{microtype.sty}{% use microtype if available
  \usepackage[]{microtype}
  \UseMicrotypeSet[protrusion]{basicmath} % disable protrusion for tt fonts
}{}
\makeatletter
\@ifundefined{KOMAClassName}{% if non-KOMA class
  \IfFileExists{parskip.sty}{%
    \usepackage{parskip}
  }{% else
    \setlength{\parindent}{0pt}
    \setlength{\parskip}{6pt plus 2pt minus 1pt}}
}{% if KOMA class
  \KOMAoptions{parskip=half}}
\makeatother
\usepackage{xcolor}
\setlength{\emergencystretch}{3em} % prevent overfull lines
\setcounter{secnumdepth}{-\maxdimen} % remove section numbering
% Make \paragraph and \subparagraph free-standing
\ifx\paragraph\undefined\else
  \let\oldparagraph\paragraph
  \renewcommand{\paragraph}[1]{\oldparagraph{#1}\mbox{}}
\fi
\ifx\subparagraph\undefined\else
  \let\oldsubparagraph\subparagraph
  \renewcommand{\subparagraph}[1]{\oldsubparagraph{#1}\mbox{}}
\fi

\usepackage{color}
\usepackage{fancyvrb}
\newcommand{\VerbBar}{|}
\newcommand{\VERB}{\Verb[commandchars=\\\{\}]}
\DefineVerbatimEnvironment{Highlighting}{Verbatim}{commandchars=\\\{\}}
% Add ',fontsize=\small' for more characters per line
\usepackage{framed}
\definecolor{shadecolor}{RGB}{241,243,245}
\newenvironment{Shaded}{\begin{snugshade}}{\end{snugshade}}
\newcommand{\AlertTok}[1]{\textcolor[rgb]{0.68,0.00,0.00}{#1}}
\newcommand{\AnnotationTok}[1]{\textcolor[rgb]{0.37,0.37,0.37}{#1}}
\newcommand{\AttributeTok}[1]{\textcolor[rgb]{0.40,0.45,0.13}{#1}}
\newcommand{\BaseNTok}[1]{\textcolor[rgb]{0.68,0.00,0.00}{#1}}
\newcommand{\BuiltInTok}[1]{\textcolor[rgb]{0.00,0.23,0.31}{#1}}
\newcommand{\CharTok}[1]{\textcolor[rgb]{0.13,0.47,0.30}{#1}}
\newcommand{\CommentTok}[1]{\textcolor[rgb]{0.37,0.37,0.37}{#1}}
\newcommand{\CommentVarTok}[1]{\textcolor[rgb]{0.37,0.37,0.37}{\textit{#1}}}
\newcommand{\ConstantTok}[1]{\textcolor[rgb]{0.56,0.35,0.01}{#1}}
\newcommand{\ControlFlowTok}[1]{\textcolor[rgb]{0.00,0.23,0.31}{#1}}
\newcommand{\DataTypeTok}[1]{\textcolor[rgb]{0.68,0.00,0.00}{#1}}
\newcommand{\DecValTok}[1]{\textcolor[rgb]{0.68,0.00,0.00}{#1}}
\newcommand{\DocumentationTok}[1]{\textcolor[rgb]{0.37,0.37,0.37}{\textit{#1}}}
\newcommand{\ErrorTok}[1]{\textcolor[rgb]{0.68,0.00,0.00}{#1}}
\newcommand{\ExtensionTok}[1]{\textcolor[rgb]{0.00,0.23,0.31}{#1}}
\newcommand{\FloatTok}[1]{\textcolor[rgb]{0.68,0.00,0.00}{#1}}
\newcommand{\FunctionTok}[1]{\textcolor[rgb]{0.28,0.35,0.67}{#1}}
\newcommand{\ImportTok}[1]{\textcolor[rgb]{0.00,0.46,0.62}{#1}}
\newcommand{\InformationTok}[1]{\textcolor[rgb]{0.37,0.37,0.37}{#1}}
\newcommand{\KeywordTok}[1]{\textcolor[rgb]{0.00,0.23,0.31}{#1}}
\newcommand{\NormalTok}[1]{\textcolor[rgb]{0.00,0.23,0.31}{#1}}
\newcommand{\OperatorTok}[1]{\textcolor[rgb]{0.37,0.37,0.37}{#1}}
\newcommand{\OtherTok}[1]{\textcolor[rgb]{0.00,0.23,0.31}{#1}}
\newcommand{\PreprocessorTok}[1]{\textcolor[rgb]{0.68,0.00,0.00}{#1}}
\newcommand{\RegionMarkerTok}[1]{\textcolor[rgb]{0.00,0.23,0.31}{#1}}
\newcommand{\SpecialCharTok}[1]{\textcolor[rgb]{0.37,0.37,0.37}{#1}}
\newcommand{\SpecialStringTok}[1]{\textcolor[rgb]{0.13,0.47,0.30}{#1}}
\newcommand{\StringTok}[1]{\textcolor[rgb]{0.13,0.47,0.30}{#1}}
\newcommand{\VariableTok}[1]{\textcolor[rgb]{0.07,0.07,0.07}{#1}}
\newcommand{\VerbatimStringTok}[1]{\textcolor[rgb]{0.13,0.47,0.30}{#1}}
\newcommand{\WarningTok}[1]{\textcolor[rgb]{0.37,0.37,0.37}{\textit{#1}}}

\providecommand{\tightlist}{%
  \setlength{\itemsep}{0pt}\setlength{\parskip}{0pt}}\usepackage{longtable,booktabs,array}
\usepackage{calc} % for calculating minipage widths
% Correct order of tables after \paragraph or \subparagraph
\usepackage{etoolbox}
\makeatletter
\patchcmd\longtable{\par}{\if@noskipsec\mbox{}\fi\par}{}{}
\makeatother
% Allow footnotes in longtable head/foot
\IfFileExists{footnotehyper.sty}{\usepackage{footnotehyper}}{\usepackage{footnote}}
\makesavenoteenv{longtable}
\usepackage{graphicx}
\makeatletter
\def\maxwidth{\ifdim\Gin@nat@width>\linewidth\linewidth\else\Gin@nat@width\fi}
\def\maxheight{\ifdim\Gin@nat@height>\textheight\textheight\else\Gin@nat@height\fi}
\makeatother
% Scale images if necessary, so that they will not overflow the page
% margins by default, and it is still possible to overwrite the defaults
% using explicit options in \includegraphics[width, height, ...]{}
\setkeys{Gin}{width=\maxwidth,height=\maxheight,keepaspectratio}
% Set default figure placement to htbp
\makeatletter
\def\fps@figure{htbp}
\makeatother

\KOMAoption{captions}{tableheading}
\makeatletter
\@ifpackageloaded{tcolorbox}{}{\usepackage[many]{tcolorbox}}
\@ifpackageloaded{fontawesome5}{}{\usepackage{fontawesome5}}
\definecolor{quarto-callout-color}{HTML}{909090}
\definecolor{quarto-callout-note-color}{HTML}{0758E5}
\definecolor{quarto-callout-important-color}{HTML}{CC1914}
\definecolor{quarto-callout-warning-color}{HTML}{EB9113}
\definecolor{quarto-callout-tip-color}{HTML}{00A047}
\definecolor{quarto-callout-caution-color}{HTML}{FC5300}
\definecolor{quarto-callout-color-frame}{HTML}{acacac}
\definecolor{quarto-callout-note-color-frame}{HTML}{4582ec}
\definecolor{quarto-callout-important-color-frame}{HTML}{d9534f}
\definecolor{quarto-callout-warning-color-frame}{HTML}{f0ad4e}
\definecolor{quarto-callout-tip-color-frame}{HTML}{02b875}
\definecolor{quarto-callout-caution-color-frame}{HTML}{fd7e14}
\makeatother
\makeatletter
\makeatother
\makeatletter
\makeatother
\makeatletter
\@ifpackageloaded{caption}{}{\usepackage{caption}}
\AtBeginDocument{%
\ifdefined\contentsname
  \renewcommand*\contentsname{Tabla de contenidos}
\else
  \newcommand\contentsname{Tabla de contenidos}
\fi
\ifdefined\listfigurename
  \renewcommand*\listfigurename{Listado de Figuras}
\else
  \newcommand\listfigurename{Listado de Figuras}
\fi
\ifdefined\listtablename
  \renewcommand*\listtablename{Listado de Tablas}
\else
  \newcommand\listtablename{Listado de Tablas}
\fi
\ifdefined\figurename
  \renewcommand*\figurename{Figura}
\else
  \newcommand\figurename{Figura}
\fi
\ifdefined\tablename
  \renewcommand*\tablename{Tabla}
\else
  \newcommand\tablename{Tabla}
\fi
}
\@ifpackageloaded{float}{}{\usepackage{float}}
\floatstyle{ruled}
\@ifundefined{c@chapter}{\newfloat{codelisting}{h}{lop}}{\newfloat{codelisting}{h}{lop}[chapter]}
\floatname{codelisting}{Listado}
\newcommand*\listoflistings{\listof{codelisting}{Listado de Listados}}
\makeatother
\makeatletter
\@ifpackageloaded{caption}{}{\usepackage{caption}}
\@ifpackageloaded{subcaption}{}{\usepackage{subcaption}}
\makeatother
\makeatletter
\@ifpackageloaded{tcolorbox}{}{\usepackage[many]{tcolorbox}}
\makeatother
\makeatletter
\@ifundefined{shadecolor}{\definecolor{shadecolor}{rgb}{.97, .97, .97}}
\makeatother
\makeatletter
\makeatother
\ifLuaTeX
\usepackage[bidi=basic]{babel}
\else
\usepackage[bidi=default]{babel}
\fi
\babelprovide[main,import]{spanish}
% get rid of language-specific shorthands (see #6817):
\let\LanguageShortHands\languageshorthands
\def\languageshorthands#1{}
\ifLuaTeX
  \usepackage{selnolig}  % disable illegal ligatures
\fi
\IfFileExists{bookmark.sty}{\usepackage{bookmark}}{\usepackage{hyperref}}
\IfFileExists{xurl.sty}{\usepackage{xurl}}{} % add URL line breaks if available
\urlstyle{same} % disable monospaced font for URLs
\hypersetup{
  pdftitle={Caso 1 · ¿Ocurre el evento? ¿Y cuándo?},
  pdfauthor={A.M.Mayoral, con asistencia de Claude Opus 4.8},
  pdflang={es},
  colorlinks=true,
  linkcolor={blue},
  filecolor={Maroon},
  citecolor={Blue},
  urlcolor={Blue},
  pdfcreator={LaTeX via pandoc}}

\title{Caso 1 · ¿Ocurre el evento? ¿Y cuándo?}
\usepackage{etoolbox}
\makeatletter
\providecommand{\subtitle}[1]{% add subtitle to \maketitle
  \apptocmd{\@title}{\par {\large #1 \par}}{}{}
}
\makeatother
\subtitle{Respuestas binarias, politómicas y tiempo hasta el evento ---
Modelos Lineales Generalizados con R}
\author{A.M.Mayoral, con asistencia de Claude Opus 4.8}
\date{3/9/26}

\begin{document}
\maketitle
\ifdefined\Shaded\renewenvironment{Shaded}{\begin{tcolorbox}[breakable, interior hidden, enhanced, sharp corners, borderline west={3pt}{0pt}{shadecolor}, frame hidden, boxrule=0pt]}{\end{tcolorbox}}\fi

\renewcommand*\contentsname{Tabla de contenidos}
{
\hypersetup{linkcolor=}
\setcounter{tocdepth}{3}
\tableofcontents
}
\hypertarget{sec-ra}{%
\subsection{Resultados de aprendizaje del Caso 1}\label{sec-ra}}

Al completar el Caso 1, el estudiante será capaz de:

\begin{itemize}
\tightlist
\item
  \textbf{RA1 · Marco GLM frente a OLS.} Reconocer por qué el modelo
  lineal (OLS) falla ante respuestas discretas y describir los tres
  componentes de un GLM: componente aleatorio, predictor lineal y
  función de enlace.
\item
  \textbf{RA2 · Formulación y estimación.} Formular un GLM en la familia
  exponencial, ajustarlo por máxima verosimilitud (IWLS) y hacer
  inferencia (deviance, Wald, razón de verosimilitudes).
\item
  \textbf{RA3 · Interpretación en la escala correcta.} Interpretar los
  efectos en la escala propia de cada respuesta: odds ratios y efectos
  marginales (binaria), odds ratios relativos (nominal), odds ratios
  acumulados (ordinal) y hazard ratios (supervivencia).
\item
  \textbf{RA4 · Diagnóstico y bondad de ajuste.} Evaluar el modelo con
  residuos, separación, calibración, el contraste χ² de bondad de ajuste
  y la sobredispersión.
\item
  \textbf{RA5 · Selección y comparación de modelos.} Comparar y
  seleccionar modelos con deviance, AIC/BIC y contrastes de razón de
  verosimilitudes.
\item
  \textbf{RA6 · Efectos aleatorios.} Ajustar e interpretar modelos
  mixtos (GLMM) con intercepto y pendiente aleatorios: componentes de
  varianza, ICC, \emph{shrinkage} y la distinción condicional vs
  marginal.
\item
  \textbf{RA7 · Tiempo-a-evento y censura.} Reformular la supervivencia
  en tiempo discreto como un GLM binomial con enlace \texttt{cloglog}
  sobre datos persona-periodo, tratando la censura y prediciendo hazard
  y supervivencia.
\item
  \textbf{RA8 · Estudio de caso reproducible.} Desarrollar, presentar y
  defender un estudio de caso reproducible que integre las lentes
  anteriores.
\end{itemize}

\begin{tcolorbox}[enhanced jigsaw, coltitle=black, rightrule=.15mm, arc=.35mm, leftrule=.75mm, colback=white, colbacktitle=quarto-callout-note-color!10!white, left=2mm, breakable, toptitle=1mm, toprule=.15mm, bottomtitle=1mm, titlerule=0mm, title=\textcolor{quarto-callout-note-color}{\faInfo}\hspace{0.5em}{Mapa del caso}, bottomrule=.15mm, opacityback=0, opacitybacktitle=0.6, colframe=quarto-callout-note-color-frame]

En esta unidad vamos a trabajar sobre respuestas categóricas. Iniciamos
con respuesta binaria, seguimos con su agrupación natural en respuestas
binomiales, y extendemos a respuesta politómica y multinomial. Por el
camino incorporamos los efectos aleatorios en la modelización, y por
último estudiamos también el abordaje como un problema de supervivencia
para ir más allá de la predición del evento, incorporando datos
censurados y prediciendo también cuándo se va a producir el evento.

Trabajaremos con la base de datos GLOW y usaremos la cohorte simulada
para el estudio de supervivencia y el estudio de caso sobre el que
trabajarás tú estudiante, para demostrar que has aprendido a modelizar
datos discretos.

\begin{longtable}[]{@{}
  >{\centering\arraybackslash}p{(\columnwidth - 4\tabcolsep) * \real{0.0455}}
  >{\raggedright\arraybackslash}p{(\columnwidth - 4\tabcolsep) * \real{0.7386}}
  >{\raggedright\arraybackslash}p{(\columnwidth - 4\tabcolsep) * \real{0.2159}}@{}}
\toprule()
\begin{minipage}[b]{\linewidth}\centering
Unidad
\end{minipage} & \begin{minipage}[b]{\linewidth}\raggedright
Contenido
\end{minipage} & \begin{minipage}[b]{\linewidth}\raggedright
Objetivo y datos
\end{minipage} \\
\midrule()
\endhead
\textbf{1.1} & \emph{«¿Ocurre el evento?»} --- Cuando la recta no llega:
OLS falla en binaria; tres componentes del GLM; tipología de respuestas;
primer \texttt{glm}. & \textbf{RA1} (y primer contacto con
RA2)\textbf{Datos:} \texttt{glow} \\
\textbf{1.2} & \emph{«¿Cómo se ajusta e interpreta una logística, y cómo
de bien clasifica?»} --- Formulación de los GLM: familia exponencial,
enlace, IWLS, \emph{deviance}; inferencia; interpretación (OR, efectos
marginales); evaluación y comparación de clasificadores. & \textbf{RA2,
RA3, RA4, RA5}\textbf{Datos:} \texttt{glow} \\
\textbf{1.3} & \emph{«¿Y si la respuesta tiene más de dos categorías?»}
--- Respuesta categórica generalizada: binomial agrupado (ajuste y
bondad de ajuste con χ²) y respuestas politómicas --- multinomial
nominal y ordinal (odds proporcionales). & \textbf{RA2--RA5} (nuevos
tipos de respuesta)\textbf{Datos:} \texttt{glow} (agrupado y
politómico) \\
\textbf{1.4} & \emph{«¿Y si los sujetos vienen agrupados?»} --- Efectos
aleatorios (GLMM logístico, intercepto y pendiente aleatorios) sobre la
cohorte y validación con datos reales (GLOW); puente a la supervivencia.
& \textbf{RA6} (y RA4 en el diagnóstico del mixto)\textbf{Datos:}
\texttt{cohorte}, \texttt{glow} \\
\textbf{1.5} & \emph{«¿Y cuándo ocurre el evento?»} --- Supervivencia:
función de supervivencia, hazard, riesgos proporcionales (Cox como
referencia, detalle en el Caso 3); el tiempo discreto como GLM
(\texttt{cloglog} persona-periodo). & \textbf{RA7} (y RA3 vía hazard
ratios)\textbf{Datos:} \texttt{cohorte} (persona-periodo) \\
\textbf{1.6} & \emph{«¿Sabes ya modelizar datos discretos?»} --- Estudio
de caso, exposiciones y evaluación (actividades, instrucciones,
rúbricas). & \textbf{RA8} (integra RA1--RA7)\textbf{Datos:} banco del
estudiante \\
\bottomrule()
\end{longtable}

\textbf{Caja de herramientas de R por unidad.} Las funciones y paquetes
que vertebran cada sesión, \textbf{enlazados a la sección donde se
explican} (el uso comentado vive en el callout «🔧 En R» de cada sección
y en los scripts de \texttt{caso1/scripts/}):

\begin{longtable}[]{@{}
  >{\centering\arraybackslash}p{(\columnwidth - 2\tabcolsep) * \real{0.0526}}
  >{\raggedright\arraybackslash}p{(\columnwidth - 2\tabcolsep) * \real{0.9474}}@{}}
\toprule()
\begin{minipage}[b]{\linewidth}\centering
Unidad
\end{minipage} & \begin{minipage}[b]{\linewidth}\raggedright
Funciones y paquetes clave
\end{minipage} \\
\midrule()
\endhead
\textbf{1.1} &
\protect\hyperlink{sec-u11-tres-componentes}{\texttt{glm(family\ =\ binomial)}
· \texttt{predict(type\ =\ "response")} ·
\texttt{broom::tidy}/\texttt{glance}/\texttt{augment} ·
\texttt{performance::binned\_residuals}} \\
\textbf{1.2} & enlace:
\protect\hyperlink{sec-u12-enlace}{\texttt{binomial("logit"/"probit"/"cloglog")}}
· interpretación:
\protect\hyperlink{sec-u12-interpretacion}{\texttt{broom::tidy} ·
\texttt{confint} ·
\texttt{marginaleffects::avg\_slopes}/\texttt{predictions}} ·
inferencia: \protect\hyperlink{sec-u12-inferencia}{\texttt{anova} ·
\texttt{AIC}/\texttt{BIC} · \texttt{MASS::stepAIC} ·
\texttt{MuMIn::dredge}} · diagnóstico:
\protect\hyperlink{sec-u12-diagnostico}{\texttt{DHARMa::simulateResiduals}
· \texttt{performance::performance\_hosmer} · \texttt{rms::val.prob} ·
\texttt{logistf} (Firth)} · evaluación:
\protect\hyperlink{sec-u12-evaluacion}{\texttt{caret::confusionMatrix} ·
\texttt{pROC::roc} · \texttt{yardstick::roc\_auc}} \\
\textbf{1.3} & binomial agrupado:
\protect\hyperlink{sec-u13-binomial}{\texttt{glm(cbind(exitos,\ fracasos),\ binomial)}}
· bondad de ajuste: \protect\hyperlink{sec-u13-gof-binomial}{deviance y
Pearson con \texttt{pchisq} · quasi-binomial} · nominal:
\protect\hyperlink{sec-u13-nominal}{\texttt{nnet::multinom} ·
\texttt{lmtest::lrtest} · Brier multiclase} · ordinal:
\protect\hyperlink{sec-u13-ordinal}{\texttt{MASS::polr} (odds
proporcionales) · \texttt{brant::brant} ·
\texttt{ordinal::nominal\_test} · \texttt{ggeffects::ggpredict}} \\
\textbf{1.4} & intercepto aleatorio:
\protect\hyperlink{sec-u14-intercepto}{\texttt{lme4::glmer} ·
\texttt{performance::icc} · \texttt{ranef}} · pendiente aleatoria:
\protect\hyperlink{sec-u14-pendiente}{\texttt{(1\ +\ x\ \textbackslash{}\textbar{}\ centro)}
· \texttt{VarCorr} · \texttt{anova}/\texttt{AIC}} · diagnóstico:
\protect\hyperlink{sec-u14-diagnostico}{\texttt{broom.mixed::tidy} ·
\texttt{DHARMa} en mixtos · \texttt{performance::check\_singularity}} ·
estructuras: \protect\hyperlink{sec-u14-extensiones}{la sintaxis
\texttt{(\ \textbackslash{}\textbar{}\ )} y \texttt{glmmTMB}} \\
\textbf{1.5} & Kaplan--Meier:
\protect\hyperlink{sec-u15-prop}{\texttt{survival::Surv} ·
\texttt{survfit}} · tiempo discreto:
\protect\hyperlink{sec-u15-lab}{\texttt{glm(binomial("cloglog"))}
persona-periodo · \texttt{lme4::glmer} con centro aleatorio} ·
\protect\hyperlink{sec-u15-cox}{Cox en tiempo continuo (aquí solo el
\textbf{concepto}: la unidad no llega a usar \texttt{coxph}; su
maquinaria, en la 3.6)} \\
\textbf{1.6} & Quarto (\texttt{.qmd}) · \texttt{renv} · Git ---
reproducibilidad y comunicación \\
\bottomrule()
\end{longtable}

\end{tcolorbox}

\hypertarget{sec-presentacion}{%
\section{Presentación del caso}\label{sec-presentacion}}

Iniciamos pues la unidad, presentando los bancos de datos con los que
vamos a trabajar. La pregunta que vertebra el caso tiene dos tiempos
básicos:

\begin{enumerate}
\def\labelenumi{\arabic{enumi}.}
\tightlist
\item
  \textbf{«¿Ocurre el evento?»}: una respuesta binaria que motiva el
  marco GLM y la regresión logística.
\item
  \textbf{«¿Cuándo?»}: el mismo evento como \textbf{tiempo hasta que
  ocurre}, con \textbf{censura}, que se reformula como un GLM sobre
  datos persona-periodo.
\end{enumerate}

En el recorrido, surgen dos matices que también trataremos:

\begin{itemize}
\tightlist
\item
  los sujetos vienen \textbf{agrupados} (por centro), e introducimos los
  \textbf{efectos aleatorios} y los \textbf{modelos mixtos},
\item
  y la respuesta puede tener \textbf{más de dos categorías}, motivando
  la extensión politómica, nominal u ordinal.
\end{itemize}

Para trabajar todas estas cuestiones utilizaremos dos fuentes de datos
con papeles distintos: un conjunto \textbf{real} para el ejemplo
logístico, y una \textbf{cohorte simulada} con proceso generador
conocido, que vertebra la extensión de la modelización a la
supervivencia.

\hypertarget{sec-datos-glow}{%
\subsection{Dato real --- GLOW}\label{sec-datos-glow}}

\texttt{glow500} procede del estudio \textbf{GLOW} (\emph{Global
Longitudinal Study of Osteoporosis in Women}) y se distribuye en el
paquete \texttt{aplore3} (licencia GPL-3), compañero del manual de
Hosmer et~al. (2013). Es una muestra de \textbf{500 mujeres} seguidas un
año en \textbf{6 centros}. La respuesta es binaria ---\texttt{fracture}:
alguna fractura en el primer año--- y el resto son factores de riesgo
basales:

\begin{longtable}[]{@{}
  >{\raggedright\arraybackslash}p{(\columnwidth - 4\tabcolsep) * \real{0.2927}}
  >{\raggedright\arraybackslash}p{(\columnwidth - 4\tabcolsep) * \real{0.5854}}
  >{\raggedright\arraybackslash}p{(\columnwidth - 4\tabcolsep) * \real{0.1220}}@{}}
\toprule()
\begin{minipage}[b]{\linewidth}\raggedright
Variable
\end{minipage} & \begin{minipage}[b]{\linewidth}\raggedright
Descripción
\end{minipage} & \begin{minipage}[b]{\linewidth}\raggedright
Tipo
\end{minipage} \\
\midrule()
\endhead
\texttt{fracture} & Fractura en el primer año (No/Yes) ---
\textbf{respuesta} & Binaria \\
\texttt{age} & Edad al ingreso (años) & Continua \\
\texttt{weight} & Peso al ingreso (kg) & Continua \\
\texttt{height} & Estatura al ingreso (cm) & Continua \\
\texttt{bmi} & Índice de masa corporal (kg/m²) & Continua \\
\texttt{priorfrac} & Antecedente de fractura previa (No/Yes) &
Binaria \\
\texttt{premeno} & Menopausia antes de los 45 (No/Yes) & Binaria \\
\texttt{momfrac} & Fractura de cadera en la madre (No/Yes) & Binaria \\
\texttt{armassist} & Necesita los brazos para levantarse de una silla
(No/Yes) & Binaria \\
\texttt{smoke} & Fumadora actual o pasada (No/Yes) & Binaria \\
\texttt{raterisk} & Riesgo de fractura autopercibido (Less/Same/Greater)
& Ordinal (3) \\
\texttt{fracscore} & Puntuación compuesta de riesgo de fractura &
Discreta \\
\texttt{sub\_id}, \texttt{site\_id}, \texttt{phy\_id} & Identificadores
de paciente, centro y médico & ID \\
\bottomrule()
\end{longtable}

Creamos también una versión 0/1 de la respuesta \texttt{fracture} para
el experimento, que utilizaremos en esta unidad.

\begin{Shaded}
\begin{Highlighting}[]
\FunctionTok{data}\NormalTok{(glow500)}
\NormalTok{glow }\OtherTok{\textless{}{-}}\NormalTok{ glow500 }\SpecialCharTok{|\textgreater{}}
  \FunctionTok{as\_tibble}\NormalTok{() }\SpecialCharTok{|\textgreater{}}
  \FunctionTok{mutate}\NormalTok{(}\AttributeTok{fractura01 =} \FunctionTok{as.integer}\NormalTok{(fracture) }\SpecialCharTok{{-}}\NormalTok{ 1L)   }\CommentTok{\# No {-}\textgreater{} 0, Yes {-}\textgreater{} 1}

\NormalTok{glow }\SpecialCharTok{|\textgreater{}} \FunctionTok{count}\NormalTok{(fracture) }\SpecialCharTok{|\textgreater{}} \FunctionTok{mutate}\NormalTok{(}\AttributeTok{prop =} \FunctionTok{round}\NormalTok{(n }\SpecialCharTok{/} \FunctionTok{sum}\NormalTok{(n), }\DecValTok{3}\NormalTok{))}
\end{Highlighting}
\end{Shaded}

\begin{verbatim}
# A tibble: 2 x 3
  fracture     n  prop
  <fct>    <int> <dbl>
1 No         375  0.75
2 Yes        125  0.25
\end{verbatim}

Visualizamos en la Figura~\ref{fig-glow-eda} los datos y las relaciones
entre variables en GLOW.

\begin{Shaded}
\begin{Highlighting}[]
\CommentTok{\# Covariables sustantivas de GLOW (dejo fuera los identificadores)}
\NormalTok{vars }\OtherTok{\textless{}{-}} \FunctionTok{c}\NormalTok{(}\StringTok{"age"}\NormalTok{, }\StringTok{"weight"}\NormalTok{, }\StringTok{"height"}\NormalTok{, }\StringTok{"bmi"}\NormalTok{, }\StringTok{"fracscore"}\NormalTok{,   }\CommentTok{\# continuas / discretas}
          \StringTok{"priorfrac"}\NormalTok{, }\StringTok{"premeno"}\NormalTok{, }\StringTok{"momfrac"}\NormalTok{,               }\CommentTok{\# factores de riesgo binarios}
          \StringTok{"armassist"}\NormalTok{, }\StringTok{"smoke"}\NormalTok{, }\StringTok{"raterisk"}\NormalTok{)                }\CommentTok{\# binarios + ordinal}

\NormalTok{glow }\SpecialCharTok{|\textgreater{}}
  \FunctionTok{select}\NormalTok{(fracture, }\FunctionTok{all\_of}\NormalTok{(vars)) }\SpecialCharTok{|\textgreater{}}
  \FunctionTok{ggpairs}\NormalTok{(}
    \AttributeTok{columns =}\NormalTok{ vars,}
    \AttributeTok{mapping =} \FunctionTok{aes}\NormalTok{(}\AttributeTok{color =}\NormalTok{ fracture, }\AttributeTok{alpha =} \FloatTok{0.6}\NormalTok{),}
    \AttributeTok{legend=}\DecValTok{1}\NormalTok{,}
    \AttributeTok{upper =} \FunctionTok{list}\NormalTok{(}\AttributeTok{continuous =} \FunctionTok{wrap}\NormalTok{(}\StringTok{"cor"}\NormalTok{, }\AttributeTok{size =} \FloatTok{2.3}\NormalTok{),     }\CommentTok{\# correlaciones (global y por grupo)}
                 \AttributeTok{combo      =} \FunctionTok{wrap}\NormalTok{(}\StringTok{"box\_no\_facet"}\NormalTok{),         }\CommentTok{\# continua \textasciitilde{} factor: boxplots}
                 \AttributeTok{discrete   =} \FunctionTok{wrap}\NormalTok{(}\StringTok{"count"}\NormalTok{)),               }\CommentTok{\# factor \textasciitilde{} factor: recuento}
    \AttributeTok{lower =} \FunctionTok{list}\NormalTok{(}\AttributeTok{continuous =} \FunctionTok{wrap}\NormalTok{(}\StringTok{"points"}\NormalTok{, }\AttributeTok{alpha =} \FloatTok{0.25}\NormalTok{, }\AttributeTok{size =} \FloatTok{0.4}\NormalTok{),}
                 \AttributeTok{combo      =} \FunctionTok{wrap}\NormalTok{(}\StringTok{"facethist"}\NormalTok{, }\AttributeTok{bins =} \DecValTok{20}\NormalTok{),}
                 \AttributeTok{discrete   =} \FunctionTok{wrap}\NormalTok{(}\StringTok{"facetbar"}\NormalTok{)),}
    \AttributeTok{diag  =} \FunctionTok{list}\NormalTok{(}\AttributeTok{continuous =} \FunctionTok{wrap}\NormalTok{(}\StringTok{"densityDiag"}\NormalTok{, }\AttributeTok{alpha =} \FloatTok{0.5}\NormalTok{),}
                 \AttributeTok{discrete   =} \FunctionTok{wrap}\NormalTok{(}\StringTok{"barDiag"}\NormalTok{)),}
    \AttributeTok{progress =} \ConstantTok{FALSE}
\NormalTok{  ) }\SpecialCharTok{+}
  \FunctionTok{theme\_bw}\NormalTok{(}\AttributeTok{base\_size =} \DecValTok{7}\NormalTok{) }\SpecialCharTok{+}
  \FunctionTok{theme}\NormalTok{(}\AttributeTok{strip.text =} \FunctionTok{element\_text}\NormalTok{(}\AttributeTok{size =} \DecValTok{6}\NormalTok{))}\SpecialCharTok{+}
  \FunctionTok{theme}\NormalTok{(}\AttributeTok{legend.position =} \StringTok{"bottom"}\NormalTok{)}
\end{Highlighting}
\end{Shaded}

\begin{figure}[H]

{\centering \includegraphics[width=1\textwidth,height=\textheight]{caso1_binaria_files/figure-pdf/fig-glow-eda-1.pdf}

}

\caption{\label{fig-glow-eda}Descripción de GLOW. En la diagonal
distribución de las variables, diferenciadas en color por
\texttt{evento}. En los cruces la relación entre ellas.}

\end{figure}

\hypertarget{sec-datos-cohorte}{%
\subsection{Cohorte simulada}\label{sec-datos-cohorte}}

Está inspirada en el problema de GLOW y simula el seguimiento de las
pacientes en el tiempo a lo largo de 6 años, para una cohorte de
mujeres, repartidas en varias clínicas (24).

El simulador \texttt{simular\_cohorte()} está disponible en
\texttt{R/dgp\_cohorte.R} y parte de unas condiciones establecidas sobre
las relaciones entre las variables consideradas (recuperables en
\texttt{str(attr(cohorte,\ "verdad"))}).

Cada paciente acude a \textbf{revisiones periódicas anuales} ---seis
como máximo--- y en cada una se constata si \textbf{ha sufrido una
fractura desde la revisión anterior}. No observamos el \textbf{instante
exacto} de la fractura, solo el \textbf{periodo} en que se detecta: por
eso \textbf{el tiempo es discreto}. Así, \texttt{tiempo}
\(\in\{1,\dots,6\}\) representa el periodo en que \textbf{termina la
observación} de cada mujer ---bien porque se le detecta una fractura
(\texttt{evento} = 1) o porque el seguimiento acaba sin ella
(\texttt{evento} = 0, \textbf{censura})---. \texttt{tiempo} = 1
significa que el desenlace ocurre en la \textbf{primera} revisión;
\texttt{tiempo} = 6, que llega a la \textbf{sexta y última}. La
expansión \textbf{persona-periodo} convierte a cada mujer en una fila
por cada revisión a la que llega sin fractura, con \texttt{y} = 1 solo
en aquella donde se detecta.

El simulador utiliza dos predictores \texttt{x1} y \texttt{x2}
abstractos, que podrían representar distintas características; a lo
largo del curso \textbf{adoptamos una caracterización fija} para
facilitar la interpretación:

\begin{itemize}
\tightlist
\item
  \texttt{x1} = \textbf{fragilidad ósea} (a partir de la densidad
  mineral, estandarizada en puntuación z): factor de riesgo
  \textbf{continuo}; a mayor valor, hueso más frágil y \textbf{más}
  riesgo de fractura (\(\beta_1>0\)).
\item
  \texttt{x2} = \textbf{tratamiento} preventivo (1 = lo recibe, 0 = no):
  factor \textbf{binario} y \textbf{protector}; reduce el riesgo
  (\(\beta_2<0\)).
\item
  \texttt{centro} = la \textbf{clínica} (24 en total); su intercepto
  aleatorio recoge que unas clínicas tienen más riesgo basal que otras
  ---la \emph{fragilidad} estadística / efecto mixto---.
\end{itemize}

Las variables simuladas son:

\begin{longtable}[]{@{}
  >{\raggedright\arraybackslash}p{(\columnwidth - 6\tabcolsep) * \real{0.1667}}
  >{\raggedright\arraybackslash}p{(\columnwidth - 6\tabcolsep) * \real{0.5000}}
  >{\raggedright\arraybackslash}p{(\columnwidth - 6\tabcolsep) * \real{0.2273}}
  >{\raggedright\arraybackslash}p{(\columnwidth - 6\tabcolsep) * \real{0.1061}}@{}}
\toprule()
\begin{minipage}[b]{\linewidth}\raggedright
Variable
\end{minipage} & \begin{minipage}[b]{\linewidth}\raggedright
Significado (lectura adoptada)
\end{minipage} & \begin{minipage}[b]{\linewidth}\raggedright
Tipo
\end{minipage} & \begin{minipage}[b]{\linewidth}\raggedright
Se usa en
\end{minipage} \\
\midrule()
\endhead
\texttt{id} & Identificador de la paciente & --- & --- \\
\texttt{centro} & \textbf{Clínica} (grupo del efecto aleatorio /
fragilidad) & Factor (24 niveles) & 1.4, 1.5 \\
\texttt{x1} & \textbf{Fragilidad ósea} (baja densidad mineral,
estandarizada) & Continua & 1.4, 1.5 \\
\texttt{x2} & \textbf{Tratamiento} preventivo (1 = sí, 0 = no) & Binaria
& 1.4, 1.5 \\
\texttt{tiempo} & \textbf{Periodo de revisión} (1--6) en que termina la
observación & Entero & 1.5 \\
\texttt{evento} & 1 = fractura \textbf{observada} en ese periodo; 0 =
\textbf{censura} & Binaria & 1.5 \\
\texttt{ever} & ¿Hubo fractura durante el seguimiento? (=
\texttt{evento}) & Binaria & 1.4 \\
\texttt{sever\_ord} & \textbf{Severidad basal ordinal} de la
osteoporosis (Leve \textless{} Moderado \textless{} Grave) & Factor
ordenado (3) & 1.4 \\
\texttt{clase\_nom} & \textbf{Endotipo nominal} de fragilidad ósea A/B/C
(sin orden) --- ilustra el contraste; el multinomial del curso se ajusta
sobre GLOW (\texttt{raterisk}) & Factor (3) & --- \\
\bottomrule()
\end{longtable}

Las respuestas se generan a partir de un modelo de \textbf{riesgo en
tiempo discreto} con enlace \texttt{cloglog} (lo estudiaremos más
adelante):
\[\operatorname{cloglog}(h_{ijt})=\alpha_t+\beta_1\,x_{1i}+\beta_2\,x_{2i}+u_j,
\qquad u_j\sim\mathcal{N}(0,\sigma_u^2),\] con \textbf{intercepto
aleatorio por clínica} \(u_j\) (fragilidad) y \textbf{censura
administrativa no informativa}. De ahí salen la respuesta binaria «¿hubo
fractura?» y el reloj de supervivencia (en qué periodo se produjo).
Construimos también una variable de \textbf{severidad ordinal}
(\texttt{sever\_ord}), que se genera aparte, con un cumulative logit más
intercepto aleatorio por clínica, disponible como ordinal con verdad
conocida para ejercicios. Las respuestas categóricas son
\textbf{independientes} del reloj de supervivencia ---comparten
covariables y clínicas, pero no se derivan de si hubo fractura, ni de
cuándo la hubo---, una decisión de diseño para no entrar en riesgos
competitivos (fuera del alcance de esta unidad).

\textbf{Las dos respuestas categóricas, en clave clínica.}

Aunque comparten las covariables \texttt{x1} y \texttt{x2} con el reloj
de supervivencia, \texttt{sever\_ord} y \texttt{clase\_nom} cuentan
historias clínicas distintas y se generan con estructuras distintas;
conviene fijar su lectura para interpretarlas con soltura.

\begin{itemize}
\item
  \texttt{sever\_ord} es una \textbf{gradación de severidad}: un único
  eje latente ordena los tres niveles ---el supuesto de \emph{odds
  proporcionales}--- y las clínicas difieren en su reparto basal
  (intercepto aleatorio por centro). Admite dos lecturas equivalentes:
  como \textbf{estadio densitométrico de la osteoporosis} (Leve
  \textless{} Moderado \textless{} Grave ≈ osteopenia \textless{}
  osteoporosis \textless{} osteoporosis establecida) o como el
  \textbf{grado semicuantitativo de fractura vertebral de Genant}
  (grados 1 \textless{} 2 \textless{} 3), que la literatura etiqueta
  literalmente leve/moderado/grave. Nos quedamos con la primera:
  ``estadio densitométrico de la osteoporosis''. Las hipótesis del
  simulador, \textbf{en tendencia}: a \textbf{mayor fragilidad ósea}
  (\texttt{x1}) la severidad tiende a \textbf{desplazarse hacia grados
  más altos}; el \textbf{tratamiento} (\texttt{x2}) se asocia a grados
  \textbf{más leves} (lectura protectora); y, a igualdad de paciente,
  \textbf{unas clínicas tienden a registrar severidades mayores que
  otras} (distinto \emph{case-mix} o criterio del lector). La relación
  con el tratamiento es \textbf{asociativa}, no causal.
\item
  \texttt{clase\_nom} es, en cambio, un \textbf{fenotipo cualitativo sin
  orden}: cada categoría sigue su propio patrón ---efectos específicos
  de categoría--- y la clínica \textbf{no} interviene (sin efecto
  aleatorio), de modo que es un rasgo del \textbf{paciente}. La leemos
  como un \textbf{endotipo de fragilidad ósea} con referencia \textbf{A}
  (baja densidad sin etiología secundaria), un fenotipo \textbf{B} de
  \textbf{caídas / fragilidad musculoesquelética} y un fenotipo
  \textbf{C} de \textbf{osteoporosis secundaria} (corticoidea o
  endocrina). En tendencia: a \textbf{mayor fragilidad ósea}
  (\texttt{x1}), más probable \textbf{alejarse de la referencia} hacia B
  y, sobre todo, hacia C; y el \textbf{tratamiento} (\texttt{x2})
  \textbf{separa} ambos fenotipos ---tiende a asociarse \textbf{menos}
  con el de caídas (que se maneja con prevención de caídas) y
  \textbf{más} con la osteoporosis secundaria (donde la profilaxis ósea
  casi se protocoliza)---. Como antes, se trata de una asociación de
  \textbf{indicación/selección}, no de causalidad.
\end{itemize}

\textbf{Simulación de los datos} Aplicamos el simulador generado y
construimos la \texttt{cohorte} completa.

\begin{Shaded}
\begin{Highlighting}[]
\FunctionTok{library}\NormalTok{(readr)}
\FunctionTok{source}\NormalTok{(}\StringTok{"R/dgp\_cohorte.R"}\NormalTok{)                        }\CommentTok{\# define simular\_cohorte(), cargar\_cohorte(), expandir\_persona\_periodo()}
\NormalTok{cohorte }\OtherTok{\textless{}{-}} \FunctionTok{cargar\_cohorte}\NormalTok{()                      }\CommentTok{\# lee datos/cohorte\_*.rds si existe; si no (o si cambió el DGP), simula y guarda}
\CommentTok{\# y guardamos las simulaciones}
\CommentTok{\#write\_csv(cohorte, "cohorte.csv")        \# nivel individuo}
\CommentTok{\# resumen}
\FunctionTok{glimpse}\NormalTok{(cohorte)}
\end{Highlighting}
\end{Shaded}

\begin{verbatim}
Rows: 877
Columns: 9
$ id        <int> 1, 2, 3, 4, 5, 6, 7, 8, 9, 10, 11, 12, 13, 14, 15, 16, 17, 1~
$ centro    <fct> 1, 1, 1, 1, 1, 1, 1, 1, 1, 1, 1, 1, 1, 1, 1, 1, 1, 1, 1, 1, ~
$ x1        <dbl> -1.67363040, 1.09293799, -0.35266071, 0.29762674, -0.3236454~
$ x2        <int> 0, 1, 1, 0, 1, 1, 1, 0, 1, 1, 0, 0, 0, 0, 1, 1, 0, 0, 1, 1, ~
$ tiempo    <int> 4, 3, 6, 4, 6, 4, 6, 3, 4, 5, 6, 5, 4, 1, 6, 4, 4, 5, 3, 6, ~
$ evento    <int> 1, 1, 0, 0, 0, 0, 0, 1, 0, 0, 0, 0, 0, 1, 1, 0, 0, 0, 0, 0, ~
$ ever      <int> 1, 1, 0, 0, 0, 0, 0, 1, 0, 0, 0, 0, 0, 1, 1, 0, 0, 0, 0, 0, ~
$ clase_nom <fct> A, B, A, C, A, A, C, A, A, B, A, A, B, B, A, C, C, A, C, B, ~
$ sever_ord <ord> Leve, Grave, Moderado, Moderado, Leve, Leve, Leve, Grave, Mo~
\end{verbatim}

\begin{Shaded}
\begin{Highlighting}[]
\FunctionTok{head}\NormalTok{(cohorte)}
\end{Highlighting}
\end{Shaded}

\begin{verbatim}
  id centro         x1 x2 tiempo evento ever clase_nom sever_ord
1  1      1 -1.6736304  0      4      1    1         A      Leve
2  2      1  1.0929380  1      3      1    1         B     Grave
3  3      1 -0.3526607  1      6      0    0         A  Moderado
4  4      1  0.2976267  0      4      0    0         C  Moderado
5  5      1 -0.3236454  1      6      0    0         A      Leve
6  6      1 -1.6879930  1      4      0    0         A      Leve
\end{verbatim}

Más tarde, para trabajar la supervivencia (Unidad 1.5), transformaremos
esta \texttt{cohorte} a un formato \textbf{persona-periodo} ---una fila
por paciente y periodo en riesgo---; esa expansión se construye y se
explica allí, donde se usa.

Visualizamos en la Figura~\ref{fig-cohorte-eda} los datos simulados en
la \texttt{cohorte}.

\begin{Shaded}
\begin{Highlighting}[]
\CommentTok{\# Preparo factores legibles (x2, tiempo y el desenlace como categóricos)}
\NormalTok{cohorte\_eda }\OtherTok{\textless{}{-}}\NormalTok{ cohorte }\SpecialCharTok{|\textgreater{}}
  \FunctionTok{mutate}\NormalTok{(}
    \AttributeTok{evento =} \FunctionTok{factor}\NormalTok{(evento, }\AttributeTok{levels =} \FunctionTok{c}\NormalTok{(}\DecValTok{0}\NormalTok{, }\DecValTok{1}\NormalTok{), }\AttributeTok{labels =} \FunctionTok{c}\NormalTok{(}\StringTok{"Sin fractura"}\NormalTok{, }\StringTok{"Fractura"}\NormalTok{)),}
    \AttributeTok{x2     =} \FunctionTok{factor}\NormalTok{(x2,     }\AttributeTok{levels =} \FunctionTok{c}\NormalTok{(}\DecValTok{0}\NormalTok{, }\DecValTok{1}\NormalTok{), }\AttributeTok{labels =} \FunctionTok{c}\NormalTok{(}\StringTok{"No"}\NormalTok{, }\StringTok{"Sí"}\NormalTok{)),}
    \AttributeTok{tiempo =} \FunctionTok{factor}\NormalTok{(tiempo, }\AttributeTok{levels =} \DecValTok{1}\SpecialCharTok{:}\DecValTok{6}\NormalTok{)}
\NormalTok{  )}

\CommentTok{\# Variables sustantivas (fuera \textquotesingle{}id\textquotesingle{}, \textquotesingle{}centro\textquotesingle{} y \textquotesingle{}ever\textquotesingle{} = duplicado de \textquotesingle{}evento\textquotesingle{})}
\NormalTok{vars }\OtherTok{\textless{}{-}} \FunctionTok{c}\NormalTok{(}\StringTok{"x1"}\NormalTok{, }\StringTok{"x2"}\NormalTok{, }\StringTok{"tiempo"}\NormalTok{, }\StringTok{"clase\_nom"}\NormalTok{, }\StringTok{"sever\_ord"}\NormalTok{)}

\NormalTok{cohorte\_eda }\SpecialCharTok{|\textgreater{}}
  \FunctionTok{select}\NormalTok{(evento, }\FunctionTok{all\_of}\NormalTok{(vars)) }\SpecialCharTok{|\textgreater{}}
  \FunctionTok{ggpairs}\NormalTok{(}
    \AttributeTok{columns =}\NormalTok{ vars,}
    \AttributeTok{mapping =} \FunctionTok{aes}\NormalTok{(}\AttributeTok{color =}\NormalTok{ evento, }\AttributeTok{alpha =} \FloatTok{0.6}\NormalTok{),}
    \AttributeTok{legend=}\DecValTok{1}\NormalTok{,}
    \AttributeTok{upper =} \FunctionTok{list}\NormalTok{(}\AttributeTok{continuous =} \FunctionTok{wrap}\NormalTok{(}\StringTok{"cor"}\NormalTok{, }\AttributeTok{size =} \FloatTok{2.5}\NormalTok{),}
                 \AttributeTok{combo      =} \FunctionTok{wrap}\NormalTok{(}\StringTok{"box\_no\_facet"}\NormalTok{),}
                 \AttributeTok{discrete   =} \FunctionTok{wrap}\NormalTok{(}\StringTok{"count"}\NormalTok{)),}
    \AttributeTok{lower =} \FunctionTok{list}\NormalTok{(}\AttributeTok{continuous =} \FunctionTok{wrap}\NormalTok{(}\StringTok{"points"}\NormalTok{, }\AttributeTok{alpha =} \FloatTok{0.25}\NormalTok{, }\AttributeTok{size =} \FloatTok{0.4}\NormalTok{),}
                 \AttributeTok{combo      =} \FunctionTok{wrap}\NormalTok{(}\StringTok{"facethist"}\NormalTok{, }\AttributeTok{bins =} \DecValTok{20}\NormalTok{),}
                 \AttributeTok{discrete   =} \FunctionTok{wrap}\NormalTok{(}\StringTok{"facetbar"}\NormalTok{)),}
    \AttributeTok{diag  =} \FunctionTok{list}\NormalTok{(}\AttributeTok{continuous =} \FunctionTok{wrap}\NormalTok{(}\StringTok{"densityDiag"}\NormalTok{, }\AttributeTok{alpha =} \FloatTok{0.5}\NormalTok{),}
                 \AttributeTok{discrete   =} \FunctionTok{wrap}\NormalTok{(}\StringTok{"barDiag"}\NormalTok{)),}
    \AttributeTok{progress =} \ConstantTok{FALSE}
\NormalTok{  ) }\SpecialCharTok{+}
  \FunctionTok{theme\_bw}\NormalTok{(}\AttributeTok{base\_size =} \DecValTok{8}\NormalTok{)}\SpecialCharTok{+}
  \FunctionTok{theme}\NormalTok{(}\AttributeTok{legend.position =} \StringTok{"bottom"}\NormalTok{)}
\end{Highlighting}
\end{Shaded}

\begin{figure}[H]

{\centering \includegraphics[width=1\textwidth,height=\textheight]{caso1_binaria_files/figure-pdf/fig-cohorte-eda-1.pdf}

}

\caption{\label{fig-cohorte-eda}Descripción de \texttt{cohorte}. En la
diagonal distribución de las variables, diferenciadas en color por
\texttt{evento}. En los cruces la relación entre ellas.}

\end{figure}

\hypertarget{sec-u11}{%
\section{1 · Cuando la recta no llega: del dato binario al marco
GLM}\label{sec-u11}}

Iniciamos nuestro curso descubriendo las limitaciones del modelo lineal
cuando la respuesta ni es normal, ni hay modo de transformarla para
hacerla lineal, como ocurre cuando tenemos una variable dicotómica por
predecir. Nos apoyamos en la base de datos GLOW para ilustrar dónde
falla el modelo lineal, y mostrar las bondades que reporta un modelo
lineal generalizado, en el que se tiene en cuenta la distribución
adecuada para la respuesta, y se busca una transformación del predictor
lineal que pueda generar predicciones viables sin salir de un modelo
lineal. Comparamos en esta sección, los resultados de un modelo lineal y
los de un modelo lineal generalizado, para fundamentar el desarrollo de
este modelo en la siguiente sección.

\begin{tcolorbox}[enhanced jigsaw, coltitle=black, rightrule=.15mm, arc=.35mm, leftrule=.75mm, colback=white, colbacktitle=quarto-callout-note-color!10!white, left=2mm, breakable, toptitle=1mm, toprule=.15mm, bottomtitle=1mm, titlerule=0mm, title=\textcolor{quarto-callout-note-color}{\faInfo}\hspace{0.5em}{Sobre esta unidad}, bottomrule=.15mm, opacityback=0, opacitybacktitle=0.6, colframe=quarto-callout-note-color-frame]

RA implicados:

\begin{itemize}
\tightlist
\item
  \textbf{RA1 · Marco GLM frente a OLS.} Reconocer por qué el modelo
  lineal (OLS) falla ante respuestas discretas y describir los tres
  componentes de un GLM (aleatorio, predictor lineal y enlace).
\item
  \textbf{RA2 · Formulación y estimación} (primer contacto). Formular y
  estimar un GLM como generalización del modelo lineal.
\end{itemize}

\end{tcolorbox}

\hypertarget{el-modelo-lineal-en-glow}{%
\subsection{1.1 El modelo lineal en
GLOW}\label{el-modelo-lineal-en-glow}}

La primera tarea antes de proponer una modelización, pasa por observar
el comportamiento de los datos. Nos interesan, de momento, las variables
\texttt{fracture}, que actúa como respuesta e indica si la paciente ha
sufrido una fractura, \texttt{age}, con la edad de la paciente, y
\texttt{priorfrac}, que indica si tiene antecedentes de haber sufrido
alguna fractura previa. Indagamos pues, cómo se relaciona la respuesta
con la edad, distinguiendo por antecedente de fractura previa (en la
Figura~\ref{fig-glow-eda-box}):

\begin{Shaded}
\begin{Highlighting}[]
\NormalTok{glow }\SpecialCharTok{|\textgreater{}}
  \FunctionTok{ggplot}\NormalTok{(}\FunctionTok{aes}\NormalTok{(}\AttributeTok{x =}\NormalTok{ fracture, }\AttributeTok{y =}\NormalTok{ age, }\AttributeTok{fill =}\NormalTok{ fracture)) }\SpecialCharTok{+}
  \FunctionTok{geom\_boxplot}\NormalTok{(}\AttributeTok{alpha =} \FloatTok{0.75}\NormalTok{, }\AttributeTok{outlier.alpha =} \FloatTok{0.3}\NormalTok{) }\SpecialCharTok{+}
  \FunctionTok{facet\_wrap}\NormalTok{(}\SpecialCharTok{\textasciitilde{}}\NormalTok{ priorfrac,}
             \AttributeTok{labeller =} \FunctionTok{labeller}\NormalTok{(}\AttributeTok{priorfrac =} \FunctionTok{c}\NormalTok{(}\AttributeTok{No =} \StringTok{"Sin fractura previa"}\NormalTok{,}
                                                \AttributeTok{Yes =} \StringTok{"Con fractura previa"}\NormalTok{))) }\SpecialCharTok{+}
  \FunctionTok{labs}\NormalTok{(}\AttributeTok{x =} \StringTok{"Fractura en el primer año"}\NormalTok{, }\AttributeTok{y =} \StringTok{"Edad (años)"}\NormalTok{) }\SpecialCharTok{+}
  \FunctionTok{guides}\NormalTok{(}\AttributeTok{fill =} \StringTok{"none"}\NormalTok{)}
\end{Highlighting}
\end{Shaded}

\begin{figure}[H]

{\centering \includegraphics[width=1\textwidth,height=\textheight]{caso1_binaria_files/figure-pdf/fig-glow-eda-box-1.pdf}

}

\caption{\label{fig-glow-eda-box}Distribución de la edad según haya
habido o no fractura en el primer año, condicionada por el antecedente
de fractura previa (priorfrac).}

\end{figure}

Observamos que las mujeres que \textbf{sí} fracturan son, en mediana,
\textbf{algo mayores} que las que no, en los dos paneles: la edad es un
\textbf{factor de riesgo}. Pero la separación es \textbf{modesta} y las
cajas se \textbf{solapan} bastante ---la edad por sí sola no discrimina
bien, de ahí que el efecto lineal de la edad resultará débil, como
veremos más tarde---. Además, el \textbf{antecedente de fractura previa}
(panel derecho) eleva el riesgo por sí mismo ---hay fracturas a
cualquier edad---, así que conviene un modelo con \textbf{ambas}
covariables, no una recta única sobre la edad.

Tiene sentido pues, plantear un modelo predictivo para \texttt{fracture}
en función de \texttt{age} y \texttt{priorfrac}. Para intentar capturar
el efecto de la edad (\texttt{age}) sobre la probabilidad de fractura (y
valorar), agrupamos la edad por franjas y calculamos la proporción de
fracturas en cada franja, diferenciando por el hecho de que la paciente
hubiera sufrido o no, una fractura previa (\texttt{priorfrac}). En la
Figura~\ref{fig-u11-eda} tenemos esta representación, y superponemos
sendas rectas, que transmiten una considerable variabilidad de los
puntos en torno a ellas, especialmente en las pacientes que han sufrido
fractura.

\begin{Shaded}
\begin{Highlighting}[]
\NormalTok{glow }\SpecialCharTok{|\textgreater{}}
  \FunctionTok{mutate}\NormalTok{(}\AttributeTok{edad\_grupo =} \FunctionTok{cut}\NormalTok{(age, }\AttributeTok{breaks =} \FunctionTok{seq}\NormalTok{(}\DecValTok{55}\NormalTok{, }\DecValTok{95}\NormalTok{, }\AttributeTok{by =} \DecValTok{5}\NormalTok{))) }\SpecialCharTok{|\textgreater{}}
  \FunctionTok{filter}\NormalTok{(}\SpecialCharTok{!}\FunctionTok{is.na}\NormalTok{(edad\_grupo)) }\SpecialCharTok{|\textgreater{}}
  \FunctionTok{group\_by}\NormalTok{(edad\_grupo, priorfrac) }\SpecialCharTok{|\textgreater{}}
  \FunctionTok{summarise}\NormalTok{(}\AttributeTok{age =} \FunctionTok{mean}\NormalTok{(age), }\AttributeTok{p\_fractura =} \FunctionTok{mean}\NormalTok{(fractura01), }\AttributeTok{n =} \FunctionTok{n}\NormalTok{(), }\AttributeTok{.groups =} \StringTok{"drop"}\NormalTok{) }\SpecialCharTok{|\textgreater{}}
  \FunctionTok{ggplot}\NormalTok{(}\FunctionTok{aes}\NormalTok{(age, p\_fractura, }\AttributeTok{color =}\NormalTok{ priorfrac)) }\SpecialCharTok{+}
  \CommentTok{\#geom\_line() + }
  \FunctionTok{geom\_point}\NormalTok{(}\FunctionTok{aes}\NormalTok{(}\AttributeTok{size =}\NormalTok{ n)) }\SpecialCharTok{+}
  \FunctionTok{geom\_smooth}\NormalTok{(}\FunctionTok{aes}\NormalTok{(}\AttributeTok{weight =}\NormalTok{ n), }\AttributeTok{method =} \StringTok{"lm"}\NormalTok{, }\AttributeTok{se =} \ConstantTok{FALSE}\NormalTok{, }\AttributeTok{linewidth =} \DecValTok{1}\NormalTok{) }\SpecialCharTok{+}
  \FunctionTok{labs}\NormalTok{(}\AttributeTok{x =} \StringTok{"Edad (media del tramo)"}\NormalTok{, }\AttributeTok{y =} \StringTok{"Proporción de fractura"}\NormalTok{,}
       \AttributeTok{color =} \StringTok{"Fractura previa"}\NormalTok{, }\AttributeTok{size =} \StringTok{"n"}\NormalTok{)}
\end{Highlighting}
\end{Shaded}

\begin{figure}[H]

{\centering \includegraphics{caso1_binaria_files/figure-pdf/fig-u11-eda-1.pdf}

}

\caption{\label{fig-u11-eda}Proporción observada de fractura por tramos
de edad y antecedente previo, con sendas rectas de regresión por grupo:
el ajuste lineal es pobre y apenas captura tendencia.}

\end{figure}

Planteamos pues, una primera aproximación a través de un modelo lineal
(sin interacción) en el que la respuesta ha de ser numérica, por lo que
consideramos la recodificación de \texttt{fracture} con valores 0/1
-\texttt{fracture01}-, y los predictores son \texttt{age} y
\texttt{priorfrac}.

\begin{verbatim}
fit_ols = lm(fractura01 ~ age + priorfrac, data = glow)
\end{verbatim}

Y comenzamos a identificar los problemas que van apareciendo.

En primer lugar, un modelo lineal no garantiza una predicción de la
respuesta dentro de un rango, y en particular como en este caso, dentro
del rango \((0,1)\) de probabilidades. Aunque para el tramo observado de
edad el modelo lineal se comporta bien a este respecto, dando
predicciones dentro del rango \((0,1)\), vemos en la
Figura~\ref{fig-u11-ols-extrapola} cómo la recta de ajuste da
predicciones por debajo de 0 para edades inferiores a 40 años.

\begin{Shaded}
\begin{Highlighting}[]
\NormalTok{fit\_ols }\OtherTok{\textless{}{-}} \FunctionTok{lm}\NormalTok{(fractura01 }\SpecialCharTok{\textasciitilde{}}\NormalTok{ age }\SpecialCharTok{+}\NormalTok{ priorfrac, }\AttributeTok{data =}\NormalTok{ glow)}
\NormalTok{rango }\OtherTok{\textless{}{-}} \FunctionTok{range}\NormalTok{(glow}\SpecialCharTok{$}\NormalTok{age)}

\NormalTok{grid\_ext }\OtherTok{\textless{}{-}} \FunctionTok{expand\_grid}\NormalTok{(}
  \AttributeTok{age       =} \FunctionTok{seq}\NormalTok{(}\DecValTok{35}\NormalTok{, }\DecValTok{110}\NormalTok{, }\AttributeTok{length.out =} \DecValTok{200}\NormalTok{),}
  \AttributeTok{priorfrac =} \FunctionTok{factor}\NormalTok{(}\FunctionTok{c}\NormalTok{(}\StringTok{"No"}\NormalTok{, }\StringTok{"Yes"}\NormalTok{), }\AttributeTok{levels =} \FunctionTok{levels}\NormalTok{(glow}\SpecialCharTok{$}\NormalTok{priorfrac)))}
\NormalTok{grid\_ext}\SpecialCharTok{$}\NormalTok{p\_hat }\OtherTok{\textless{}{-}} \FunctionTok{predict}\NormalTok{(fit\_ols, }\AttributeTok{newdata =}\NormalTok{ grid\_ext)}

\FunctionTok{ggplot}\NormalTok{(grid\_ext, }\FunctionTok{aes}\NormalTok{(age, p\_hat, }\AttributeTok{color =}\NormalTok{ priorfrac)) }\SpecialCharTok{+}
  \FunctionTok{annotate}\NormalTok{(}\StringTok{"rect"}\NormalTok{, }\AttributeTok{xmin =}\NormalTok{ rango[}\DecValTok{1}\NormalTok{], }\AttributeTok{xmax =}\NormalTok{ rango[}\DecValTok{2}\NormalTok{], }\AttributeTok{ymin =} \SpecialCharTok{{-}}\ConstantTok{Inf}\NormalTok{, }\AttributeTok{ymax =} \ConstantTok{Inf}\NormalTok{,}
           \AttributeTok{alpha =} \FloatTok{0.12}\NormalTok{, }\AttributeTok{fill =} \StringTok{"grey40"}\NormalTok{) }\SpecialCharTok{+}
  \FunctionTok{geom\_line}\NormalTok{(}\AttributeTok{linewidth =} \DecValTok{1}\NormalTok{) }\SpecialCharTok{+}
  \FunctionTok{geom\_hline}\NormalTok{(}\AttributeTok{yintercept =} \FunctionTok{c}\NormalTok{(}\DecValTok{0}\NormalTok{, }\DecValTok{1}\NormalTok{), }\AttributeTok{linetype =} \StringTok{"dashed"}\NormalTok{) }\SpecialCharTok{+}
  \FunctionTok{labs}\NormalTok{(}\AttributeTok{x =} \StringTok{"Edad (rango extrapolado; sombreado = rango observado)"}\NormalTok{,}
       \AttributeTok{y =} \StringTok{"Probabilidad predicha (OLS)"}\NormalTok{, }\AttributeTok{color =} \StringTok{"Fractura previa"}\NormalTok{)}
\end{Highlighting}
\end{Shaded}

\begin{figure}[H]

{\centering \includegraphics{caso1_binaria_files/figure-pdf/fig-u11-ols-extrapola-1.pdf}

}

\caption{\label{fig-u11-ols-extrapola}Predicción del modelo lineal
extrapolada fuera del rango de edad observado (franja sombreada). Más
allá de los datos, OLS predice probabilidades por debajo de 0:
imposibles para una probabilidad.}

\end{figure}

Otro fallo se da en el diagnóstico de los residuos, como vemos en la
Figura~\ref{fig-u11-diagnostico-ols}: a la izquierda apreciamos que los
residuos, para un x dado, toman solo dos valores posibles, y se mueven
con la media ajustada; se rompe la independencia de los residuos y la
homocedasticidad. En el gráfico a la derecha se muestran los residuos
versus los percentiles de una distribución normal, y de nuevo se aprecia
la ruptura de la hipótesis de normalidad. Esto ya era previsible desde
el inicio, puesto que la respuesta es Bernouilli y no Normal.

\begin{Shaded}
\begin{Highlighting}[]
\NormalTok{fit\_ols }\OtherTok{\textless{}{-}} \FunctionTok{lm}\NormalTok{(fractura01 }\SpecialCharTok{\textasciitilde{}}\NormalTok{ age }\SpecialCharTok{+}\NormalTok{ priorfrac, }\AttributeTok{data =}\NormalTok{ glow)}

\NormalTok{g\_res }\OtherTok{\textless{}{-}} \FunctionTok{tibble}\NormalTok{(}\AttributeTok{ajustado =} \FunctionTok{fitted}\NormalTok{(fit\_ols),}
                \AttributeTok{residuo  =}\NormalTok{ glow}\SpecialCharTok{$}\NormalTok{fractura01 }\SpecialCharTok{{-}} \FunctionTok{fitted}\NormalTok{(fit\_ols),}
                \AttributeTok{y =} \FunctionTok{factor}\NormalTok{(glow}\SpecialCharTok{$}\NormalTok{fractura01, }\AttributeTok{labels =} \FunctionTok{c}\NormalTok{(}\StringTok{"No fractura (y=0)"}\NormalTok{, }\StringTok{"Fractura (y=1)"}\NormalTok{))) }\SpecialCharTok{|\textgreater{}}
  \FunctionTok{ggplot}\NormalTok{(}\FunctionTok{aes}\NormalTok{(ajustado, residuo, }\AttributeTok{color =}\NormalTok{ y)) }\SpecialCharTok{+}
  \FunctionTok{geom\_point}\NormalTok{(}\AttributeTok{alpha =} \FloatTok{0.3}\NormalTok{) }\SpecialCharTok{+}
  \FunctionTok{geom\_hline}\NormalTok{(}\AttributeTok{yintercept =} \DecValTok{0}\NormalTok{, }\AttributeTok{linetype =} \StringTok{"dashed"}\NormalTok{) }\SpecialCharTok{+}
  \FunctionTok{labs}\NormalTok{(}\AttributeTok{x =} \StringTok{"Valores ajustados"}\NormalTok{, }\AttributeTok{y =} \StringTok{"Residuo de respuesta (y − μ̂)"}\NormalTok{,}
       \AttributeTok{title =} \StringTok{"Residuos vs ajustados (OLS)"}\NormalTok{, }\AttributeTok{color =} \ConstantTok{NULL}\NormalTok{)}\SpecialCharTok{+}
  \FunctionTok{theme}\NormalTok{(}\AttributeTok{legend.position=}\StringTok{"bottom"}\NormalTok{)}

\NormalTok{g\_qplot }\OtherTok{\textless{}{-}} \FunctionTok{tibble}\NormalTok{(}\AttributeTok{residuo =} \FunctionTok{residuals}\NormalTok{(fit\_ols)) }\SpecialCharTok{|\textgreater{}}
  \FunctionTok{ggplot}\NormalTok{(}\FunctionTok{aes}\NormalTok{(}\AttributeTok{sample =}\NormalTok{ residuo)) }\SpecialCharTok{+}
  \FunctionTok{stat\_qq}\NormalTok{(}\AttributeTok{alpha =} \FloatTok{0.4}\NormalTok{, }\AttributeTok{size =} \FloatTok{0.6}\NormalTok{) }\SpecialCharTok{+}
  \FunctionTok{stat\_qq\_line}\NormalTok{(}\AttributeTok{color =} \StringTok{"firebrick"}\NormalTok{) }\SpecialCharTok{+}
  \FunctionTok{labs}\NormalTok{(}\AttributeTok{x =} \StringTok{"Cuantiles teóricos (Normal)"}\NormalTok{, }\AttributeTok{y =} \StringTok{"Cuantiles de los residuos de OLS"}\NormalTok{)}

\NormalTok{patchwork}\SpecialCharTok{::}\FunctionTok{wrap\_plots}\NormalTok{(g\_res, g\_qplot, }\AttributeTok{ncol =} \DecValTok{2}\NormalTok{)}
\end{Highlighting}
\end{Shaded}

\begin{figure}[H]

{\centering \includegraphics[width=1\textwidth,height=\textheight]{caso1_binaria_files/figure-pdf/fig-u11-diagnostico-ols-1.pdf}

}

\caption{\label{fig-u11-diagnostico-ols}Diagnóstico del ajuste OLS sobre
respuesta binaria. Izquierda: residuos frente a ajustados (no hay nube
de puntos aleatoria en torno al cero). Derecha: distribución de los
residuos frente a la Normal de igual media y desviación (fallo de
normalidad).}

\end{figure}

\hypertarget{sec-u11-ols-falla}{%
\subsection{1.2 Por qué OLS no sirve para una respuesta
binaria}\label{sec-u11-ols-falla}}

El modelo OLS postula
\(Y = \beta_0 + \beta_1 x_1 + \dots + \varepsilon\) con
\(\varepsilon \sim \mathcal{N}(0,\sigma^2)\). Como \(Y\in\{0,1\}\), su
esperanza condicional es directamente la probabilidad de éxito,
\(\mathbb{E}[Y\mid x]=\Pr(Y=1\mid x)=\pi(x)\): OLS estaría modelando una
\textbf{probabilidad como función lineal}. Cuatro problemas qurgen todos
ellos críticos (Dobson \& Barnett, 2018, caps. 1--3; Fox, 2016, 14):

\begin{enumerate}
\def\labelenumi{\arabic{enumi}.}
\tightlist
\item
  \textbf{Predicciones fuera de rango.} El predictor lineal recorre toda
  la recta real, pero una probabilidad vive en \([0,1]\). Para valores
  extremos de las covariables, OLS predice probabilidades
  \textbf{negativas o mayores que 1}.
\item
  \textbf{Varianza no constante.} Si \(Y\) es Bernoulli con media
  \(\pi(x)\), entonces \(\operatorname{Var}(Y\mid x)=\pi(x)[1-\pi(x)]\):
  la varianza \textbf{depende de la media}. El supuesto de
  homocedasticidad se viola por construcción.
\item
  \textbf{Error no normal.} Fijado un \(x\) concreto, \(\pi(x)\) es una
  \textbf{constante} e \(Y\in\{0,1\}\), así que \(\varepsilon=Y-\pi(x)\)
  solo puede tomar \textbf{dos valores}: \(-\pi(x)\) con probabilidad
  \(1-\pi(x)\), o \(1-\pi(x)\) con probabilidad \(\pi(x)\). La
  distribución \textbf{condicional} del error (dado \(x\)) es de
  \textbf{dos puntos}, no una campana gaussiana. Al variar \(x\) esos
  dos valores se desplazan ---de ahí las dos bandas de residuos del
  laboratorio---, pero el supuesto de normalidad de OLS recae sobre la
  distribución condicional, que es la que aquí es falsa.
\item
  \textbf{Linealidad implausible.} OLS impone que el efecto de una
  covariable sobre la probabilidad sea \textbf{constante}: la misma
  pendiente \(\beta\) tanto si \(\pi\) ronda \(0{,}5\) como si ya está
  pegada a 0 o a 1. Eso no es creíble, porque en los extremos apenas
  queda margen para moverse: si una paciente tiene una probabilidad de
  fractura casi nula (o casi segura), un año más de edad casi no puede
  cambiarla; en cambio, en la zona central ese mismo año podría
  desplazar bastante la probabilidad. Lo razonable es entonces una
  relación \textbf{en forma de S}: pronunciada en el centro y
  \textbf{aplanada (saturada) en los extremos}. La inversa del logit
  dibuja exactamente esa curva y, de paso, mantiene \(\pi\) dentro de
  \((0,1)\).
\end{enumerate}

\hypertarget{sec-u11-tres-componentes}{%
\subsection{1.3 Los tres componentes de un
GLM}\label{sec-u11-tres-componentes}}

La alternativa que plantearon Nelder \& Wedderburn (1972) unificó la
regresión lineal, logística, de Poisson, etc., bajo una estructura de
\textbf{tres piezas}:

\begin{enumerate}
\def\labelenumi{\arabic{enumi}.}
\item
  \textbf{Componente aleatoria.} \(Y\) sigue una distribución de la
  \textbf{familia exponencial}.

  \begin{quote}
  Para GLOW, \(Y\mid x \sim \text{Bernoulli}(\pi(x))\). Esta pieza fija
  la relación media--varianza y absorbe los problemas 2 y 3.
  \end{quote}
\item
  \textbf{Componente sistemática.} Mantiene un \textbf{predictor lineal}
  \(\eta=\mathbf{x}^\top \boldsymbol\beta\), lineal en los parámetros.
\item
  Introduce una \textbf{función de enlace} o \emph{link} \(g\), monótona
  y diferenciable, que liga la media con el predictor lineal,
  \(g(\mu)=\eta\) y \(\mu=g^{-1}(\eta)\).

  \begin{quote}
  Para binaria, el enlace por defecto es el \textbf{logit}: \[
  g(\pi)=\log\frac{\pi}{1-\pi}, \qquad \pi=g^{-1}(\eta)=\frac{e^{\eta}}{1+e^{\eta}},
  \] que resuelve los problemas 1 y 4: el logit lleva \((0,1)\) a toda
  la recta real (nunca fuera de rango).
  \end{quote}
\end{enumerate}

En la Figura~\ref{fig-u11-logit-inversa} se muestra cómo la función
logit transforma una probabilidad cualquiera en el rango \([0,1]\) en un
valor sobre toda la recta real. Su inversa, la sigmoide, a la izquierda,
permite recuperar un valor real en un rango de probabilidades, y se
utilizará para estimar la probabilidad a partir del predictor lineal.

\begin{Shaded}
\begin{Highlighting}[]
\FunctionTok{library}\NormalTok{(patchwork)}
\NormalTok{g1 }\OtherTok{\textless{}{-}} \FunctionTok{tibble}\NormalTok{(}\AttributeTok{p =} \FunctionTok{seq}\NormalTok{(}\FloatTok{0.001}\NormalTok{, }\FloatTok{0.999}\NormalTok{, }\AttributeTok{by =} \FloatTok{0.001}\NormalTok{)) }\SpecialCharTok{|\textgreater{}}
  \FunctionTok{mutate}\NormalTok{(}\AttributeTok{y =} \FunctionTok{qlogis}\NormalTok{(p)) }\SpecialCharTok{|\textgreater{}}
  \FunctionTok{ggplot}\NormalTok{(}\FunctionTok{aes}\NormalTok{(p, y)) }\SpecialCharTok{+}
  \FunctionTok{geom\_line}\NormalTok{(}\AttributeTok{linewidth =} \DecValTok{1}\NormalTok{, }\AttributeTok{color =} \StringTok{"steelblue"}\NormalTok{) }\SpecialCharTok{+}
  \FunctionTok{geom\_hline}\NormalTok{(}\AttributeTok{yintercept =} \DecValTok{0}\NormalTok{, }\AttributeTok{linetype =} \StringTok{"dashed"}\NormalTok{) }\SpecialCharTok{+}
  \FunctionTok{labs}\NormalTok{(}\AttributeTok{title =} \StringTok{"logit:}\SpecialCharTok{\textbackslash{}n}\StringTok{ estira (0,1) sobre la recta real"}\NormalTok{,}
       \AttributeTok{x =} \FunctionTok{expression}\NormalTok{(pi), }\AttributeTok{y =} \FunctionTok{expression}\NormalTok{(}\FunctionTok{log}\NormalTok{(pi}\SpecialCharTok{/}\NormalTok{(}\DecValTok{1} \SpecialCharTok{{-}}\NormalTok{ pi))))}

\NormalTok{g2 }\OtherTok{\textless{}{-}} \FunctionTok{tibble}\NormalTok{(}\AttributeTok{eta =} \FunctionTok{seq}\NormalTok{(}\SpecialCharTok{{-}}\DecValTok{6}\NormalTok{, }\DecValTok{6}\NormalTok{, }\AttributeTok{by =} \FloatTok{0.05}\NormalTok{)) }\SpecialCharTok{|\textgreater{}}
  \FunctionTok{mutate}\NormalTok{(}\AttributeTok{p =} \FunctionTok{plogis}\NormalTok{(eta)) }\SpecialCharTok{|\textgreater{}}
  \FunctionTok{ggplot}\NormalTok{(}\FunctionTok{aes}\NormalTok{(eta, p)) }\SpecialCharTok{+}
  \FunctionTok{geom\_line}\NormalTok{(}\AttributeTok{linewidth =} \DecValTok{1}\NormalTok{, }\AttributeTok{color =} \StringTok{"firebrick"}\NormalTok{) }\SpecialCharTok{+}
  \FunctionTok{geom\_hline}\NormalTok{(}\AttributeTok{yintercept =} \FunctionTok{c}\NormalTok{(}\DecValTok{0}\NormalTok{, }\DecValTok{1}\NormalTok{), }\AttributeTok{linetype =} \StringTok{"dotted"}\NormalTok{) }\SpecialCharTok{+}
  \FunctionTok{labs}\NormalTok{(}\AttributeTok{title =} \StringTok{"sigmoide:}\SpecialCharTok{\textbackslash{}n}\StringTok{ devuelve el predictor a (0,1)"}\NormalTok{,}
       \AttributeTok{x =} \FunctionTok{expression}\NormalTok{(eta), }\AttributeTok{y =} \FunctionTok{expression}\NormalTok{(e}\SpecialCharTok{\^{}}\NormalTok{eta}\SpecialCharTok{/}\NormalTok{(}\DecValTok{1} \SpecialCharTok{+}\NormalTok{ e}\SpecialCharTok{\^{}}\NormalTok{eta)))}

\NormalTok{g1 }\SpecialCharTok{+}\NormalTok{ g2}
\end{Highlighting}
\end{Shaded}

\begin{figure}[H]

{\centering \includegraphics{caso1_binaria_files/figure-pdf/fig-u11-logit-inversa-1.pdf}

}

\caption{\label{fig-u11-logit-inversa}Función logit y sigmoide: las dos
direcciones. El logit lleva (0,1) a toda la recta real; su inversa (la
sigmoide) devuelve cualquier η a (0,1).}

\end{figure}

\begin{tcolorbox}[enhanced jigsaw, coltitle=black, rightrule=.15mm, arc=.35mm, leftrule=.75mm, colback=white, colbacktitle=quarto-callout-important-color!10!white, left=2mm, breakable, toptitle=1mm, toprule=.15mm, bottomtitle=1mm, titlerule=0mm, title=\textcolor{quarto-callout-important-color}{\faExclamation}\hspace{0.5em}{OLS es un GLM}, bottomrule=.15mm, opacityback=0, opacitybacktitle=0.6, colframe=quarto-callout-important-color-frame]

La regresión lineal es un modelo lineal generalizado con la familia
Gaussiana y la función de enlace identidad. No abandonamos OLS: lo
integramos en los GLM. Por eso \texttt{lm(y\ \textasciitilde{}\ x)} y
\texttt{glm(y\ \textasciitilde{}\ x)} se parecen tanto.

\end{tcolorbox}

En esta unidad trabajaremos con:

\begin{longtable}[]{@{}
  >{\raggedright\arraybackslash}p{(\columnwidth - 8\tabcolsep) * \real{0.2000}}
  >{\raggedright\arraybackslash}p{(\columnwidth - 8\tabcolsep) * \real{0.2000}}
  >{\raggedright\arraybackslash}p{(\columnwidth - 8\tabcolsep) * \real{0.2000}}
  >{\raggedright\arraybackslash}p{(\columnwidth - 8\tabcolsep) * \real{0.2000}}
  >{\raggedright\arraybackslash}p{(\columnwidth - 8\tabcolsep) * \real{0.2000}}@{}}
\toprule()
\begin{minipage}[b]{\linewidth}\raggedright
Tipo de respuesta
\end{minipage} & \begin{minipage}[b]{\linewidth}\raggedright
Ejemplo
\end{minipage} & \begin{minipage}[b]{\linewidth}\raggedright
Familia (por defecto)
\end{minipage} & \begin{minipage}[b]{\linewidth}\raggedright
Enlace canónico
\end{minipage} & \begin{minipage}[b]{\linewidth}\raggedright
Unidad
\end{minipage} \\
\midrule()
\endhead
Binaria & fractura sí/no & Binomial & logit & Unidades
\protect\hyperlink{sec-u11}{1} y \protect\hyperlink{sec-u12}{2} \\
Binomial & nº fracturas por edad & Binomial & logit &
\protect\hyperlink{sec-u13}{Unidad 3} \\
Politómica nominal & perfil A/B/C & Multinomial & logit (gen.) &
\protect\hyperlink{sec-u13}{Unidad 3} \\
Politómica ordinal & severidad Leve/Mod/Grave & Multinomial ordinal &
logit acumulado & \protect\hyperlink{sec-u13}{Unidad 3} \\
\bottomrule()
\end{longtable}

\hypertarget{glm-en-datos-binarios}{%
\subsection{1.4 GLM en datos binarios}\label{glm-en-datos-binarios}}

Para modelizar con un modelo lineal generalizado la predicción de
\texttt{fracture}, mantenemos el carácter dicotómico de la respuesta, y
su modelización con una variable Bernouilli, con \(p\) su valor esperado
y la probabilidad de fractura. Buscamos pues, un modelo a través del que
relacionar esta probabilidad con un predictor lineal construido en
función de las variables \texttt{age} y \texttt{priorfrac}. Como hemos
visto antes, la función \texttt{logit} nos permite manejar este cambio
de escala de una probabilidad a toda la recta real, y por otro lado, la
familia binomial integra a la distribución Bernouilli, por lo que
planteamos lo que popularmente se conoce como \textbf{el modelo logit},
que aplicado a GLOW para predecir la fractura en función de la edad y el
haber tenido una fractura previa, se formularía así:

\[Y=\begin{cases} 1 \ \text{fractura, con probabilidad } \pi \\
0 \ \text{ no fractura, con probabilidad } 1-\pi
\end{cases} \ \rightarrow \ logit(\pi)=\theta + \beta \cdot \text{age} + \gamma \cdot \mathbf{1}(\text{priorfrac}=\text{Yes})
\]

\begin{tcolorbox}[enhanced jigsaw, coltitle=black, rightrule=.15mm, arc=.35mm, leftrule=.75mm, colback=white, colbacktitle=quarto-callout-tip-color!10!white, left=2mm, breakable, toptitle=1mm, toprule=.15mm, bottomtitle=1mm, titlerule=0mm, title=\textcolor{quarto-callout-tip-color}{\faLightbulb}\hspace{0.5em}{🔧 En R. El \texttt{glm()} con datos binarios}, bottomrule=.15mm, opacityback=0, opacitybacktitle=0.6, colframe=quarto-callout-tip-color-frame]

La función
\texttt{glm(y\ \textasciitilde{}\ x1+x2+x1:x2+...,family=binomial(link\ =\ "logit"),data)}
ajusta un modelo lineal generalizado en R, integrando los tres
componentes del GLM:

\begin{itemize}
\tightlist
\item
  \textbf{Respuesta (componente aleatoria):} Si se trata de una
  respuesta binaria, esta puede venir etiquetada (No/Yes) o recodificada
  (0/1), siempre vinculada a la familia binomial →
  \texttt{family\ =\ binomial}. La categoría de \textbf{``éxito''} -y
  esto es relevante porque define la probabilidad que se estima- siempre
  es ``1'' en la codificación 0/1, y la que alfabéticamente esté
  posicionada la última (en No/Yes, sería Yes, pues se asume
  No=0/Yes=1).
\item
  \textbf{Predictor lineal (componente sistemática):} Se especifica
  igual que en un \texttt{lm()} -lineal en los parámetros:
  \texttt{x1+x2+x1:x2+...}.
\item
  \textbf{Enlace:} \texttt{link\ =\ "logit"} conecta la media ---la
  probabilidad--- con el predictor y mantiene \(\hat\pi\) en \((0,1)\).
\end{itemize}

Los principales descriptivos los obtenemos con: -
\texttt{summary(fit\_glm)}, el resultado del ajuste -
\texttt{tidy(fit\_glm)} formatea el resultado -\texttt{glance()} da la
\textbf{deviance} y el \textbf{AIC}, con las que evaluamos la calidad
del ajuste. - \texttt{fitted(fit\_glm)} da la predicción de
probabilidades en los registros observados. -
\texttt{broom::augment(fit\_glm,\ type.predict\ =\ "response")} da los
ajustes ordenados en \texttt{{[}0,1{]}} - \texttt{predict()} devuelve la
predicción en la \textbf{escala del predictor} (log-odds) para cualquier
valor de los predictores. Para obtenerla en \textbf{probabilidad} hay
que pedir la escala de respuesta:

\begin{Shaded}
\begin{Highlighting}[]
\FunctionTok{predict}\NormalTok{(fit\_glm, }\AttributeTok{newdata =}\NormalTok{ ..., }\AttributeTok{type =} \StringTok{"response"}\NormalTok{)    }\CommentTok{\# probabilidad en [0,1]}
\FunctionTok{predict}\NormalTok{(fit\_glm, }\AttributeTok{newdata =}\NormalTok{ ...)                       }\CommentTok{\# por defecto: log{-}odds (type = "link")}
\end{Highlighting}
\end{Shaded}

\end{tcolorbox}

Ajustamos pues el modelo, y mostramos el resultado.

\begin{Shaded}
\begin{Highlighting}[]
\NormalTok{fit\_glm }\OtherTok{\textless{}{-}} \FunctionTok{glm}\NormalTok{(fracture }\SpecialCharTok{\textasciitilde{}}\NormalTok{ age }\SpecialCharTok{+}\NormalTok{ priorfrac, }\AttributeTok{family =}\NormalTok{ binomial, }\AttributeTok{data =}\NormalTok{ glow)}
\CommentTok{\# Resultado del ajuste}
\FunctionTok{summary}\NormalTok{(fit\_glm)}
\end{Highlighting}
\end{Shaded}

\begin{verbatim}

Call:
glm(formula = fracture ~ age + priorfrac, family = binomial, 
    data = glow)

Coefficients:
             Estimate Std. Error z value Pr(>|z|)    
(Intercept)  -4.21429    0.84784  -4.971 6.67e-07 ***
age           0.04119    0.01218   3.382 0.000719 ***
priorfracYes  0.83884    0.23416   3.582 0.000340 ***
---
Signif. codes:  0 '***' 0.001 '**' 0.01 '*' 0.05 '.' 0.1 ' ' 1

(Dispersion parameter for binomial family taken to be 1)

    Null deviance: 562.34  on 499  degrees of freedom
Residual deviance: 528.52  on 497  degrees of freedom
AIC: 534.52

Number of Fisher Scoring iterations: 4
\end{verbatim}

Este ajuste proporciona, como el \texttt{lm()}, coeficientes, medidas de
bondad de ajuste, que podemos mostrar con formato con la librería
\texttt{broom}.

\begin{Shaded}
\begin{Highlighting}[]
\FunctionTok{tidy}\NormalTok{(fit\_glm)                                     }\CommentTok{\# coeficientes (escala log{-}odds)}
\end{Highlighting}
\end{Shaded}

\begin{verbatim}
# A tibble: 3 x 5
  term         estimate std.error statistic     p.value
  <chr>           <dbl>     <dbl>     <dbl>       <dbl>
1 (Intercept)   -4.21      0.848      -4.97 0.000000667
2 age            0.0412    0.0122      3.38 0.000719   
3 priorfracYes   0.839     0.234       3.58 0.000340   
\end{verbatim}

\begin{Shaded}
\begin{Highlighting}[]
\FunctionTok{glance}\NormalTok{(fit\_glm)                                   }\CommentTok{\# deviance, AIC, BIC, gl...}
\end{Highlighting}
\end{Shaded}

\begin{verbatim}
# A tibble: 1 x 8
  null.deviance df.null logLik   AIC   BIC deviance df.residual  nobs
          <dbl>   <int>  <dbl> <dbl> <dbl>    <dbl>       <int> <int>
1          562.     499  -264.  535.  547.     529.         497   500
\end{verbatim}

\begin{Shaded}
\begin{Highlighting}[]
\FunctionTok{augment}\NormalTok{(fit\_glm, }\AttributeTok{type.predict =} \StringTok{"response"}\NormalTok{) }\SpecialCharTok{|\textgreater{}}    \CommentTok{\# ajustes en escala probabilidad}
\NormalTok{  dplyr}\SpecialCharTok{::}\FunctionTok{select}\NormalTok{(fracture, age, priorfrac, .fitted) }\SpecialCharTok{|\textgreater{}}
  \FunctionTok{slice\_head}\NormalTok{(}\AttributeTok{n =} \DecValTok{5}\NormalTok{)}
\end{Highlighting}
\end{Shaded}

\begin{verbatim}
# A tibble: 5 x 4
  fracture   age priorfrac .fitted
  <fct>    <int> <fct>       <dbl>
1 No          62 No          0.160
2 No          65 No          0.177
3 No          88 Yes         0.562
4 No          82 No          0.302
5 No          61 No          0.154
\end{verbatim}

La Figura~\ref{fig-u11-ols-vs-logistica} muestra las predicciones que se
obtienen en los dos modelos considerados hasta el momento: el lineal y
el lineal generalizado.

\begin{Shaded}
\begin{Highlighting}[]
\CommentTok{\# Rejilla de predicción de los modelos }
\NormalTok{grid }\OtherTok{\textless{}{-}} \FunctionTok{expand\_grid}\NormalTok{(}
  \AttributeTok{age       =} \FunctionTok{seq}\NormalTok{(}\FunctionTok{min}\NormalTok{(glow}\SpecialCharTok{$}\NormalTok{age), }\FunctionTok{max}\NormalTok{(glow}\SpecialCharTok{$}\NormalTok{age), }\AttributeTok{length.out =} \DecValTok{100}\NormalTok{),}
  \AttributeTok{priorfrac =} \FunctionTok{factor}\NormalTok{(}\FunctionTok{c}\NormalTok{(}\StringTok{"No"}\NormalTok{, }\StringTok{"Yes"}\NormalTok{), }\AttributeTok{levels =} \FunctionTok{levels}\NormalTok{(glow}\SpecialCharTok{$}\NormalTok{priorfrac))}
\NormalTok{)}
\NormalTok{grid}\SpecialCharTok{$}\NormalTok{OLS }\OtherTok{\textless{}{-}} \FunctionTok{predict}\NormalTok{(fit\_ols, }\AttributeTok{newdata =}\NormalTok{ grid)                    }\CommentTok{\# escala probabilidad}
\NormalTok{grid}\SpecialCharTok{$}\NormalTok{GLM }\OtherTok{\textless{}{-}} \FunctionTok{predict}\NormalTok{(fit\_glm, }\AttributeTok{newdata =}\NormalTok{ grid, }\AttributeTok{type =} \StringTok{"response"}\NormalTok{) }\CommentTok{\# escala probabilidad}
\NormalTok{grid\_long }\OtherTok{\textless{}{-}} \FunctionTok{pivot\_longer}\NormalTok{(grid, }\FunctionTok{c}\NormalTok{(OLS, GLM),}
                          \AttributeTok{names\_to =} \StringTok{"modelo"}\NormalTok{, }\AttributeTok{values\_to =} \StringTok{"p\_hat"}\NormalTok{) }\SpecialCharTok{|\textgreater{}}
  \FunctionTok{mutate}\NormalTok{(}\AttributeTok{modelo =} \FunctionTok{factor}\NormalTok{(modelo, }\AttributeTok{levels =} \FunctionTok{c}\NormalTok{(}\StringTok{"OLS"}\NormalTok{, }\StringTok{"GLM"}\NormalTok{)))}

\CommentTok{\# Proporciones observadas por grupo de edad y priorfrac (las de la EDA),}
\CommentTok{\# situadas en la edad media de cada grupo}
\NormalTok{grupos }\OtherTok{\textless{}{-}}\NormalTok{ glow }\SpecialCharTok{|\textgreater{}}
  \FunctionTok{mutate}\NormalTok{(}\AttributeTok{edad\_grupo =} \FunctionTok{cut}\NormalTok{(age, }\AttributeTok{breaks =} \FunctionTok{seq}\NormalTok{(}\DecValTok{55}\NormalTok{, }\DecValTok{95}\NormalTok{, }\AttributeTok{by =} \DecValTok{5}\NormalTok{))) }\SpecialCharTok{|\textgreater{}}
  \FunctionTok{filter}\NormalTok{(}\SpecialCharTok{!}\FunctionTok{is.na}\NormalTok{(edad\_grupo)) }\SpecialCharTok{|\textgreater{}}
  \FunctionTok{group\_by}\NormalTok{(edad\_grupo, priorfrac) }\SpecialCharTok{|\textgreater{}}
  \FunctionTok{summarise}\NormalTok{(}\AttributeTok{age =} \FunctionTok{mean}\NormalTok{(age), }\AttributeTok{p\_fractura =} \FunctionTok{mean}\NormalTok{(fractura01), }\AttributeTok{n =} \FunctionTok{n}\NormalTok{(), }\AttributeTok{.groups =} \StringTok{"drop"}\NormalTok{)}

\FunctionTok{ggplot}\NormalTok{() }\SpecialCharTok{+}
  \FunctionTok{geom\_point}\NormalTok{(}\AttributeTok{data =}\NormalTok{ grupos,}
             \FunctionTok{aes}\NormalTok{(age, p\_fractura, }\AttributeTok{color =}\NormalTok{ priorfrac, }\AttributeTok{size =}\NormalTok{ n)) }\SpecialCharTok{+}
  \FunctionTok{geom\_line}\NormalTok{(}\AttributeTok{data =}\NormalTok{ grid\_long,}
            \FunctionTok{aes}\NormalTok{(age, p\_hat, }\AttributeTok{color =}\NormalTok{ priorfrac), }\AttributeTok{linewidth =} \DecValTok{1}\NormalTok{) }\SpecialCharTok{+}
  \FunctionTok{geom\_hline}\NormalTok{(}\AttributeTok{yintercept =} \FunctionTok{c}\NormalTok{(}\DecValTok{0}\NormalTok{, }\DecValTok{1}\NormalTok{), }\AttributeTok{linetype =} \StringTok{"dashed"}\NormalTok{) }\SpecialCharTok{+}
  \FunctionTok{facet\_wrap}\NormalTok{(}\SpecialCharTok{\textasciitilde{}}\NormalTok{ modelo) }\SpecialCharTok{+}
  \FunctionTok{labs}\NormalTok{(}\AttributeTok{x =} \StringTok{"Edad"}\NormalTok{, }\AttributeTok{y =} \StringTok{"Proporción / probabilidad de fractura"}\NormalTok{,}
       \AttributeTok{color =} \StringTok{"Fractura previa"}\NormalTok{, }\AttributeTok{size =} \StringTok{"n"}\NormalTok{)}
\end{Highlighting}
\end{Shaded}

\begin{figure}[H]

{\centering \includegraphics{caso1_binaria_files/figure-pdf/fig-u11-ols-vs-logistica-1.pdf}

}

\caption{\label{fig-u11-ols-vs-logistica}Predicciones de los modelos
(\texttt{age} + \texttt{priorfrac}) frente a las proporciones de
fractura observadas por grupo de edad (el tamaño de los puntos
representa el nº de casos). Ajustes OLS (izquierda) y logístico
(derecha) deberían pasar cerca de los puntos del mismo color. En el
rango observado ambos se mantienen en {[}0,1{]}.}

\end{figure}

Se aprecia en la Figura~\ref{fig-u11-ols-vs-logistica} la diferencia en
la predicción de la probabilidad de fractura, asimilada en el OLS por
sendas rectas paralelas y en el GLM con dos curvas que intentan adaptar
mejor la forma de los datos.

Ya con el glm ajustado, al comparar los residuos medios por tramos
(\emph{binned residuals}, en la
\protect\hyperlink{sec-u12-diagnostico}{Unidad 2} para el ajuste OLS y
el GLM (Figura~\ref{fig-u11-residuos-ols-glm}), apreciamos mejor
rendimiento en este último, con la mayoría de residuos dentro de la
banda de confianza, a diferencia de los residuos OLS, que fluyen fuera
de la banda.

\begin{Shaded}
\begin{Highlighting}[]
\NormalTok{p\_ols }\OtherTok{\textless{}{-}} \FunctionTok{plot}\NormalTok{(performance}\SpecialCharTok{::}\FunctionTok{binned\_residuals}\NormalTok{(fit\_ols)) }\SpecialCharTok{+}
  \FunctionTok{labs}\NormalTok{(}\AttributeTok{title =} \StringTok{"OLS (modelo lineal de probabilidad)"}\NormalTok{)}
\NormalTok{p\_glm }\OtherTok{\textless{}{-}} \FunctionTok{plot}\NormalTok{(performance}\SpecialCharTok{::}\FunctionTok{binned\_residuals}\NormalTok{(fit\_glm)) }\SpecialCharTok{+}
  \FunctionTok{labs}\NormalTok{(}\AttributeTok{title =} \StringTok{"GLM logístico"}\NormalTok{)}
\NormalTok{patchwork}\SpecialCharTok{::}\FunctionTok{wrap\_plots}\NormalTok{( p\_ols, p\_glm, }\AttributeTok{nrow =} \DecValTok{2}\NormalTok{) }\SpecialCharTok{+}
\NormalTok{  patchwork}\SpecialCharTok{::}\FunctionTok{plot\_annotation}\NormalTok{(}\AttributeTok{title =} \StringTok{"Residuos medios por tramos: OLS vs GLM"}\NormalTok{)}
\end{Highlighting}
\end{Shaded}

\begin{figure}[H]

{\centering \includegraphics[width=1\textwidth,height=\textheight]{caso1_binaria_files/figure-pdf/fig-u11-residuos-ols-glm-1.pdf}

}

\caption{\label{fig-u11-residuos-ols-glm}Residuos medios por tramos para
los ajustes OLS y GLM.}

\end{figure}

\begin{tcolorbox}[enhanced jigsaw, coltitle=black, rightrule=.15mm, arc=.35mm, leftrule=.75mm, colback=white, colbacktitle=quarto-callout-important-color!10!white, left=2mm, breakable, toptitle=1mm, toprule=.15mm, bottomtitle=1mm, titlerule=0mm, title=\textcolor{quarto-callout-important-color}{\faExclamation}\hspace{0.5em}{Tarea del estudiante (1.1) --- arranca la Pregunta 1 del caso}, bottomrule=.15mm, opacityback=0, opacitybacktitle=0.6, colframe=quarto-callout-important-color-frame]

Empiezas a responder la \textbf{Pregunta 1} del cliente
(\protect\hyperlink{sec-u16}{Unidad 6}) ---\emph{¿quién se fractura y
cuánto influyen la fragilidad (\texttt{x1}) y el tratamiento
(\texttt{x2})?}--- sobre la \textbf{cohorte de tu equipo}. Aquí solo la
dejas \textbf{planteada y ajustada}; interpretarla, decidirla y
evaluarla es cosa de la Unidad 2.

Tareas:

\begin{enumerate}
\def\labelenumi{\arabic{enumi}.}
\tightlist
\item
  \textbf{Identifica el objetivo y el tipo de respuesta.} Argumenta que
  «¿ocurre la fractura?» es una respuesta \textbf{binaria}
  (\texttt{ever}) y que el \textbf{modelo lineal (OLS) no sirve}: ajusta
  \texttt{lm(ever\ \textasciitilde{}\ x1)} y construye \textbf{un}
  gráfico que evidencie a la vez predicciones fuera de \([0,1]\) y los
  residuos en dos bandas.
\item
  \textbf{Del fallo al GLM.} En una tabla, asocia cada fallo de OLS con
  el \textbf{componente del GLM} (aleatoria / predictor lineal / enlace)
  que lo resuelve, con una frase de porqué.
\item
  \textbf{Primer ajuste.} Ajusta
  \texttt{glm(ever\ \textasciitilde{}\ x1\ +\ x2,\ family\ =\ binomial)}
  y guárdalo en el \texttt{.qmd} del equipo: es el punto de partida de
  la Pregunta 1, que completarás en la Unidad 2.
\end{enumerate}

\end{tcolorbox}

\hypertarget{sec-u11-puente}{%
\subsection{Puente a la Unidad 2}\label{sec-u11-puente}}

Hoy dimos por buenos, sin justificarlos, varios ingredientes: que
Bernoulli pertenece a la familia exponencial, que el logit es el enlace
canónico, que \texttt{glm()} estima por máxima verosimilitud y que la
\texttt{deviance} mide el ajuste. La \protect\hyperlink{sec-u12}{Unidad
2} abre esa caja negra para desarrollarla.

\hypertarget{sec-u12}{%
\section{2 · Formulación de los GLM: estimación, inferencia,
interpretación y evaluación}\label{sec-u12}}

En la \protect\hyperlink{sec-u11}{Unidad 1} dimos por buenos, sin
justificarlos, cuatro ingredientes: que la Bernoulli pertenece a la
\textbf{familia exponencial}, que el \textbf{logit} es su enlace
canónico, que \texttt{glm()} estima por \textbf{máxima verosimilitud} y
que la \textbf{deviance} mide el ajuste. Esta unidad abre esa caja negra
y la cierra con lo que de verdad se hace en la práctica: estimar, hacer
inferencia, comparar modelos, interpretarlos en la escala correcta y
evaluarlos como clasificadores.

\begin{tcolorbox}[enhanced jigsaw, coltitle=black, rightrule=.15mm, arc=.35mm, leftrule=.75mm, colback=white, colbacktitle=quarto-callout-note-color!10!white, left=2mm, breakable, toptitle=1mm, toprule=.15mm, bottomtitle=1mm, titlerule=0mm, title=\textcolor{quarto-callout-note-color}{\faInfo}\hspace{0.5em}{Sobre esta unidad}, bottomrule=.15mm, opacityback=0, opacitybacktitle=0.6, colframe=quarto-callout-note-color-frame]

RA implicados:

\begin{itemize}
\tightlist
\item
  \textbf{RA2 · Formulación y estimación.} Formular en la familia
  exponencial, estimar por MV (IWLS) y hacer inferencia (deviance, Wald,
  razón de verosimilitudes).
\item
  \textbf{RA3 · Interpretación en la escala correcta.} Leer los efectos
  como odds ratios y efectos marginales.
\item
  \textbf{RA4 · Diagnóstico y bondad de ajuste.} Separación,
  calibración, residuos y evaluación como clasificador.
\item
  \textbf{RA5 · Selección y comparación de modelos.} Deviance, AIC/BIC y
  contrastes de razón de verosimilitudes.
\end{itemize}

(Agresti, 2015; Dobson \& Barnett, 2018; Fox, 2016; Hosmer et~al., 2013;
McCullagh \& Nelder, 1989; Nelder \& Wedderburn, 1972)

\end{tcolorbox}

\hypertarget{sec-u12-familia}{%
\subsection{2.1 Familia exponencial y relación
media--varianza}\label{sec-u12-familia}}

Un GLM exige que la respuesta siga una distribución de la
\textbf{familia exponencial de dispersión}, cuya densidad (o función de
masa) se escribe en la forma canónica

\[
f(y;\theta,\phi)=\exp\!\left\{\frac{y\theta-b(\theta)}{a(\phi)}+c(y,\phi)\right\},
\]

donde \(\theta\) es el \textbf{parámetro natural} (o canónico),
\(b(\theta)\) la \textbf{función cumulante}, \(\phi\) el
\textbf{parámetro de dispersión} y \(a(\phi)\), \(c(y,\phi)\) funciones
conocidas (Dobson \& Barnett, 2018; McCullagh \& Nelder, 1989). Esta
forma no es un capricho notacional: de ella se derivan, derivando la
cumulante, los dos momentos que gobiernan todo el marco.

\begin{tcolorbox}[enhanced jigsaw, coltitle=black, rightrule=.15mm, arc=.35mm, leftrule=.75mm, colback=white, colbacktitle=quarto-callout-note-color!10!white, left=2mm, breakable, toptitle=1mm, toprule=.15mm, bottomtitle=1mm, titlerule=0mm, title=\textcolor{quarto-callout-note-color}{\faInfo}\hspace{0.5em}{Quién vive en esta familia (y para qué sirve)}, bottomrule=.15mm, opacityback=0, opacitybacktitle=0.6, colframe=quarto-callout-note-color-frame]

La misma forma canónica agrupa distribuciones muy distintas; cada una es
la herramienta natural para un tipo de fenómeno del mundo real:

\begin{itemize}
\tightlist
\item
  \textbf{Bernoulli / Binomial} --- sucesos sí/no: ¿hay fractura?
  (GLOW), abandono de cliente (\emph{churn}), impago de un crédito,
  resultado de un test diagnóstico.
\item
  \textbf{Poisson} --- recuentos por unidad de tiempo o espacio:
  reclamaciones de un asegurado al año, visitas a urgencias, fallos de
  una máquina, llamadas por hora.
\item
  \textbf{Exponencial} --- tiempo continuo hasta un evento con riesgo
  constante (sin memoria): vida útil de un componente; abre el hilo de
  supervivencia (Caso 3).
\item
  \textbf{Geométrica} --- número de intentos hasta el primer éxito:
  envíos hasta la primera respuesta, ciclos hasta el primer fallo.
\item
  \textbf{Gaussiana} --- magnitudes continuas sin restricción de signo:
  estatura, error de medida. Es el OLS que ya conoces, subsumido en el
  marco.
\item
  \textbf{Gamma} --- magnitudes positivas y asimétricas con variabilidad
  proporcional a la media: coste de un siniestro, gasto sanitario por
  paciente, duración de una sesión (Caso 3).
\item
  \textbf{Inversa gaussiana / Tweedie} --- duraciones de cola pesada, o
  datos con masa en cero y parte continua positiva a la vez (coste
  agregado de siniestros, precipitación diaria).
\end{itemize}

\end{tcolorbox}

Partiendo de que \(\int f\,\mathrm{d}y=1\) y derivando bajo el signo
integral se obtiene el resultado central:

\[
\mathbb{E}[Y]=\mu=b'(\theta),\qquad
\operatorname{Var}(Y)=b''(\theta)\,a(\phi).
\]

La media depende solo de \(\theta\) y la varianza se \textbf{factoriza}
en dos piezas: \(a(\phi)\), común a la familia, y \(b''(\theta)\), que
reescrita en función de la media define la \textbf{función de varianza}
\(V(\mu)=b''(\theta(\mu))\). Esa \(V(\mu)\) es la \textbf{huella
digital} de cada familia: es lo que distingue una Bernoulli de una
Poisson o una Gamma, y lo que en 1.1 «absorbía los problemas 2 y 3» de
OLS (varianza no constante y error no normal).

\begin{tcolorbox}[enhanced jigsaw, coltitle=black, rightrule=.15mm, arc=.35mm, leftrule=.75mm, colback=white, colbacktitle=quarto-callout-note-color!10!white, left=2mm, breakable, toptitle=1mm, toprule=.15mm, bottomtitle=1mm, titlerule=0mm, title=\textcolor{quarto-callout-note-color}{\faInfo}\hspace{0.5em}{El otro factor de la varianza: \(a(\phi)\) y los pesos a priori}, bottomrule=.15mm, opacityback=0, opacitybacktitle=0.6, colframe=quarto-callout-note-color-frame]

En los modelos GLM estándar el primer factor de
\(\operatorname{Var}(Y_i)=a(\phi)\,V(\mu_i)\) toma
\textbf{habitualmente} la forma \[a(\phi)=\frac{\phi}{m_i},\] con
\(\phi\) el \textbf{parámetro de dispersión} (común a la familia) y
\(m_i\) un \textbf{peso a priori conocido} que registra cuántas unidades
elementales hay detrás de cada fila. En general \(a(\phi)\) es solo
\emph{alguna} función de \(\phi\); la forma \(\phi/m_i\) es la
convención que surge cuando cada observación resume \(m_i\) unidades
i.i.d.

Con datos \textbf{individuales} ---una fila, una observación--- se tiene
\(m_i=1\) y por tanto \(a(\phi)=\phi\). Es la situación de todo el Caso
1: cada mujer de GLOW aporta un único ensayo Bernoulli.

El peso \(m_i\neq1\) aparece cuando el dato llega \textbf{agregado}:

\begin{itemize}
\tightlist
\item
  \textbf{Binomial como proporción.} Si \(y_i\) es la fracción de éxitos
  sobre \(m_i\) ensayos, \(m_i\) es el denominador binomial y
  \(\operatorname{Var}(\bar Y_i)=\mu_i(1-\mu_i)/m_i\). En R: respuesta
  \texttt{cbind(éxitos,\ fracasos)} o bien \texttt{weights\ =\ m}.
\item
  \textbf{Gaussiana o Gamma sobre medias de grupo.} Si \(y_i\) es el
  promedio de \(m_i\) mediciones, \(m_i\) es el tamaño del grupo.
\end{itemize}

No conviene confundir las dos piezas: \(\phi\) es la \textbf{escala
libre} de la familia (recuadro siguiente), mientras que \(m_i\) refleja
la \textbf{estructura de los datos} (réplicas o agregación), no la
distribución.

\end{tcolorbox}

\begin{tcolorbox}[enhanced jigsaw, coltitle=black, rightrule=.15mm, arc=.35mm, leftrule=.75mm, colback=white, colbacktitle=quarto-callout-note-color!10!white, left=2mm, breakable, toptitle=1mm, toprule=.15mm, bottomtitle=1mm, titlerule=0mm, title=\textcolor{quarto-callout-note-color}{\faInfo}\hspace{0.5em}{¿Uno o dos parámetros? El papel de \(\phi\)}, bottomrule=.15mm, opacityback=0, opacitybacktitle=0.6, colframe=quarto-callout-note-color-frame]

En \textbf{Bernoulli, Poisson, exponencial y geométrica} el parámetro de
dispersión es \textbf{conocido y fijo, \(\phi=1\)}: no hay escala libre
que estimar, y un \textbf{único parámetro \(\theta\)} (equivalentemente,
la media \(\mu\)) gobierna a la vez la media y la varianza. Por eso
\(V(\mu)\) las liga de forma rígida ---\(\mu(1-\mu)\), \(\mu\)\ldots---.
Son familias exponenciales \textbf{lineales} (un solo parámetro).

El resto de variantes ---\textbf{Gaussiana, Gamma, inversa gaussiana,
Tweedie}--- añaden un parámetro de \textbf{dispersión propio} que se
estima de los datos: el \textbf{escalamiento se modeliza con
independencia de la media}, lo que permite describir situaciones más
ricas (varianza constante en la Gaussiana, coeficiente de variación
constante en la Gamma\ldots). Son los \textbf{modelos de dispersión
exponencial}.

Ese mando \(\phi\) no es ornamental: reaparece como
\textbf{sobredispersión} en los conteos del Caso 2 (quasi-Poisson,
binomial negativa) y como las \textbf{continuas positivas} del Caso 3.
En el Caso 1, al ser todo Bernoulli/Binomial, \(\phi=1\) y no hay nada
que estimar.

\end{tcolorbox}

\hypertarget{tbl-familias}{}
\begin{longtable}[]{@{}
  >{\raggedright\arraybackslash}p{(\columnwidth - 8\tabcolsep) * \real{0.2000}}
  >{\raggedright\arraybackslash}p{(\columnwidth - 8\tabcolsep) * \real{0.2000}}
  >{\raggedright\arraybackslash}p{(\columnwidth - 8\tabcolsep) * \real{0.2000}}
  >{\raggedright\arraybackslash}p{(\columnwidth - 8\tabcolsep) * \real{0.2000}}
  >{\raggedright\arraybackslash}p{(\columnwidth - 8\tabcolsep) * \real{0.2000}}@{}}
\caption{\label{tbl-familias}Familias de la exponencial: parametrización
e ingredientes canónicos.}\tabularnewline
\toprule()
\begin{minipage}[b]{\linewidth}\raggedright
\end{minipage} & \begin{minipage}[b]{\linewidth}\raggedright
Bernoulli
\end{minipage} & \begin{minipage}[b]{\linewidth}\raggedright
Poisson
\end{minipage} & \begin{minipage}[b]{\linewidth}\raggedright
Gaussiana
\end{minipage} & \begin{minipage}[b]{\linewidth}\raggedright
Gamma
\end{minipage} \\
\midrule()
\endfirsthead
\toprule()
\begin{minipage}[b]{\linewidth}\raggedright
\end{minipage} & \begin{minipage}[b]{\linewidth}\raggedright
Bernoulli
\end{minipage} & \begin{minipage}[b]{\linewidth}\raggedright
Poisson
\end{minipage} & \begin{minipage}[b]{\linewidth}\raggedright
Gaussiana
\end{minipage} & \begin{minipage}[b]{\linewidth}\raggedright
Gamma
\end{minipage} \\
\midrule()
\endhead
\textbf{Parametrización} & \(\mathrm{Br}(\mu)\) & \(\mathrm{Pois}(\mu)\)
& \(\mathcal{N}(\mu,\sigma^2)\) &
\(\mathrm{Ga}\!\left(\nu,\tfrac{\mu}{\nu}\right)\) \\
\(\theta(\mu)\) & \(\log\frac{\mu}{1-\mu}\) & \(\log\mu\) & \(\mu\) &
\(-1/\mu\) \\
\(b(\theta)\) & \(\log(1+e^\theta)\) & \(e^\theta\) & \(\theta^2/2\) &
\(-\log(-\theta)\) \\
\(b'(\theta)\) & \(\dfrac{e^\theta}{1+e^\theta}\) & \(e^\theta\) &
\(\theta\) & \(-1/\theta\) \\
\(b''(\theta)\) & \(\dfrac{e^\theta}{(1+e^\theta)^2}\) & \(e^\theta\) &
\(1\) & \(1/\theta^2\) \\
\(a(\phi)\) & \(1\) & \(1\) & \(\phi=\sigma^2\) & \(\phi=\nu^{-1}\) \\
\(c(y,\phi)\) & \(0\) & \(-\log y!\) &
\(-\tfrac{1}{2}\!\left(\tfrac{y^2}{\phi}+\log(2\pi\phi)\right)\) &
\(\tfrac{1-\phi}{\phi}\log y-\tfrac{\log\phi}{\phi}-\log\Gamma\!\left(\tfrac{1}{\phi}\right)\) \\
\(V(\mu)\) & \(\mu(1-\mu)\) & \(\mu\) & \(1\) & \(\mu^2\) \\
\textbf{Enlace canónico} & logit & log & identidad & inversa \\
\bottomrule()
\end{longtable}

La fila \(V(\mu)\) explica de un vistazo por qué OLS (Gaussiana,
\(V\equiv 1\)) fracasa en binaria: para la Bernoulli,
\(V(\mu)=\mu(1-\mu)\) \textbf{no} es constante. El objeto
\texttt{family()} de R expone exactamente estas piezas, y cada familia
permite recuperar sus ingredientes canónicos:

\begin{Shaded}
\begin{Highlighting}[]
\FunctionTok{family}\NormalTok{(fit\_glm)              }\CommentTok{\# familia y enlace del ajuste de 1.1}
\end{Highlighting}
\end{Shaded}

\begin{verbatim}

Family: binomial 
Link function: logit 
\end{verbatim}

\begin{Shaded}
\begin{Highlighting}[]
\FunctionTok{binomial}\NormalTok{()}\SpecialCharTok{$}\NormalTok{variance          }\CommentTok{\# función de varianza V(mu) = mu(1{-}mu)}
\end{Highlighting}
\end{Shaded}

\begin{verbatim}
function (mu) 
mu * (1 - mu)
<bytecode: 0x7ff63e311b30>
<environment: 0x7ff6380e9208>
\end{verbatim}

\begin{Shaded}
\begin{Highlighting}[]
\FunctionTok{binomial}\NormalTok{()}\SpecialCharTok{$}\FunctionTok{linkfun}\NormalTok{(}\FloatTok{0.25}\NormalTok{)     }\CommentTok{\# logit(0.25): de la media al predictor}
\end{Highlighting}
\end{Shaded}

\begin{verbatim}
[1] -1.098612
\end{verbatim}

\begin{Shaded}
\begin{Highlighting}[]
\FunctionTok{binomial}\NormalTok{()}\SpecialCharTok{$}\FunctionTok{linkinv}\NormalTok{(}\SpecialCharTok{{-}}\FloatTok{1.0986}\NormalTok{)  }\CommentTok{\# sigmoide: del predictor a la media}
\end{Highlighting}
\end{Shaded}

\begin{verbatim}
[1] 0.2500023
\end{verbatim}

La varianza de una sola respuesta binaria no se observa: un 0/1 no tiene
dispersión propia. Para \emph{verla}, agrupamos GLOW en diez tramos
(deciles) según la probabilidad ajustada \(\hat\mu\); dentro de cada
tramo el riesgo es parecido, así que podemos medir allí la varianza
empírica de la respuesta. Confrontamos ese valor con la parábola
\(\mu(1-\mu)\): si los puntos la siguen, confirmamos ---solo con
datos--- que en la Bernoulli la varianza la \textbf{dicta la media}
(máxima en \(\mu=0{,}5\), casi nula cerca de 0 o 1), no una constante
como supondría OLS. La Figura~\ref{fig-u12-varianza-mu} muestra que los
puntos siguen la parábola.

\begin{Shaded}
\begin{Highlighting}[]
\NormalTok{glow }\SpecialCharTok{|\textgreater{}}
  \FunctionTok{mutate}\NormalTok{(}\AttributeTok{p\_hat =} \FunctionTok{fitted}\NormalTok{(fit\_glm),}
         \AttributeTok{tramo =} \FunctionTok{ntile}\NormalTok{(p\_hat, }\DecValTok{10}\NormalTok{)) }\SpecialCharTok{|\textgreater{}}
  \FunctionTok{group\_by}\NormalTok{(tramo) }\SpecialCharTok{|\textgreater{}}
  \FunctionTok{summarise}\NormalTok{(}\AttributeTok{mu =} \FunctionTok{mean}\NormalTok{(p\_hat),}
            \AttributeTok{var\_emp =} \FunctionTok{var}\NormalTok{(fractura01), }\AttributeTok{.groups =} \StringTok{"drop"}\NormalTok{) }\SpecialCharTok{|\textgreater{}}
  \FunctionTok{ggplot}\NormalTok{(}\FunctionTok{aes}\NormalTok{(mu, var\_emp)) }\SpecialCharTok{+}
  \FunctionTok{geom\_function}\NormalTok{(}\AttributeTok{fun =}\NormalTok{ \textbackslash{}(m) m }\SpecialCharTok{*}\NormalTok{ (}\DecValTok{1} \SpecialCharTok{{-}}\NormalTok{ m), }\AttributeTok{linewidth =} \DecValTok{1}\NormalTok{, }\AttributeTok{color =} \StringTok{"firebrick"}\NormalTok{) }\SpecialCharTok{+}
  \FunctionTok{geom\_point}\NormalTok{(}\AttributeTok{size =} \DecValTok{2}\NormalTok{) }\SpecialCharTok{+}
  \FunctionTok{labs}\NormalTok{(}\AttributeTok{x =} \FunctionTok{expression}\NormalTok{(}\FunctionTok{hat}\NormalTok{(mu)), }\AttributeTok{y =} \StringTok{"Varianza empírica de Y por tramo"}\NormalTok{,}
       \AttributeTok{caption =} \FunctionTok{expression}\NormalTok{(}\StringTok{"Curva: V("} \SpecialCharTok{*}\NormalTok{ mu }\SpecialCharTok{*} \StringTok{") = "} \SpecialCharTok{*}\NormalTok{ mu }\SpecialCharTok{*} \StringTok{"(1 {-} "} \SpecialCharTok{*}\NormalTok{ mu }\SpecialCharTok{*} \StringTok{")"}\NormalTok{))}
\end{Highlighting}
\end{Shaded}

\begin{figure}[H]

{\centering \includegraphics{caso1_binaria_files/figure-pdf/fig-u12-varianza-mu-1.pdf}

}

\caption{\label{fig-u12-varianza-mu}Relación media--varianza en GLOW. La
varianza empírica de la respuesta por tramos de probabilidad ajustada
(puntos) sigue la parábola teórica \(V(\mu)=\mu(1-\mu)\) (línea), no una
constante: justo lo que OLS no puede reproducir.}

\end{figure}

\hypertarget{sec-u12-enlace}{%
\subsection{2.2 El enlace: canónico y
alternativas}\label{sec-u12-enlace}}

El \textbf{enlace} o \emph{función link} \(g\) conecta la media con el
predictor lineal, \(g(\mu)=\eta=\mathbf{x}^\top \boldsymbol\beta\). Una
misma familia puede utilizar cualquier enlaces \(g\) que sea una función
monótona y diferenciable, pero el \textbf{enlace canónico} es aquel que
iguala el parámetro natural al predictor lineal, esto es,
\(\theta=\eta\). Este enlace canónico simplemente aporta ventajas
formales en las matemáticas, como estadísticos suficientes de dimensión
mínima o simplificación de las ecuaciones de verosimilitud.

En el caso binario de la familia binomial, las tres alternativas que
re-escalan un predictor lineal -teóricamente definido en toda la recta
real- a la escala de una probabilidad -en el rango {[}0,1{]}- son el
enlace logit, el probit y el complementario log-log -cloglog-, que
reproducen la forma de S invertida a la hora de relacionar el predictor
lineal con la probabilidad, como se aprecia en la
Figura~\ref{fig-u12-enlaces}. Logit y probit son \textbf{simétricos}
(casi intercambiables tras reescalar: el coeficiente probit ≈
logit/1,6); el \textbf{cloglog} (complementario log-log) es
\textbf{asimétrico} y será el enlace natural de la supervivencia en
tiempo discreto en la \protect\hyperlink{sec-u15}{Unidad 5} (Agresti,
2015).

\[
\text{logit: }\log\tfrac{\pi}{1-\pi}=\eta,\qquad
\text{probit: }\Phi^{-1}(\pi)=\eta,\qquad
\text{cloglog: }\log(-\log(1-\pi))=\eta .
\]

\begin{tcolorbox}[enhanced jigsaw, coltitle=black, rightrule=.15mm, arc=.35mm, leftrule=.75mm, colback=white, colbacktitle=quarto-callout-tip-color!10!white, left=2mm, breakable, toptitle=1mm, toprule=.15mm, bottomtitle=1mm, titlerule=0mm, title=\textcolor{quarto-callout-tip-color}{\faLightbulb}\hspace{0.5em}{🔧 En R. Especificar el enlace (logit, probit, cloglog)}, bottomrule=.15mm, opacityback=0, opacitybacktitle=0.6, colframe=quarto-callout-tip-color-frame]

En R, la forma de especificar los distintos logit se hace a través del
argumento
\texttt{link=\{\textquotesingle{}logit\textquotesingle{},\textquotesingle{}probit\textquotesingle{},\textquotesingle{}cloglog\textquotesingle{}\}},
o bien directamente como un argumento de la familia,
\texttt{family=\{binomial("logit"),\ binomial("probit"),\ binomial("cloglog")\}}.

\end{tcolorbox}

\begin{Shaded}
\begin{Highlighting}[]
\FunctionTok{tibble}\NormalTok{(}\AttributeTok{eta =} \FunctionTok{seq}\NormalTok{(}\SpecialCharTok{{-}}\DecValTok{5}\NormalTok{, }\DecValTok{5}\NormalTok{, }\AttributeTok{by =} \FloatTok{0.02}\NormalTok{)) }\SpecialCharTok{|\textgreater{}}
  \FunctionTok{mutate}\NormalTok{(}\AttributeTok{logit   =} \FunctionTok{plogis}\NormalTok{(eta),}
         \AttributeTok{probit  =} \FunctionTok{pnorm}\NormalTok{(eta),}
         \AttributeTok{cloglog =} \DecValTok{1} \SpecialCharTok{{-}} \FunctionTok{exp}\NormalTok{(}\SpecialCharTok{{-}}\FunctionTok{exp}\NormalTok{(eta))) }\SpecialCharTok{|\textgreater{}}
  \FunctionTok{pivot\_longer}\NormalTok{(}\FunctionTok{c}\NormalTok{(logit, probit, cloglog), }\AttributeTok{names\_to =} \StringTok{"enlace"}\NormalTok{, }\AttributeTok{values\_to =} \StringTok{"p"}\NormalTok{) }\SpecialCharTok{|\textgreater{}}
  \FunctionTok{ggplot}\NormalTok{(}\FunctionTok{aes}\NormalTok{(eta, p, }\AttributeTok{color =}\NormalTok{ enlace)) }\SpecialCharTok{+}
  \FunctionTok{geom\_line}\NormalTok{(}\AttributeTok{linewidth =} \DecValTok{1}\NormalTok{) }\SpecialCharTok{+}
  \FunctionTok{geom\_hline}\NormalTok{(}\AttributeTok{yintercept =} \FunctionTok{c}\NormalTok{(}\DecValTok{0}\NormalTok{, }\DecValTok{1}\NormalTok{), }\AttributeTok{linetype =} \StringTok{"dotted"}\NormalTok{) }\SpecialCharTok{+}
  \FunctionTok{labs}\NormalTok{(}\AttributeTok{x =} \FunctionTok{expression}\NormalTok{(eta), }\AttributeTok{y =} \FunctionTok{expression}\NormalTok{(p }\SpecialCharTok{==}\NormalTok{ g}\SpecialCharTok{\^{}}\NormalTok{\{}\SpecialCharTok{{-}}\DecValTok{1}\NormalTok{\}(eta)), }\AttributeTok{color =} \StringTok{"Enlace"}\NormalTok{)}
\end{Highlighting}
\end{Shaded}

\begin{figure}[H]

{\centering \includegraphics{caso1_binaria_files/figure-pdf/fig-u12-enlaces-1.pdf}

}

\caption{\label{fig-u12-enlaces}Funciones enlace en la familia binomial.
Las tres inversas devuelven η a (0,1). logit y probit son simétricas
(casi intercambiables tras reescalar); el cloglog es asimétrico.}

\end{figure}

Para ver el efecto de utilizar diferentes enlaces, ajustamos GLOW con
los tres y comparamos:

\begin{Shaded}
\begin{Highlighting}[]
\NormalTok{ajustes }\OtherTok{\textless{}{-}} \FunctionTok{list}\NormalTok{(}
  \AttributeTok{logit   =} \FunctionTok{glm}\NormalTok{(fracture }\SpecialCharTok{\textasciitilde{}}\NormalTok{ age }\SpecialCharTok{+}\NormalTok{ priorfrac, }\FunctionTok{binomial}\NormalTok{(}\StringTok{"logit"}\NormalTok{),   glow),}
  \AttributeTok{probit  =} \FunctionTok{glm}\NormalTok{(fracture }\SpecialCharTok{\textasciitilde{}}\NormalTok{ age }\SpecialCharTok{+}\NormalTok{ priorfrac, }\FunctionTok{binomial}\NormalTok{(}\StringTok{"probit"}\NormalTok{),  glow),}
  \AttributeTok{cloglog =} \FunctionTok{glm}\NormalTok{(fracture }\SpecialCharTok{\textasciitilde{}}\NormalTok{ age }\SpecialCharTok{+}\NormalTok{ priorfrac, }\FunctionTok{binomial}\NormalTok{(}\StringTok{"cloglog"}\NormalTok{), glow))}

\FunctionTok{map\_dfr}\NormalTok{(ajustes, \textbackslash{}(m) }\FunctionTok{glance}\NormalTok{(m) }\SpecialCharTok{|\textgreater{}} \FunctionTok{select}\NormalTok{(deviance, AIC), }\AttributeTok{.id =} \StringTok{"enlace"}\NormalTok{)}
\end{Highlighting}
\end{Shaded}

\begin{verbatim}
# A tibble: 3 x 3
  enlace  deviance   AIC
  <chr>      <dbl> <dbl>
1 logit       529.  535.
2 probit      528.  534.
3 cloglog     530.  536.
\end{verbatim}

\begin{Shaded}
\begin{Highlighting}[]
\FunctionTok{map\_dfr}\NormalTok{(ajustes, tidy, }\AttributeTok{.id =} \StringTok{"enlace"}\NormalTok{) }\SpecialCharTok{|\textgreater{}}
  \FunctionTok{select}\NormalTok{(enlace, term, estimate) }\SpecialCharTok{|\textgreater{}}
  \FunctionTok{pivot\_wider}\NormalTok{(}\AttributeTok{names\_from =}\NormalTok{ enlace, }\AttributeTok{values\_from =}\NormalTok{ estimate)}
\end{Highlighting}
\end{Shaded}

\begin{verbatim}
# A tibble: 3 x 4
  term           logit  probit cloglog
  <chr>          <dbl>   <dbl>   <dbl>
1 (Intercept)  -4.21   -2.54   -3.78  
2 age           0.0412  0.0248  0.0332
3 priorfracYes  0.839   0.508   0.680 
\end{verbatim}

Los coeficientes estimados \textbf{no} son directamente comparables
entre enlaces (viven en escalas distintas y como veremos, se interpretan
de forma distinta), pero el ajuste global, calibrado en términos del AIC
-que ya conocemos del modelo lineal- y la deviance -que presentaremos
más tarde-, es muy parecido, con lo cual en este caso concluiríamos que
\textbf{la elección de enlace apenas cambia las conclusiones}. Esto es
lo habitual salvo en comportamientos extremos en las colas o en datos
con mucha asimetría. Por interpretabilidad solemos utilizar el logit en
la familia binomial.

Esta equivalencia entre enlaces se aprecia en la escala de la
\textbf{probabilidad}, en la Figura~\ref{fig-u12-enlaces-prob}, que
superpone las curvas edad--probabilidad de los tres enlaces, separando
en dos paneles las predicciones según el antecedente de fractura previa.
Con el fin de ver la variación que se produce en las colas, y a pesar de
que no es pertinente la extrapolación de la predicción fuera del rango
observado, se han punteado las predicciones prolongadas fuera de dicho
rango, para visualizar la forma de S invertida completa, y las
diferencias en dichas colas.

\begin{Shaded}
\begin{Highlighting}[]
\NormalTok{rango\_obs }\OtherTok{\textless{}{-}} \FunctionTok{range}\NormalTok{(glow}\SpecialCharTok{$}\NormalTok{age)}
\NormalTok{grid }\OtherTok{\textless{}{-}} \FunctionTok{expand\_grid}\NormalTok{(}
  \AttributeTok{age       =} \FunctionTok{seq}\NormalTok{(}\DecValTok{40}\NormalTok{, }\DecValTok{105}\NormalTok{, }\AttributeTok{length.out =} \DecValTok{300}\NormalTok{),}
  \AttributeTok{priorfrac =} \FunctionTok{factor}\NormalTok{(}\FunctionTok{c}\NormalTok{(}\StringTok{"No"}\NormalTok{, }\StringTok{"Yes"}\NormalTok{), }\AttributeTok{levels =} \FunctionTok{levels}\NormalTok{(glow}\SpecialCharTok{$}\NormalTok{priorfrac)))}

\CommentTok{\# Predicción en probabilidad de cada uno de los tres modelos de \textasciigrave{}ajustes\textasciigrave{}}
\NormalTok{pred\_enlaces }\OtherTok{\textless{}{-}} \FunctionTok{map\_dfr}\NormalTok{(ajustes, \textbackslash{}(m) \{}
\NormalTok{  grid }\SpecialCharTok{|\textgreater{}} \FunctionTok{mutate}\NormalTok{(}\AttributeTok{p\_hat =} \FunctionTok{predict}\NormalTok{(m, }\AttributeTok{newdata =}\NormalTok{ grid, }\AttributeTok{type =} \StringTok{"response"}\NormalTok{))}
\NormalTok{\}, }\AttributeTok{.id =} \StringTok{"enlace"}\NormalTok{) }\SpecialCharTok{|\textgreater{}}
  \FunctionTok{mutate}\NormalTok{(}\AttributeTok{enlace =} \FunctionTok{factor}\NormalTok{(enlace, }\AttributeTok{levels =} \FunctionTok{c}\NormalTok{(}\StringTok{"logit"}\NormalTok{, }\StringTok{"probit"}\NormalTok{, }\StringTok{"cloglog"}\NormalTok{)))}
\NormalTok{pred\_obs }\OtherTok{\textless{}{-}} \FunctionTok{filter}\NormalTok{(pred\_enlaces, age }\SpecialCharTok{\textgreater{}=}\NormalTok{ rango\_obs[}\DecValTok{1}\NormalTok{], age }\SpecialCharTok{\textless{}=}\NormalTok{ rango\_obs[}\DecValTok{2}\NormalTok{])}

\NormalTok{etiq\_pf }\OtherTok{\textless{}{-}} \FunctionTok{c}\NormalTok{(}\AttributeTok{No =} \StringTok{"Sin fractura previa"}\NormalTok{, }\AttributeTok{Yes =} \StringTok{"Con fractura previa"}\NormalTok{)}

\FunctionTok{ggplot}\NormalTok{(pred\_enlaces, }\FunctionTok{aes}\NormalTok{(age, p\_hat, }\AttributeTok{color =}\NormalTok{ enlace)) }\SpecialCharTok{+}
  \FunctionTok{annotate}\NormalTok{(}\StringTok{"rect"}\NormalTok{, }\AttributeTok{xmin =}\NormalTok{ rango\_obs[}\DecValTok{1}\NormalTok{], }\AttributeTok{xmax =}\NormalTok{ rango\_obs[}\DecValTok{2}\NormalTok{],}
           \AttributeTok{ymin =} \SpecialCharTok{{-}}\ConstantTok{Inf}\NormalTok{, }\AttributeTok{ymax =} \ConstantTok{Inf}\NormalTok{, }\AttributeTok{alpha =} \FloatTok{0.05}\NormalTok{) }\SpecialCharTok{+}
  \FunctionTok{geom\_line}\NormalTok{(}\AttributeTok{linewidth =} \FloatTok{0.7}\NormalTok{, }\AttributeTok{linetype =} \StringTok{"22"}\NormalTok{) }\SpecialCharTok{+}          \CommentTok{\# extrapolación}
  \FunctionTok{geom\_line}\NormalTok{(}\AttributeTok{data =}\NormalTok{ pred\_obs, }\AttributeTok{linewidth =} \FloatTok{1.1}\NormalTok{) }\SpecialCharTok{+}          \CommentTok{\# rango observado}
  \FunctionTok{facet\_wrap}\NormalTok{(}\SpecialCharTok{\textasciitilde{}}\NormalTok{ priorfrac, }\AttributeTok{labeller =} \FunctionTok{labeller}\NormalTok{(}\AttributeTok{priorfrac =}\NormalTok{ etiq\_pf)) }\SpecialCharTok{+}
  \FunctionTok{scale\_y\_continuous}\NormalTok{(}\AttributeTok{limits =} \FunctionTok{c}\NormalTok{(}\DecValTok{0}\NormalTok{, }\DecValTok{1}\NormalTok{), }\AttributeTok{labels =}\NormalTok{ scales}\SpecialCharTok{::}\FunctionTok{label\_percent}\NormalTok{()) }\SpecialCharTok{+}
  \FunctionTok{labs}\NormalTok{(}\AttributeTok{x =} \StringTok{"Edad (años)"}\NormalTok{, }\AttributeTok{y =} \StringTok{"P(fractura)"}\NormalTok{, }\AttributeTok{color =} \StringTok{"Enlace"}\NormalTok{) }\SpecialCharTok{+}
  \FunctionTok{theme}\NormalTok{(}\AttributeTok{legend.position =} \StringTok{"bottom"}\NormalTok{)}
\end{Highlighting}
\end{Shaded}

\begin{figure}[H]

{\centering \includegraphics[width=1\textwidth,height=\textheight]{caso1_binaria_files/figure-pdf/fig-u12-enlaces-prob-1.pdf}

}

\caption{\label{fig-u12-enlaces-prob}Probabilidad de fractura predicha
por los tres enlaces (logit, probit, cloglog) sobre
\texttt{age\ +\ priorfrac}, en paneles por antecedente de fractura
previa. Trazo sólido y franja sombreada: rango observado en GLOW (55--90
años); punteado: extrapolación para apreciar la forma en S.}

\end{figure}

En el rango de edad \textbf{realmente observado} (zona sombreada, 55--90
años) las tres curvas son prácticamente \textbf{indistinguibles}: la
elección de enlace no mueve ni las probabilidades predichas ni, por
tanto, las conclusiones, en línea con la deviance y el AIC casi
idénticos de la tabla anterior. Las diferencias solo asoman al
\textbf{extrapolar} hacia las colas, donde el \texttt{cloglog}
---asimétrico--- se despega un poco de logit y probit (simétricos, casi
superpuestos entre sí). Es el comportamiento típico: el enlace importa
en las colas y bajo fuerte asimetría, no en el grueso de los datos, lo
que respalda quedarnos con el logit por interpretabilidad.

\hypertarget{sec-u12-interpretacion}{%
\subsection{2.3 Interpretación: del coeficiente a la
probabilidad}\label{sec-u12-interpretacion}}

\textbf{El enlace rompió la línea recta.} En OLS, \(\beta\) era el
cambio en la respuesta por cada unidad de \(x\): una lectura directa.
Aquí ya no. El modelo es lineal en el \textbf{predictor}
\(\eta=\beta_0+\beta x\), pero la respuesta es \(\mu=g^{-1}(\eta)\), una
transformación \textbf{no lineal} de \(\eta\). Por eso \(\beta\) ya
\textbf{no} mide el cambio en la probabilidad, sino el cambio en
\(\eta\): la misma diferencia en \(x\) mueve mucho la probabilidad en el
centro de la curva y casi nada en las colas. Para interpretar, leeremos
los coeficientes en la escala donde el modelo \textbf{sí} es lineal ---y
le pondremos nombre.

\textbf{Lo que da el modelo:} \texttt{summary()} (o
\texttt{broom::tidy}) devuelve los coeficientes en la escala del
predictor lineal:

Recordemos el modelo que estamos ajustando sobre GLOW:

\[\pi = P(Y=1) \ \rightarrow \ logit(\pi)=\theta + \beta \cdot \text{age} + \gamma \cdot \mathbf{1}(\text{priorfrac}=\text{Yes})\]

\begin{Shaded}
\begin{Highlighting}[]
\FunctionTok{tidy}\NormalTok{(fit\_glm)   }\CommentTok{\# coeficientes en escala log{-}odds: lo que da summary() por defecto}
\end{Highlighting}
\end{Shaded}

\begin{verbatim}
# A tibble: 3 x 5
  term         estimate std.error statistic     p.value
  <chr>           <dbl>     <dbl>     <dbl>       <dbl>
1 (Intercept)   -4.21      0.848      -4.97 0.000000667
2 age            0.0412    0.0122      3.38 0.000719   
3 priorfracYes   0.839     0.234       3.58 0.000340   
\end{verbatim}

Su \textbf{signo} ya informa ---positivo sube \(\eta\) y con él la
probabilidad; negativo la baja---, pero su \textbf{magnitud} no es un
cambio en probabilidad. ¿Qué es, entonces, esa escala?

\hypertarget{interpretaciuxf3n-en-la-escala-de-los-odds-y-log-odds}{%
\subsubsection{Interpretación en la escala de los odds y
log-odds}\label{interpretaciuxf3n-en-la-escala-de-los-odds-y-log-odds}}

Aparecen en la ecuación los \textbf{odds y log-odds}, que provienen del
enlace canónico de la Bernoulli.

\begin{itemize}
\tightlist
\item
  El \emph{odds} de un suceso es su \textbf{ventaja}, la probabilidad de
  que ocurra frente a la de que no: un odds de 2 dice que el evento es
  el doble de probable que su ausencia.
  \[\text{odds}(x)=\frac{\pi(x)}{1-\pi(x)}\]
\item
  Su logaritmo ---el \textbf{log-odds} o \emph{logit}--- es precisamente
  el enlace canónico, y por tanto \textbf{es} el predictor lineal:
  \[\eta(x)=\log\frac{\pi(x)}{1-\pi(x)}=\log.\text{odds}(x).\] La escala
  de los log.odds es donde viven los coeficientes y donde el modelo es
  una recta.
  \[ \log.\text{odds}(x)=\eta(x) \rightarrow \text{odds}(x)=e^{\eta(x)}\]
\end{itemize}

\textbf{Qué dice \(\beta\) ahí, y el odds ratio:}

\begin{itemize}
\tightlist
\item
  Subir \(x\) una unidad \textbf{suma} \(\beta\) al log-odds: el efecto
  es \textbf{aditivo}
  \[\beta=\eta(x+1)-\eta(x) \rightarrow \log.odds(x+1)= \beta +\log.odds(x)\]
\item
  Al deshacer el logaritmo surge el \textbf{odds ratio} (\(OR\)), como
  un cociente de odds:
  \[\beta=\log\frac{\text{odds}(x+1)}{\text{odds}(x)} = log(OR) \rightarrow OR=e^{\beta}\]
  de modo que el efecto en los \textbf{odds} es multiplicativo, y cada
  unidad de \(x\) multiplica los odds por el odds ratio:
  \[\text{odds}(x+1)=e^{\eta(x)+\beta} = \text{odds}(x) \cdot e^{\beta}= \text{odds}(x) \cdot OR\]
\end{itemize}

Por eso, si bien el log-odds depende de \(x\), el OR es un factor
\textbf{constante} que aplica a toda la curva, y de ahí su
interpretabilidad: cada movimiento de una unidad en \(x\), sea cual sea
el \(x\), multiplica los odds por el mismo factor \(e^{\beta}\). La
relación entre el signo de los coeficientes, el OR y el efecto que
genera en las probabilidades, es ahora más patente:

\begin{itemize}
\tightlist
\item
  \(\beta>0\Leftrightarrow\text{OR}>1\): aumenta la ventaja del evento;
\item
  \(\beta<0\Leftrightarrow\text{OR}<1\): reduce la ventaja del evento.
\end{itemize}

La Figura~\ref{fig-u12-tres-escalas} lo resume. Utiliza los valores
\(\beta=1\) (en azul) y \(\beta=-1\) (en rojo), y muestra el efecto de
variar el predictor \(x\) (en el eje X) sobre el log-odds, el odds y la
probabilidad del evento. A la izquierda, el log-odds η = β₀ + βx
identifica una recta cuya pendiente es β (efecto aditivo, constante). En
el centro, se aprecia la curva exponencial del odds = e\^{}η: cada
unidad de x multiplica las odds por e\^{}β. A la derecha, la relación
entre el predictor x y la probabilidad = 1/(1+e\^{}(−η)) viene
representada en forma de S, creciente o decreciente, con efecto máximo
en el centro y casi nulo en las colas. Esa no linealidad de la derecha
es justo lo que impide leer β como un cambio de probabilidad.

\begin{Shaded}
\begin{Highlighting}[]
\NormalTok{betas }\OtherTok{\textless{}{-}} \FunctionTok{c}\NormalTok{(}\StringTok{"β = +1 (positivo)"} \OtherTok{=} \DecValTok{1}\NormalTok{, }\StringTok{"β = −1 (negativo)"} \OtherTok{=} \SpecialCharTok{{-}}\DecValTok{1}\NormalTok{)}
\NormalTok{demo }\OtherTok{\textless{}{-}} \FunctionTok{expand\_grid}\NormalTok{(}\AttributeTok{x =} \FunctionTok{seq}\NormalTok{(}\SpecialCharTok{{-}}\DecValTok{3}\NormalTok{, }\DecValTok{3}\NormalTok{, }\AttributeTok{length.out =} \DecValTok{300}\NormalTok{), }\AttributeTok{signo =} \FunctionTok{names}\NormalTok{(betas)) }\SpecialCharTok{|\textgreater{}}
  \FunctionTok{mutate}\NormalTok{(}\AttributeTok{beta =}\NormalTok{ betas[signo],}
         \StringTok{\textasciigrave{}}\AttributeTok{log{-}odds}\StringTok{\textasciigrave{}} \OtherTok{=}\NormalTok{ beta }\SpecialCharTok{*}\NormalTok{ x,          }\CommentTok{\# β0 = 0}
         \AttributeTok{odds       =} \FunctionTok{exp}\NormalTok{(beta }\SpecialCharTok{*}\NormalTok{ x),}
         \AttributeTok{prob       =} \FunctionTok{plogis}\NormalTok{(beta }\SpecialCharTok{*}\NormalTok{ x))}

\NormalTok{col }\OtherTok{\textless{}{-}} \FunctionTok{c}\NormalTok{(}\StringTok{"β = +1 (positivo)"} \OtherTok{=} \StringTok{"\#2c7fb8"}\NormalTok{, }\StringTok{"β = −1 (negativo)"} \OtherTok{=} \StringTok{"\#d7301f"}\NormalTok{)}
\NormalTok{escala }\OtherTok{\textless{}{-}} \ControlFlowTok{function}\NormalTok{(y, ylab, titulo, y0)}
  \FunctionTok{ggplot}\NormalTok{(demo, }\FunctionTok{aes}\NormalTok{(x, .data[[y]], }\AttributeTok{color =}\NormalTok{ signo)) }\SpecialCharTok{+}
    \FunctionTok{geom\_hline}\NormalTok{(}\AttributeTok{yintercept =}\NormalTok{ y0, }\AttributeTok{linetype =} \StringTok{"dotted"}\NormalTok{, }\AttributeTok{color =} \StringTok{"grey60"}\NormalTok{) }\SpecialCharTok{+}
    \FunctionTok{geom\_line}\NormalTok{(}\AttributeTok{linewidth =} \FloatTok{1.1}\NormalTok{) }\SpecialCharTok{+}
    \FunctionTok{scale\_color\_manual}\NormalTok{(}\AttributeTok{values =}\NormalTok{ col) }\SpecialCharTok{+}
    \FunctionTok{labs}\NormalTok{(}\AttributeTok{x =} \StringTok{"predictor x"}\NormalTok{, }\AttributeTok{y =}\NormalTok{ ylab, }\AttributeTok{color =} \ConstantTok{NULL}\NormalTok{, }\AttributeTok{title =}\NormalTok{ titulo)}

\NormalTok{p1 }\OtherTok{\textless{}{-}} \FunctionTok{escala}\NormalTok{(}\StringTok{"log{-}odds"}\NormalTok{, }\StringTok{"log{-}odds (η)"}\NormalTok{, }\StringTok{"η = β₀ + βx  ·  recta (pendiente β)"}\NormalTok{, }\DecValTok{0}\NormalTok{)}
\NormalTok{p2 }\OtherTok{\textless{}{-}} \FunctionTok{escala}\NormalTok{(}\StringTok{"odds"}\NormalTok{,     }\StringTok{"odds = e\^{}η"}\NormalTok{,   }\StringTok{"odds  ·  multiplica e\^{}β por unidad"}\NormalTok{,            }\DecValTok{1}\NormalTok{)}
\NormalTok{p3 }\OtherTok{\textless{}{-}} \FunctionTok{escala}\NormalTok{(}\StringTok{"prob"}\NormalTok{,     }\StringTok{"P(evento)"}\NormalTok{,    }\StringTok{"probabilidad  ·  forma en S"}\NormalTok{,        }\FloatTok{0.5}\NormalTok{) }\SpecialCharTok{+}
        \FunctionTok{scale\_y\_continuous}\NormalTok{(}\AttributeTok{limits =} \FunctionTok{c}\NormalTok{(}\DecValTok{0}\NormalTok{, }\DecValTok{1}\NormalTok{), }\AttributeTok{labels =}\NormalTok{ scales}\SpecialCharTok{::}\FunctionTok{label\_percent}\NormalTok{())}

\NormalTok{(p1 }\SpecialCharTok{+}\NormalTok{ p2 }\SpecialCharTok{+}\NormalTok{ p3) }\SpecialCharTok{+} \FunctionTok{plot\_layout}\NormalTok{(}\AttributeTok{guides =} \StringTok{"collect"}\NormalTok{) }\SpecialCharTok{\&}
  \FunctionTok{theme}\NormalTok{(}\AttributeTok{legend.position =} \StringTok{"bottom"}\NormalTok{,}
        \AttributeTok{plot.title =} \FunctionTok{element\_text}\NormalTok{(}\AttributeTok{size =} \FloatTok{9.5}\NormalTok{, }\AttributeTok{face =} \StringTok{"bold"}\NormalTok{))}
\end{Highlighting}
\end{Shaded}

\begin{figure}[H]

{\centering \includegraphics[width=1\textwidth,height=\textheight]{caso1_binaria_files/figure-pdf/fig-u12-tres-escalas-1.pdf}

}

\caption{\label{fig-u12-tres-escalas}Efecto de la variación en un
predictor genérico x (β₀ = 0) sobre el log-odds, el odds y la
probabilidad del evento. En azul un coeficiente positivo (β = +1); en
rojo uno negativo (β = −1).}

\end{figure}

\begin{tcolorbox}[enhanced jigsaw, coltitle=black, rightrule=.15mm, arc=.35mm, leftrule=.75mm, colback=white, colbacktitle=quarto-callout-warning-color!10!white, left=2mm, breakable, toptitle=1mm, toprule=.15mm, bottomtitle=1mm, titlerule=0mm, title=\textcolor{quarto-callout-warning-color}{\faExclamationTriangle}\hspace{0.5em}{Odds ratio y riesgo relativo}, bottomrule=.15mm, opacityback=0, opacitybacktitle=0.6, colframe=quarto-callout-warning-color-frame]

Puntualizar que odds-ratio OR no es lo mismo que riesgo relativo RR. El
riesgo relativo es útil para comparar probabilidades, luego en el
contexto del modelo lineal cabría valorar el efecto, sobre la
probabilidad, de añadir una unidad a la \(x\)
\[\text{RR}=\pi(x+1)/\pi(x)\] En general \(\text{OR}\neq\text{RR}\), si
bien ambos se relacionan por \[
\text{OR}=\text{RR}\cdot\frac{1-\pi(x)}{1-\pi(x+1)},
\] de modo que el OR \textbf{se aleja de 1 más que el RR} y exagera el
efecto cuando el evento es frecuente. Solo si el evento es \textbf{raro}
(\(\pi\) pequeña) se cumple \(1-\pi\approx 1\) y entonces
\(\text{OR}\approx\text{RR}\).

Por eso decir «el riesgo se triplica» al leer un \(\text{OR}=3\) es una
aproximación cómoda pero informal: la lectura exacta es en términos de
\emph{odds} (ventaja), o bien en probabilidad mediante los efectos
marginales (más abajo).

\end{tcolorbox}

\textbf{Intervalos de confianza.} Como \(\text{OR}=e^{\beta}\) y la
exponencial es \textbf{monótona creciente}, el IC del OR es la
\textbf{exponencial de los extremos} del IC del coeficiente, sin
recalcular nada:
\[\text{IC}_{95\%}(\text{OR})=\big[e^{\beta_{\text{inf}}},\,e^{\beta_{\text{sup}}}\big].\]
Y como \(e^{0}=1\), que el IC de \(\beta\) \textbf{no contenga 0}
equivale a que el del OR \textbf{no contenga 1}: la misma señal de
efecto en las dos escalas.

\textbf{En GLOW.} Aplicado a las dos predictoras:

\begin{Shaded}
\begin{Highlighting}[]
\FunctionTok{exp}\NormalTok{(}\FunctionTok{cbind}\NormalTok{(}\AttributeTok{OR =} \FunctionTok{coef}\NormalTok{(fit\_glm), }\FunctionTok{confint}\NormalTok{(fit\_glm)))   }\CommentTok{\# OR e IC: exponencial del IC de beta}
\end{Highlighting}
\end{Shaded}

\begin{verbatim}
                     OR       2.5 %     97.5 %
(Intercept)  0.01478275 0.002737803 0.07645453
age          1.04205301 1.017574674 1.06743089
priorfracYes 2.31367006 1.458786254 3.65820301
\end{verbatim}

La fractura previa multiplica las odds de fractura por
\(\text{OR}\approx 2{,}31\) (IC 95 \%: 1,46--3,66) y cada año de edad
por \(\approx 1{,}04\) (IC 95 \%: 1,02--1,07); ninguno cruza 1, de modo
que ambos efectos son claros. (El IC de \(\beta\) procede del perfil de
verosimilitud, \protect\hyperlink{sec-u12-inferencia}{Sección 2.5}.)

\textbf{La misma idea en otros enlaces.} El predictor lineal \(\eta\) es
\textbf{idéntico} en los tres enlaces; cambia qué \emph{representa} y,
con ello, cómo se lee el coeficiente:

\hypertarget{tbl-enlaces-eta-mu}{}
\begin{longtable}[]{@{}
  >{\raggedright\arraybackslash}p{(\columnwidth - 6\tabcolsep) * \real{0.2500}}
  >{\raggedright\arraybackslash}p{(\columnwidth - 6\tabcolsep) * \real{0.2500}}
  >{\raggedright\arraybackslash}p{(\columnwidth - 6\tabcolsep) * \real{0.2500}}
  >{\raggedright\arraybackslash}p{(\columnwidth - 6\tabcolsep) * \real{0.2500}}@{}}
\caption{\label{tbl-enlaces-eta-mu}Relación entre el predictor lineal
\(\eta\) ---donde actúan los coeficientes--- y la media \(\mu\) en los
tres enlaces. En el logit, \(\eta\) es exactamente el log-odds; cada
enlace «lee» \(\eta\) en una escala distinta.}\tabularnewline
\toprule()
\begin{minipage}[b]{\linewidth}\raggedright
Enlace
\end{minipage} & \begin{minipage}[b]{\linewidth}\raggedright
\(\eta\) es\ldots{}
\end{minipage} & \begin{minipage}[b]{\linewidth}\raggedright
\(\eta=g(\mu)\)
\end{minipage} & \begin{minipage}[b]{\linewidth}\raggedright
\(\mu=g^{-1}(\eta)\)
\end{minipage} \\
\midrule()
\endfirsthead
\toprule()
\begin{minipage}[b]{\linewidth}\raggedright
Enlace
\end{minipage} & \begin{minipage}[b]{\linewidth}\raggedright
\(\eta\) es\ldots{}
\end{minipage} & \begin{minipage}[b]{\linewidth}\raggedright
\(\eta=g(\mu)\)
\end{minipage} & \begin{minipage}[b]{\linewidth}\raggedright
\(\mu=g^{-1}(\eta)\)
\end{minipage} \\
\midrule()
\endhead
\textbf{logit} & el \textbf{log-odds} & \(\log\dfrac{\mu}{1-\mu}\) &
\(\dfrac{e^{\eta}}{1+e^{\eta}}\) \\
\textbf{probit} & el cuantil normal (\(z\)) & \(\Phi^{-1}(\mu)\) &
\(\Phi(\eta)\) \\
\textbf{cloglog} & el log-riesgo (log del \emph{hazard} acumulado) &
\(\log\!\big(-\log(1-\mu)\big)\) & \(1-e^{-e^{\eta}}\) \\
\bottomrule()
\end{longtable}

De ahí que el mismo \(\beta\) signifique cosas distintas aunque el
ajuste apenas varíe:

\begin{itemize}
\tightlist
\item
  \textbf{logit:} \(e^{\beta}\) = \textbf{odds ratio} (multiplica las
  odds); la lectura habitual en clínica y epidemiología.
\item
  \textbf{probit:} \(\beta\) \textbf{desplaza la escala latente normal}
  (\(z\)), sin lectura de odds; como guía,
  \(\beta_{\text{probit}}\approx\beta_{\text{logit}}/1{,}6\text{–}1{,}8\).
\item
  \textbf{cloglog:} \(e^{\beta}\) = \textbf{razón de riesgos}
  (\emph{hazard ratio}); lectura natural en supervivencia
  (\protect\hyperlink{sec-u15}{Unidad 5}).
\end{itemize}

\begin{Shaded}
\begin{Highlighting}[]
\CommentTok{\# Mismos datos, distinta escala del coeficiente según el enlace}
\FunctionTok{map\_dfr}\NormalTok{(ajustes, \textbackslash{}(m) }\FunctionTok{tidy}\NormalTok{(m) }\SpecialCharTok{|\textgreater{}} \FunctionTok{select}\NormalTok{(term, estimate), }\AttributeTok{.id =} \StringTok{"enlace"}\NormalTok{) }\SpecialCharTok{|\textgreater{}}
  \FunctionTok{pivot\_wider}\NormalTok{(}\AttributeTok{names\_from =}\NormalTok{ enlace, }\AttributeTok{values\_from =}\NormalTok{ estimate) }\SpecialCharTok{|\textgreater{}}
  \FunctionTok{mutate}\NormalTok{(}\AttributeTok{razon\_logit\_probit =}\NormalTok{ logit }\SpecialCharTok{/}\NormalTok{ probit)   }\CommentTok{\# \textasciitilde{}1,6{-}1,8 (salvo intercepto)}
\end{Highlighting}
\end{Shaded}

\begin{verbatim}
# A tibble: 3 x 5
  term           logit  probit cloglog razon_logit_probit
  <chr>          <dbl>   <dbl>   <dbl>              <dbl>
1 (Intercept)  -4.21   -2.54   -3.78                 1.66
2 age           0.0412  0.0248  0.0332               1.66
3 priorfracYes  0.839   0.508   0.680                1.65
\end{verbatim}

\hypertarget{interpretaciuxf3n-en-la-escala-de-la-probabilidad-efectos-marginales}{%
\subsubsection{Interpretación en la escala de la probabilidad: efectos
marginales}\label{interpretaciuxf3n-en-la-escala-de-la-probabilidad-efectos-marginales}}

Cuando queremos interpretar un ajuste directamente en términos de las
probabilidades, acudimos a los \textbf{efectos marginales}. El odds
ratio es constante, pero el cambio en \textbf{probabilidad} no: un año
más de \texttt{age} apenas mueve nada en las colas y mucho en el centro
de la S. Para resumirlo en un número, el \textbf{efecto marginal medio
(AME)} responde a la pregunta que de verdad importa:

\begin{quote}
``¿Cuántos puntos de probabilidad sube el evento, en promedio, por una
unidad más de la covariable?
\end{quote}

En GLOW, el AME de \texttt{age} es ≈ \textbf{0,007}: unos \textbf{0,7
puntos porcentuales de riesgo de fractura por año}.

\textbf{¿De dónde sale ese número?}

Para calcular el AME o efecto medio marginal asociado a un predictor
continuo \(x_k\) (con coeficiente \(\beta_k\) en el predictor lineal),
se considera el efecto de cada individuo a través de la pendiente de su
propia curva, y dado que la derivada del logit es \(\mu(1-\mu)\), esta
pendiente es:
\[\frac{\partial\mu_i}{\partial x_k}=\hat\pi_i(1-\hat\pi_i)\,\beta_k.\]
El AME es el promedio sobre las pendientes para las \(n\) personas
observadas: literalmente, el coeficiente por la pendiente media de la
sigmoide.
\[\text{AME}_k=\beta_k\cdot\frac{1}{n}\sum_i \hat\pi_i(1-\hat\pi_i)\]

\begin{Shaded}
\begin{Highlighting}[]
\CommentTok{\# AME de cada covariable, en puntos de probabilidad}
\FunctionTok{avg\_slopes}\NormalTok{(fit\_glm)                       }
\end{Highlighting}
\end{Shaded}

\begin{verbatim}

      Term Contrast Estimate Std. Error    z Pr(>|z|)    S   2.5 % 97.5 %
 age       dY/dX     0.00717    0.00205 3.50   <0.001 11.1 0.00316 0.0112
 priorfrac Yes - No  0.16353    0.04900 3.34   <0.001 10.2 0.06749 0.2596

Type: response
\end{verbatim}

\begin{Shaded}
\begin{Highlighting}[]
\CommentTok{\# Reproducción "a mano" del cálculo del AME para \textquotesingle{}age\textquotesingle{}:}
\NormalTok{p }\OtherTok{\textless{}{-}} \FunctionTok{fitted}\NormalTok{(fit\_glm)                      }\CommentTok{\# pi\_hat de cada individuo}
\FunctionTok{mean}\NormalTok{(p }\SpecialCharTok{*}\NormalTok{ (}\DecValTok{1} \SpecialCharTok{{-}}\NormalTok{ p)) }\SpecialCharTok{*} \FunctionTok{coef}\NormalTok{(fit\_glm)[}\StringTok{"age"}\NormalTok{]  }\CommentTok{\# = AME de age}
\end{Highlighting}
\end{Shaded}

\begin{verbatim}
        age 
0.007174101 
\end{verbatim}

Cuando se trata de evaluar el AME para un factor como
\texttt{priorfrac}, el coeficiente no es una pendiente (derivada) sino
una \textbf{diferencia contrafactual}. El AME entonces se calcula con el
efecto diferencial sobre la categoría de referencia (en
\texttt{priorfrac}, \(\hat\pi\) con \texttt{Yes} menos con \texttt{No}),
promediada entre individuos, es decir: \[
\text{AME}_{\text{priorfrac}}=\frac{1}{n}\sum_{i=1}^{n}\Big[\hat\pi\big(\texttt{Yes},\,x_i\big)-\hat\pi\big(\texttt{No},\,x_i\big)\Big],
\] donde \(\hat\pi(\texttt{nivel},x_i)\) se evalúa para todos los
sujetos en la muestra en todos los niveles (en \texttt{priorfrac}, Yes y
No). \[
\hat\pi(\text{nivel},x_i)=\operatorname{logit}^{-1}\!\big(\hat\theta+\hat\beta\,\text{age}_i+\hat\delta_{\text{pf}}^{\,\text{nivel}}\big),\qquad \hat\delta_{\text{pf}}^{\,\text{ref}}=0.
\]

\begin{Shaded}
\begin{Highlighting}[]
\CommentTok{\# Reproducción "a mano" del cálculo de AME para \textquotesingle{}priorfrac\textquotesingle{}}
\NormalTok{g\_yes }\OtherTok{\textless{}{-}} \FunctionTok{transform}\NormalTok{(glow, }\AttributeTok{priorfrac =} \FunctionTok{factor}\NormalTok{(}\StringTok{"Yes"}\NormalTok{, }\AttributeTok{levels =} \FunctionTok{levels}\NormalTok{(glow}\SpecialCharTok{$}\NormalTok{priorfrac)))}
\NormalTok{g\_no  }\OtherTok{\textless{}{-}} \FunctionTok{transform}\NormalTok{(glow, }\AttributeTok{priorfrac =} \FunctionTok{factor}\NormalTok{(}\StringTok{"No"}\NormalTok{,  }\AttributeTok{levels =} \FunctionTok{levels}\NormalTok{(glow}\SpecialCharTok{$}\NormalTok{priorfrac)))}

\FunctionTok{mean}\NormalTok{(}\FunctionTok{predict}\NormalTok{(fit\_glm, g\_yes, }\AttributeTok{type =} \StringTok{"response"}\NormalTok{) }\SpecialCharTok{{-}}
     \FunctionTok{predict}\NormalTok{(fit\_glm, g\_no,  }\AttributeTok{type =} \StringTok{"response"}\NormalTok{))   }\CommentTok{\# = AME de priorfrac}
\end{Highlighting}
\end{Shaded}

\begin{verbatim}
[1] 0.1635282
\end{verbatim}

El AME en probabilidad es, además, \textbf{casi invariante al enlace}
(los tres ajustes lo dan parecido pese a coeficientes distintos) y se
prefiere frente a otros indicadore por promediar sobre la población
observada (Arel-Bundock et~al., 2024). En una frase: \textgreater El
\textbf{OR} comunica la fuerza del efecto; el \textbf{AME}, su impacto
absoluto en probabilidad.

En la Figura~\ref{fig-u12-dos-escalas} se visualiza el efecto de la edad
en dos escalas. A la izquierda, en escala log-odds, los dos perfiles son
\textbf{rectas paralelas}: la edad entra de forma lineal y
\texttt{priorfrac} solo desplaza la recta una cantidad
\textbf{constante}, \(\log(\text{OR})\approx 0{,}84\) ---ahí vive el
odds ratio---. A la derecha, en probabilidad, esas mismas rectas se
vuelven \textbf{curvas en S}: el efecto de \texttt{priorfrac} (la
distancia entre curvas) es máximo en el centro y se desvanece en los
extremos, \textbf{aunque el OR sea constante}. Claramente se aprecia la
diferencia constante en los odds y variable en las probabilidades.

\begin{Shaded}
\begin{Highlighting}[]
\NormalTok{rango\_obs }\OtherTok{\textless{}{-}} \FunctionTok{range}\NormalTok{(glow}\SpecialCharTok{$}\NormalTok{age)        }\CommentTok{\# \textasciitilde{}55–90: edad observada en GLOW}
\NormalTok{grid }\OtherTok{\textless{}{-}} \FunctionTok{expand\_grid}\NormalTok{(}
  \AttributeTok{age       =} \FunctionTok{seq}\NormalTok{(}\DecValTok{40}\NormalTok{, }\DecValTok{105}\NormalTok{, }\AttributeTok{length.out =} \DecValTok{300}\NormalTok{),}
  \AttributeTok{priorfrac =} \FunctionTok{factor}\NormalTok{(}\FunctionTok{c}\NormalTok{(}\StringTok{"No"}\NormalTok{, }\StringTok{"Yes"}\NormalTok{), }\AttributeTok{levels =} \FunctionTok{levels}\NormalTok{(glow}\SpecialCharTok{$}\NormalTok{priorfrac)))}
\NormalTok{pr   }\OtherTok{\textless{}{-}} \FunctionTok{predict}\NormalTok{(fit\_glm, }\AttributeTok{newdata =}\NormalTok{ grid, }\AttributeTok{type =} \StringTok{"link"}\NormalTok{, }\AttributeTok{se.fit =} \ConstantTok{TRUE}\NormalTok{)}
\NormalTok{grid }\OtherTok{\textless{}{-}}\NormalTok{ grid }\SpecialCharTok{|\textgreater{}}
  \FunctionTok{mutate}\NormalTok{(}\AttributeTok{eta =}\NormalTok{ pr}\SpecialCharTok{$}\NormalTok{fit, }\AttributeTok{se =}\NormalTok{ pr}\SpecialCharTok{$}\NormalTok{se.fit,}
         \AttributeTok{lo\_eta =}\NormalTok{ eta }\SpecialCharTok{{-}} \FloatTok{1.96} \SpecialCharTok{*}\NormalTok{ se, }\AttributeTok{hi\_eta =}\NormalTok{ eta }\SpecialCharTok{+} \FloatTok{1.96} \SpecialCharTok{*}\NormalTok{ se,   }\CommentTok{\# IC 95 \% en log{-}odds}
         \AttributeTok{p\_hat  =} \FunctionTok{plogis}\NormalTok{(eta),}
         \AttributeTok{lo\_p   =} \FunctionTok{plogis}\NormalTok{(lo\_eta), }\AttributeTok{hi\_p =} \FunctionTok{plogis}\NormalTok{(hi\_eta))       }\CommentTok{\# IC 95 \% en probabilidad}
\NormalTok{grid\_obs }\OtherTok{\textless{}{-}} \FunctionTok{filter}\NormalTok{(grid, age }\SpecialCharTok{\textgreater{}=}\NormalTok{ rango\_obs[}\DecValTok{1}\NormalTok{], age }\SpecialCharTok{\textless{}=}\NormalTok{ rango\_obs[}\DecValTok{2}\NormalTok{])}

\NormalTok{pA }\OtherTok{\textless{}{-}} \FunctionTok{ggplot}\NormalTok{(grid, }\FunctionTok{aes}\NormalTok{(age, eta, }\AttributeTok{color =}\NormalTok{ priorfrac, }\AttributeTok{fill =}\NormalTok{ priorfrac)) }\SpecialCharTok{+}
  \FunctionTok{annotate}\NormalTok{(}\StringTok{"rect"}\NormalTok{, }\AttributeTok{xmin =}\NormalTok{ rango\_obs[}\DecValTok{1}\NormalTok{], }\AttributeTok{xmax =}\NormalTok{ rango\_obs[}\DecValTok{2}\NormalTok{],}
           \AttributeTok{ymin =} \SpecialCharTok{{-}}\ConstantTok{Inf}\NormalTok{, }\AttributeTok{ymax =} \ConstantTok{Inf}\NormalTok{, }\AttributeTok{alpha =} \FloatTok{0.06}\NormalTok{) }\SpecialCharTok{+}
  \FunctionTok{geom\_ribbon}\NormalTok{(}\FunctionTok{aes}\NormalTok{(}\AttributeTok{ymin =}\NormalTok{ lo\_eta, }\AttributeTok{ymax =}\NormalTok{ hi\_eta), }\AttributeTok{alpha =} \FloatTok{0.15}\NormalTok{, }\AttributeTok{color =} \ConstantTok{NA}\NormalTok{) }\SpecialCharTok{+}
  \FunctionTok{geom\_line}\NormalTok{(}\AttributeTok{linewidth =} \FloatTok{0.7}\NormalTok{, }\AttributeTok{linetype =} \StringTok{"22"}\NormalTok{) }\SpecialCharTok{+}          \CommentTok{\# extrapolación}
  \FunctionTok{geom\_line}\NormalTok{(}\AttributeTok{data =}\NormalTok{ grid\_obs, }\AttributeTok{linewidth =} \FloatTok{1.2}\NormalTok{) }\SpecialCharTok{+}          \CommentTok{\# rango observado}
  \FunctionTok{labs}\NormalTok{(}\AttributeTok{x =} \StringTok{"Edad (años)"}\NormalTok{, }\AttributeTok{y =} \StringTok{"log{-}odds (η)"}\NormalTok{,}
       \AttributeTok{color =} \StringTok{"Fractura previa"}\NormalTok{, }\AttributeTok{fill =} \StringTok{"Fractura previa"}\NormalTok{,}
       \AttributeTok{title =} \StringTok{"Escala log{-}odds"}\NormalTok{, }\AttributeTok{subtitle =} \StringTok{"rectas paralelas (separación ≈ 0,84)"}\NormalTok{)}

\NormalTok{pB }\OtherTok{\textless{}{-}} \FunctionTok{ggplot}\NormalTok{(grid, }\FunctionTok{aes}\NormalTok{(age, p\_hat, }\AttributeTok{color =}\NormalTok{ priorfrac, }\AttributeTok{fill =}\NormalTok{ priorfrac)) }\SpecialCharTok{+}
  \FunctionTok{annotate}\NormalTok{(}\StringTok{"rect"}\NormalTok{, }\AttributeTok{xmin =}\NormalTok{ rango\_obs[}\DecValTok{1}\NormalTok{], }\AttributeTok{xmax =}\NormalTok{ rango\_obs[}\DecValTok{2}\NormalTok{],}
           \AttributeTok{ymin =} \SpecialCharTok{{-}}\ConstantTok{Inf}\NormalTok{, }\AttributeTok{ymax =} \ConstantTok{Inf}\NormalTok{, }\AttributeTok{alpha =} \FloatTok{0.06}\NormalTok{) }\SpecialCharTok{+}
  \FunctionTok{geom\_ribbon}\NormalTok{(}\FunctionTok{aes}\NormalTok{(}\AttributeTok{ymin =}\NormalTok{ lo\_p, }\AttributeTok{ymax =}\NormalTok{ hi\_p), }\AttributeTok{alpha =} \FloatTok{0.15}\NormalTok{, }\AttributeTok{color =} \ConstantTok{NA}\NormalTok{) }\SpecialCharTok{+}
  \FunctionTok{geom\_line}\NormalTok{(}\AttributeTok{linewidth =} \FloatTok{0.7}\NormalTok{, }\AttributeTok{linetype =} \StringTok{"22"}\NormalTok{) }\SpecialCharTok{+}
  \FunctionTok{geom\_line}\NormalTok{(}\AttributeTok{data =}\NormalTok{ grid\_obs, }\AttributeTok{linewidth =} \FloatTok{1.2}\NormalTok{) }\SpecialCharTok{+}
  \FunctionTok{scale\_y\_continuous}\NormalTok{(}\AttributeTok{limits =} \FunctionTok{c}\NormalTok{(}\DecValTok{0}\NormalTok{, }\DecValTok{1}\NormalTok{), }\AttributeTok{labels =}\NormalTok{ scales}\SpecialCharTok{::}\FunctionTok{label\_percent}\NormalTok{()) }\SpecialCharTok{+}
  \FunctionTok{labs}\NormalTok{(}\AttributeTok{x =} \StringTok{"Edad (años)"}\NormalTok{, }\AttributeTok{y =} \StringTok{"P(fractura)"}\NormalTok{,}
       \AttributeTok{color =} \StringTok{"Fractura previa"}\NormalTok{, }\AttributeTok{fill =} \StringTok{"Fractura previa"}\NormalTok{,}
       \AttributeTok{title =} \StringTok{"Escala probabilidad"}\NormalTok{, }\AttributeTok{subtitle =} \StringTok{"curvas en S (separación variable)"}\NormalTok{)}

\NormalTok{(pA }\SpecialCharTok{+}\NormalTok{ pB) }\SpecialCharTok{+}
  \FunctionTok{plot\_layout}\NormalTok{(}\AttributeTok{guides =} \StringTok{"collect"}\NormalTok{) }\SpecialCharTok{\&}
  \FunctionTok{theme}\NormalTok{(}\AttributeTok{legend.position =} \StringTok{"bottom"}\NormalTok{,}
        \AttributeTok{plot.title    =} \FunctionTok{element\_text}\NormalTok{(}\AttributeTok{size =} \DecValTok{11}\NormalTok{, }\AttributeTok{face =} \StringTok{"bold"}\NormalTok{),}
        \AttributeTok{plot.subtitle =} \FunctionTok{element\_text}\NormalTok{(}\AttributeTok{size =} \FloatTok{8.5}\NormalTok{))}
\end{Highlighting}
\end{Shaded}

\begin{figure}[H]

{\centering \includegraphics[width=1\textwidth,height=\textheight]{caso1_binaria_files/figure-pdf/fig-u12-dos-escalas-1.pdf}

}

\caption{\label{fig-u12-dos-escalas}El mismo modelo en dos escalas.
Izquierda (log-odds): rectas paralelas, separación constante ≈ 0,84
(logaritmo del OR). Derecha (probabilidad): curvas en S, con separación
variable. Trazo sólido y franja sombreada: rango de edad observado en
GLOW (55--90 años); fuera, extrapolación (punteado). Bandas: IC 95 \%.}

\end{figure}

\begin{tcolorbox}[enhanced jigsaw, coltitle=black, rightrule=.15mm, arc=.35mm, leftrule=.75mm, colback=white, colbacktitle=quarto-callout-tip-color!10!white, left=2mm, breakable, toptitle=1mm, toprule=.15mm, bottomtitle=1mm, titlerule=0mm, title=\textcolor{quarto-callout-tip-color}{\faLightbulb}\hspace{0.5em}{🔧 En R. Odds ratios, probabilidades y efectos marginales}, bottomrule=.15mm, opacityback=0, opacitybacktitle=0.6, colframe=quarto-callout-tip-color-frame]

Tres escalas de lectura, cada una con su herramienta:

\begin{itemize}
\tightlist
\item
  \textbf{Odds ratios con intervalo} ---
  \texttt{broom::tidy(m,\ exponentiate\ =\ TRUE,\ conf.int\ =\ TRUE)}, o
  a mano \texttt{exp(cbind(OR\ =\ coef(m),\ confint(m)))}.
  \texttt{confint()} da el IC por \textbf{perfil de verosimilitud} (más
  fiable que Wald); que el IC del OR no contenga el \textbf{1} equivale
  a que el de \(\beta\) no contenga el 0.
\item
  \textbf{Probabilidades predichas} ---
  \texttt{predict(m,\ newdata,\ type\ =\ "response")} da \(\hat\pi\);
  con \texttt{type\ =\ "link",\ se.fit\ =\ TRUE} se obtiene \(\hat\eta\)
  y su error, y el IC se construye en la escala del enlace y \textbf{se
  pasa por la sigmoide} (\texttt{plogis}) ---nunca sumando ± en
  probabilidad---.
\item
  \textbf{Efecto marginal medio (AME)} ---
  \texttt{marginaleffects::avg\_slopes(m)}: el efecto en \textbf{puntos
  de probabilidad}, casi invariante al enlace.
  \texttt{marginaleffects::predictions()} / \texttt{emmeans} dan
  probabilidades por perfil. El \textbf{OR} comunica la fuerza del
  efecto; el \textbf{AME}, su impacto absoluto.
\end{itemize}

\end{tcolorbox}

\hypertarget{sec-u12-iwls}{%
\subsection{2.4 Estimación: máxima verosimilitud e
IWLS}\label{sec-u12-iwls}}

A diferencia de OLS, \textbf{no hay fórmula cerrada} para
\(\hat{\boldsymbol\beta}\): se maximiza la log-verosimilitud
numéricamente. Pero ese ``numéricamente'' no es una caja negra; al
abrirlo aparece algo muy reconocible ---una sucesión de regresiones
lineales ponderadas---. Veámoslo en tres pasos: las ecuaciones que hay
que resolver, el algoritmo que las resuelve y su reproducción a mano
sobre GLOW.

\textbf{Las ecuaciones de score: ¿dónde están los \(\boldsymbol\beta\)?}

La estimación de los coeficientes \(\beta\) del predictor lineal pasa
por maximizar la verosimilitud, o lo que es lo mismo, la
log-verosimilitud \(l(\beta)\). Maximizar
\(\ell(\boldsymbol\beta)=\sum_i \ell_i\) exige anular su gradiente. Como
los coeficientes entran en la log-verosimilitud a través del predictor
lineal, \(\eta_i=\mathbf{x}_i^\top\boldsymbol\beta\), aplicando la regla
de la cadena en la derivación,
\(\ell_i\to\mu_i\to\eta_i\to\boldsymbol\beta\), y dado que que
\(\partial\eta_i/\partial\boldsymbol\beta=\mathbf{x}_i\), llegamos a:

\[
\frac{\partial \ell}{\partial\boldsymbol\beta}
=\sum_{i=1}^{n}
\underbrace{\frac{\partial \ell_i}{\partial\mu_i}}_{\textstyle \frac{y_i-\mu_i}{a(\phi)V(\mu_i)}}\;
\underbrace{\frac{\partial\mu_i}{\partial\eta_i}}_{\text{enlace}}\;
\underbrace{\frac{\partial\eta_i}{\partial\boldsymbol\beta}}_{\textstyle \mathbf{x}_i}
=\sum_{i=1}^{n}\frac{(y_i-\mu_i)}{a(\phi)\,V(\mu_i)}\,
\frac{\mathrm{d}\mu_i}{\mathrm{d}\eta_i}\,\mathbf{x}_i=\mathbf{0}.
\]

Resulta pues un sistema que iguala a cero una \textbf{suma ponderada de
residuos} \((y_i-\mu_i)\), que no tiene solución cerrada y requiere de
un procedimiento iterativo de optimización.

\textbf{Del sistema al algoritmo: Fisher scoring = IWLS (\emph{Iterative
Weighted Least Squares})}

El método de \textbf{Newton--Raphson}, usando la \textbf{información
esperada} (\emph{Fisher scoring}), resuelve ese sistema, y tiene una
lectura ``interesante'': \textbf{cada iteración equivale a una regresión
por mínimos cuadrados ponderados} (IWLS) (McCullagh \& Nelder, 1989;
Nelder \& Wedderburn, 1972). En cada paso se fabrican: una
\textbf{respuesta de trabajo} \(z_i\) (una linealización del enlace
alrededor del ajuste actual) y unos \textbf{pesos} \(w_i\) (que dan más
voz a las observaciones más precisas):

\[
z_i=\eta_i+(y_i-\mu_i)\frac{\mathrm{d}\eta_i}{\mathrm{d}\mu_i},
\qquad
w_i=\frac{(\mathrm{d}\mu_i/\mathrm{d}\eta_i)^2}{V(\mu_i)},
\]

y se actualiza
\(\boldsymbol\beta\leftarrow(\mathbf{X}^\top\mathbf{W}\mathbf{X})^{-1} \mathbf{X}^\top\mathbf{W}\mathbf{z}\)
---la fórmula de WLS--- hasta que deja de cambiar. Así, \textbf{estimar
un GLM es resolver una sucesión de regresiones lineales ponderadas}.

\begin{tcolorbox}[enhanced jigsaw, coltitle=black, rightrule=.15mm, arc=.35mm, leftrule=.75mm, colback=white, colbacktitle=quarto-callout-note-color!10!white, left=2mm, breakable, toptitle=1mm, toprule=.15mm, bottomtitle=1mm, titlerule=0mm, title=\textcolor{quarto-callout-note-color}{\faInfo}\hspace{0.5em}{Recordatorio: Newton--Raphson y Fisher scoring}, bottomrule=.15mm, opacityback=0, opacitybacktitle=0.6, colframe=quarto-callout-note-color-frame]

Para hallar la raíz de la ecuación de score
\(\partial\ell/\partial\boldsymbol\beta=\mathbf{0}\),
\textbf{Newton--Raphson} itera
\(\boldsymbol\beta\leftarrow\boldsymbol\beta-\mathbf{H}^{-1}\mathbf{s}\),
con \(\mathbf{s}\) el score (gradiente) y \(\mathbf{H}\) el hessiano
(segunda derivada). \textbf{Fisher scoring} sustituye \(\mathbf{H}\) por
su valor esperado, la \textbf{información esperada}
\(\mathcal{I}=\mathbb{E}[-\mathbf{H}]\): más estable y, en un GLM, lo
que convierte cada paso en una regresión ponderada.

\end{tcolorbox}

\begin{tcolorbox}[enhanced jigsaw, coltitle=black, rightrule=.15mm, arc=.35mm, leftrule=.75mm, colback=white, colbacktitle=quarto-callout-note-color!10!white, left=2mm, breakable, toptitle=1mm, toprule=.15mm, bottomtitle=1mm, titlerule=0mm, title=\textcolor{quarto-callout-note-color}{\faInfo}\hspace{0.5em}{Recordatorio: mínimos cuadrados ponderados (WLS)}, bottomrule=.15mm, opacityback=0, opacitybacktitle=0.6, colframe=quarto-callout-note-color-frame]

OLS minimiza \(\sum_i (z_i-\mathbf{x}_i^\top\boldsymbol\beta)^2\)
tratando todas las observaciones por igual. \textbf{WLS} asigna a cada
una un peso \(w_i\) y minimiza
\(\sum_i w_i\,(z_i-\mathbf{x}_i^\top\boldsymbol\beta)^2\), dando más
influencia a las de menor varianza; su solución es
\(\hat{\boldsymbol\beta}=(\mathbf{X}^\top\mathbf{W}\mathbf{X})^{-1} \mathbf{X}^\top\mathbf{W}\mathbf{z}\),
con \(\mathbf{W}=\operatorname{diag}(w_i)\). El IWLS no es más que
aplicarla una y otra vez, recalculando \(z\) y \(w\) en cada vuelta.

\end{tcolorbox}

\textbf{Reproducción a mano sobre GLOW.}

Reproducimos ahora el bucle con las piezas que el propio objeto
\texttt{family()} proporciona (\texttt{linkinv}, \texttt{mu.eta},
\texttt{variance}) y comprobamos que coincide con la solución de
\texttt{glm()}:

\begin{Shaded}
\begin{Highlighting}[]
\NormalTok{X   }\OtherTok{\textless{}{-}} \FunctionTok{model.matrix}\NormalTok{(fit\_glm)          }\CommentTok{\# matriz de diseño: columnas (Intercept), age, priorfracYes}
\NormalTok{y   }\OtherTok{\textless{}{-}}\NormalTok{ glow}\SpecialCharTok{$}\NormalTok{fractura01                }\CommentTok{\# respuesta binaria 0/1}
\NormalTok{fam }\OtherTok{\textless{}{-}} \FunctionTok{binomial}\NormalTok{()                     }\CommentTok{\# la familia expone las piezas: linkinv, mu.eta, variance}

\NormalTok{beta  }\OtherTok{\textless{}{-}} \FunctionTok{rep}\NormalTok{(}\DecValTok{0}\NormalTok{, }\FunctionTok{ncol}\NormalTok{(X))              }\CommentTok{\# arranque: todos los coeficientes a 0}
\NormalTok{traza }\OtherTok{\textless{}{-}} \FunctionTok{list}\NormalTok{()                       }\CommentTok{\# aquí guardaremos beta de cada iteración (para la figura)}

\ControlFlowTok{for}\NormalTok{ (it }\ControlFlowTok{in} \DecValTok{1}\SpecialCharTok{:}\DecValTok{25}\NormalTok{) \{}
\NormalTok{  eta }\OtherTok{\textless{}{-}} \FunctionTok{as.vector}\NormalTok{(X }\SpecialCharTok{\%*\%}\NormalTok{ beta)        }\CommentTok{\# predictor lineal      eta = X beta}
\NormalTok{  mu  }\OtherTok{\textless{}{-}}\NormalTok{ fam}\SpecialCharTok{$}\FunctionTok{linkinv}\NormalTok{(eta)             }\CommentTok{\# media ajustada        mu  = g\^{}\{{-}1\}(eta) = plogis(eta)}
\NormalTok{  mep }\OtherTok{\textless{}{-}}\NormalTok{ fam}\SpecialCharTok{$}\FunctionTok{mu.eta}\NormalTok{(eta)              }\CommentTok{\# derivada del enlace    dmu/deta = mu(1{-}mu)}
\NormalTok{  z   }\OtherTok{\textless{}{-}}\NormalTok{ eta }\SpecialCharTok{+}\NormalTok{ (y }\SpecialCharTok{{-}}\NormalTok{ mu) }\SpecialCharTok{/}\NormalTok{ mep         }\CommentTok{\# respuesta de trabajo  (linealización de un paso)}
\NormalTok{  w   }\OtherTok{\textless{}{-}}\NormalTok{ mep}\SpecialCharTok{\^{}}\DecValTok{2} \SpecialCharTok{/}\NormalTok{ fam}\SpecialCharTok{$}\FunctionTok{variance}\NormalTok{(mu)     }\CommentTok{\# pesos IWLS             w = (dmu/deta)\^{}2 / V(mu)}
\NormalTok{  WX  }\OtherTok{\textless{}{-}}\NormalTok{ X }\SpecialCharTok{*}\NormalTok{ w                        }\CommentTok{\# cada fila de X escalada por su peso}
  \CommentTok{\# paso de mínimos cuadrados ponderados: beta = (X\textquotesingle{}WX)\^{}\{{-}1\} X\textquotesingle{}W z}
\NormalTok{  beta\_new }\OtherTok{\textless{}{-}} \FunctionTok{as.vector}\NormalTok{(}\FunctionTok{solve}\NormalTok{(}\FunctionTok{crossprod}\NormalTok{(WX, X), }\FunctionTok{crossprod}\NormalTok{(WX, z)))}
\NormalTok{  traza[[it]] }\OtherTok{\textless{}{-}}\NormalTok{ beta\_new             }\CommentTok{\# registramos el beta de esta vuelta}
  \ControlFlowTok{if}\NormalTok{ (}\FunctionTok{max}\NormalTok{(}\FunctionTok{abs}\NormalTok{(beta\_new }\SpecialCharTok{{-}}\NormalTok{ beta)) }\SpecialCharTok{\textless{}} \FloatTok{1e{-}10}\NormalTok{) \{ beta }\OtherTok{\textless{}{-}}\NormalTok{ beta\_new; }\ControlFlowTok{break}\NormalTok{ \}  }\CommentTok{\# ¿convergió?}
\NormalTok{  beta }\OtherTok{\textless{}{-}}\NormalTok{ beta\_new                    }\CommentTok{\# si no, otra iteración partiendo del nuevo beta}
\NormalTok{\}}
\FunctionTok{cat}\NormalTok{(}\StringTok{"Convergió en"}\NormalTok{, it, }\StringTok{"iteraciones}\SpecialCharTok{\textbackslash{}n}\StringTok{"}\NormalTok{)}
\end{Highlighting}
\end{Shaded}

\begin{verbatim}
Convergió en 6 iteraciones
\end{verbatim}

\begin{Shaded}
\begin{Highlighting}[]
\FunctionTok{cbind}\NormalTok{(}\AttributeTok{IWLS\_a\_mano =}\NormalTok{ beta, }\AttributeTok{glm =} \FunctionTok{coef}\NormalTok{(fit\_glm))   }\CommentTok{\# deben coincidir con glm()}
\end{Highlighting}
\end{Shaded}

\begin{verbatim}
             IWLS_a_mano         glm
(Intercept)  -4.21429456 -4.21429456
age           0.04119282  0.04119282
priorfracYes  0.83883503  0.83883503
\end{verbatim}

\textbf{La trayectoria hasta la convergencia.}

En cada vuelta hemos ido guardando el vector de coeficientes en la lista
\texttt{traza}; ahora \texttt{do.call(rbind,\ traza)} \textbf{apila}
esos vectores en una \textbf{matriz} ---una fila por iteración, una
columna por coeficiente---, que pasamos a formato largo para dibujar
cómo evoluciona cada coeficiente. La
Figura~\ref{fig-u12-iwls-convergencia} lo muestra: partiendo de cero,
los tres coeficientes alcanzan su valor de máxima verosimilitud ---las
líneas punteadas de \texttt{glm()}--- en apenas un puñado de
iteraciones. Esa convergencia casi inmediata es la marca del Fisher
scoring cerca del óptimo.

\begin{Shaded}
\begin{Highlighting}[]
\FunctionTok{do.call}\NormalTok{(rbind, traza) }\SpecialCharTok{|\textgreater{}}          \CommentTok{\# apila la lista en matriz: filas = iteraciones, columnas = coeficientes}
  \FunctionTok{as.data.frame}\NormalTok{() }\SpecialCharTok{|\textgreater{}}
  \FunctionTok{setNames}\NormalTok{(}\FunctionTok{colnames}\NormalTok{(X)) }\SpecialCharTok{|\textgreater{}}        \CommentTok{\# nombra las columnas como los coeficientes}
  \FunctionTok{mutate}\NormalTok{(}\AttributeTok{iter =} \FunctionTok{row\_number}\NormalTok{()) }\SpecialCharTok{|\textgreater{}}  \CommentTok{\# índice de iteración}
  \FunctionTok{pivot\_longer}\NormalTok{(}\SpecialCharTok{{-}}\NormalTok{iter, }\AttributeTok{names\_to =} \StringTok{"coef"}\NormalTok{, }\AttributeTok{values\_to =} \StringTok{"valor"}\NormalTok{) }\SpecialCharTok{|\textgreater{}}   \CommentTok{\# formato largo para ggplot}
  \FunctionTok{ggplot}\NormalTok{(}\FunctionTok{aes}\NormalTok{(iter, valor, }\AttributeTok{color =}\NormalTok{ coef)) }\SpecialCharTok{+}
  \FunctionTok{geom\_line}\NormalTok{() }\SpecialCharTok{+} \FunctionTok{geom\_point}\NormalTok{(}\AttributeTok{size =} \FloatTok{1.5}\NormalTok{) }\SpecialCharTok{+}
  \FunctionTok{geom\_hline}\NormalTok{(}\AttributeTok{data =} \FunctionTok{enframe}\NormalTok{(}\FunctionTok{coef}\NormalTok{(fit\_glm), }\StringTok{"coef"}\NormalTok{, }\StringTok{"obj"}\NormalTok{),         }\CommentTok{\# valores de glm() como referencia}
             \FunctionTok{aes}\NormalTok{(}\AttributeTok{yintercept =}\NormalTok{ obj, }\AttributeTok{color =}\NormalTok{ coef), }\AttributeTok{linetype =} \StringTok{"dashed"}\NormalTok{) }\SpecialCharTok{+}
  \FunctionTok{labs}\NormalTok{(}\AttributeTok{x =} \StringTok{"Iteración IWLS"}\NormalTok{, }\AttributeTok{y =} \StringTok{"Coeficiente"}\NormalTok{, }\AttributeTok{color =} \ConstantTok{NULL}\NormalTok{)}
\end{Highlighting}
\end{Shaded}

\begin{figure}[H]

{\centering \includegraphics{caso1_binaria_files/figure-pdf/fig-u12-iwls-convergencia-1.pdf}

}

\caption{\label{fig-u12-iwls-convergencia}Convergencia del IWLS «a
mano». Cada coeficiente se estabiliza en su valor de máxima
verosimilitud (líneas punteadas, las de \texttt{glm()}) en muy pocas
iteraciones.}

\end{figure}

\hypertarget{sec-u12-inferencia}{%
\subsection{2.5 Inferencia y selección de
modelos}\label{sec-u12-inferencia}}

Una vez tenemos los coeficientes estimados, \(\hat{\boldsymbol\beta}\),
toca interrogar al modelo en tres niveles, de lo particular a lo global:
1. ¿es real el efecto de cada coeficiente? 1. ¿aporta conocimiento un
bloque de variables? y 1. ¿qué modelo elijo entre varias
especificaciones?

Las tres preguntas se responden con la misma moneda ---la
verosimilitud--- y su sombra operativa, la \textbf{deviance}.

\hypertarget{la-deviance}{%
\subsubsection{La Deviance}\label{la-deviance}}

Para medir cuánto se aleja un ajuste del mejor posible, se compara con
el \textbf{modelo saturado}, que es el que tiene un parámetro por
observación, y por lo tanto reproduce los datos exactamente
(\(\hat\mu_i=y_i\)) y proporciona el valor de log-verosimilitud máxima
\(\ell_{\text{sat}}\). La \textbf{deviance} es esa distancia, en
unidades de log-verosimilitud, convertido en un valor positivo: \[
D=-2\bigl[\ell(\hat{\boldsymbol\theta})-\ell_{\text{sat}}\bigr].
\] La deviance viene a ser el análogo GLM de la suma de cuadrados
residual: cuanto menor valor toma, mejor ajuste.

La \textbf{deviance nula} \(D_0\) es la del modelo solo con
interceptación (el más simple que no asume ningún efecto de los
predictores); la \textbf{deviance residual} \(D\), representa la
deviance del modelo ajustado.

\begin{tcolorbox}[enhanced jigsaw, coltitle=black, rightrule=.15mm, arc=.35mm, leftrule=.75mm, colback=white, colbacktitle=quarto-callout-tip-color!10!white, left=2mm, breakable, toptitle=1mm, toprule=.15mm, bottomtitle=1mm, titlerule=0mm, title=\textcolor{quarto-callout-tip-color}{\faLightbulb}\hspace{0.5em}{🔧 En R. Deviances y grados de libertad de un \texttt{glm}}, bottomrule=.15mm, opacityback=0, opacitybacktitle=0.6, colframe=quarto-callout-tip-color-frame]

En un modelo \texttt{glm} ajustado, \texttt{fit\_glm}, las deviances
nula, del modelo saturado y del modelo ajustado, junto con sus grados de
libertad, se obtienen así:

\begin{Shaded}
\begin{Highlighting}[]
\NormalTok{fit\_glm}\SpecialCharTok{$}\NormalTok{null.deviance     }\CommentTok{\# deviance NULA (modelo solo con intercepto)}
\NormalTok{fit\_glm}\SpecialCharTok{$}\NormalTok{df.null           }\CommentTok{\#   sus g.l.: n {-} 1}

\FunctionTok{deviance}\NormalTok{(fit\_glm)         }\CommentTok{\# deviance del modelo AJUSTADO (= residual); igual a fit\_glm$deviance}
\FunctionTok{df.residual}\NormalTok{(fit\_glm)      }\CommentTok{\#   sus g.l.: n {-} p; igual a fit\_glm$df.residual}
\end{Highlighting}
\end{Shaded}

La deviance del modelo \textbf{saturado es 0 por definición}: es el
modelo de referencia ---un parámetro por observación, con g.l.
residuales = 0--- frente al que se mide \emph{toda} deviance, de modo
que no hay una función que la devuelva (la ``residual'' ya es la
distancia del modelo ajustado al saturado).

Para mostrarlas todas conjuntamente, podemos utilizar \texttt{glance()},
de la librería \texttt{broom}, que devuelve en una sola fila la
\texttt{null.deviance} con \texttt{df.null}, la \texttt{deviance}
residual con \texttt{df.residual}, y además \texttt{logLik},
\texttt{AIC}, \texttt{BIC} y \texttt{nobs}:

\begin{Shaded}
\begin{Highlighting}[]
\FunctionTok{glance}\NormalTok{(fit\_glm)}
\end{Highlighting}
\end{Shaded}

\end{tcolorbox}

\begin{tcolorbox}[enhanced jigsaw, coltitle=black, rightrule=.15mm, arc=.35mm, leftrule=.75mm, colback=white, colbacktitle=quarto-callout-tip-color!10!white, left=2mm, breakable, toptitle=1mm, toprule=.15mm, bottomtitle=1mm, titlerule=0mm, title=\textcolor{quarto-callout-tip-color}{\faLightbulb}\hspace{0.5em}{Tip}, bottomrule=.15mm, opacityback=0, opacitybacktitle=0.6, colframe=quarto-callout-tip-color-frame]

Detalle útil: con datos binarios \textbf{individuales},
\(\ell_{\text{sat}}=0\), de modo que
\(D=-2\,\ell(\hat{\boldsymbol\beta})\).

\end{tcolorbox}

Mostramos las deviances y otros indicadores como el AIC y BIC, que nos
serán útiles más adelante para la comparación de modelos:

\begin{Shaded}
\begin{Highlighting}[]
\FunctionTok{glance}\NormalTok{(fit\_glm) }\SpecialCharTok{|\textgreater{}}
  \FunctionTok{select}\NormalTok{(null.deviance, df.null, deviance, df.residual, logLik, AIC, BIC)}
\end{Highlighting}
\end{Shaded}

\begin{verbatim}
# A tibble: 1 x 7
  null.deviance df.null deviance df.residual logLik   AIC   BIC
          <dbl>   <int>    <dbl>       <int>  <dbl> <dbl> <dbl>
1          562.     499     529.         497  -264.  535.  547.
\end{verbatim}

En GLOW, la deviance baja de \textbf{562,3} (nula) a \textbf{528,5}
(residual) al añadir \texttt{age} y \texttt{priorfrac}: esa reducción de
≈ 33,8 puntos es la materia prima de los contrastes que siguen.

La deviance residual \textbf{no} sirve como test \emph{absoluto} de
ajuste con datos individuales (lo veremos en la
\protect\hyperlink{sec-u12-diagnostico}{Sección 2.6}). De momento la
usaremos solo para \textbf{comparar} modelos.

\hypertarget{inferencia-sobre-los-coeficientes}{%
\subsubsection{Inferencia sobre los
coeficientes}\label{inferencia-sobre-los-coeficientes}}

Para contestar la pregunta ``¿es real el efecto de cada coeficiente?'',
se utilizan los intervalos de confianza de Wald y de perfil.

\textbf{Test de Wald}

El test de Wald es el análogo directo de los tests \(t\) del modelo
lineal. Se construye con el cociente entre el estimador del coeficiente
\(\beta_k\) (asociado al predictor \(x_k\)) y su error estándar:

\[
z_k=\frac{\hat\beta_k}{\widehat{\text{EE}}(\hat\beta_k)}.
\] Este estadístico de Wald se compara con la normal asintótica (de ahí
\(z\) y no \(t\)), y es el que imprime
\texttt{summary()}/\texttt{tidy()}. Es un test cómodo, pero \textbf{poco
fiable} con muestras pequeñas o coeficientes grandes, y \textbf{engañoso
bajo separación} (da por no significativo un efecto que en realidad es
muy fuerte, debido a un inflado del error estándar -efecto (Hauck \&
Donner, 1977)).

\textbf{Intervalos de confianza: Wald frente a perfil.}

Por la misma razón, los \textbf{intervalos de confianza} se prefieren
por \textbf{perfil de verosimilitud} antes que por Wald (Agresti, 2015).
El de \textbf{Wald}, \(\hat\beta_k\pm1{,}96\,\widehat{\text{EE}}\),
aproxima la log-verosimilitud por una \textbf{parábola} y es siempre
simétrico. El de \textbf{perfil} se obtiene \textbf{invirtiendo la razón
de verosimilitudes}: es el conjunto de valores \(b\) que el LRT no
rechazaría, es decir, aquellos en los que la \textbf{log-verosimilitud
perfil} \[
\ell_p(b)=\max_{\boldsymbol\beta_{-k}}\ \ell(b,\boldsymbol\beta_{-k})
\] ---la verosimilitud maximizada sobre el \textbf{resto} de parámetros
con \(\beta_k=b\) fijo--- cae como mucho
\(1{,}92=\tfrac12\chi^2_{1,\,0{,}95}\) por debajo de su máximo: \[
\text{IC}_{95\%}(\beta_k)=\Bigl\{\,b:\ \ell(\hat{\boldsymbol\beta})-\ell_p(b)\le 1{,}92\,\Bigr\}.
\] Geométricamente, se corta la curva de verosimilitud perfil a
\(1{,}92\) del pico y se toman las dos abscisas del cruce. Como usa la
\textbf{forma real} ---a menudo \textbf{asimétrica}--- de la
verosimilitud, no asume normalidad de \(\hat\beta\) y \textbf{no tiene
por qué ser simétrico}: es más fiable justo cuando la parábola de Wald
aproxima mal (muestras pequeñas, coeficientes grandes, cerca de la
separación).

\begin{tcolorbox}[enhanced jigsaw, coltitle=black, rightrule=.15mm, arc=.35mm, leftrule=.75mm, colback=white, colbacktitle=quarto-callout-tip-color!10!white, left=2mm, breakable, toptitle=1mm, toprule=.15mm, bottomtitle=1mm, titlerule=0mm, title=\textcolor{quarto-callout-tip-color}{\faLightbulb}\hspace{0.5em}{🔧 En R. Intervalos de confianza sobre los coeficientes}, bottomrule=.15mm, opacityback=0, opacitybacktitle=0.6, colframe=quarto-callout-tip-color-frame]

En R,

\begin{itemize}
\tightlist
\item
  \texttt{confint()} devuelve el intervalo de confianza del perfil de
  verosimilitud ---el método de perfilado implementado en \texttt{MASS}
  (Venables \& Ripley, 2002)---
\item
  \texttt{confint.default()}, el de Wald.
\end{itemize}

\end{tcolorbox}

Vemos que en GLOW, que se comporta bien, casi coinciden ambos
intervalos:

\begin{Shaded}
\begin{Highlighting}[]
\FunctionTok{list}\NormalTok{(}\AttributeTok{perfil =} \FunctionTok{confint}\NormalTok{(fit\_glm),}
     \AttributeTok{wald   =} \FunctionTok{confint.default}\NormalTok{(fit\_glm)) }\SpecialCharTok{|\textgreater{}}
  \FunctionTok{map}\NormalTok{(\textbackslash{}(m) }\FunctionTok{round}\NormalTok{(m, }\DecValTok{4}\NormalTok{))}
\end{Highlighting}
\end{Shaded}

\begin{verbatim}
$perfil
               2.5 %  97.5 %
(Intercept)  -5.9006 -2.5711
age           0.0174  0.0653
priorfracYes  0.3776  1.2970

$wald
               2.5 %  97.5 %
(Intercept)  -5.8760 -2.5526
age           0.0173  0.0651
priorfracYes  0.3799  1.2978
\end{verbatim}

\hypertarget{inferencia-sobre-bloques-el-contraste-de-razuxf3n-de-verosimilitudes-lrt}{%
\subsubsection{Inferencia sobre bloques: el contraste de razón de
verosimilitudes
(LRT)}\label{inferencia-sobre-bloques-el-contraste-de-razuxf3n-de-verosimilitudes-lrt}}

Toda la inferencia que \textbf{compara modelos} se apoya en un mismo
motor: el \textbf{contraste de razón de verosimilitudes}
(\emph{likelihood ratio test}, LRT). Mide \textbf{cuánto mejora la
verosimilitud} al pasar de un modelo \textbf{pequeño} a otro \textbf{más
grande que lo contiene} (anidado). El estadístico es \[
\Lambda=2\big(\ell_{\text{grande}}-\ell_{\text{pequeño}}\big)=D_{\text{pequeño}}-D_{\text{grande}}\;\sim\;\chi^2_{\Delta\text{gl}},
\] es decir, \textbf{la diferencia de deviances} ---recuerda que
\(D=2(\ell_{\text{sat}}-\ell_{\text{modelo}})\), así que el término del
saturado se cancela---. Bajo \(H_0\) (``el modelo pequeño basta; los
parámetros extra no aportan''), \(\Lambda\) sigue una \(\chi^2\) con
tantos grados de libertad como \textbf{parámetros añadidos}. Un
\(\Lambda\) \textbf{grande} (\(p\) bajo) indica que esos parámetros
\textbf{sí} explican señal real: nos quedamos con el modelo grande.

Es \textbf{un único contraste} que, según qué par anidado compares,
responde a preguntas distintas:

\begin{itemize}
\tightlist
\item
  \textbf{pequeño = nulo, grande = ajustado} → \emph{«¿aporta algo el
  modelo?»} (el \textbf{ómnibus}, justo abajo).
\item
  \textbf{pequeño = reducido, grande = completo} → \emph{«¿aporta este
  bloque de variables?»} (la escalera de modelos que recorreremos
  enseguida).
\item
  \textbf{grande = saturado} → \emph{bondad de ajuste}: la
  \textbf{deviance residual} es el LRT contra el ajuste perfecto
  (\protect\hyperlink{sec-u12-diagnostico}{Sección 2.6}).
\end{itemize}

Dos requisitos: los modelos han de estar \textbf{anidados} y ajustados
sobre los \textbf{mismos datos}. Para especificaciones \textbf{no
anidadas} no hay LRT y se recurre a \textbf{AIC/BIC} (más abajo).

\textbf{El ómnibus} ---``¿aporta \emph{algo} el modelo?''--- es el caso
\emph{nulo vs ajustado}, análogo al \textbf{test \(F\) global} del
modelo lineal: \[D_0-D\sim\chi^2_{p}.\] Para valorar el aporte de
\emph{cada término ajustando por el resto}, se echa mano del
\texttt{Anova} de tipo II:

\begin{Shaded}
\begin{Highlighting}[]
\CommentTok{\# Wald en escala log{-}odds: aporte de cada término independientemente}
\FunctionTok{tidy}\NormalTok{(fit\_glm)  }\CommentTok{\# proporcionado por summary(fit\_glm)}
\end{Highlighting}
\end{Shaded}

\begin{verbatim}
# A tibble: 3 x 5
  term         estimate std.error statistic     p.value
  <chr>           <dbl>     <dbl>     <dbl>       <dbl>
1 (Intercept)   -4.21      0.848      -4.97 0.000000667
2 age            0.0412    0.0122      3.38 0.000719   
3 priorfracYes   0.839     0.234       3.58 0.000340   
\end{verbatim}

\begin{Shaded}
\begin{Highlighting}[]
\CommentTok{\# Anova tipo II: aporte de cada término ajustando por el resto}
\NormalTok{car}\SpecialCharTok{::}\FunctionTok{Anova}\NormalTok{(fit\_glm, }\AttributeTok{type =} \StringTok{"II"}\NormalTok{)}
\end{Highlighting}
\end{Shaded}

\begin{verbatim}
Analysis of Deviance Table (Type II tests)

Response: fracture
          LR Chisq Df Pr(>Chisq)    
age         11.551  1  0.0006772 ***
priorfrac   12.544  1  0.0003975 ***
---
Signif. codes:  0 '***' 0.001 '**' 0.01 '*' 0.05 '.' 0.1 ' ' 1
\end{verbatim}

\begin{Shaded}
\begin{Highlighting}[]
\CommentTok{\# LRT ómnibus: aporte del bloque completo frente al modelo nulo}
\FunctionTok{anova}\NormalTok{(}\FunctionTok{glm}\NormalTok{(fracture }\SpecialCharTok{\textasciitilde{}} \DecValTok{1}\NormalTok{, binomial, glow), fit\_glm, }\AttributeTok{test =} \StringTok{"LRT"}\NormalTok{)}
\end{Highlighting}
\end{Shaded}

\begin{verbatim}
Analysis of Deviance Table

Model 1: fracture ~ 1
Model 2: fracture ~ age + priorfrac
  Resid. Df Resid. Dev Df Deviance  Pr(>Chi)    
1       499     562.34                          
2       497     528.52  2   33.818 4.535e-08 ***
---
Signif. codes:  0 '***' 0.001 '**' 0.01 '*' 0.05 '.' 0.1 ' ' 1
\end{verbatim}

\begin{tcolorbox}[enhanced jigsaw, coltitle=black, rightrule=.15mm, arc=.35mm, leftrule=.75mm, colback=white, colbacktitle=quarto-callout-tip-color!10!white, left=2mm, breakable, toptitle=1mm, toprule=.15mm, bottomtitle=1mm, titlerule=0mm, title=\textcolor{quarto-callout-tip-color}{\faLightbulb}\hspace{0.5em}{🔧 En R. Contrastes de razón de verosimilitudes (LRT)}, bottomrule=.15mm, opacityback=0, opacitybacktitle=0.6, colframe=quarto-callout-tip-color-frame]

El LRT compara dos modelos \textbf{anidados} (ajustados sobre los
\textbf{mismos datos}) por su diferencia de deviances, con
\texttt{anova(...,\ test\ =\ "LRT")} (equivale a
\texttt{test\ =\ "Chisq"} en \texttt{glm}):

\begin{Shaded}
\begin{Highlighting}[]
\CommentTok{\# Ómnibus: ¿aporta algo el modelo frente al nulo (solo intercepto)?}
\FunctionTok{anova}\NormalTok{(}\FunctionTok{glm}\NormalTok{(fracture }\SpecialCharTok{\textasciitilde{}} \DecValTok{1}\NormalTok{, }\AttributeTok{family =}\NormalTok{ binomial, }\AttributeTok{data =}\NormalTok{ glow), fit\_glm, }\AttributeTok{test =} \StringTok{"LRT"}\NormalTok{)}

\CommentTok{\# LRT entre dos modelos anidados cualesquiera (reducido vs completo)}
\FunctionTok{anova}\NormalTok{(mod\_reducido, mod\_completo, }\AttributeTok{test =} \StringTok{"LRT"}\NormalTok{)}

\CommentTok{\# Aporte de CADA término ajustando por el resto (LRT por término)}
\NormalTok{car}\SpecialCharTok{::}\FunctionTok{Anova}\NormalTok{(fit\_glm, }\AttributeTok{type =} \StringTok{"II"}\NormalTok{)     }\CommentTok{\# type = "III" si hay interacciones}
\end{Highlighting}
\end{Shaded}

\texttt{summary(fit\_glm)} / \texttt{tidy(fit\_glm)} dan el contraste
\textbf{de Wald} por coeficiente (aproximado); el \textbf{LRT} suele ser
más fiable. Recuerda: el LRT \textbf{solo} vale para modelos
\textbf{anidados}.

\end{tcolorbox}

\hypertarget{comparar-modelos-no-anidados-aic-y-bic}{%
\subsubsection{Comparar modelos no anidados: AIC y
BIC}\label{comparar-modelos-no-anidados-aic-y-bic}}

El LRT solo vale para comparar modelos anidados. Para comparar
especificaciones cualesquiera se usan criterios que penalizan la
verosimilitud por complejidad, como el AIC y el BIC (Akaike, 1974;
Schwarz, 1978), que tienen en cuenta el número \(k\) parámetros y \(n\)
observaciones:

\[
\text{AIC}=-2\ell+2k,\qquad \text{BIC}=-2\ell+k\log n,
\]

Menor es mejor, si bien el \textbf{BIC penaliza más} y por tanto
\textbf{favorece la parsimonia}. Es importante remarcar que son medidas
\textbf{relativas}, útiles para ordenar modelos, no para concluir si
alguno ajusta bien en términos absolutos.

\begin{tcolorbox}[enhanced jigsaw, coltitle=black, rightrule=.15mm, arc=.35mm, leftrule=.75mm, colback=white, colbacktitle=quarto-callout-tip-color!10!white, left=2mm, breakable, toptitle=1mm, toprule=.15mm, bottomtitle=1mm, titlerule=0mm, title=\textcolor{quarto-callout-tip-color}{\faLightbulb}\hspace{0.5em}{🔧 En R. AIC y BIC de un modelo}, bottomrule=.15mm, opacityback=0, opacitybacktitle=0.6, colframe=quarto-callout-tip-color-frame]

Sobre un modelo ajustado \texttt{fit\_glm}:

\begin{Shaded}
\begin{Highlighting}[]
\FunctionTok{AIC}\NormalTok{(fit\_glm)              }\CommentTok{\# AIC}
\FunctionTok{BIC}\NormalTok{(fit\_glm)              }\CommentTok{\# BIC}
\FunctionTok{AIC}\NormalTok{(m\_1, m\_2, m\_3)        }\CommentTok{\# tabla con los AIC de varios modelos (BIC(...) para el BIC)}
\FunctionTok{glance}\NormalTok{(fit\_glm)           }\CommentTok{\# broom: AIC, BIC, deviance, logLik, gl... en una fila}
\end{Highlighting}
\end{Shaded}

Menor es mejor; recuerda que son \textbf{relativos} (sirven para
ordenar, no para juzgar el ajuste absoluto).

\end{tcolorbox}

\textbf{La escalera de modelos.}

Seleccionar el mejor modelo no es elegir entre dos concretos, sino
\textbf{recorrer una escalera o jerarquía} de modelos, de menos a más
compleja (o viceversa), decidiendo en cada peldaño si la complejidad
añadida (o excluida) se paga ``de modo suficiente'' con ajuste. Para
realizar esta selección el analista propone un criterio y procede; los
más habituales son LRT, AIC y BIC. La decisión final es \textbf{del
analista}, y debe argumentarse.

A la hora de ajustar un mejor modelo sobre GLOW, consideramos diversos
modelos de menos a más complejos, con las variables \texttt{age},
\texttt{priorfrac}, \texttt{momfrac} (fractura de cadera en la madre) y
\texttt{bmi} (índice de masa corporal), teniendo en cuenta posibles
interacciones entre factores o entre factores y covariables (los dos
factores son \texttt{priorfrac} y \texttt{momfrac}).

\begin{Shaded}
\begin{Highlighting}[]
\NormalTok{mod\_0 }\OtherTok{\textless{}{-}} \FunctionTok{glm}\NormalTok{(fracture }\SpecialCharTok{\textasciitilde{}}\NormalTok{ age }\SpecialCharTok{+}\NormalTok{ priorfrac,                     }\AttributeTok{family =}\NormalTok{ binomial, }\AttributeTok{data =}\NormalTok{ glow)}
\NormalTok{mod\_1 }\OtherTok{\textless{}{-}} \FunctionTok{glm}\NormalTok{(fracture }\SpecialCharTok{\textasciitilde{}}\NormalTok{ age }\SpecialCharTok{+}\NormalTok{ priorfrac }\SpecialCharTok{+}\NormalTok{ momfrac }\SpecialCharTok{+}\NormalTok{ bmi,     }\AttributeTok{family =}\NormalTok{ binomial, }\AttributeTok{data =}\NormalTok{ glow)}
\NormalTok{mod\_2 }\OtherTok{\textless{}{-}} \FunctionTok{glm}\NormalTok{(fracture }\SpecialCharTok{\textasciitilde{}}\NormalTok{ age }\SpecialCharTok{*}\NormalTok{ priorfrac }\SpecialCharTok{+}\NormalTok{ bmi }\SpecialCharTok{+}\NormalTok{ momfrac, }\AttributeTok{family =}\NormalTok{ binomial, }\AttributeTok{data =}\NormalTok{ glow)}
\NormalTok{mod\_3 }\OtherTok{\textless{}{-}} \FunctionTok{glm}\NormalTok{(fracture }\SpecialCharTok{\textasciitilde{}}\NormalTok{ (age }\SpecialCharTok{+}\NormalTok{ bmi) }\SpecialCharTok{*}\NormalTok{ priorfrac }\SpecialCharTok{+}\NormalTok{ momfrac, }\AttributeTok{family =}\NormalTok{ binomial, }\AttributeTok{data =}\NormalTok{ glow)}
\NormalTok{mod\_4 }\OtherTok{\textless{}{-}} \FunctionTok{glm}\NormalTok{(fracture }\SpecialCharTok{\textasciitilde{}}\NormalTok{ (age }\SpecialCharTok{+}\NormalTok{ bmi) }\SpecialCharTok{*}\NormalTok{ (priorfrac }\SpecialCharTok{+}\NormalTok{ momfrac), }\AttributeTok{family =}\NormalTok{ binomial, }\AttributeTok{data =}\NormalTok{ glow)}
\FunctionTok{tibble}\NormalTok{(}\AttributeTok{modelo =} \FunctionTok{paste0}\NormalTok{(}\StringTok{"mod\_"}\NormalTok{, }\DecValTok{1}\SpecialCharTok{:}\DecValTok{4}\NormalTok{)) }\SpecialCharTok{|\textgreater{}}
  \FunctionTok{bind\_cols}\NormalTok{(}\FunctionTok{map\_dfr}\NormalTok{(}\FunctionTok{list}\NormalTok{(mod\_1, mod\_2, mod\_3, mod\_4),}
\NormalTok{                    \textbackslash{}(m) }\FunctionTok{glance}\NormalTok{(m) }\SpecialCharTok{|\textgreater{}}\NormalTok{ dplyr}\SpecialCharTok{::}\FunctionTok{select}\NormalTok{(}\AttributeTok{df =}\NormalTok{ df.residual, deviance, AIC, BIC)))}
\end{Highlighting}
\end{Shaded}

\begin{verbatim}
# A tibble: 4 x 5
  modelo    df deviance   AIC   BIC
  <chr>  <int>    <dbl> <dbl> <dbl>
1 mod_1    495     522.  532.  553.
2 mod_2    494     516.  528.  553.
3 mod_3    493     514.  528.  558.
4 mod_4    491     514.  532.  570.
\end{verbatim}

Los efectos principales \textbf{bajan la deviance con claridad}; las
interacciones, menos. Lo contrastamos peldaño a peldaño con LRT:

\begin{Shaded}
\begin{Highlighting}[]
\FunctionTok{anova}\NormalTok{(mod\_4, mod\_3,  }\AttributeTok{test =} \StringTok{"LRT"}\NormalTok{)   }
\end{Highlighting}
\end{Shaded}

\begin{verbatim}
Analysis of Deviance Table

Model 1: fracture ~ (age + bmi) * (priorfrac + momfrac)
Model 2: fracture ~ (age + bmi) * priorfrac + momfrac
  Resid. Df Resid. Dev Df Deviance Pr(>Chi)
1       491     513.80                     
2       493     514.11 -2 -0.31228   0.8554
\end{verbatim}

\begin{Shaded}
\begin{Highlighting}[]
\FunctionTok{anova}\NormalTok{(mod\_3, mod\_2, }\AttributeTok{test =} \StringTok{"LRT"}\NormalTok{)   }
\end{Highlighting}
\end{Shaded}

\begin{verbatim}
Analysis of Deviance Table

Model 1: fracture ~ (age + bmi) * priorfrac + momfrac
Model 2: fracture ~ age * priorfrac + bmi + momfrac
  Resid. Df Resid. Dev Df Deviance Pr(>Chi)
1       493     514.11                     
2       494     516.09 -1  -1.9787   0.1595
\end{verbatim}

\begin{Shaded}
\begin{Highlighting}[]
\FunctionTok{anova}\NormalTok{(mod\_2, mod\_1, }\AttributeTok{test =} \StringTok{"LRT"}\NormalTok{) }
\end{Highlighting}
\end{Shaded}

\begin{verbatim}
Analysis of Deviance Table

Model 1: fracture ~ age * priorfrac + bmi + momfrac
Model 2: fracture ~ age + priorfrac + momfrac + bmi
  Resid. Df Resid. Dev Df Deviance Pr(>Chi)  
1       494     516.09                       
2       495     521.87 -1  -5.7753  0.01625 *
---
Signif. codes:  0 '***' 0.001 '**' 0.01 '*' 0.05 '.' 0.1 ' ' 1
\end{verbatim}

\begin{Shaded}
\begin{Highlighting}[]
\FunctionTok{anova}\NormalTok{(mod\_1, mod\_0, }\AttributeTok{test =} \StringTok{"LRT"}\NormalTok{) }
\end{Highlighting}
\end{Shaded}

\begin{verbatim}
Analysis of Deviance Table

Model 1: fracture ~ age + priorfrac + momfrac + bmi
Model 2: fracture ~ age + priorfrac
  Resid. Df Resid. Dev Df Deviance Pr(>Chi)  
1       495     521.87                       
2       497     528.52 -2  -6.6485    0.036 *
---
Signif. codes:  0 '***' 0.001 '**' 0.01 '*' 0.05 '.' 0.1 ' ' 1
\end{verbatim}

El primer contraste contesta si son significativas (en bloque) las
interacciones de edad y bmi con \texttt{momfrac}, y la respuesta es que
no (p-valor 0.8554), por lo que cabe descartarlas para simplificar el
modelo. La siguiente tabla Anova resuelve si tiene sentido la
interacción de \texttt{bmi} con \texttt{priorfrac}, y la respuesta es no
(p-valor 0.1595). A continuación se comparan los modelos \texttt{mod\_2}
y \texttt{mod\_1}, esto es, si es significativa la interacción
\texttt{age:priorfrac} y el resultado es sí, con un p-valor de 0.0162.
Llegado a este punto, pararíamos y nos quedaríamos con el
\texttt{mod\_2}. La última comparación valora si merece la pena incluir
las variables \texttt{momfrac} y \texttt{bmi} en el modelo, y la
respuesta es nuevamente que sí (p-valor 0.036), pues amplían la
información que dan exclusivamente \texttt{age} y \texttt{priorfrac}.

Interpretar una interacción exige cuidado: el coeficiente de
\texttt{priorfrac} ya \textbf{no} es ``el efecto de la fractura
previa'', sino su efecto \textbf{a edad 0} (sin sentido sin centrar).
Por eso, con interacciones, la lectura se hace en \textbf{efectos
marginales o predicciones} (ya vista en la
\protect\hyperlink{sec-u12-interpretacion}{Sección 2.3}), no en
coeficientes crudos:

\begin{Shaded}
\begin{Highlighting}[]
\FunctionTok{tidy}\NormalTok{(mod\_2) }
\end{Highlighting}
\end{Shaded}

\begin{verbatim}
# A tibble: 6 x 5
  term             estimate std.error statistic     p.value
  <chr>               <dbl>     <dbl>     <dbl>       <dbl>
1 (Intercept)       -6.84      1.33       -5.15 0.000000263
2 age                0.0678    0.0159      4.25 0.0000214  
3 priorfracYes       5.21      1.84        2.83 0.00464    
4 bmi                0.0248    0.0189      1.31 0.189      
5 momfracYes         0.748     0.299       2.50 0.0125     
6 age:priorfracYes  -0.0612    0.0254     -2.40 0.0162     
\end{verbatim}

\begin{tcolorbox}[enhanced jigsaw, coltitle=black, rightrule=.15mm, arc=.35mm, leftrule=.75mm, colback=white, colbacktitle=quarto-callout-tip-color!10!white, left=2mm, breakable, toptitle=1mm, toprule=.15mm, bottomtitle=1mm, titlerule=0mm, title=\textcolor{quarto-callout-tip-color}{\faLightbulb}\hspace{0.5em}{🔧 En R. Automatizar la escalera: \texttt{dredge} (exhaustiva) y
\texttt{stepAIC} (paso a paso)}, bottomrule=.15mm, opacityback=0, opacitybacktitle=0.6, colframe=quarto-callout-tip-color-frame]

Con varias variables e interacciones el número de modelos crece deprisa.
Dos funciones lo automatizan, con filosofías opuestas; ambas parten de
un \textbf{modelo global} ajustado con \texttt{na.action\ =\ na.fail}
(requisito de \texttt{dredge}):

\begin{Shaded}
\begin{Highlighting}[]
\NormalTok{global }\OtherTok{\textless{}{-}} \FunctionTok{glm}\NormalTok{(fracture }\SpecialCharTok{\textasciitilde{}}\NormalTok{ (age }\SpecialCharTok{+}\NormalTok{ bmi) }\SpecialCharTok{*}\NormalTok{ (priorfrac }\SpecialCharTok{+}\NormalTok{ momfrac),}
              \AttributeTok{family =}\NormalTok{ binomial, }\AttributeTok{data =}\NormalTok{ glow, }\AttributeTok{na.action =}\NormalTok{ na.fail)}
\end{Highlighting}
\end{Shaded}

\textbf{\texttt{MuMIn::dredge} --- exhaustiva} (Bartoń, 2024): ajusta
\textbf{todos} los submodelos (respetando la marginalidad) y los ordena
por el criterio elegido:

\begin{Shaded}
\begin{Highlighting}[]
\NormalTok{sel   }\OtherTok{\textless{}{-}} \FunctionTok{dredge}\NormalTok{(global, }\AttributeTok{rank =} \StringTok{"AIC"}\NormalTok{)        }\CommentTok{\# rank = "BIC" para favorecer la parsimonia}
\FunctionTok{head}\NormalTok{(sel, }\DecValTok{3}\NormalTok{)                                 }\CommentTok{\# mejores modelos}
\NormalTok{mejor }\OtherTok{\textless{}{-}} \FunctionTok{get.models}\NormalTok{(sel, }\AttributeTok{subset =} \DecValTok{1}\NormalTok{)[[}\DecValTok{1}\NormalTok{]]    }\CommentTok{\# el modelo ganador, ya ajustado}
\FunctionTok{formula}\NormalTok{(mejor)                               }\CommentTok{\# solo su fórmula}
\end{Highlighting}
\end{Shaded}

\textbf{\texttt{MASS::stepAIC} --- voraz / paso a paso} (Venables \&
Ripley, 2002): añade o quita un término cada vez siguiendo el AIC, sin
recorrer todo el espacio (más rápida, pero puede quedarse en un óptimo
local):

\begin{Shaded}
\begin{Highlighting}[]
\NormalTok{m\_step }\OtherTok{\textless{}{-}}\NormalTok{ MASS}\SpecialCharTok{::}\FunctionTok{stepAIC}\NormalTok{(global, }\AttributeTok{direction =} \StringTok{"both"}\NormalTok{, }\AttributeTok{trace =} \ConstantTok{FALSE}\NormalTok{)  }\CommentTok{\# \textquotesingle{}both\textquotesingle{}/\textquotesingle{}forward\textquotesingle{}/\textquotesingle{}backward\textquotesingle{}; devuelve YA el modelo ganador}
\FunctionTok{formula}\NormalTok{(m\_step)                                                     }\CommentTok{\# solo su fórmula}
\end{Highlighting}
\end{Shaded}

\texttt{dredge} recorre \textbf{todo} el espacio; \texttt{stepAIC},
\textbf{un solo camino}: si coinciden, refuerza la elección; si
discrepan, hay varios modelos casi indistinguibles.

\end{tcolorbox}

Las comparamos entre sí:

\begin{Shaded}
\begin{Highlighting}[]
\CommentTok{\#library(MuMIn)}
\NormalTok{global }\OtherTok{\textless{}{-}} \FunctionTok{glm}\NormalTok{(fracture }\SpecialCharTok{\textasciitilde{}}\NormalTok{ (age }\SpecialCharTok{+}\NormalTok{ bmi) }\SpecialCharTok{*}\NormalTok{ (priorfrac }\SpecialCharTok{+}\NormalTok{ momfrac),}
              \AttributeTok{family =}\NormalTok{ binomial, }\AttributeTok{data =}\NormalTok{ glow, }\AttributeTok{na.action =}\NormalTok{ na.fail)  }\CommentTok{\# na.fail: requisito de dredge}

\NormalTok{sel }\OtherTok{\textless{}{-}} \FunctionTok{dredge}\NormalTok{(global, }\AttributeTok{rank =} \StringTok{"AIC"}\NormalTok{)       }\CommentTok{\# todos los submodelos, ordenados por AIC (respeta marginalidad)}
\FunctionTok{head}\NormalTok{(sel, }\DecValTok{3}\NormalTok{)                              }\CommentTok{\# los cinco mejores}
\end{Highlighting}
\end{Shaded}

\begin{verbatim}
Global model call: glm(formula = fracture ~ (age + bmi) * (priorfrac + momfrac), 
    family = binomial, data = glow, na.action = na.fail)
---
Model selection table 
     (Int)     age      bmi mmf prr age:prr bmi:prr df   logLik   AIC delta
46  -5.896 0.06400            +   +       +          5 -258.900 527.8  0.00
48  -6.845 0.06777 0.024850   +   +       +          6 -258.047 528.1  0.29
176 -6.079 0.06475 0.004593   +   +       +       +  7 -257.057 528.1  0.31
    weight
46   0.368
48   0.318
176  0.314
Models ranked by AIC(x) 
\end{verbatim}

\begin{Shaded}
\begin{Highlighting}[]
\FunctionTok{dredge}\NormalTok{(global, }\AttributeTok{rank =} \StringTok{"BIC"}\NormalTok{) }\SpecialCharTok{|\textgreater{}} \FunctionTok{head}\NormalTok{(}\DecValTok{3}\NormalTok{)   }\CommentTok{\# el BIC, más parsimonioso, suele elegir otro}
\end{Highlighting}
\end{Shaded}

\begin{verbatim}
Global model call: glm(formula = fracture ~ (age + bmi) * (priorfrac + momfrac), 
    family = binomial, data = glow, na.action = na.fail)
---
Model selection table 
    (Int)     age mmf prr age:prr df   logLik   BIC delta weight
10 -4.214 0.04119       +          3 -264.259 547.2  0.00  0.485
42 -5.689 0.06251       +       +  4 -261.636 548.1  0.97  0.299
14 -4.277 0.04075   +   +          4 -261.963 548.8  1.62  0.216
Models ranked by BIC(x) 
\end{verbatim}

\begin{Shaded}
\begin{Highlighting}[]
\CommentTok{\# Mejor por AIC (exhaustivo) frente a la búsqueda paso a paso:}
\FunctionTok{formula}\NormalTok{(}\FunctionTok{get.models}\NormalTok{(sel, }\AttributeTok{subset =} \DecValTok{1}\NormalTok{)[[}\DecValTok{1}\NormalTok{]])                          }\CommentTok{\# ganador de dredge}
\end{Highlighting}
\end{Shaded}

\begin{verbatim}
fracture ~ age + momfrac + priorfrac + age:priorfrac + 1
\end{verbatim}

\begin{Shaded}
\begin{Highlighting}[]
\FunctionTok{formula}\NormalTok{(MASS}\SpecialCharTok{::}\FunctionTok{stepAIC}\NormalTok{(global, }\AttributeTok{direction =} \StringTok{"both"}\NormalTok{, }\AttributeTok{trace =} \ConstantTok{FALSE}\NormalTok{))  }\CommentTok{\# ganador de stepAIC}
\end{Highlighting}
\end{Shaded}

\begin{verbatim}
fracture ~ age + priorfrac + momfrac + age:priorfrac
\end{verbatim}

En este caso, ambos métodos han llegado al mismo óptimo.

\begin{tcolorbox}[enhanced jigsaw, coltitle=black, rightrule=.15mm, arc=.35mm, leftrule=.75mm, colback=white, colbacktitle=quarto-callout-tip-color!10!white, left=2mm, breakable, toptitle=1mm, toprule=.15mm, bottomtitle=1mm, titlerule=0mm, title=\textcolor{quarto-callout-tip-color}{\faLightbulb}\hspace{0.5em}{🔧 En R. Generar modelos anidados con \texttt{MuMIn} (\texttt{dredge})}, bottomrule=.15mm, opacityback=0, opacitybacktitle=0.6, colframe=quarto-callout-tip-color-frame]

En cualquier caso, cabe tener como referencia que a partir de un modelo
más complejo, podemos extraer todos los posibles modelos (anidados)
viables con la librería \texttt{MuMIn}, para compararlos por cualquier
criterio.

\begin{verbatim}
library(MuMIn)
# Definimos el modelo más complejo de partida
global <- glm(fracture ~ (age + bmi )* (priorfrac + momfrac ) + priorfrac:momfrac,
              family = binomial, data = glow, na.action = na.fail)  # na.fail: requisito de dredge

calls <- dredge(global, evaluate = FALSE)     # NO ajusta: solo genera las llamadas/fórmulas
length(calls)                                 # cuántos modelos hay
calls[[10]]                                     # una llamada concreta
formulas = lapply(calls, function(cl) formula(eval(cl))) # extrae las fórmulas, si las quieres como tales
formulas[[10]]
\end{verbatim}

\end{tcolorbox}

Conviene leer las tres herramientas \textbf{juntas}, no como árbitros
que se sustituyan: el \textbf{LRT} dice qué términos aportan señal real
(aquí, \texttt{age:priorfrac} sí; \texttt{priorfrac:momfrac} no);
\textbf{dredge} y \textbf{stepAIC}, al optimizar AIC, tienden a un
modelo de buena predicción, pero pueden discrepar entre sí ---dredge ve
todo el espacio; \texttt{step}, solo un camino--- y, sobre todo, del
\textbf{BIC}, más parsimonioso. Que dos algoritmos converjan o no es ya
información: si discrepan, varios modelos son casi indistinguibles y la
elección no debe delegarse al criterio.

\begin{tcolorbox}[enhanced jigsaw, coltitle=black, rightrule=.15mm, arc=.35mm, leftrule=.75mm, colback=white, colbacktitle=quarto-callout-warning-color!10!white, left=2mm, breakable, toptitle=1mm, toprule=.15mm, bottomtitle=1mm, titlerule=0mm, title=\textcolor{quarto-callout-warning-color}{\faExclamationTriangle}\hspace{0.5em}{Sobre la selección automática (\texttt{step}, \texttt{dredge})}, bottomrule=.15mm, opacityback=0, opacitybacktitle=0.6, colframe=quarto-callout-warning-color-frame]

La búsqueda automática ---paso a paso (\texttt{MASS::stepAIC}) o
exhaustiva (\texttt{MuMIn::dredge})--- es útil como
\textbf{exploración}, pero no como árbitro. Ambas invalidan los
\(p\)-valores posteriores (inferencia \textbf{post-selección}) y
favorecen el sobreajuste; la exhaustiva, además, \textbf{agrava la
multiplicidad} al comparar muchísimos modelos. En este curso las usamos
para mapear el espacio, y la decisión final se justifica con
LRT/AIC/BIC, validación y conocimiento del problema (Harrell, 2015).

\end{tcolorbox}

\begin{tcolorbox}[enhanced jigsaw, coltitle=black, rightrule=.15mm, arc=.35mm, leftrule=.75mm, colback=white, colbacktitle=quarto-callout-tip-color!10!white, left=2mm, breakable, toptitle=1mm, toprule=.15mm, bottomtitle=1mm, titlerule=0mm, title=\textcolor{quarto-callout-tip-color}{\faLightbulb}\hspace{0.5em}{Protocolo de inferencia y selección}, bottomrule=.15mm, opacityback=0, opacitybacktitle=0.6, colframe=quarto-callout-tip-color-frame]

\begin{enumerate}
\def\labelenumi{\arabic{enumi}.}
\tightlist
\item
  \textbf{¿Aporta algo el modelo?} Contraste ómnibus (LRT nulo vs
  completo).
\item
  \textbf{¿Aporta cada término?} LRT / \texttt{Anova} tipo II por
  término, con Wald e IC (mejor por perfil) como lectura individual.
\item
  \textbf{¿Qué especificación?} LRT entre anidados; AIC/BIC entre
  cualesquiera (AIC → predicción, BIC → parsimonia).
\item
  \textbf{Selección de variables.} Recorre la escalera \textbf{guiada
  por el problema}, no a ciegas; combina LRT/AIC/BIC; evita la búsqueda
  automática como árbitro; \textbf{valida y argumenta} la elección.
\end{enumerate}

\end{tcolorbox}

\hypertarget{sec-u12-diagnostico}{%
\subsection{2.6 Bondad de ajuste y
diagnóstico}\label{sec-u12-diagnostico}}

En 2.5 preguntamos si el modelo \textbf{aporta} y \textbf{cuál} elegir.
Ahora la pregunta es otra: el modelo elegido, \textbf{¿ajusta bien los
datos?} Es la \textbf{bondad de ajuste}, y la bibliografía la trata de
la mano del \textbf{análisis de residuos} (Agresti, 2015; Faraway, 2016;
Hosmer et~al., 2013), porque ambos responden a lo mismo desde dos
ángulos: un \textbf{contraste global} y una \textbf{mirada caso a caso}.

\hypertarget{el-contraste-global-aportar-o-ajustar-y-por-quuxe9-falla-en-binaria}{%
\subsubsection{El contraste global: ¿aportar o ajustar? (y por qué falla
en
binaria)}\label{el-contraste-global-aportar-o-ajustar-y-por-quuxe9-falla-en-binaria}}

Conviene situar tres modelos en una \textbf{escalera}, del peor al mejor
ajuste posible: el \textbf{nulo} (solo intercepto, deviance \(D_0\)); el
\textbf{ajustado}, el nuestro (deviance residual \(D\)); y el
\textbf{saturado}, con un parámetro por observación, que \textbf{ajusta
los datos a la perfección} (deviance \(=0\)). Con ellos se responden
\textbf{dos preguntas distintas que no hay que mezclar}:

\begin{itemize}
\tightlist
\item
  \textbf{¿Aporta el modelo?} (utilidad, la de la
  \protect\hyperlink{sec-u12-inferencia}{Sección 2.5}). Se compara
  \textbf{ajustado vs nulo}: \(D_0-D\sim\chi^2_{p}\). Un \(p\) bajo
  \textbf{rechaza el nulo} a favor del nuestro ---las covariables
  explican algo---. Es la pregunta \emph{«¿mejor que nada?»}.
\item
  \textbf{¿Ajusta bien?} (bondad de ajuste, la de aquí). Se compara
  \textbf{ajustado vs saturado}: la propia \textbf{deviance residual}
  \(D=2(\ell_{\text{sat}}-\ell_{\text{ajustado}})\) \textbf{es} el LRT
  frente al mejor ajuste posible y, bajo un modelo correcto, se
  distribuiría \(\chi^2_{n-p}\). Mide \emph{«¿cuánto queda sin explicar
  respecto a un ajuste perfecto?»}; un \(p\) bajo señala \textbf{falta
  de ajuste}.
\end{itemize}

¿Por qué el \textbf{saturado} y no el nulo es la referencia de la
\emph{bondad de ajuste}? Porque ``ajustar bien'' es \textbf{acercarse al
mejor ajuste alcanzable}, no solo superar a ``no ajustar nada'': un
modelo puede \textbf{batir al nulo} y, a la vez, quedar \textbf{lejos
del saturado} (forma funcional pobre, enlace equivocado, no linealidad
ignorada). Aportaría, pero no ajustaría. Por eso son contrastes
\textbf{complementarios}, no intercambiables.

\textbf{Y aquí está el problema en binaria.} El \(\chi^2_{n-p}\) exige
un régimen asintótico concreto: que el \textbf{número de celdas} ---los
parámetros del saturado--- quede \textbf{fijo} mientras la
\textbf{información dentro de cada celda crece}. Con datos
\textbf{agrupados} (binomial, \(m_i\) ensayos por celda) se cumple: hay
un número fijo de celdas y, al aumentar \(m_i\), cada proporción
\(\hat\pi_i\) se estima con precisión. Con respuesta \textbf{binaria
individual} (\(m_i=1\), una fila por sujeto) ocurre lo contrario: cada
``celda'' es \textbf{una sola} observación, así que al crecer \(n\)
\textbf{crece también el saturado} ---un parámetro por dato, estimado
siempre con un único \(0/1\)---; nunca se acumula información por
parámetro, \(D\) \textbf{no} sigue \(\chi^2_{n-p}\) y el contraste
\textbf{carece de valor}. De hecho, en GLOW \(\ell_{\text{sat}}=0\) y
\(D=-2\ell(\hat{\boldsymbol\beta})\) depende \textbf{solo} de las
probabilidades ajustadas, sin informar del ajuste absoluto (Agresti,
2015; Hosmer et~al., 2013).

\textbf{Cómo se resuelve, entonces.} Dos vías: (i) \textbf{agrupar}
---si las covariables son categóricas, se pasa a binomial: tanto \(D\)
como el \textbf{estadístico de Pearson} \(X^2\) \textbf{recuperan su
sentido} como contraste \(\chi^2\), y allí los definimos con detalle
(\protect\hyperlink{sec-u13-gof-binomial}{Sección 3.2})---; (ii) cuando
no se puede (covariables continuas, como en GLOW), se sustituye el
contraste global por herramientas que \textbf{agrupan implícitamente} o
\textbf{simulan}: el \textbf{test de Hosmer--Lemeshow} (deciles de
riesgo), la \textbf{curva de calibración} y los \textbf{residuos por
tramos}, que vemos a continuación.

\hypertarget{residuos}{%
\subsubsection{Residuos}\label{residuos}}

Para diagnosticar un GLM \textbf{observación a observación} se examinan
los \textbf{residuos}. Los estándar de un GLM son los de
\textbf{Pearson} y los de \textbf{deviance} (McCullagh \& Nelder, 1989);
son informativos cuando hay \textbf{suficiente información por
observación} ---datos agrupados, conteos altos---, y los desarrollamos
en propiedad (fórmulas, gráfico y variantes ---estandarizados, leverage,
distancia de Cook---) en la
\protect\hyperlink{sec-u13-gof-binomial}{Sección 3.2}, donde \textbf{sí}
funcionan.

\textbf{El problema en binaria.} Con respuesta \textbf{binaria
individual}, \(y_i\in\{0,1\}\): para cada \(\hat\mu_i\) el residuo crudo
solo toma dos valores y se ordena en \textbf{dos bandas}, y ni Pearson
ni deviance llegan a ser \(\approx N(0,1)\). Hacen falta herramientas
que \textbf{rescaten} el diagnóstico, y conviene elegirlas \textbf{según
para qué}:

\begin{itemize}
\tightlist
\item
  \textbf{Diagnóstico general} (mala especificación de cualquier tipo:
  forma funcional, dispersión, \emph{outliers}) → los \textbf{residuos
  simulados / de cuantil} (DHARMa) (Dunn \& Smyth, 1996; Hartig, 2024).
  Se simula muchas veces del modelo ajustado y se sitúa cada observación
  en su distribución simulada; bajo un modelo correcto son
  \textbf{uniformes} \((0,1)\), comparables con un Q--Q y un test de
  Kolmogorov--Smirnov. No son ``de binaria'': valen para
  \textbf{cualquier} GLM ---son, simplemente, el rescate cuando los
  estándar fallan---.
\item
  \textbf{Diagnóstico de calibración} (¿son fiables las probabilidades,
  y en qué tramo fallan?) → los \textbf{residuos por tramos}
  (\emph{binned}), que promedian observado − predicho por tramos de
  riesgo. Como responden a la pregunta de la calibración, los tratamos
  \textbf{donde se usan}: en la subsección siguiente.
\end{itemize}

\begin{tcolorbox}[enhanced jigsaw, coltitle=black, rightrule=.15mm, arc=.35mm, leftrule=.75mm, colback=white, colbacktitle=quarto-callout-tip-color!10!white, left=2mm, breakable, toptitle=1mm, toprule=.15mm, bottomtitle=1mm, titlerule=0mm, title=\textcolor{quarto-callout-tip-color}{\faLightbulb}\hspace{0.5em}{🔧 En R. Residuos simulados (DHARMa)}, bottomrule=.15mm, opacityback=0, opacitybacktitle=0.6, colframe=quarto-callout-tip-color-frame]

\begin{Shaded}
\begin{Highlighting}[]
\NormalTok{sim }\OtherTok{\textless{}{-}}\NormalTok{ DHARMa}\SpecialCharTok{::}\FunctionTok{simulateResiduals}\NormalTok{(fit\_glm)   }\CommentTok{\# plot = TRUE por defecto}
\FunctionTok{plot}\NormalTok{(sim)                                   }\CommentTok{\# Q–Q (con test KS) + residuos vs predicción}
\NormalTok{DHARMa}\SpecialCharTok{::}\FunctionTok{testResiduals}\NormalTok{(sim)                  }\CommentTok{\# tests formales: uniformidad, dispersión, outliers}
\end{Highlighting}
\end{Shaded}

Los de Pearson y deviance, con
\texttt{residuals(fit\_glm,\ type\ =\ "pearson"\ \textbar{}\ "deviance")}.

\end{tcolorbox}

\begin{Shaded}
\begin{Highlighting}[]
\NormalTok{sim }\OtherTok{\textless{}{-}}\NormalTok{ DHARMa}\SpecialCharTok{::}\FunctionTok{simulateResiduals}\NormalTok{(fit\_glm, }\AttributeTok{plot =} \ConstantTok{FALSE}\NormalTok{)}
\FunctionTok{plot}\NormalTok{(sim)}
\end{Highlighting}
\end{Shaded}

\begin{figure}[H]

{\centering \includegraphics{caso1_binaria_files/figure-pdf/fig-u12-dharma-1.pdf}

}

\caption{\label{fig-u12-dharma}Residuos simulados (DHARMa) del modelo
logístico en GLOW. Izquierda: Q--Q de los residuos, que deberían seguir
la diagonal uniforme. Derecha: residuos frente a predicción, sin patrón
si el ajuste es adecuado.}

\end{figure}

En GLOW, los residuos simulados \textbf{no muestran ningún problema}: el
Q--Q sigue la diagonal uniforme (test de Kolmogorov--Smirnov
\(p\approx0{,}81\)) y no hay señal de \textbf{sobredispersión}
(\(p\approx0{,}87\)) ni de \textbf{outliers} (\(p\approx1\)); a la
derecha, los residuos se reparten \textbf{sin patrón} frente a la
predicción. Ninguna evidencia, por tanto, de mala especificación ---ni
en la forma, ni en la dispersión, ni en casos atípicos---.

\hypertarget{calibraciuxf3n-son-fiables-las-probabilidades}{%
\subsubsection{Calibración: ¿son fiables las
probabilidades?}\label{calibraciuxf3n-son-fiables-las-probabilidades}}

Confronta la probabilidad \textbf{predicha} con la frecuencia
\textbf{observada}. El contraste de referencia es el \textbf{test de
Hosmer--Lemeshow} (HL) (Hosmer \& Lemeshow, 1980): se ordenan las \(n\)
observaciones por su probabilidad predicha \(\hat\pi_i\) y se reparten
en \(G\) grupos de igual tamaño (habitualmente \(G=10\), deciles de
riesgo). En cada grupo \(g\) se cuentan los eventos \textbf{observados}
\(O_g=\sum_{i\in g}y_i\) y los \textbf{esperados}
\(E_g=\sum_{i\in g}\hat\pi_i=n_g\bar\pi_g\) (con \(\bar\pi_g\) la
probabilidad media predicha del grupo), y el estadístico los confronta:
\[
\hat C=\sum_{g=1}^{G}\frac{(O_g-E_g)^2}{n_g\,\bar\pi_g(1-\bar\pi_g)}.
\] Bajo \(H_0\) («el modelo calibra bien»), \(\hat C\) se distribuye
aproximadamente \(\chi^2_{G-2}\) ---los \(G-2\) grados de libertad,
establecidos por simulación, descuentan dos por el agrupamiento y la
estimación---. Un \(\hat C\) \textbf{grande} (\(p\) bajo) indica que
observado y esperado \textbf{discrepan} → mala calibración; un \(p\)
\textbf{alto} señala \textbf{ausencia de evidencia} de mala calibración
(no «prueba» que calibre bien).

\begin{Shaded}
\begin{Highlighting}[]
\NormalTok{performance}\SpecialCharTok{::}\FunctionTok{performance\_hosmer}\NormalTok{(fit\_glm, }\AttributeTok{n\_bins =} \DecValTok{10}\NormalTok{)}
\end{Highlighting}
\end{Shaded}

\begin{verbatim}
# Hosmer-Lemeshow Goodness-of-Fit Test

  Chi-squared: 10.816
           df:  8    
      p-value:  0.212
\end{verbatim}

En GLOW, \(p\approx 0{,}21\): sin evidencia de mala calibración. Pero el
HL tiene \textbf{limitaciones conocidas} ---sensible al número de grupos
y poco potente---, así que la práctica moderna lo \textbf{acompaña de un
gráfico de calibración} y de los residuos por tramos (Agresti, 2015;
Harrell, 2015):

\begin{Shaded}
\begin{Highlighting}[]
\NormalTok{diag }\OtherTok{\textless{}{-}} \FunctionTok{tibble}\NormalTok{(}\AttributeTok{p\_hat =} \FunctionTok{fitted}\NormalTok{(fit\_glm), }\AttributeTok{y =}\NormalTok{ glow}\SpecialCharTok{$}\NormalTok{fractura01)}

\CommentTok{\# {-}{-}{-} Panel A: residuos por tramos (estilo binnedplot, en ggplot) {-}{-}{-}}
\NormalTok{binned }\OtherTok{\textless{}{-}}\NormalTok{ diag }\SpecialCharTok{|\textgreater{}}
  \FunctionTok{mutate}\NormalTok{(}\AttributeTok{tramo =} \FunctionTok{ntile}\NormalTok{(p\_hat, }\DecValTok{20}\NormalTok{)) }\SpecialCharTok{|\textgreater{}}
  \FunctionTok{group\_by}\NormalTok{(tramo) }\SpecialCharTok{|\textgreater{}}
  \FunctionTok{summarise}\NormalTok{(}\AttributeTok{p\_med =} \FunctionTok{mean}\NormalTok{(p\_hat),}
            \AttributeTok{r\_med =} \FunctionTok{mean}\NormalTok{(y }\SpecialCharTok{{-}}\NormalTok{ p\_hat),}
            \AttributeTok{se    =} \FunctionTok{sqrt}\NormalTok{(}\FunctionTok{mean}\NormalTok{(p\_hat }\SpecialCharTok{*}\NormalTok{ (}\DecValTok{1} \SpecialCharTok{{-}}\NormalTok{ p\_hat)) }\SpecialCharTok{/} \FunctionTok{n}\NormalTok{()), }\AttributeTok{.groups =} \StringTok{"drop"}\NormalTok{)}

\NormalTok{pA }\OtherTok{\textless{}{-}} \FunctionTok{ggplot}\NormalTok{(binned, }\FunctionTok{aes}\NormalTok{(p\_med, r\_med)) }\SpecialCharTok{+}
  \FunctionTok{geom\_ribbon}\NormalTok{(}\FunctionTok{aes}\NormalTok{(}\AttributeTok{ymin =} \SpecialCharTok{{-}}\DecValTok{2} \SpecialCharTok{*}\NormalTok{ se, }\AttributeTok{ymax =} \DecValTok{2} \SpecialCharTok{*}\NormalTok{ se), }\AttributeTok{alpha =} \FloatTok{0.15}\NormalTok{) }\SpecialCharTok{+}
  \FunctionTok{geom\_hline}\NormalTok{(}\AttributeTok{yintercept =} \DecValTok{0}\NormalTok{, }\AttributeTok{linetype =} \StringTok{"dashed"}\NormalTok{) }\SpecialCharTok{+}
  \FunctionTok{geom\_point}\NormalTok{(}\AttributeTok{size =} \FloatTok{1.8}\NormalTok{) }\SpecialCharTok{+}
  \FunctionTok{labs}\NormalTok{(}\AttributeTok{x =} \StringTok{"Probabilidad ajustada"}\NormalTok{, }\AttributeTok{y =} \StringTok{"Residuo medio por tramo"}\NormalTok{,}
       \AttributeTok{title =} \StringTok{"Residuos por tramos"}\NormalTok{)}

\CommentTok{\# {-}{-}{-} Panel B: calibración por deciles {-}{-}{-}}
\NormalTok{calib }\OtherTok{\textless{}{-}}\NormalTok{ diag }\SpecialCharTok{|\textgreater{}}
  \FunctionTok{mutate}\NormalTok{(}\AttributeTok{decil =} \FunctionTok{ntile}\NormalTok{(p\_hat, }\DecValTok{10}\NormalTok{)) }\SpecialCharTok{|\textgreater{}}
  \FunctionTok{group\_by}\NormalTok{(decil) }\SpecialCharTok{|\textgreater{}}
  \FunctionTok{summarise}\NormalTok{(}\AttributeTok{pred =} \FunctionTok{mean}\NormalTok{(p\_hat), }\AttributeTok{obs =} \FunctionTok{mean}\NormalTok{(y), }\AttributeTok{.groups =} \StringTok{"drop"}\NormalTok{)}

\NormalTok{pB }\OtherTok{\textless{}{-}} \FunctionTok{ggplot}\NormalTok{(calib, }\FunctionTok{aes}\NormalTok{(pred, obs)) }\SpecialCharTok{+}
  \FunctionTok{geom\_abline}\NormalTok{(}\AttributeTok{slope =} \DecValTok{1}\NormalTok{, }\AttributeTok{intercept =} \DecValTok{0}\NormalTok{, }\AttributeTok{linetype =} \StringTok{"dashed"}\NormalTok{) }\SpecialCharTok{+}   \CommentTok{\# calibración perfecta}
  \FunctionTok{geom\_line}\NormalTok{() }\SpecialCharTok{+} \FunctionTok{geom\_point}\NormalTok{(}\AttributeTok{size =} \FloatTok{1.8}\NormalTok{) }\SpecialCharTok{+}
  \FunctionTok{coord\_equal}\NormalTok{(}\AttributeTok{xlim =} \FunctionTok{c}\NormalTok{(}\DecValTok{0}\NormalTok{, }\DecValTok{1}\NormalTok{), }\AttributeTok{ylim =} \FunctionTok{c}\NormalTok{(}\DecValTok{0}\NormalTok{, }\DecValTok{1}\NormalTok{)) }\SpecialCharTok{+}
  \FunctionTok{labs}\NormalTok{(}\AttributeTok{x =} \StringTok{"Riesgo predicho (media por decil)"}\NormalTok{, }\AttributeTok{y =} \StringTok{"Frecuencia observada"}\NormalTok{,}
       \AttributeTok{title =} \StringTok{"Calibración"}\NormalTok{)}

\NormalTok{pA }\SpecialCharTok{+}\NormalTok{ pB        }\CommentTok{\# patchwork: lado a lado}
\end{Highlighting}
\end{Shaded}

\begin{figure}[H]

{\centering \includegraphics[width=1\textwidth,height=\textheight]{caso1_binaria_files/figure-pdf/fig-u12-calibracion-1.pdf}

}

\caption{\label{fig-u12-calibracion}Diagnóstico visual. Izquierda:
residuos por tramos (deberían caer dentro de las bandas de ±2 EE).
Derecha: curva de calibración por deciles de riesgo (debería seguir la
diagonal).}

\end{figure}

Ambos paneles convierten el ``pasa / no pasa'' del HL en un
\textbf{diagnóstico localizado}: dicen no solo \emph{si} hay problema,
sino \emph{dónde}, \emph{en qué sentido} y \emph{qué} podría corregirlo.

\begin{tcolorbox}[enhanced jigsaw, coltitle=black, rightrule=.15mm, arc=.35mm, leftrule=.75mm, colback=white, colbacktitle=quarto-callout-tip-color!10!white, left=2mm, breakable, toptitle=1mm, toprule=.15mm, bottomtitle=1mm, titlerule=0mm, title=\textcolor{quarto-callout-tip-color}{\faLightbulb}\hspace{0.5em}{Cómo leer los dos gráficos}, bottomrule=.15mm, opacityback=0, opacitybacktitle=0.6, colframe=quarto-callout-tip-color-frame]

\textbf{Curva de calibración} (probabilidad predicha en X, frecuencia
observada en Y; la referencia es la diagonal de 45°):

\begin{itemize}
\tightlist
\item
  \textbf{Sobre la diagonal} → bien calibrado: entre los casos con
  riesgo predicho \(p\), una fracción \(\approx p\) tiene el evento.
\item
  \textbf{Por debajo de la diagonal} (sobre todo en riesgo alto) → el
  modelo \textbf{sobreestima} el riesgo en esa zona.
\item
  \textbf{Por encima de la diagonal} (sobre todo en riesgo bajo) →
  \textbf{subestima}.
\item
  \textbf{Forma en S} alrededor de la diagonal → exceso o defecto de
  confianza en los extremos.
\item
  Lo informativo es \textbf{en qué tramo} de riesgo aparece la
  desviación, no solo que exista.
\end{itemize}

\textbf{Residuos por tramos} (residuo medio por tramo en Y, probabilidad
media en X, con bandas de \(\pm 2\) EE \(\approx 95\,\%\)):

\begin{itemize}
\tightlist
\item
  \textbf{Dispersos en torno a 0 y dentro de las bandas} → sin falta de
  ajuste sistemática.
\item
  \textbf{Patrón sistemático} (p.~ej. negativos en riesgo bajo y
  positivos en alto, o una curva) → el modelo \textbf{yerra la media} en
  ese rango.
\item
  \textbf{Más de \textasciitilde5 \% de puntos fuera de las bandas} →
  señal de mal ajuste.
\item
  Si se representa el residuo medio por tramos frente a una
  \textbf{covariable} (la edad, por ejemplo) en lugar de frente a la
  probabilidad ajustada, un patrón indica que \textbf{esa variable pide
  un término no lineal} (un cuadrático) o una interacción: orienta sobre
  \textbf{qué} cambiar.
\end{itemize}

\end{tcolorbox}

En GLOW (Figura~\ref{fig-u12-calibracion}), coherente con el HL
(\(p\approx0{,}21\)): los residuos por tramos quedan dentro de las
bandas y la calibración se ajusta a la diagonal, sin señales de
descalibración ni de falta de ajuste en el rango observado.

\begin{tcolorbox}[enhanced jigsaw, coltitle=black, rightrule=.15mm, arc=.35mm, leftrule=.75mm, colback=white, colbacktitle=quarto-callout-tip-color!10!white, left=2mm, breakable, toptitle=1mm, toprule=.15mm, bottomtitle=1mm, titlerule=0mm, title=\textcolor{quarto-callout-tip-color}{\faLightbulb}\hspace{0.5em}{🔧 En R. Hosmer--Lemeshow y curva de calibración}, bottomrule=.15mm, opacityback=0, opacitybacktitle=0.6, colframe=quarto-callout-tip-color-frame]

\textbf{Test de Hosmer--Lemeshow:}

\begin{Shaded}
\begin{Highlighting}[]
\NormalTok{performance}\SpecialCharTok{::}\FunctionTok{performance\_hosmer}\NormalTok{(fit\_glm, }\AttributeTok{n\_bins =} \DecValTok{10}\NormalTok{)}
\NormalTok{ResourceSelection}\SpecialCharTok{::}\FunctionTok{hoslem.test}\NormalTok{(glow}\SpecialCharTok{$}\NormalTok{fractura01, }\FunctionTok{fitted}\NormalTok{(fit\_glm), }\AttributeTok{g =} \DecValTok{10}\NormalTok{)   }\CommentTok{\# alternativa}
\end{Highlighting}
\end{Shaded}

\textbf{Residuos por tramos (binned):}

\begin{Shaded}
\begin{Highlighting}[]
\FunctionTok{plot}\NormalTok{(performance}\SpecialCharTok{::}\FunctionTok{binned\_residuals}\NormalTok{(fit\_glm))   }\CommentTok{\# observado − predicho promediado por tramo, ±2 EE}
\end{Highlighting}
\end{Shaded}

\textbf{Curva de calibración} (riesgo predicho vs frecuencia observada,
por deciles):

\begin{Shaded}
\begin{Highlighting}[]
\FunctionTok{tibble}\NormalTok{(}\AttributeTok{p =} \FunctionTok{fitted}\NormalTok{(fit\_glm), }\AttributeTok{y =}\NormalTok{ glow}\SpecialCharTok{$}\NormalTok{fractura01) }\SpecialCharTok{|\textgreater{}}
  \FunctionTok{mutate}\NormalTok{(}\AttributeTok{grupo =} \FunctionTok{ntile}\NormalTok{(p, }\DecValTok{10}\NormalTok{)) }\SpecialCharTok{|\textgreater{}}
  \FunctionTok{group\_by}\NormalTok{(grupo) }\SpecialCharTok{|\textgreater{}} \FunctionTok{summarise}\NormalTok{(}\AttributeTok{pred =} \FunctionTok{mean}\NormalTok{(p), }\AttributeTok{obs =} \FunctionTok{mean}\NormalTok{(y)) }\SpecialCharTok{|\textgreater{}}
  \FunctionTok{ggplot}\NormalTok{(}\FunctionTok{aes}\NormalTok{(pred, obs)) }\SpecialCharTok{+}
  \FunctionTok{geom\_abline}\NormalTok{(}\AttributeTok{linetype =} \StringTok{"dashed"}\NormalTok{) }\SpecialCharTok{+} \FunctionTok{geom\_point}\NormalTok{() }\SpecialCharTok{+} \FunctionTok{geom\_line}\NormalTok{()}
\end{Highlighting}
\end{Shaded}

(Curva suavizada e índices: \texttt{rms::val.prob()},
\texttt{performance::check\_model()}.)

\end{tcolorbox}

\hypertarget{separaciuxf3n-cuando-el-ajuste-perfecto-rompe-la-inferencia.}{%
\subsubsection{Separación: cuando el ajuste ``perfecto'' rompe la
inferencia.}\label{separaciuxf3n-cuando-el-ajuste-perfecto-rompe-la-inferencia.}}

Es el patológico que vuelve inútil al Wald. Hablamos de ``separación''
cuando una combinación de covariables \textbf{predice perfectamente} el
resultado en algún subgrupo: la verosimilitud sigue creciendo al alejar
el coeficiente hacia \(\pm\infty\), así que \texttt{glm()} devuelve
estimaciones \textbf{enormes con errores estándar gigantescos} (Albert
\& Anderson, 1984).

GLOW lo exhibe de forma natural: entre las fumadoras que autoperciben su
riesgo como «menor» (\texttt{smoke\ =\ Yes}, \texttt{raterisk\ =\ Less})
\textbf{no hay ninguna fractura} (0 de 12), lo que separa ese grupo.

\begin{Shaded}
\begin{Highlighting}[]
\NormalTok{glow }\SpecialCharTok{|\textgreater{}} \FunctionTok{count}\NormalTok{(smoke, raterisk, fracture) }\SpecialCharTok{|\textgreater{}}
  \FunctionTok{filter}\NormalTok{(smoke }\SpecialCharTok{==} \StringTok{"Yes"}\NormalTok{, raterisk }\SpecialCharTok{==} \StringTok{"Less"}\NormalTok{)        }\CommentTok{\# celda con 0 fracturas}
\end{Highlighting}
\end{Shaded}

\begin{verbatim}
# A tibble: 1 x 4
  smoke raterisk fracture     n
  <fct> <fct>    <fct>    <int>
1 Yes   Less     No          12
\end{verbatim}

\begin{Shaded}
\begin{Highlighting}[]
\NormalTok{fit\_sep }\OtherTok{\textless{}{-}} \FunctionTok{glm}\NormalTok{(fracture }\SpecialCharTok{\textasciitilde{}}\NormalTok{ smoke }\SpecialCharTok{*}\NormalTok{ raterisk, binomial, glow)}
\FunctionTok{tidy}\NormalTok{(fit\_sep) }\SpecialCharTok{|\textgreater{}} \FunctionTok{filter}\NormalTok{(}\FunctionTok{str\_detect}\NormalTok{(term, }\StringTok{"smoke"}\NormalTok{))  }\CommentTok{\# coef enorme, EE descomunal}
\end{Highlighting}
\end{Shaded}

\begin{verbatim}
# A tibble: 3 x 5
  term                     estimate std.error statistic p.value
  <chr>                       <dbl>     <dbl>     <dbl>   <dbl>
1 smokeYes                    -15.1      693.   -0.0217   0.983
2 smokeYes:rateriskSame        14.5      693.    0.0209   0.983
3 smokeYes:rateriskGreater     15.6      693.    0.0225   0.982
\end{verbatim}

El coeficiente de \texttt{smokeYes} ronda \textbf{−15 con un EE de
varios cientos}: síntoma inequívoco de cuasi-separación, no un efecto
«protector» real. El Wald carece de sentido aquí.

La solución es detectarla
(\texttt{detectseparation::detect\_separation}) y, si aparece, estimar
por \textbf{verosimilitud penalizada de Firth} (Firth, 1993; Heinze \&
Schemper, 2002), que produce estimaciones finitas:

\begin{Shaded}
\begin{Highlighting}[]
\CommentTok{\# install.packages("logistf")  \# o brglm2::brglm\_fit}
\CommentTok{\# La verosimilitud penalizada de Firth devuelve estimaciones FINITAS e inferencia con sentido:}
\NormalTok{logistf}\SpecialCharTok{::}\FunctionTok{logistf}\NormalTok{(fracture }\SpecialCharTok{\textasciitilde{}}\NormalTok{ smoke }\SpecialCharTok{*}\NormalTok{ raterisk, }\AttributeTok{data =}\NormalTok{ glow)}
\end{Highlighting}
\end{Shaded}

\begin{verbatim}
logistf::logistf(formula = fracture ~ smoke * raterisk, data = glow)
Model fitted by Penalized ML
Confidence intervals and p-values by Profile Likelihood 

Coefficients:
             (Intercept)                 smokeYes             rateriskSame 
              -1.4982123               -1.7206635                0.4817357 
         rateriskGreater    smokeYes:rateriskSame smokeYes:rateriskGreater 
               0.7664928                1.3020556                2.2853289 

Likelihood ratio test=15.12048 on 5 df, p=0.009859754, n=500
\end{verbatim}

Se ve inmediatamente la corrección: el término de \texttt{smoke} pasa de
\(\approx-15\) con EE de cientos y \(p\approx1\) (Wald, roto) a una
estimación \textbf{finita}, con IC y \(p\) razonables (Firth).

\begin{tcolorbox}[enhanced jigsaw, coltitle=black, rightrule=.15mm, arc=.35mm, leftrule=.75mm, colback=white, colbacktitle=quarto-callout-tip-color!10!white, left=2mm, breakable, toptitle=1mm, toprule=.15mm, bottomtitle=1mm, titlerule=0mm, title=\textcolor{quarto-callout-tip-color}{\faLightbulb}\hspace{0.5em}{🔧 En R. Detectar la separación y estimar por verosimilitud penalizada
(Firth)}, bottomrule=.15mm, opacityback=0, opacitybacktitle=0.6, colframe=quarto-callout-tip-color-frame]

\textbf{Detección} (antes de fiarte de un \texttt{glm} con coeficientes
enormes y EE descomunales; también avisa el warning \emph{``fitted
probabilities numerically 0 or 1 occurred''}):

\begin{Shaded}
\begin{Highlighting}[]
\NormalTok{detectseparation}\SpecialCharTok{::}\FunctionTok{detect\_separation}\NormalTok{(fracture }\SpecialCharTok{\textasciitilde{}}\NormalTok{ smoke }\SpecialCharTok{*}\NormalTok{ raterisk,}
                                    \AttributeTok{family =}\NormalTok{ binomial, }\AttributeTok{data =}\NormalTok{ glow)   }\CommentTok{\# marca los términos separados}
\end{Highlighting}
\end{Shaded}

\textbf{Estimación penalizada de Firth} ---penaliza la verosimilitud
(corrección de Jeffreys), de modo que los coeficientes quedan
\textbf{finitos} incluso bajo separación, con inferencia por
verosimilitud de perfil---:

\begin{Shaded}
\begin{Highlighting}[]
\NormalTok{logistf}\SpecialCharTok{::}\FunctionTok{logistf}\NormalTok{(fracture }\SpecialCharTok{\textasciitilde{}}\NormalTok{ smoke }\SpecialCharTok{*}\NormalTok{ raterisk, }\AttributeTok{data =}\NormalTok{ glow)            }\CommentTok{\# coef, IC y p finitos}
\FunctionTok{glm}\NormalTok{(fracture }\SpecialCharTok{\textasciitilde{}}\NormalTok{ smoke }\SpecialCharTok{*}\NormalTok{ raterisk, }\AttributeTok{family =}\NormalTok{ binomial, }\AttributeTok{data =}\NormalTok{ glow,}
    \AttributeTok{method =}\NormalTok{ brglm2}\SpecialCharTok{::}\NormalTok{brglm\_fit)                                       }\CommentTok{\# alternativa (brglm2)}
\end{Highlighting}
\end{Shaded}

\end{tcolorbox}

\hypertarget{sec-u12-evaluacion}{%
\subsection{2.7 Evaluación como clasificador}\label{sec-u12-evaluacion}}

El modelo entrega \textbf{probabilidades}; convertirlo en un
\textbf{clasificador} (predecir sí/no) exige un \textbf{umbral} \(c\) y
abre un tercer eje de evaluación del modelo: la \textbf{discriminación}
---¿separa eventos de no-eventos?---. Este enfoque de evaluación es
distinta al que se realiza sobre la bondad del ajuste (en la
\protect\hyperlink{sec-u12-diagnostico}{Sección 2.6} y sobre la utilidad
inferencial (en la \protect\hyperlink{sec-u12-inferencia}{Sección 2.5}).

\hypertarget{del-umbral-a-la-matriz-de-confusiuxf3n.}{%
\subsubsection{Del umbral a la matriz de
confusión.}\label{del-umbral-a-la-matriz-de-confusiuxf3n.}}

Fijado un umbral \(c \in [0,1]\), por defecto \(c=0{,}5\), clasificamos
con una observación en función de su predicción \(\hat\pi_i\) con:

\[\hat y_i=\mathbf{1}[\hat\pi_i\ge c]\]

Para evaluar la calidad de la clasificación, tenemos fácil la
comparación entre lo observado (los datos iniciales 0/1) y lo predicho
al clasificar. El cruce de lo predicho con lo observado define la
\textbf{matriz de confusión}, punto de partida de todas las métricas.

\textbf{Las cuatro celdas.}

Con un umbral dado \(c\), cada caso cae en una de cuatro celdas según
coincidan predicción y realidad: \textbf{verdadero positivo} (TP,
\emph{true positive}), \textbf{falso positivo} (FP, \emph{false
positive}), \textbf{falso negativo} (FN, \emph{false negative}) y
\textbf{verdadero negativo} (TN, \emph{true negative}). De ahí salen las
métricas más habituales (nombradas en inglés porque es así como las
nombra casi todo el software):

\begin{itemize}
\tightlist
\item
  \textbf{Exactitud} --- \emph{accuracy}:
  \(\dfrac{TP+TN}{\text{total}}\). Proporción global de aciertos.
\item
  \textbf{Sensibilidad} o \emph{recall} --- \emph{sensitivity / true
  positive rate}: \(\dfrac{TP}{TP+FN}\). De los eventos reales, cuántos
  detecta.
\item
  \textbf{Especificidad} --- \emph{specificity / true negative rate}:
  \(\dfrac{TN}{TN+FP}\). De los no-eventos reales, cuántos descarta.
\item
  \textbf{Precisión} o valor predictivo positivo --- \emph{precision /
  positive predictive value (PPV)}: \(\dfrac{TP}{TP+FP}\). De los que
  predice como eventos, cuántos lo son de verdad.
\end{itemize}

Sensibilidad y precisión suelen estar \textbf{en tensión} (bajar el
umbral detecta más eventos pero genera más falsas alarmas); cuando ambas
importan se utiliza el

\begin{itemize}
\tightlist
\item
  \textbf{F1} (\emph{F1-score}), la media armónica de sensibilidad y
  precisión,
  \(2\,\dfrac{\text{precisión}\cdot\text{recall}}{\text{precisión}+\text{recall}}\).
\end{itemize}

\textbf{Cuál mirar.}

A la hora de valorar la calidad de la clasificación, cabe tener en
cuenta estas consideraciones:

\begin{itemize}
\tightlist
\item
  La \textbf{exactitud engaña con clases desbalanceadas}: si una clase
  es rara, ``clasificar todo como la mayoritaria'' logra una exactitud
  alta y es inútil. Conviene entonces mirar \textbf{sensibilidad,
  especificidad y precisión por separado} y elegir según el
  \textbf{coste de cada error}: la sensibilidad prima cuando un falso
  negativo es caro (no detectar un evento); la precisión, cuando lo caro
  es la falsa alarma.
\item
  Además, sensibilidad y especificidad son \textbf{invariantes a la
  prevalencia} (son propiedades del clasificador dada la clase real),
  mientras que exactitud y precisión \textbf{sí dependen} de ella; por
  eso, y para resumir todos los umbrales a la vez, conviene
  complementarlas con el \textbf{AUC}, también invariante a la
  prevalencia.
\end{itemize}

Calculamos a mano la matriz de confusión para GLOW.

\begin{Shaded}
\begin{Highlighting}[]
\CommentTok{\# {-}{-}{-} A mano (para ver qué es cada cosa) {-}{-}{-}}
\NormalTok{umbral }\OtherTok{\textless{}{-}} \FloatTok{0.5}
\NormalTok{pred   }\OtherTok{\textless{}{-}} \FunctionTok{factor}\NormalTok{(}\FunctionTok{ifelse}\NormalTok{(}\FunctionTok{fitted}\NormalTok{(fit\_glm) }\SpecialCharTok{\textgreater{}=}\NormalTok{ umbral, }\StringTok{"Yes"}\NormalTok{, }\StringTok{"No"}\NormalTok{),}
                 \AttributeTok{levels =} \FunctionTok{c}\NormalTok{(}\StringTok{"No"}\NormalTok{, }\StringTok{"Yes"}\NormalTok{))}
\NormalTok{cm }\OtherTok{\textless{}{-}} \FunctionTok{table}\NormalTok{(}\AttributeTok{Predicho =}\NormalTok{ pred, }\AttributeTok{Observado =}\NormalTok{ glow}\SpecialCharTok{$}\NormalTok{fracture)}
\NormalTok{cm}
\end{Highlighting}
\end{Shaded}

\begin{verbatim}
        Observado
Predicho  No Yes
     No  360 111
     Yes  15  14
\end{verbatim}

\begin{Shaded}
\begin{Highlighting}[]
\NormalTok{TP }\OtherTok{\textless{}{-}}\NormalTok{ cm[}\StringTok{"Yes"}\NormalTok{, }\StringTok{"Yes"}\NormalTok{]; FP }\OtherTok{\textless{}{-}}\NormalTok{ cm[}\StringTok{"Yes"}\NormalTok{, }\StringTok{"No"}\NormalTok{]}
\NormalTok{FN }\OtherTok{\textless{}{-}}\NormalTok{ cm[}\StringTok{"No"}\NormalTok{,  }\StringTok{"Yes"}\NormalTok{]; TN }\OtherTok{\textless{}{-}}\NormalTok{ cm[}\StringTok{"No"}\NormalTok{,  }\StringTok{"No"}\NormalTok{]}
\FunctionTok{c}\NormalTok{(}\AttributeTok{exactitud     =}\NormalTok{ (TP }\SpecialCharTok{+}\NormalTok{ TN) }\SpecialCharTok{/} \FunctionTok{sum}\NormalTok{(cm),   }\CommentTok{\# accuracy}
  \AttributeTok{sensibilidad  =}\NormalTok{ TP }\SpecialCharTok{/}\NormalTok{ (TP }\SpecialCharTok{+}\NormalTok{ FN),        }\CommentTok{\# sensitivity / recall}
  \AttributeTok{especificidad =}\NormalTok{ TN }\SpecialCharTok{/}\NormalTok{ (TN }\SpecialCharTok{+}\NormalTok{ FP),        }\CommentTok{\# specificity}
  \AttributeTok{precision     =}\NormalTok{ TP }\SpecialCharTok{/}\NormalTok{ (TP }\SpecialCharTok{+}\NormalTok{ FP)) }\SpecialCharTok{|\textgreater{}} \FunctionTok{round}\NormalTok{(}\DecValTok{3}\NormalTok{)   }\CommentTok{\# precision / PPV}
\end{Highlighting}
\end{Shaded}

\begin{verbatim}
    exactitud  sensibilidad especificidad     precision 
        0.748         0.112         0.960         0.483 
\end{verbatim}

En la práctica no se calcula a mano: \texttt{yardstick} (Kuhn et~al.,
2026) devuelve la matriz y todas las métricas en una llamada. El único
cuidado es declarar la \textbf{clase positiva}; como aquí \texttt{Yes}
es el \textbf{segundo} nivel del factor, se indica con
\texttt{event\_level\ =\ "second"} (si no, sensibilidad y especificidad
saldrían intercambiadas):

\begin{Shaded}
\begin{Highlighting}[]
\FunctionTok{library}\NormalTok{(yardstick)}
\NormalTok{eval\_df }\OtherTok{\textless{}{-}} \FunctionTok{tibble}\NormalTok{(}\AttributeTok{obs =}\NormalTok{ glow}\SpecialCharTok{$}\NormalTok{fracture, }\AttributeTok{pred =}\NormalTok{ pred)}

\FunctionTok{conf\_mat}\NormalTok{(eval\_df, }\AttributeTok{truth =}\NormalTok{ obs, }\AttributeTok{estimate =}\NormalTok{ pred)        }\CommentTok{\# misma matriz de confusión}
\end{Highlighting}
\end{Shaded}

\begin{verbatim}
          Truth
Prediction  No Yes
       No  360 111
       Yes  15  14
\end{verbatim}

\begin{Shaded}
\begin{Highlighting}[]
\FunctionTok{metric\_set}\NormalTok{(accuracy, sensitivity, specificity, precision, recall, f\_meas)(}
\NormalTok{  eval\_df, }\AttributeTok{truth =}\NormalTok{ obs, }\AttributeTok{estimate =}\NormalTok{ pred, }\AttributeTok{event\_level =} \StringTok{"second"}\NormalTok{)}
\end{Highlighting}
\end{Shaded}

\begin{verbatim}
# A tibble: 6 x 3
  .metric     .estimator .estimate
  <chr>       <chr>          <dbl>
1 accuracy    binary         0.748
2 sensitivity binary         0.112
3 specificity binary         0.96 
4 precision   binary         0.483
5 recall      binary         0.112
6 f_meas      binary         0.182
\end{verbatim}

\begin{tcolorbox}[enhanced jigsaw, coltitle=black, rightrule=.15mm, arc=.35mm, leftrule=.75mm, colback=white, colbacktitle=quarto-callout-tip-color!10!white, left=2mm, breakable, toptitle=1mm, toprule=.15mm, bottomtitle=1mm, titlerule=0mm, title=\textcolor{quarto-callout-tip-color}{\faLightbulb}\hspace{0.5em}{🔧 En R. Matriz de confusión y métricas}, bottomrule=.15mm, opacityback=0, opacitybacktitle=0.6, colframe=quarto-callout-tip-color-frame]

A partir de la respuesta observada y la predicha (factores con los
\textbf{mismos niveles}), \texttt{yardstick} da todo en una llamada; el
único cuidado es declarar la \textbf{clase positiva} con
\texttt{event\_level}:

\begin{Shaded}
\begin{Highlighting}[]
\FunctionTok{library}\NormalTok{(yardstick)}
\NormalTok{pred    }\OtherTok{\textless{}{-}} \FunctionTok{factor}\NormalTok{(}\FunctionTok{ifelse}\NormalTok{(}\FunctionTok{fitted}\NormalTok{(fit\_glm) }\SpecialCharTok{\textgreater{}=} \FloatTok{0.5}\NormalTok{, }\StringTok{"Yes"}\NormalTok{, }\StringTok{"No"}\NormalTok{), }\AttributeTok{levels =} \FunctionTok{c}\NormalTok{(}\StringTok{"No"}\NormalTok{, }\StringTok{"Yes"}\NormalTok{))}
\NormalTok{eval\_df }\OtherTok{\textless{}{-}} \FunctionTok{tibble}\NormalTok{(}\AttributeTok{obs =}\NormalTok{ glow}\SpecialCharTok{$}\NormalTok{fracture, }\AttributeTok{pred =}\NormalTok{ pred)}

\FunctionTok{conf\_mat}\NormalTok{(eval\_df, }\AttributeTok{truth =}\NormalTok{ obs, }\AttributeTok{estimate =}\NormalTok{ pred)                       }\CommentTok{\# matriz de confusión}
\FunctionTok{metric\_set}\NormalTok{(accuracy, sensitivity, specificity, precision, recall, f\_meas)(}
\NormalTok{  eval\_df, }\AttributeTok{truth =}\NormalTok{ obs, }\AttributeTok{estimate =}\NormalTok{ pred, }\AttributeTok{event\_level =} \StringTok{"second"}\NormalTok{)      }\CommentTok{\# \textquotesingle{}second\textquotesingle{} = positivo en el 2º nivel}
\end{Highlighting}
\end{Shaded}

Alternativa:
\texttt{caret::confusionMatrix(pred,\ obs,\ positive\ =\ "Yes")}.

\end{tcolorbox}

\hypertarget{roc-y-auc}{%
\subsubsection{ROC y AUC}\label{roc-y-auc}}

Todas las métricas anteriores dependen de un \textbf{umbral fijo} para
clasificar a los sujetos en Yes/No (\(c=0{,}5\) por defecto): cambiar el
corte cambia sensibilidad, especificidad y precisión. Para evaluar el
clasificador \textbf{con independencia del umbral} ---y, como acabamos
de ver, de la prevalencia--- se recorre el corte de 0 a 1 y se resume el
compromiso entre acertar eventos y evitar falsas alarmas. Eso es la
curva \textbf{ROC} y su área, el \textbf{AUC} (Robin et~al., 2011).

La \textbf{curva ROC} (\emph{receiver operating characteristic}) recorre
\textbf{todos los umbrales} \(c\in[0,1]\) y, para cada uno, dibuja la
\textbf{sensibilidad} (tasa de verdaderos positivos, TPR) frente a
\textbf{1 − especificidad} (tasa de falsos positivos, FPR):

\begin{eqnarray}
\text{TPR}(c)&=&\frac{TP(c)}{TP(c)+FN(c)}=\Pr(\hat\pi\ge c\mid Y=1),\\
\text{FPR}(c)&=&\frac{FP(c)}{FP(c)+TN(c)}=\Pr(\hat\pi\ge c\mid Y=0).
\end{eqnarray}

La curva es el conjunto de puntos
\(\{(\text{FPR}(c),\,\text{TPR}(c)):c\in[0,1]\}\); cada punto es un
clasificador posible. La \textbf{diagonal} es el azar (TPR = FPR), y
cuanto más se arrima la curva a la \textbf{esquina superior izquierda},
mejor separa (ver Figura~\ref{fig-u12-roc}).

El \textbf{AUC} (\emph{area under the curve}) es el área bajo esa curva,
\[
\text{AUC}=\int_0^1 \text{TPR}\,\mathrm{d}(\text{FPR}),
\] con \(0{,}5\) = azar y \(1\) = perfecto. Tiene una lectura
\textbf{probabilística} elegante (equivale al estadístico de
Mann--Whitney): es la \textbf{probabilidad de que el modelo asigne más
riesgo a un evento que a un no-evento} tomados al azar, \[
\text{AUC}=\Pr(\hat\pi_i>\hat\pi_j\mid Y_i=1,\,Y_j=0),
\] que se estima como la \textbf{proporción de pares (evento, no-evento)
bien ordenados} (los empates cuentan media). Al promediar sobre todos
los umbrales, mide la \textbf{discriminación} con independencia del
corte y de la prevalencia.

\begin{Shaded}
\begin{Highlighting}[]
\NormalTok{roc\_obj }\OtherTok{\textless{}{-}}\NormalTok{ pROC}\SpecialCharTok{::}\FunctionTok{roc}\NormalTok{(glow}\SpecialCharTok{$}\NormalTok{fracture, }\FunctionTok{fitted}\NormalTok{(fit\_glm), }\AttributeTok{quiet =} \ConstantTok{TRUE}\NormalTok{)}
\NormalTok{pROC}\SpecialCharTok{::}\FunctionTok{auc}\NormalTok{(roc\_obj)}
\end{Highlighting}
\end{Shaded}

\begin{verbatim}
Area under the curve: 0.6807
\end{verbatim}

\begin{Shaded}
\begin{Highlighting}[]
\NormalTok{pROC}\SpecialCharTok{::}\FunctionTok{ggroc}\NormalTok{(roc\_obj) }\SpecialCharTok{+}
  \FunctionTok{geom\_abline}\NormalTok{(}\AttributeTok{slope =} \DecValTok{1}\NormalTok{, }\AttributeTok{intercept =} \DecValTok{1}\NormalTok{, }\AttributeTok{linetype =} \StringTok{"dashed"}\NormalTok{) }\SpecialCharTok{+}
  \FunctionTok{labs}\NormalTok{(}\AttributeTok{x =} \StringTok{"Especificidad"}\NormalTok{, }\AttributeTok{y =} \StringTok{"Sensibilidad"}\NormalTok{)}
\end{Highlighting}
\end{Shaded}

\begin{figure}[H]

{\centering \includegraphics{caso1_binaria_files/figure-pdf/fig-u12-roc-1.pdf}

}

\caption{\label{fig-u12-roc}Curva ROC del modelo logístico en GLOW. El
AUC resume la capacidad de discriminación agregada sobre todos los
umbrales.}

\end{figure}

El AUC ≈ \textbf{0,68}: discriminación \textbf{modesta}, coherente con
que la fractura depende de muchos factores no incluidos. Un AUC alto no
garantiza buenas probabilidades: eso es \textbf{calibración}, ya vista
en la \protect\hyperlink{sec-u12-diagnostico}{Sección 2.6}.

\textbf{Elegir el umbral es una decisión, no un dato.}

El \(0{,}5\) por defecto rara vez es óptimo: depende del \textbf{coste
relativo} de un falso positivo frente a un falso negativo. Un criterio
habitual y \textbf{simétrico} (trata ambos errores por igual) es el
\textbf{índice de Youden}, \[
J(c)=\text{sensibilidad}(c)+\text{especificidad}(c)-1=\text{TPR}(c)-\text{FPR}(c),
\] con \(J\in[0,1]\) (\(0\) = no discrimina, sobre la diagonal; \(1\) =
perfecto). El umbral elegido es el que \textbf{maximiza} \(J\), \[
c^\star=\arg\max_{c}\big[\text{TPR}(c)-\text{FPR}(c)\big],
\] que geométricamente es el punto de la curva ROC \textbf{más alejado
verticalmente de la diagonal}. Cuando los costes de los dos errores
\textbf{no} son simétricos, Youden deja de ser apropiado y se elige el
umbral que \textbf{minimiza el coste esperado} (ponderando FP y FN por
su coste).

\begin{Shaded}
\begin{Highlighting}[]
\NormalTok{pROC}\SpecialCharTok{::}\FunctionTok{coords}\NormalTok{(roc\_obj, }\StringTok{"best"}\NormalTok{, }\AttributeTok{best.method =} \StringTok{"youden"}\NormalTok{,}
             \AttributeTok{ret =} \FunctionTok{c}\NormalTok{(}\StringTok{"threshold"}\NormalTok{, }\StringTok{"sensitivity"}\NormalTok{, }\StringTok{"specificity"}\NormalTok{))}
\end{Highlighting}
\end{Shaded}

\begin{verbatim}
  threshold sensitivity specificity
1 0.2265844        0.68   0.6293333
\end{verbatim}

\begin{tcolorbox}[enhanced jigsaw, coltitle=black, rightrule=.15mm, arc=.35mm, leftrule=.75mm, colback=white, colbacktitle=quarto-callout-tip-color!10!white, left=2mm, breakable, toptitle=1mm, toprule=.15mm, bottomtitle=1mm, titlerule=0mm, title=\textcolor{quarto-callout-tip-color}{\faLightbulb}\hspace{0.5em}{🔧 En R. Curva ROC, AUC y umbral de Youden}, bottomrule=.15mm, opacityback=0, opacitybacktitle=0.6, colframe=quarto-callout-tip-color-frame]

Con \texttt{pROC}, a partir de la respuesta observada y las
probabilidades predichas:

\begin{Shaded}
\begin{Highlighting}[]
\FunctionTok{library}\NormalTok{(pROC)}
\NormalTok{roc\_obj }\OtherTok{\textless{}{-}} \FunctionTok{roc}\NormalTok{(glow}\SpecialCharTok{$}\NormalTok{fracture, }\FunctionTok{fitted}\NormalTok{(fit\_glm))     }\CommentTok{\# (respuesta, predictor)}
\FunctionTok{auc}\NormalTok{(roc\_obj)                                        }\CommentTok{\# AUC ; ci.auc(roc\_obj) para su IC 95\%}
\FunctionTok{ggroc}\NormalTok{(roc\_obj)                                      }\CommentTok{\# la curva ROC (ggplot)}

\FunctionTok{coords}\NormalTok{(roc\_obj, }\StringTok{"best"}\NormalTok{, }\AttributeTok{best.method =} \StringTok{"youden"}\NormalTok{,     }\CommentTok{\# umbral óptimo de Youden}
       \AttributeTok{ret =} \FunctionTok{c}\NormalTok{(}\StringTok{"threshold"}\NormalTok{, }\StringTok{"sensitivity"}\NormalTok{, }\StringTok{"specificity"}\NormalTok{))}
\end{Highlighting}
\end{Shaded}

Alternativa con \texttt{yardstick}:
\texttt{roc\_auc(eval\_df,\ truth,\ prob,\ event\_level\ =\ "second")} y
\texttt{roc\_curve(...)}.

\end{tcolorbox}

\begin{tcolorbox}[enhanced jigsaw, coltitle=black, rightrule=.15mm, arc=.35mm, leftrule=.75mm, colback=white, colbacktitle=quarto-callout-note-color!10!white, left=2mm, breakable, toptitle=1mm, toprule=.15mm, bottomtitle=1mm, titlerule=0mm, title=\textcolor{quarto-callout-note-color}{\faInfo}\hspace{0.5em}{Tres preguntas, tres herramientas}, bottomrule=.15mm, opacityback=0, opacitybacktitle=0.6, colframe=quarto-callout-note-color-frame]

Con el fin de resumir lo trabajado en esta unidad al respecto de cómo
evaluar y analizar qué nos proporciona un ajuste, hemos respondido una
serie de preguntas, y proporcionado una serie de herramientas con las
que contestarlas, y son:

\begin{itemize}
\tightlist
\item
  \textbf{¿Aporta? ¿qué modelo?} (utilidad e inferencia) → Wald, LRT
  ómnibus, AIC/BIC (\protect\hyperlink{sec-u12-inferencia}{2.5}).
\item
  \textbf{¿Ajusta bien?} (bondad de ajuste) → residuos, calibración,
  Hosmer--Lemeshow (\protect\hyperlink{sec-u12-diagnostico}{2.6}).
\item
  \textbf{¿Discrimina bien?} (clasificación) → matriz de confusión,
  ROC/AUC, umbral (\protect\hyperlink{sec-u12-evaluacion}{2.7}).
\end{itemize}

Un modelo puede ser \textbf{significativo} y aun así \textbf{calibrar
mal} o \textbf{discriminar poco}: son preguntas independientes.

\end{tcolorbox}

\begin{tcolorbox}[enhanced jigsaw, coltitle=black, rightrule=.15mm, arc=.35mm, leftrule=.75mm, colback=white, colbacktitle=quarto-callout-important-color!10!white, left=2mm, breakable, toptitle=1mm, toprule=.15mm, bottomtitle=1mm, titlerule=0mm, title=\textcolor{quarto-callout-important-color}{\faExclamation}\hspace{0.5em}{Tarea del estudiante (1.2) --- resuelve la Pregunta 1 del caso}, bottomrule=.15mm, opacityback=0, opacitybacktitle=0.6, colframe=quarto-callout-important-color-frame]

Completa la \textbf{Pregunta 1} del cliente
(\protect\hyperlink{sec-u16}{Unidad 6}) sobre la cohorte de tu equipo:
partiendo del \texttt{glm(ever\ \textasciitilde{}\ x1\ +\ x2)} de la
Unidad 1, interprétalo, decídelo y evalúalo.

Tareas:

\begin{enumerate}
\def\labelenumi{\arabic{enumi}.}
\tightlist
\item
  \textbf{Interpreta (RA3).} Da el efecto de la fragilidad y del
  tratamiento en \textbf{dos escalas} ---odds ratio (con IC) y efecto
  marginal medio (AME)---, en lenguaje clínico; explica por qué el
  coeficiente crudo no es un efecto sobre la probabilidad. Estima el
  \textbf{riesgo de una paciente nueva} con un perfil concreto.
\item
  \textbf{¿Protege por igual a todas? (RA5).} El cliente pregunta si el
  efecto del tratamiento depende de la fragilidad: añade la
  \textbf{interacción \texttt{x1:x2}}, contrástala (\textbf{LRT}) dentro
  de una escalera de modelos y compáralos por \textbf{AIC y BIC};
  documenta dónde discrepan y \textbf{justifica tu modelo final} (no
  basta el AIC mínimo).
\item
  \textbf{Evalúa (RA4).} Reporta \textbf{AUC}, evalúa la
  \textbf{calibración} (gráfico + Hosmer--Lemeshow, con su salvedad) y
  comprueba si hay \textbf{separación} en algún subgrupo; si la hubiera,
  explica cómo la tratarías.
\item
  \textbf{Responde al cliente.} Con lo anterior, redacta la respuesta a
  la Pregunta 1 y una recomendación accionable.
\end{enumerate}

\emph{Formato:} fragmento reproducible en el \texttt{.qmd} del equipo;
es el borrador de la Pregunta 1 del informe (ver
\protect\hyperlink{sec-u16}{Unidad 6}).

\end{tcolorbox}

\hypertarget{sec-u12-puente}{%
\subsection{Puente a la Unidad 3}\label{sec-u12-puente}}

Hasta aquí, una respuesta \textbf{binaria} modelada con \textbf{efectos
fijos}. Las próximas unidades generalizan en dos frentes. La
\protect\hyperlink{sec-u13}{Unidad 3} amplía el \textbf{tipo de
respuesta}: datos \textbf{agrupados} (binomial) y \textbf{más de dos
categorías} (politómica nominal y ordinal). La
\protect\hyperlink{sec-u14}{Unidad 4} introduce los \textbf{efectos
aleatorios} (GLMM) para el agrupamiento por centro ---GLOW tiene
pacientes anidadas en 6 centros y la cohorte simulada, en 24---, donde
las observaciones de un mismo centro se parecen más de lo que el azar
explicaría. Y por último la \protect\hyperlink{sec-u15}{Unidad 5} aborda
la supervivencia desde la perspectiva de los datos binarios.

\hypertarget{sec-u13}{%
\section{3 · Extensión de la respuesta binaria: binomial y
politómica}\label{sec-u13}}

La respuesta binaria de las unidades anteriores admite \textbf{dos
extensiones naturales}, y esta unidad recorre ambas.

La primera es \textbf{agrupar}. Cuando las predictoras son
\textbf{categóricas}, dos sujetos con el mismo perfil de covariables son
\textbf{intercambiables} ---para el modelo, ensayos idénticos---, así
que podemos \textbf{colapsar} cada perfil en una fila que cuente cuántos
casos hay en el grupo. Con respuesta \textbf{binaria}, agrupar convierte
la Bernoulli individual (\(y_i\in\{0,1\}\)) en una \textbf{binomial},
\(Y_i\sim\mathrm{Bin}(m_i,\pi_i)\) con \(m_i\) el tamaño del grupo: es
\textbf{el mismo modelo logístico}, solo cambia la disposición del dato.
Lo vemos en las Secciones \protect\hyperlink{sec-u13-binomial}{3.1} y
\protect\hyperlink{sec-u13-gof-binomial}{3.2}, donde agrupar además
\textbf{recupera los diagnósticos de ajuste}.

La segunda es admitir \textbf{más de dos categorías}. Si la respuesta es
\textbf{politómica} ---categórica con tres o más niveles---, el logit
binario se generaliza al \textbf{logit de categoría base} cuando no hay
orden (\protect\hyperlink{sec-u13-nominal}{Sección 3.3}) y al de
\textbf{odds proporcionales} cuando sí lo hay
(\protect\hyperlink{sec-u13-ordinal}{Sección 3.4}). Aquí las tratamos
sobre datos \textbf{individuales} de GLOW (\texttt{raterisk}).

\begin{tcolorbox}[enhanced jigsaw, coltitle=black, rightrule=.15mm, arc=.35mm, leftrule=.75mm, colback=white, colbacktitle=quarto-callout-tip-color!10!white, left=2mm, breakable, toptitle=1mm, toprule=.15mm, bottomtitle=1mm, titlerule=0mm, title=\textcolor{quarto-callout-tip-color}{\faLightbulb}\hspace{0.5em}{Tip}, bottomrule=.15mm, opacityback=0, opacitybacktitle=0.6, colframe=quarto-callout-tip-color-frame]

Las respuestas politómicas también podrían agruparse en recuentos,
generando una \textbf{multinomial},
\(Y_i\sim\mathrm{Mult}(m_i,\boldsymbol{\pi}_i)\), pero su ajuste se
apoya en la equivalencia con un modelo \textbf{log-lineal de Poisson},
que desarrollamos en el Caso 2, donde se introduce la Poisson, de modo
que lo estudiaremos allí.

\end{tcolorbox}

\begin{tcolorbox}[enhanced jigsaw, coltitle=black, rightrule=.15mm, arc=.35mm, leftrule=.75mm, colback=white, colbacktitle=quarto-callout-note-color!10!white, left=2mm, breakable, toptitle=1mm, toprule=.15mm, bottomtitle=1mm, titlerule=0mm, title=\textcolor{quarto-callout-note-color}{\faInfo}\hspace{0.5em}{Sobre esta unidad}, bottomrule=.15mm, opacityback=0, opacitybacktitle=0.6, colframe=quarto-callout-note-color-frame]

RA implicados (el marco de 1.1--1.2 aplicado a nuevos tipos de
respuesta):

\begin{itemize}
\tightlist
\item
  \textbf{RA2 · Formulación y estimación.} Ajustar el binomial agrupado
  y los modelos politómicos (nominal, ordinal).
\item
  \textbf{RA3 · Interpretación en la escala correcta.} Odds ratios
  relativos (nominal) y acumulados (ordinal).
\item
  \textbf{RA4 · Diagnóstico y bondad de ajuste.} Contraste χ²,
  calibración y sobredispersión.
\item
  \textbf{RA5 · Selección y comparación de modelos.} Nominal frente a
  ordinal (parsimonia).
\end{itemize}

\end{tcolorbox}

\hypertarget{sec-u13-binomial}{%
\subsection{3.1 Del dato individual al agrupado: respuesta
binomial}\label{sec-u13-binomial}}

Para ejemplificar la implementación de esta versión del modelo binomial,
appliquemos una agrupación a GLOW, considerando ahora como predictoras
las variables \texttt{priorfrac}, \texttt{momfrac} y \texttt{raterisk}.
Como las tres predictoras son \textbf{factores}, cada combinación de
niveles reúne a mujeres intercambiables. Al colapsarlas, la respuesta
pasa de un \(0/1\) por sujeto a un \textbf{recuento de fracturas sobre
el total del grupo} ---una binomial---. Lo llamativo es que \textbf{el
ajuste no cambia} respecto a la logística individual; lo que sí cambia
---y es la ganancia de agrupar--- lo veremos en la bondad de ajuste
(\protect\hyperlink{sec-u13-gof-binomial}{Sección 3.2}).

Agrupados los datos en \(k\) celdas (perfiles de covariables), en la
celda \(i\)

\[Y_i \sim \mathrm{Bi}(m_i,\pi_i), \qquad i=1,\dots,k,\quad \textstyle\sum_i m_i=n,\]

donde \(Y_i\) es el número de fracturas entre las \(m_i\) mujeres del
perfil y \(\pi_i\) su probabilidad de fractura.

La parte sistemática es común a todos los perfiles, así que la
escribimos sin índice ---con la respuesta
\((\text{éxitos},\text{fracasos})=(Y,\,m-Y)\)---, y así el modelo
binomial plantea estimar:

\[\operatorname{logit}(\pi)=\theta+\beta\,\text{priorfrac}+\gamma\,\text{momfrac}+\delta_{\text{Same}}\mathbf 1(\text{raterisk}=\text{Same})+\delta_{\text{Greater}}\mathbf 1(\text{raterisk}=\text{Greater}),\]

con \texttt{raterisk\ =\ Less} como referencia.

\begin{tcolorbox}[enhanced jigsaw, coltitle=black, rightrule=.15mm, arc=.35mm, leftrule=.75mm, colback=white, colbacktitle=quarto-callout-tip-color!10!white, left=2mm, breakable, toptitle=1mm, toprule=.15mm, bottomtitle=1mm, titlerule=0mm, title=\textcolor{quarto-callout-tip-color}{\faLightbulb}\hspace{0.5em}{🔧 En R. El \texttt{glm()} con datos binomiales}, bottomrule=.15mm, opacityback=0, opacitybacktitle=0.6, colframe=quarto-callout-tip-color-frame]

El ajuste de datos binomiales es similar al de datos binarios, con la
salvedad de cómo se especifica la respuesta aquí, a partir de una matriz
con dos columnas: éxitos y fracasos:

\begin{verbatim}
glm(cbind(exitos,fracasos) ~ x1+x2+x1:x2+...,family=binomial(link = "logit"),data)
\end{verbatim}

\end{tcolorbox}

Agrupamos entonces GLOW por tres factores (\texttt{priorfrac},
\texttt{momfrac}, \texttt{raterisk}) y contamos fracturas:

\begin{Shaded}
\begin{Highlighting}[]
\NormalTok{glow\_agg }\OtherTok{\textless{}{-}}\NormalTok{ glow }\SpecialCharTok{|\textgreater{}}
  \FunctionTok{count}\NormalTok{(priorfrac, momfrac, raterisk, fracture) }\SpecialCharTok{|\textgreater{}}
\NormalTok{  tidyr}\SpecialCharTok{::}\FunctionTok{pivot\_wider}\NormalTok{(}\AttributeTok{names\_from =}\NormalTok{ fracture, }\AttributeTok{values\_from =}\NormalTok{ n, }\AttributeTok{values\_fill =} \DecValTok{0}\NormalTok{) }\SpecialCharTok{|\textgreater{}}
  \FunctionTok{mutate}\NormalTok{(}\AttributeTok{total =}\NormalTok{ No }\SpecialCharTok{+}\NormalTok{ Yes)            }\CommentTok{\# ensayos por celda; Yes = nº de fracturas}
\NormalTok{glow\_agg}
\end{Highlighting}
\end{Shaded}

\begin{verbatim}
# A tibble: 12 x 6
   priorfrac momfrac raterisk    No   Yes total
   <fct>     <fct>   <fct>    <int> <int> <int>
 1 No        No      Less       109    18   127
 2 No        No      Same       101    23   124
 3 No        No      Greater     60    16    76
 4 No        Yes     Less         9     2    11
 5 No        Yes     Same        11     8    19
 6 No        Yes     Greater     11     6    17
 7 Yes       No      Less        19     8    27
 8 Yes       No      Same        23    14    37
 9 Yes       No      Greater     22    22    44
10 Yes       Yes     Less         2     0     2
11 Yes       Yes     Same         3     3     6
12 Yes       Yes     Greater      5     5    10
\end{verbatim}

Antes de modelar conviene \textbf{ver} los datos agrupados. La
Figura~\ref{fig-u13-barras} muestra, para cada perfil de covariables, la
\textbf{proporción observada de fractura} \((y_i/m_i)\); la cifra sobre
cada barra es el número de mujeres de la celda (\(m_i\)), que recuerda
cuánta información sostiene cada proporción.

\begin{Shaded}
\begin{Highlighting}[]
\NormalTok{glow\_agg }\SpecialCharTok{|\textgreater{}}
  \FunctionTok{mutate}\NormalTok{(}\AttributeTok{prop =}\NormalTok{ Yes }\SpecialCharTok{/}\NormalTok{ total) }\SpecialCharTok{|\textgreater{}}
  \FunctionTok{ggplot}\NormalTok{(}\FunctionTok{aes}\NormalTok{(}\AttributeTok{x =}\NormalTok{ priorfrac, }\AttributeTok{y =}\NormalTok{ prop, }\AttributeTok{fill =}\NormalTok{ momfrac)) }\SpecialCharTok{+}
  \FunctionTok{geom\_col}\NormalTok{(}\AttributeTok{position =} \FunctionTok{position\_dodge}\NormalTok{(}\AttributeTok{width =} \FloatTok{0.8}\NormalTok{), }\AttributeTok{width =} \FloatTok{0.7}\NormalTok{) }\SpecialCharTok{+}
  \FunctionTok{geom\_text}\NormalTok{(}\FunctionTok{aes}\NormalTok{(}\AttributeTok{label =}\NormalTok{ total), }\AttributeTok{position =} \FunctionTok{position\_dodge}\NormalTok{(}\AttributeTok{width =} \FloatTok{0.8}\NormalTok{),}
            \AttributeTok{vjust =} \SpecialCharTok{{-}}\FloatTok{0.3}\NormalTok{, }\AttributeTok{size =} \DecValTok{3}\NormalTok{) }\SpecialCharTok{+}
  \FunctionTok{facet\_wrap}\NormalTok{(}\SpecialCharTok{\textasciitilde{}}\NormalTok{ raterisk, }\AttributeTok{labeller =}\NormalTok{ label\_both) }\SpecialCharTok{+}
  \FunctionTok{labs}\NormalTok{(}\AttributeTok{x =} \StringTok{"Fractura previa (priorfrac)"}\NormalTok{, }\AttributeTok{y =} \StringTok{"Proporción de fractura (Yes/total)"}\NormalTok{,}
       \AttributeTok{fill =} \StringTok{"Antecedente materno}\SpecialCharTok{\textbackslash{}n}\StringTok{(momfrac)"}\NormalTok{) }\SpecialCharTok{+}
  \FunctionTok{ylim}\NormalTok{(}\DecValTok{0}\NormalTok{, }\ConstantTok{NA}\NormalTok{)}
\end{Highlighting}
\end{Shaded}

\begin{figure}[H]

{\centering \includegraphics{caso1_binaria_files/figure-pdf/fig-u13-barras-1.pdf}

}

\caption{\label{fig-u13-barras}Proporción observada de fractura por
perfil de covariables en GLOW. Cada panel es un nivel de riesgo
autopercibido (raterisk); dentro de cada panel, las barras cruzan
fractura previa (priorfrac) y antecedente materno (momfrac). La cifra
sobre la barra es el nº de mujeres de la celda.}

\end{figure}

En la Figura~\ref{fig-u13-barras} se aprecian ya las tendencias que el
modelo recogerá: la proporción de fractura \textbf{sube con la fractura
previa} (\texttt{priorfrac}), \textbf{con el antecedente materno}
(\texttt{momfrac}) y al pasar de riesgo autopercibido \texttt{Less} a
\texttt{Same} y \texttt{Greater} ---el orden esperado---, lo que
respalda incluir las tres covariables. Conviene, eso sí, leer cada barra
junto al \textbf{tamaño de su celda} (la cifra encima): las tres celdas
basales (\texttt{priorfrac\ =\ No}, \texttt{momfrac\ =\ No}) concentran
la mayor parte de los datos (\(n = 127, 124, 76\)), mientras que las que
acumulan ambos antecedentes son \textbf{diminutas} (\(n = 2, 6, 10\)).
Sus proporciones son muy inestables ---la celda \texttt{Less} con ambos
antecedentes y \(n=2\) cae a 0, y su análoga en \texttt{Same} con
\(n=6\) salta a 0,5---: es \textbf{ruido de muestreo} que no conviene
sobreinterpretar y que, además, es justo el que degrada la aproximación
\(\chi^2\) del contraste de bondad de ajuste cuando hay celdas pequeñas
(la cautela de la \protect\hyperlink{sec-u13-gof-binomial}{Sección
3.2}).

Ahora ajustamos el binomial con la sintaxis
\texttt{cbind(éxitos,\ fracasos)} y comprobamos que \textbf{recupera
exactamente} los mismos coeficientes que la logística individual con
esas predictoras ---misma verosimilitud, distinta presentación---:

\begin{Shaded}
\begin{Highlighting}[]
\NormalTok{fit\_bin }\OtherTok{\textless{}{-}} \FunctionTok{glm}\NormalTok{(}\FunctionTok{cbind}\NormalTok{(Yes, No) }\SpecialCharTok{\textasciitilde{}}\NormalTok{ priorfrac }\SpecialCharTok{+}\NormalTok{ momfrac }\SpecialCharTok{+}\NormalTok{ raterisk,}
               \AttributeTok{family =}\NormalTok{ binomial, }\AttributeTok{data =}\NormalTok{ glow\_agg)}
\NormalTok{fit\_ind }\OtherTok{\textless{}{-}} \FunctionTok{glm}\NormalTok{(fracture }\SpecialCharTok{\textasciitilde{}}\NormalTok{ priorfrac }\SpecialCharTok{+}\NormalTok{ momfrac }\SpecialCharTok{+}\NormalTok{ raterisk,}
               \AttributeTok{family =}\NormalTok{ binomial, }\AttributeTok{data =}\NormalTok{ glow)}

\FunctionTok{cbind}\NormalTok{(}\AttributeTok{agrupado =} \FunctionTok{coef}\NormalTok{(fit\_bin), }\AttributeTok{individual =} \FunctionTok{coef}\NormalTok{(fit\_ind))   }\CommentTok{\# idénticos}
\end{Highlighting}
\end{Shaded}

\begin{verbatim}
                  agrupado individual
(Intercept)     -1.8711967 -1.8711967
priorfracYes     0.9797192  0.9797192
momfracYes       0.5750477  0.5750477
rateriskSame     0.4647529  0.4647529
rateriskGreater  0.6717294  0.6717294
\end{verbatim}

La interpretación de los coeficientes (odds ratios,
\protect\hyperlink{sec-u12-interpretacion}{Sección 2.3}) y la selección
de modelos (LRT/AIC/BIC, \protect\hyperlink{sec-u12-inferencia}{Sección
2.5}) son \textbf{idénticas} a las del caso individual. Lo que cambia
---y es la ganancia de agrupar--- es la bondad de ajuste.

\hypertarget{sec-u13-gof-binomial}{%
\subsection{3.2 Bondad de ajuste y diagnóstico en datos
binomiales}\label{sec-u13-gof-binomial}}

En la \protect\hyperlink{sec-u12-diagnostico}{Sección 2.6} vimos que,
con datos \textbf{individuales}, ni la deviance ni el estadístico de
Pearson sirven como test absoluto de ajuste: el saturado crece con \(n\)
y la asintótica \(\chi^2\) no se sostiene. Con datos \textbf{agrupados}
(y \(m_i\) suficientemente grandes) esa asintótica \textbf{se recupera}
(Agresti, 2015; Faraway, 2016). También cobran relevancia ciertos
residuos comunes en los GLM que no funcionaban en el caso binario.
Presentamos en primer lugar estos residuos y construimos, a partir de
ellos y con su \textbf{suma de cuadrados}, los estadísticos globales de
bondad de ajuste.

\hypertarget{residuos-1}{%
\subsubsection{Residuos}\label{residuos-1}}

Al agrupar los datos, cada celda reúne \(m_i\) ensayos, de modo que los
residuos \textbf{estándar} ya son informativos ---aproximadamente
\(N(0,1)\) con \(m_i\) grande (McCullagh \& Nelder, 1989)---. Estos
residuos, genéricos en los GLM, son:

\begin{itemize}
\item
  \textbf{Pearson:}, la diferencia entre lo observado y lo esperado en
  la celda, tipificada por su desviación estándar,
  \[r_{P,i}=\dfrac{y_i-m_i\hat\pi_i}{\sqrt{m_i\hat\pi_i(1-\hat\pi_i)}}\]
\item
  \textbf{Deviance:}, la raíz con signo de la contribución \(d_i\) de
  cada celda a la deviance,
  \[r_{D,i}=\operatorname{sign}(y_i-m_i\hat\pi_i)\sqrt{d_i}\]
\item
  Para compararlos en igualdad de condiciones se usan sus versiones
  \textbf{estandarizadas} (o \emph{estudentizadas}), que dividen por
  \(\sqrt{1-h_{ii}}\), \[
  r^{\,s}_{P,i}=\frac{r_{P,i}}{\sqrt{1-h_{ii}}}, \qquad r^{\,s}_{D,i}=\frac{r_{D,i}}{\sqrt{1-h_{ii}}}
  \] donde \(h_{ii}\) es el \textbf{leverage} (apalancamiento) de la
  celda ---su influencia \emph{potencial}, dada por el elemento diagonal
  de la matriz gorro \(\hat H\), con \(\sum_i h_{ii}=p\)---: un
  \(h_{ii}\) alto marca una celda con un perfil de covariables extremo,
  esto es, sujetos con características alejadas de la mayoría de los
  datos.
\end{itemize}

\begin{tcolorbox}[enhanced jigsaw, coltitle=black, rightrule=.15mm, arc=.35mm, leftrule=.75mm, colback=white, colbacktitle=quarto-callout-note-color!10!white, left=2mm, breakable, toptitle=1mm, toprule=.15mm, bottomtitle=1mm, titlerule=0mm, title=\textcolor{quarto-callout-note-color}{\faInfo}\hspace{0.5em}{La matriz sombrero en un GLM}, bottomrule=.15mm, opacityback=0, opacitybacktitle=0.6, colframe=quarto-callout-note-color-frame]

En regresión lineal \(\hat H = X(X^{\top}X)^{-1}X^{\top}\) proyecta
\(\mathbf y\) sobre los valores ajustados. En un GLM el ajuste es por
mínimos cuadrados \textbf{ponderados} (IWLS), así que la matriz sombrero
incorpora los pesos de la última iteración
\(W=\operatorname{diag}(w_i)\): \[
\hat H = W^{1/2}X\,(X^{\top}WX)^{-1}X^{\top}W^{1/2}.
\] Es \textbf{simétrica e idempotente}, sus diagonales son los
\textbf{leverages} \(h_{ii}\) y suman el número de parámetros,
\(\sum_i h_{ii}=\operatorname{tr}(\hat H)=p\). En el binomial/logístico
los pesos de trabajo son \(w_i=m_i\hat\pi_i(1-\hat\pi_i)\), de modo que
\(h_{ii}\) combina lo extremo del perfil de covariables con su varianza.
Es el \(h_{ii}\) que aparece en
\(r^{\,s}_{P,i}=r_{P,i}/\sqrt{1-h_{ii}}\) y en la distancia de Cook.

\end{tcolorbox}

\begin{itemize}
\tightlist
\item
  La \textbf{distancia de Cook} combina residuo y leverage, y mide
  \textbf{cuánto cambiarían los coeficientes} si se eliminara esa celda;
  valores grandes señalan \textbf{celdas influyentes} que conviene
  revisar. \[
  D_i=\frac{(r^{\,s}_{P,i})^2}{p}\,\frac{h_{ii}}{1-h_{ii}}
  \] Calculamos y representamos los residuos estandarizados y las
  distancias de Cook para GLOW en la Figura~\ref{fig-u13-residuos}.
\end{itemize}

\begin{Shaded}
\begin{Highlighting}[]
\CommentTok{\#library(patchwork)}

\NormalTok{diag\_bin }\OtherTok{\textless{}{-}}\NormalTok{ glow\_agg }\SpecialCharTok{|\textgreater{}}
  \FunctionTok{mutate}\NormalTok{(}\AttributeTok{pred  =} \FunctionTok{fitted}\NormalTok{(fit\_bin),}
         \AttributeTok{r\_std =} \FunctionTok{rstandard}\NormalTok{(fit\_bin, }\AttributeTok{type =} \StringTok{"pearson"}\NormalTok{),}
         \AttributeTok{cook  =} \FunctionTok{cooks.distance}\NormalTok{(fit\_bin),}
         \AttributeTok{cell  =} \FunctionTok{row\_number}\NormalTok{())}

\NormalTok{p\_res }\OtherTok{\textless{}{-}} \FunctionTok{ggplot}\NormalTok{(diag\_bin, }\FunctionTok{aes}\NormalTok{(pred, r\_std, }\AttributeTok{size =}\NormalTok{ total)) }\SpecialCharTok{+}
  \FunctionTok{geom\_hline}\NormalTok{(}\AttributeTok{yintercept =} \FunctionTok{c}\NormalTok{(}\SpecialCharTok{{-}}\DecValTok{2}\NormalTok{, }\DecValTok{0}\NormalTok{, }\DecValTok{2}\NormalTok{), }\AttributeTok{linetype =} \FunctionTok{c}\NormalTok{(}\StringTok{"dotted"}\NormalTok{, }\StringTok{"dashed"}\NormalTok{, }\StringTok{"dotted"}\NormalTok{)) }\SpecialCharTok{+}
  \FunctionTok{geom\_point}\NormalTok{(}\AttributeTok{alpha =} \FloatTok{0.7}\NormalTok{) }\SpecialCharTok{+}
  \FunctionTok{labs}\NormalTok{(}\AttributeTok{x =} \StringTok{"Probabilidad ajustada"}\NormalTok{, }\AttributeTok{y =} \StringTok{"Residuo de Pearson estandarizado"}\NormalTok{,}
       \AttributeTok{size =} \StringTok{"n por celda"}\NormalTok{)}

\NormalTok{p\_cook }\OtherTok{\textless{}{-}} \FunctionTok{ggplot}\NormalTok{(diag\_bin, }\FunctionTok{aes}\NormalTok{(cell, cook, }\AttributeTok{size =}\NormalTok{ total)) }\SpecialCharTok{+}
  \FunctionTok{geom\_hline}\NormalTok{(}\AttributeTok{yintercept =} \DecValTok{4} \SpecialCharTok{/} \FunctionTok{nrow}\NormalTok{(diag\_bin), }\AttributeTok{linetype =} \StringTok{"dotted"}\NormalTok{) }\SpecialCharTok{+}
  \FunctionTok{geom\_segment}\NormalTok{(}\FunctionTok{aes}\NormalTok{(}\AttributeTok{xend =}\NormalTok{ cell, }\AttributeTok{yend =} \DecValTok{0}\NormalTok{), }\AttributeTok{linewidth =} \FloatTok{0.3}\NormalTok{) }\SpecialCharTok{+}
  \FunctionTok{geom\_point}\NormalTok{(}\AttributeTok{alpha =} \FloatTok{0.7}\NormalTok{) }\SpecialCharTok{+}
  \FunctionTok{scale\_x\_continuous}\NormalTok{(}\AttributeTok{breaks =}\NormalTok{ diag\_bin}\SpecialCharTok{$}\NormalTok{cell) }\SpecialCharTok{+}
  \FunctionTok{labs}\NormalTok{(}\AttributeTok{x =} \StringTok{"Celda"}\NormalTok{, }\AttributeTok{y =} \StringTok{"Distancia de Cook"}\NormalTok{, }\AttributeTok{size =} \StringTok{"n por celda"}\NormalTok{)}

\NormalTok{p\_res }\SpecialCharTok{+}\NormalTok{ p\_cook }\SpecialCharTok{+} \FunctionTok{plot\_layout}\NormalTok{(}\AttributeTok{guides =} \StringTok{"collect"}\NormalTok{)}
\end{Highlighting}
\end{Shaded}

\begin{figure}[H]

{\centering \includegraphics{caso1_binaria_files/figure-pdf/fig-u13-residuos-1.pdf}

}

\caption{\label{fig-u13-residuos}Diagnósticos por celda (el tamaño del
punto es el nº de ensayos). Izquierda: residuos de Pearson
estandarizados frente a la probabilidad ajustada; con un buen ajuste se
reparten en torno a 0 dentro de la banda de ±2. Derecha: distancia de
Cook por celda, con línea de referencia orientativa en 4/n.}

\end{figure}

En GLOW, vemos, a la izquierda en la Figura~\ref{fig-u13-residuos}, que
los \textbf{12 perfiles} (celdas) dejan todos los residuos
estandarizados (de Pearson) dentro de la banda de ±2, y se reparten sin
patrón alrededor de 0 (rango ≈ −1,3 a +1,4); las celdas con más ensayos
---las de menor riesgo--- caen muy cerca de 0. No hay, pues,
\textbf{celdas mal ajustadas ni influyentes}, en línea con el χ² de
bondad de ajuste y la calibración. En el panel derecho, las distancias
de Cook son pequeñas en términos absolutos ---la mayor (celda 9) ronda
0,33 y apenas roza la referencia orientativa de \(4/n\)---, muy por
debajo del umbral clásico \(D\gtrsim 1\). Las celdas con más ensayos
(1--3) presentan un Cook moderado por su peso en el ajuste, no por
desajuste. Junto con la ausencia de residuos fuera de la banda de ±2, se
confirma que ninguna celda resulta influyente.

\begin{tcolorbox}[enhanced jigsaw, coltitle=black, rightrule=.15mm, arc=.35mm, leftrule=.75mm, colback=white, colbacktitle=quarto-callout-tip-color!10!white, left=2mm, breakable, toptitle=1mm, toprule=.15mm, bottomtitle=1mm, titlerule=0mm, title=\textcolor{quarto-callout-tip-color}{\faLightbulb}\hspace{0.5em}{🔧 En R. Residuos e influencia}, bottomrule=.15mm, opacityback=0, opacitybacktitle=0.6, colframe=quarto-callout-tip-color-frame]

\begin{Shaded}
\begin{Highlighting}[]
\FunctionTok{residuals}\NormalTok{(fit\_bin, }\AttributeTok{type =} \StringTok{"deviance"}\NormalTok{)     }\CommentTok{\# o "deviance", "response"}
\FunctionTok{rstandard}\NormalTok{(fit\_bin, }\AttributeTok{type =} \StringTok{"pearson"}\NormalTok{)     }\CommentTok{\# estandarizados (ajustan por el leverage)}
\FunctionTok{hatvalues}\NormalTok{(fit\_bin)                        }\CommentTok{\# leverage h\_ii de cada celda}
\FunctionTok{cooks.distance}\NormalTok{(fit\_bin)                   }\CommentTok{\# influencia de cada celda}
\end{Highlighting}
\end{Shaded}

\end{tcolorbox}

\hypertarget{bondad-de-ajuste}{%
\subsubsection{Bondad de ajuste}\label{bondad-de-ajuste}}

Sumando los cuadrados de los residuos deviance y Pearson se obtienen los
\textbf{dos estadísticos globales}, que comparan el modelo
\textbf{ajustado con el saturado} ---el mejor ajuste posible---:

\begin{itemize}
\tightlist
\item
  \textbf{Pearson:} \(X^2=\sum_i r_{P,i}^2\), la suma directa de los
  residuos de Pearson al cuadrado.
\item
  \textbf{Deviance:}
  \(D=\sum_i r_{D,i}^2=2(\ell_{\text{sat}}-\ell_{\text{ajustado}})\),
  que es además el \textbf{LRT frente al saturado}.
\end{itemize}

Bajo un modelo correcto y con celdas no demasiado pequeñas, ambos se
distribuyen \(\chi^2_{n-p}\) (Agresti, 2015; Faraway, 2016), con lo que
\textbf{ya podemos contrastar el ajuste de forma absoluta} contra el
modelo saturado y responder a la pregunta \textbf{¿ajusta bien nuestro
modelo?}.

\begin{Shaded}
\begin{Highlighting}[]
\NormalTok{gl }\OtherTok{\textless{}{-}} \FunctionTok{df.residual}\NormalTok{(fit\_bin)                       }\CommentTok{\# n {-} p = nº de celdas {-} nº de parámetros}
\FunctionTok{c}\NormalTok{(}\AttributeTok{deviance =} \FunctionTok{deviance}\NormalTok{(fit\_bin),}
  \AttributeTok{pearson  =} \FunctionTok{sum}\NormalTok{(}\FunctionTok{residuals}\NormalTok{(fit\_bin, }\AttributeTok{type =} \StringTok{"pearson"}\NormalTok{)}\SpecialCharTok{\^{}}\DecValTok{2}\NormalTok{),}
  \AttributeTok{gl       =}\NormalTok{ gl)}
\end{Highlighting}
\end{Shaded}

\begin{verbatim}
deviance  pearson       gl 
4.732757 4.072651 7.000000 
\end{verbatim}

\begin{Shaded}
\begin{Highlighting}[]
\CommentTok{\# p{-}valores del contraste de bondad de ajuste (H0: el modelo ajusta)}
\FunctionTok{c}\NormalTok{(}\AttributeTok{p\_deviance =} \FunctionTok{pchisq}\NormalTok{(}\FunctionTok{deviance}\NormalTok{(fit\_bin), gl, }\AttributeTok{lower.tail =} \ConstantTok{FALSE}\NormalTok{),}
  \AttributeTok{p\_pearson  =} \FunctionTok{pchisq}\NormalTok{(}\FunctionTok{sum}\NormalTok{(}\FunctionTok{residuals}\NormalTok{(fit\_bin, }\StringTok{"pearson"}\NormalTok{)}\SpecialCharTok{\^{}}\DecValTok{2}\NormalTok{), gl, }\AttributeTok{lower.tail =} \ConstantTok{FALSE}\NormalTok{))}
\end{Highlighting}
\end{Shaded}

\begin{verbatim}
p_deviance  p_pearson 
 0.6925382  0.7713725 
\end{verbatim}

Al conseguir un p-valor \textbf{alto} (muy superior a la
significatividad 0.05) no evidenciamos falta de ajuste (no podemos
rechazar el modelo ajustado frente al saturado), y así concluimos que el
modelo reproduce razonablemente las proporciones por celda. Es
exactamente el contraste que en la
\protect\hyperlink{sec-u12-diagnostico}{Sección 2.6} no podíamos usar y
que ahora, al agrupar, recupera su sentido.

\begin{tcolorbox}[enhanced jigsaw, coltitle=black, rightrule=.15mm, arc=.35mm, leftrule=.75mm, colback=white, colbacktitle=quarto-callout-important-color!10!white, left=2mm, breakable, toptitle=1mm, toprule=.15mm, bottomtitle=1mm, titlerule=0mm, title=\textcolor{quarto-callout-important-color}{\faExclamation}\hspace{0.5em}{Importante}, bottomrule=.15mm, opacityback=0, opacitybacktitle=0.6, colframe=quarto-callout-important-color-frame]

\textbf{Cautela:} la aproximación exige celdas no demasiado pequeñas;
con muchos \(m_i\) pequeños vuelve a degradarse.

\end{tcolorbox}

\begin{tcolorbox}[enhanced jigsaw, coltitle=black, rightrule=.15mm, arc=.35mm, leftrule=.75mm, colback=white, colbacktitle=quarto-callout-tip-color!10!white, left=2mm, breakable, toptitle=1mm, toprule=.15mm, bottomtitle=1mm, titlerule=0mm, title=\textcolor{quarto-callout-tip-color}{\faLightbulb}\hspace{0.5em}{🔧 En R. El contraste χ² de bondad de ajuste (deviance y Pearson)}, bottomrule=.15mm, opacityback=0, opacitybacktitle=0.6, colframe=quarto-callout-tip-color-frame]

Con el modelo agrupado ajustado, los grados de libertad son
\textbf{celdas menos parámetros} (\texttt{df.residual}) y cada
estadístico se compara con una \(\chi^2_{n-p}\) mediante
\texttt{pchisq(...,\ lower.tail\ =\ FALSE)}:

\begin{Shaded}
\begin{Highlighting}[]
\NormalTok{gl }\OtherTok{\textless{}{-}} \FunctionTok{df.residual}\NormalTok{(fit\_bin)                       }\CommentTok{\# n {-} p = nº de celdas {-} nº de parámetros}

\CommentTok{\# Deviance}
\FunctionTok{pchisq}\NormalTok{(}\FunctionTok{deviance}\NormalTok{(fit\_bin), gl, }\AttributeTok{lower.tail =} \ConstantTok{FALSE}\NormalTok{)}

\CommentTok{\# Pearson}
\NormalTok{X2 }\OtherTok{\textless{}{-}} \FunctionTok{sum}\NormalTok{(}\FunctionTok{residuals}\NormalTok{(fit\_bin, }\AttributeTok{type =} \StringTok{"pearson"}\NormalTok{)}\SpecialCharTok{\^{}}\DecValTok{2}\NormalTok{)}
\FunctionTok{pchisq}\NormalTok{(X2, gl, }\AttributeTok{lower.tail =} \ConstantTok{FALSE}\NormalTok{)}
\end{Highlighting}
\end{Shaded}

Un \(p\) \textbf{alto} = no hay evidencia de falta de ajuste. ¿Cuál
usar?

\begin{itemize}
\tightlist
\item
  \textbf{Pearson (\(X^2\))} para el \textbf{test absoluto} de bondad de
  ajuste: aproxima mejor la \(\chi^2\) con recuentos moderados y es,
  además, la base de la \textbf{sobredispersión}
  (\(\hat\phi = X^2/\text{gl}\)).
\item
  \textbf{Deviance} para \textbf{comparar modelos anidados} (su
  diferencia es el LRT) y como medida ligada a la verosimilitud; como
  test \textbf{absoluto} es algo menos fiable que Pearson cuando los
  recuentos son pequeños.
\item
  \textbf{Ambos} exigen celdas \textbf{no pequeñas} (regla práctica:
  frecuencias esperadas \(\gtrsim 5\)); con muchos \(m_i\) pequeños la
  aproximación \(\chi^2\) se degrada ---y la deviance, antes---.
\end{itemize}

\end{tcolorbox}

\hypertarget{calibraciuxf3n}{%
\subsubsection{Calibración}\label{calibraciuxf3n}}

Al agrupar, la \textbf{calibración} es inmediata ---algo que el caso
individual no permitía sin recurrir a deciles---: cada celda aporta su
\textbf{probabilidad predicha} y su \textbf{proporción observada}, y
basta enfrentarlas. Si el modelo está bien calibrado, los puntos se
alinean sobre la diagonal:

\begin{Shaded}
\begin{Highlighting}[]
\CommentTok{\# Calibración DIRECTA: cada celda = (probabilidad predicha, proporción observada)}
\NormalTok{calib\_bin }\OtherTok{\textless{}{-}}\NormalTok{ glow\_agg }\SpecialCharTok{|\textgreater{}}
  \FunctionTok{mutate}\NormalTok{(}\AttributeTok{pred =} \FunctionTok{fitted}\NormalTok{(fit\_bin), }\AttributeTok{obs =}\NormalTok{ Yes }\SpecialCharTok{/}\NormalTok{ total)}

\FunctionTok{ggplot}\NormalTok{(calib\_bin, }\FunctionTok{aes}\NormalTok{(pred, obs, }\AttributeTok{size =}\NormalTok{ total)) }\SpecialCharTok{+}
  \FunctionTok{geom\_abline}\NormalTok{(}\AttributeTok{linetype =} \StringTok{"dashed"}\NormalTok{) }\SpecialCharTok{+}                 \CommentTok{\# calibración perfecta}
  \FunctionTok{geom\_point}\NormalTok{(}\AttributeTok{alpha =} \FloatTok{0.7}\NormalTok{) }\SpecialCharTok{+}
  \FunctionTok{labs}\NormalTok{(}\AttributeTok{x =} \StringTok{"Probabilidad predicha"}\NormalTok{, }\AttributeTok{y =} \StringTok{"Proporción observada (Yes/total)"}\NormalTok{,}
       \AttributeTok{size =} \StringTok{"n por celda"}\NormalTok{)}
\end{Highlighting}
\end{Shaded}

\begin{figure}[H]

{\centering \includegraphics{caso1_binaria_files/figure-pdf/fig-u13-calibracion-1.pdf}

}

\caption{\label{fig-u13-calibracion}Curva de calibración por celda en
GLOW: probabilidad predicha frente a proporción observada de fractura
(el tamaño del punto es el nº de mujeres de la celda). La diagonal
discontinua marca la calibración perfecta; los puntos próximos a ella,
sin sesgo sistemático, indican un modelo bien calibrado.}

\end{figure}

Como se aprecia en Figura~\ref{fig-u13-calibracion}, en GLOW los 12
puntos caen cerca de la diagonal, sin sesgo sistemático por tramos de
riesgo: el modelo está \textbf{bien calibrado}, en línea con el χ² de
bondad de ajuste y los residuos.

\begin{tcolorbox}[enhanced jigsaw, coltitle=black, rightrule=.15mm, arc=.35mm, leftrule=.75mm, colback=white, colbacktitle=quarto-callout-tip-color!10!white, left=2mm, breakable, toptitle=1mm, toprule=.15mm, bottomtitle=1mm, titlerule=0mm, title=\textcolor{quarto-callout-tip-color}{\faLightbulb}\hspace{0.5em}{🔧 En R. El data frame de observados vs predichos para la calibración}, bottomrule=.15mm, opacityback=0, opacitybacktitle=0.6, colframe=quarto-callout-tip-color-frame]

La calibración solo necesita, por celda, la probabilidad
\textbf{predicha} y la proporción \textbf{observada}. La forma directa
es añadirlas con \texttt{mutate()}:

\begin{Shaded}
\begin{Highlighting}[]
\NormalTok{calib\_bin }\OtherTok{\textless{}{-}}\NormalTok{ glow\_agg }\SpecialCharTok{|\textgreater{}}
  \FunctionTok{mutate}\NormalTok{(}\AttributeTok{pred =} \FunctionTok{fitted}\NormalTok{(fit\_bin),     }\CommentTok{\# probabilidad predicha por el modelo}
         \AttributeTok{obs  =}\NormalTok{ Yes }\SpecialCharTok{/}\NormalTok{ total)         }\CommentTok{\# proporción observada en la celda}
\end{Highlighting}
\end{Shaded}

Con \texttt{broom} se obtiene lo mismo de forma compacta
---\texttt{type.predict\ =\ "response"} devuelve la predicción en escala
de probabilidad en la columna \texttt{.fitted}---:

\begin{Shaded}
\begin{Highlighting}[]
\FunctionTok{library}\NormalTok{(broom)}
\NormalTok{calib\_bin }\OtherTok{\textless{}{-}} \FunctionTok{augment}\NormalTok{(fit\_bin, }\AttributeTok{data =}\NormalTok{ glow\_agg, }\AttributeTok{type.predict =} \StringTok{"response"}\NormalTok{) }\SpecialCharTok{|\textgreater{}}
  \FunctionTok{mutate}\NormalTok{(}\AttributeTok{obs =}\NormalTok{ Yes }\SpecialCharTok{/}\NormalTok{ total)}
\end{Highlighting}
\end{Shaded}

Sobre ese data frame, \texttt{geom\_point(aes(pred,\ obs))} dibuja la
calibración por celda y \texttt{geom\_abline()} la diagonal de
referencia; un \texttt{geom\_smooth(aes(pred,\ obs))} traza además la
\textbf{curva de calibración} suavizada.

\end{tcolorbox}

\hypertarget{sobredispersiuxf3n}{%
\subsubsection{Sobredispersión}\label{sobredispersiuxf3n}}

El modelo binomial \(y_i \sim Bi(m_i,\pi_i)\), impone que la varianza de
cada celda sea exactamente \(m_i\hat\pi_i(1-\hat\pi_i)=Var(y_i)\).
Cuando la varianza real en los datos \textbf{excede} ese valor hay
\textbf{sobredispersión}, y al agrupar pasa a ser \textbf{estimable}
mediante el cociente de Pearson:

\[
\hat\phi=\frac{X^2}{n-p}=\frac{1}{n-p}\sum_i r_{P,i}^2 .
\]

Un \(\hat\phi\approx 1\) indica que el binomial basta; \(\hat\phi>1\)
indica sobredispersión (variación \emph{extra-binomial}); \(\hat\phi<1\)
indica infradispersión (es menos frecuente).

La sobredispersión puede ser debida a cierta \textbf{correlación entre
los ensayos de una misma celda} (sujetos agrupados que no son del todo
independientes), a \textbf{heterogeneidad no modelada} o por una
\textbf{covariable omitida}. Lo que importa es su efecto: los
coeficientes \(\hat\beta\) siguen siendo \textbf{consistentes}, pero los
\textbf{errores estándar se subestiman}, de modo que los contrastes de
Wald y LRT resultan \textbf{demasiado liberales}, y los intervalos
\textbf{demasiado estrechos}.

El remedio inmediato es el modelo \textbf{quasi-binomial}: no cambia las
estimaciones puntuales, pero \textbf{estima \(\phi\) de los datos} e
\textbf{infla los errores estándar} por \(\sqrt{\hat\phi}\) (y sustituye
los contrastes \(\chi^2\) por contrastes \textbf{F}). Cuando la
sobredispersión procede de un efecto de agrupamiento identificable, la
alternativa estructural pasa por incluir \textbf{efectos aleatorios}
(\protect\hyperlink{sec-u14}{Unidad 4}).

\begin{Shaded}
\begin{Highlighting}[]
\CommentTok{\# Índice de sobredispersión (cociente de Pearson)}
\FunctionTok{sum}\NormalTok{(}\FunctionTok{residuals}\NormalTok{(fit\_bin, }\AttributeTok{type =} \StringTok{"pearson"}\NormalTok{)}\SpecialCharTok{\^{}}\DecValTok{2}\NormalTok{) }\SpecialCharTok{/} \FunctionTok{df.residual}\NormalTok{(fit\_bin)}
\end{Highlighting}
\end{Shaded}

\begin{verbatim}
[1] 0.5818072
\end{verbatim}

En GLOW, \(\hat\phi\approx 0{,}58\): \textbf{por debajo de 1}, así que
\textbf{no hay rastro de sobredispersión} ---el binomial es de sobra
suficiente--- y el quasi-binomial no haría falta. Lo mostramos
igualmente como herramienta para cuando \(\hat\phi\) supere con claridad
la unidad.

\begin{tcolorbox}[enhanced jigsaw, coltitle=black, rightrule=.15mm, arc=.35mm, leftrule=.75mm, colback=white, colbacktitle=quarto-callout-tip-color!10!white, left=2mm, breakable, toptitle=1mm, toprule=.15mm, bottomtitle=1mm, titlerule=0mm, title=\textcolor{quarto-callout-tip-color}{\faLightbulb}\hspace{0.5em}{🔧 En R. Sobredispersión y quasi-binomial}, bottomrule=.15mm, opacityback=0, opacitybacktitle=0.6, colframe=quarto-callout-tip-color-frame]

\begin{Shaded}
\begin{Highlighting}[]
\CommentTok{\# Índice de sobredispersión (cociente de Pearson): \textasciitilde{}1 si basta el binomial}
\FunctionTok{sum}\NormalTok{(}\FunctionTok{residuals}\NormalTok{(fit\_bin, }\AttributeTok{type =} \StringTok{"pearson"}\NormalTok{)}\SpecialCharTok{\^{}}\DecValTok{2}\NormalTok{) }\SpecialCharTok{/} \FunctionTok{df.residual}\NormalTok{(fit\_bin)}

\CommentTok{\# Ajuste robusto a sobredispersión: mismas estimaciones, SE inflados por sqrt(phi)}
\NormalTok{fit\_q }\OtherTok{\textless{}{-}} \FunctionTok{glm}\NormalTok{(}\FunctionTok{cbind}\NormalTok{(Yes, No) }\SpecialCharTok{\textasciitilde{}}\NormalTok{ priorfrac }\SpecialCharTok{+}\NormalTok{ momfrac }\SpecialCharTok{+}\NormalTok{ raterisk,}
             \AttributeTok{family =}\NormalTok{ quasibinomial, }\AttributeTok{data =}\NormalTok{ glow\_agg)}
\FunctionTok{summary}\NormalTok{(fit\_q)             }\CommentTok{\# \textquotesingle{}Dispersion parameter\textquotesingle{} = phi estimado; contrastes tipo t/F}
\FunctionTok{drop1}\NormalTok{(fit\_q, }\AttributeTok{test =} \StringTok{"F"}\NormalTok{)   }\CommentTok{\# comparación de modelos con test F (no chi\^{}2)}
\end{Highlighting}
\end{Shaded}

\end{tcolorbox}

\begin{tcolorbox}[enhanced jigsaw, coltitle=black, rightrule=.15mm, arc=.35mm, leftrule=.75mm, colback=white, colbacktitle=quarto-callout-note-color!10!white, left=2mm, breakable, toptitle=1mm, toprule=.15mm, bottomtitle=1mm, titlerule=0mm, title=\textcolor{quarto-callout-note-color}{\faInfo}\hspace{0.5em}{¿Y tratar la sobredispersión binomial con quasi-Poisson?}, bottomrule=.15mm, opacityback=0, opacitybacktitle=0.6, colframe=quarto-callout-note-color-frame]

Hay una excepción a tratar la sobredispersión en el modelo binomial con
un \textbf{quasi-Poisson}, y es el régimen de \textbf{sucesos raros}
(\(m_i\) grande, \(\pi_i\) pequeña y \(m_i\pi_i\) moderado), donde la
binomial se aproxima a una Poisson. Ahí sí tiene sentido modelar el
\textbf{conteo} con una Poisson de enlace log y un \emph{offset}
\(\log(m_i)\) ---es decir, modelar una \textbf{tasa}---, y un
quasi-Poisson o una binomial negativa absorberían la variación extra.
Pero eso es porque has entrado \textbf{deliberadamente} en el dominio
Poisson, no porque el quasi-Poisson arregle una binomial. Estos modelos
los veremos en el Caso 2.

\end{tcolorbox}

En \textbf{todo lo demás}, el binomial agrupado se ajusta e interpreta
\textbf{exactamente igual} que el logístico individual de la
\protect\hyperlink{sec-u12}{Unidad 2}: misma estimación (máxima
verosimilitud / IWLS), misma inferencia (Wald, LRT, IC por perfil) y
misma selección de modelos (deviance, AIC/BIC, la escalera de la
\protect\hyperlink{sec-u12-inferencia}{Sección 2.5}). Y la
\textbf{lectura de los resultados no cambia}: los coeficientes siguen
siendo log-odds, \(e^{\beta}\) es un odds ratio y las predicciones se
interpretan en probabilidad igual que en la
\protect\hyperlink{sec-u12-interpretacion}{Sección 2.3}. La única
diferencia es la que acabamos de explotar: al agrupar, recuperamos los
\textbf{diagnósticos de ajuste} (el \(\chi^2\) de bondad de ajuste, los
residuos, la calibración y la sobredispersión por celda) que en el caso
individual no eran válidos.

\hypertarget{sec-u13-nominal}{%
\subsection{3.3 Más de dos categorías sin orden: politómica
nominal}\label{sec-u13-nominal}}

Cuando seguimos teniendo respuestas individualizadas y la respuesta
tiene \textbf{más de dos categorías sin orden}, el logit binario se
generaliza al \textbf{logit de categoría base}: se fija una categoría de
referencia y se modela, frente a ella, el log-odds de cada una de las
demás.

Para ilustrar este modelo con GLOW, usamos \texttt{raterisk} ---el
riesgo de fractura \textbf{autopercibido} como variable respuesta \(Y\),
observada en tres niveles
\texttt{Less}/\texttt{Same}/\texttt{Greater}---, y usamos la categoría
\texttt{Same} como referencia.

\[
 \log\frac{\Pr(Y=k)}{\Pr(Y=\text{Same})}=\eta_{k}= \theta_{k} + \beta_{k} \cdot \text{age} + \gamma_{k} \cdot \mathbf{1}(\text{priorfrac}=\text{Yes}), \qquad k\in\{\text{Less},\text{Greater}\}
 \]

Conviene leerla como \textbf{dos rectas, una por cada categoría
no-referencia}. Fijada \texttt{Same} como base, a cada nivel
\(k\in\{\text{Less},\text{Greater}\}\) le asignamos su \textbf{propio
predictor lineal} \(\eta_{k}\), con \textbf{su propia} interceptación y
\textbf{sus propias} pendientes: una recta describe cómo varían, con la
edad y la fractura previa, los log-odds de percibirse en \textbf{menor}
riesgo frente a «igual» (\texttt{Less} vs \texttt{Same}); la otra, los
de percibirse en \textbf{mayor} riesgo frente a «igual»
(\texttt{Greater} vs \texttt{Same}). La categoría base no necesita recta
propia: hace de \textbf{vara de medir}, con \(\eta_{\text{Same}}=0\) por
convención.

Ajustar el modelo nominal equivale, pues, a estimar \textbf{a la vez}
\(K−1=2\) regresiones ---\texttt{Less\textbar{}Same} y
\texttt{Greater\textbar{}Same}---, cada una con su intercepto
\(\theta_k+\gamma_k\mathbf{1}(\text{priorfrac}=\text{Yes})\) y su
pendiente \(\beta_k\). Que cada categoría tenga \textbf{su propia recta}
es lo que da al modelo nominal su flexibilidad ---cada nivel puede
responder a la edad o a la fractura previa de manera distinta, incluso
en \textbf{direcciones opuestas}---, frente al ordinal de la
\protect\hyperlink{sec-u13-ordinal}{Sección 3.4}, que impondrá una
\textbf{pendiente común} a todos los cortes.

Al exponenciar, \(e^{\beta_k}\) es un \textbf{odds ratio relativo}: el
factor por el que se multiplican las \textbf{odds de pertenecer a la
categoría \(k\) en lugar de a \texttt{Same}} cuando \texttt{age} aumenta
una unidad, manteniendo el resto constante. Por ejemplo,
\(e^{\beta_{\text{Greater}}}\) es cuánto cambia la ventaja de percibirse
en \textbf{mayor} riesgo (\texttt{Greater}) frente a \textbf{igual}
(\texttt{Same}) por cada año más de edad. Cuando se estiman los
coeficientes del modelo, su \textbf{signo} ya orienta: \(\beta_{k}>0\)
aumenta el riesgo de \(k\) sobre \texttt{Same} al aumentar la edad, y
\(\beta_{k}<0\), disminuye el riesgo en \(k\) respecto de \texttt{Same}.

Es siempre una comparación \textbf{por pares contra la referencia}: el
odds de \texttt{Less} frente a \texttt{Greater} no aparece directamente
---se obtiene combinando ambos juegos
(\(e^{\eta_{\text{Less}}-\eta_{\text{Greater}}}\)) o reeligiendo la
categoría base---. Cambiar la referencia cambia los coeficientes, pero
\textbf{no} las probabilidades ajustadas.

\begin{tcolorbox}[enhanced jigsaw, coltitle=black, rightrule=.15mm, arc=.35mm, leftrule=.75mm, colback=white, colbacktitle=quarto-callout-note-color!10!white, left=2mm, breakable, toptitle=1mm, toprule=.15mm, bottomtitle=1mm, titlerule=0mm, title=\textcolor{quarto-callout-note-color}{\faInfo}\hspace{0.5em}{Individual o agrupado: qué cambia y qué no}, bottomrule=.15mm, opacityback=0, opacitybacktitle=0.6, colframe=quarto-callout-note-color-frame]

Tanto si la respuesta categórica es \textbf{individual} (una fila por
sujeto, como aquí con \texttt{raterisk}) como \textbf{agrupada}
(recuentos por perfil, multinomial), el \textbf{modelo y sus
estimaciones son los mismos}: misma verosimilitud, misma interpretación
de coeficientes, y la \textbf{selección de modelos} (LRT, AIC/BIC)
funciona igual. Lo que \textbf{no} vale con datos individuales es el
contraste \textbf{\(\chi^2\) de bondad de ajuste} ni el análisis de
\textbf{residuos}: como en el paso Bernoulli→binomial
(\protect\hyperlink{sec-u13-gof-binomial}{Sección 3.2}), solo recuperan
su sentido al \textbf{agrupar}. Por eso aquí valoramos el ajuste mirando
las \textbf{probabilidades predichas}, no un \(\chi^2\).

\end{tcolorbox}

\textbf{Un cambio de motor: esto ya no se ajusta con \texttt{glm()}.} La
función \texttt{glm()} está programada para respuestas de la familia
exponencial \textbf{univariante} ---un predictor lineal para una
respuesta escalar---; el multinomial tiene respuesta \textbf{vectorial}
(\(K\) probabilidades que suman 1) y necesita \textbf{\(K-1\)
predictores lineales a la vez}, de modo que se sale de lo que cubre
\texttt{glm()} (aunque sigue siendo un GLM en sentido amplio). En R se
ajusta con \texttt{nnet::multinom()} ---lo que usamos aquí--- o con
\texttt{VGAM::vglm(...,\ family\ =\ multinomial)}; ambas estiman por
máxima verosimilitud el \textbf{mismo} logit de categoría base.

\begin{tcolorbox}[enhanced jigsaw, coltitle=black, rightrule=.15mm, arc=.35mm, leftrule=.75mm, colback=white, colbacktitle=quarto-callout-tip-color!10!white, left=2mm, breakable, toptitle=1mm, toprule=.15mm, bottomtitle=1mm, titlerule=0mm, title=\textcolor{quarto-callout-tip-color}{\faLightbulb}\hspace{0.5em}{🔧 En R. Ajustar e inspeccionar el multinomial nominal}, bottomrule=.15mm, opacityback=0, opacitybacktitle=0.6, colframe=quarto-callout-tip-color-frame]

La respuesta debe ser un \textbf{factor} en el que hay que identificar
la categoría de referencia.

\textbf{Ajuste (dos motores equivalentes).}

\begin{Shaded}
\begin{Highlighting}[]
\FunctionTok{library}\NormalTok{(nnet)}
\NormalTok{m\_nom  }\OtherTok{\textless{}{-}} \FunctionTok{multinom}\NormalTok{(}\FunctionTok{relevel}\NormalTok{(raterisk, }\AttributeTok{ref =} \StringTok{"Same"}\NormalTok{) }\SpecialCharTok{\textasciitilde{}}\NormalTok{ age }\SpecialCharTok{+}\NormalTok{ priorfrac, }\AttributeTok{data =}\NormalTok{ glow, }\AttributeTok{trace =} \ConstantTok{FALSE}\NormalTok{)}

\FunctionTok{library}\NormalTok{(VGAM)}
\NormalTok{m\_vglm }\OtherTok{\textless{}{-}} \FunctionTok{vglm}\NormalTok{(raterisk }\SpecialCharTok{\textasciitilde{}}\NormalTok{ age }\SpecialCharTok{+}\NormalTok{ priorfrac,}
               \AttributeTok{family =} \FunctionTok{multinomial}\NormalTok{(}\AttributeTok{refLevel =} \StringTok{"Same"}\NormalTok{), }\AttributeTok{data =}\NormalTok{ glow)}
\end{Highlighting}
\end{Shaded}

\textbf{Resúmenes.}

En \texttt{nnet}:

\begin{itemize}
\tightlist
\item
  \texttt{summary(m\_nom)} da un bloque de coeficientes por categoría
  no-referencia (\texttt{Less\textbar{}Same},
  \texttt{Greater\textbar{}Same}) con sus errores estándar;
\item
  \texttt{multinom} no imprime \(p\)-valores, así que se sacan con
  \texttt{broom::tidy(m\_nom)} (añade \(z\) y \(p\)), y
\item
  \texttt{exp(coef(m\_nom))} da los \textbf{odds ratios relativos}.
\end{itemize}

En \texttt{vglm}:

\begin{itemize}
\tightlist
\item
  \texttt{summary(m\_vglm)} y
\item
  \texttt{coef(m\_vglm,\ matrix\ =\ TRUE)}.
\end{itemize}

\textbf{Bondad de ajuste.} Con datos individuales no hay \(\chi^2\) de
bondad de ajuste (\protect\hyperlink{sec-u12-diagnostico}{Sección 2.6});
se usan la verosimilitud y los criterios de información:
\texttt{logLik()}, \texttt{AIC()}, \texttt{BIC()}, \texttt{deviance()}.

\textbf{Comparación de modelos.} Modelos \textbf{anidados} con el LRT
---\texttt{lmtest::lrtest(m0,\ m\_nom)} (o \texttt{anova()} en
\texttt{vglm})---; \textbf{no anidados}, con AIC/BIC.

\textbf{Ajustes y predicciones.}

En \texttt{nnet}:

\begin{verbatim}
fitted(m_nom) # matriz de probabilidades** (una columna por nivel) 
predict(m_nom, type = "probs") #  predice probabilidades para datos nuevos
predict(m_nom, type ="class") # da la categoría más probable. 
ggeffects::ggpredict(m_nom, terms = "age [all]") # para visualizarlas frente a una covariable
\end{verbatim}

En \texttt{vglm}:

\begin{verbatim}
fitted(m_vglm)` # equivale a `predict(m_vglm, type = "response")`
nd <- data.frame(age = 60:90, priorfrac = "No") # datos nuevos
predict(m_vglm, newdata = nd, type = "response")   # predicción de probabilidades para datos nuevos, matriz n x K
predict(m_vglm, newdata = nd, type = "link") # predicción de logit frente a referencia
colnames(p)[max.col(p, ties.method = "first")] # da la categoría más probable
\end{verbatim}

\end{tcolorbox}

Ajustamos a continuación este modelo sobre GLOW, utilizando
\texttt{nnet}.

\begin{Shaded}
\begin{Highlighting}[]
\FunctionTok{library}\NormalTok{(nnet)}
\NormalTok{m\_nom }\OtherTok{\textless{}{-}} \FunctionTok{multinom}\NormalTok{(}\FunctionTok{relevel}\NormalTok{(raterisk, }\AttributeTok{ref =} \StringTok{"Same"}\NormalTok{) }\SpecialCharTok{\textasciitilde{}}\NormalTok{ age }\SpecialCharTok{+}\NormalTok{ priorfrac,}
                  \AttributeTok{data =}\NormalTok{ glow, }\AttributeTok{trace =} \ConstantTok{FALSE}\NormalTok{)}
\FunctionTok{summary}\NormalTok{(m\_nom) }\CommentTok{\# dos bloques de coeficientes: Less|Same y Greater|Same}
\end{Highlighting}
\end{Shaded}

\begin{verbatim}
Call:
multinom(formula = relevel(raterisk, ref = "Same") ~ age + priorfrac, 
    data = glow, trace = FALSE)

Coefficients:
        (Intercept)          age priorfracYes
Less     -0.7074585  0.009952521   -0.4240676
Greater   1.0758421 -0.022686702    0.8022770

Std. Errors:
        (Intercept)        age priorfracYes
Less      0.8471448 0.01242217    0.2809105
Greater   0.8876437 0.01325035    0.2598469

Residual Deviance: 1072.149 
AIC: 1084.149 
\end{verbatim}

\begin{Shaded}
\begin{Highlighting}[]
\FunctionTok{tidy}\NormalTok{(m\_nom) }\CommentTok{\# visualización ordenada}
\end{Highlighting}
\end{Shaded}

\begin{verbatim}
# A tibble: 6 x 6
  y.level term         estimate std.error statistic p.value
  <chr>   <chr>           <dbl>     <dbl>     <dbl>   <dbl>
1 Less    (Intercept)  -0.707      0.847     -0.835 0.404  
2 Less    age           0.00995    0.0124     0.801 0.423  
3 Less    priorfracYes -0.424      0.281     -1.51  0.131  
4 Greater (Intercept)   1.08       0.888      1.21  0.226  
5 Greater age          -0.0227     0.0133    -1.71  0.0869 
6 Greater priorfracYes  0.802      0.260      3.09  0.00202
\end{verbatim}

\begin{Shaded}
\begin{Highlighting}[]
\FunctionTok{exp}\NormalTok{(}\FunctionTok{coef}\NormalTok{(m\_nom))      }\CommentTok{\# odds ratios relativos}
\end{Highlighting}
\end{Shaded}

\begin{verbatim}
        (Intercept)       age priorfracYes
Less      0.4928953 1.0100022    0.6543797
Greater   2.9324614 0.9775687    2.2306142
\end{verbatim}

\texttt{multinom} devuelve \textbf{dos bloques} de odds ratios
relativos, \texttt{Less\textbar{}Same} y \texttt{Greater\textbar{}Same}.
El motor es la \textbf{fractura previa}: tener una multiplica por
\textbf{2,23} las odds de percibirse en \textbf{mayor} riesgo
(\texttt{Greater}) frente a «igual» (\texttt{Same}) y, a la vez, las
reduce a \textbf{0,65} (−35 \%) para percibirse en \textbf{menor} riesgo
(\texttt{Less}); es decir, empuja la percepción fuera de «Same» hacia
«Greater» ---clínicamente coherente y buena señal de que el modelo capta
algo real---. La \textbf{edad}, en cambio, apenas influye (OR ≈ 1: 1,01
para \texttt{Less}, 0,98 para \texttt{Greater}). Dos matices: son odds
\textbf{relativas a \texttt{Same}}, no frente a todas las categorías; y
los interceptos (OR de \texttt{Less} ≈ 0,49 y \texttt{Greater} ≈ 2,93)
corresponden a \texttt{age\ =\ 0}, una extrapolación sin sentido
sustantivo ---para leer en probabilidad por edades, la
Figura~\ref{fig-u13-pred-nominal}---.

Con los coeficientes estimados, la \textbf{probabilidad predicha} sale
de combinar los dos predictores lineales (con
\(\mathbf1(\text{priorfrac}=\text{Yes})\) y \(\text{age}\) en años):

\[
\eta_{\text{Less}}    = -0{,}71 + 0{,}010\,\text{age} - 0{,}42\cdot\mathbf1(\text{priorfrac}=\text{Yes}),\qquad
\eta_{\text{Greater}} =  1{,}08 - 0{,}023\,\text{age} + 0{,}80\cdot\mathbf1(\text{priorfrac}=\text{Yes}).
\]

Como \texttt{Same} es la referencia y las tres probabilidades suman 1,
\[
\Pr(Y=\text{Less})=\frac{e^{\eta_{\text{Less}}}}{1+e^{\eta_{\text{Less}}}+e^{\eta_{\text{Greater}}}},
\qquad
\Pr(Y=\text{Greater})=\frac{e^{\eta_{\text{Greater}}}}{1+e^{\eta_{\text{Less}}}+e^{\eta_{\text{Greater}}}},
\] y
\(\Pr(Y=\text{Same})=1\big/\!\left(1+e^{\eta_{\text{Less}}}+e^{\eta_{\text{Greater}}}\right)\).

Lo más comunicable es la \textbf{probabilidad predicha} de cada nivel.
La Figura~\ref{fig-u13-pred-nominal} la traza frente a la edad:

\begin{Shaded}
\begin{Highlighting}[]
\FunctionTok{library}\NormalTok{(ggeffects)}
\FunctionTok{plot}\NormalTok{(}\FunctionTok{ggpredict}\NormalTok{(m\_nom, }\AttributeTok{terms =} \StringTok{"age [all]"}\NormalTok{))}
\end{Highlighting}
\end{Shaded}

\begin{figure}[H]

{\centering \includegraphics{caso1_binaria_files/figure-pdf/fig-u13-pred-nominal-1.pdf}

}

\caption{\label{fig-u13-pred-nominal}Probabilidad predicha de cada nivel
de raterisk según la edad (multinomial nominal, priorfrac = No). Cada
categoría sigue su propia curva, sin restricción entre ellas.}

\end{figure}

La Figura~\ref{fig-u13-pred-nominal} traduce los coeficientes a
\textbf{probabilidad} (con \texttt{priorfrac\ =\ No}) y revela el patrón
con la edad: la de percibirse en \textbf{menor} riesgo (\texttt{Less})
\textbf{crece} de ≈32 \% a los 55 años a ≈47 \% a los 90, mientras que
la de \textbf{mayor} riesgo (\texttt{Greater}) \textbf{cae} de ≈31 \% a
≈14 \%; la referencia \texttt{Same} se mantiene casi plana en torno al
38 \% (con una leve curvatura: en la escala de probabilidad las curvas
no tienen por qué ser monótonas aunque los logit sean lineales). Las
tres se mueven \textbf{con independencia} ---sin restricción de orden,
propio del modelo nominal--- pero \textbf{suman 1} a cada edad. El
desplazamiento concuerda con los OR por año (\texttt{Less} 1,01;
\texttt{Greater} 0,98): minúsculos uno a uno, pero apreciables
\textbf{acumulados} sobre 35 años. Sustantivamente llama la atención
---y es clínicamente relevante--- que las mujeres mayores tiendan a
\textbf{autopercibir menos riesgo} justo cuando el objetivo es mayor,
una brecha de percepción que el modelo capta. (Con
\texttt{priorfrac\ =\ Yes} el perfil se desplazaría hacia
\texttt{Greater}, acorde con su OR de 2,23.)

\textbf{Bondad de ajuste, diagnóstico y comparación.}

\begin{itemize}
\tightlist
\item
  Con la respuesta \textbf{individual} no disponemos del \(\chi^2\) de
  bondad de ajuste ni de los residuos estándar
  (\protect\hyperlink{sec-u13-gof-binomial}{Sección 3.2}): el ajuste se
  valora con la \textbf{verosimilitud}.
\item
  El contraste \textbf{global} ---¿aporta el modelo frente al nulo?---
  es un LRT, \(D_0-D\sim\chi^2\) con \((K-1)\,p\) grados de libertad
  (aquí, los dos bloques de pendientes a la vez).
\item
  El \textbf{diagnóstico} se apoya en las \textbf{probabilidades
  predichas}
\item
  La \textbf{calibración} se examina \textbf{por categoría}
  ---probabilidad media predicha frente a frecuencia observada en grupos
  de riesgo, una \emph{reliability} por nivel---.
\item
  El acierto clasificatorio se concluye con la \textbf{matriz de
  confusión} (\texttt{predict(m\_nom,\ type\ =\ "class")}) o el
  \emph{Brier score} multiclase.
\item
  Para residuos formales hay que recurrir a residuos
  \textbf{simulados/surrogate}, no a los clásicos.
\item
  La \textbf{comparación de modelos} sigue la lógica habitual:

  \begin{itemize}
  \tightlist
  \item
    \textbf{anidados} por LRT (\texttt{lmtest::lrtest}, \texttt{anova})
    ---por ejemplo, contrastar si \texttt{priorfrac} aporta--- y
  \item
    \textbf{no anidados} por \textbf{AIC/BIC}.
  \end{itemize}
\end{itemize}

Una comparación natural es con el modelo \textbf{ordinal} de la
\protect\hyperlink{sec-u13-ordinal}{Sección 3.4}, que aprovecha el orden
de las categorías y gana parsimonia: un solo \(\beta\) por covariable en
vez de uno por categoría.

El \emph{Brier score} multiclase, que acabamos de mencionar, merece su
definición. Es el \textbf{error cuadrático medio} entre el vector de
probabilidades predichas y la clase observada codificada como
\emph{one-hot}: \[
\text{BS}=\frac{1}{n}\sum_{i=1}^{n}\sum_{k=1}^{K}\bigl(\hat p_{ik}-y_{ik}\bigr)^2,
\qquad y_{ik}=\mathbf 1\{\text{la clase observada de }i\text{ es }k\},
\] donde \(\hat p_{ik}\) es la probabilidad predicha de que la
observación \(i\) pertenezca a la categoría \(k\). Es un \emph{scoring
rule} \textbf{propio} ---penaliza a la vez la mala \textbf{calibración}
y la falta de \textbf{nitidez}---: vale \textbf{0} con predicción
perfecta y, sumando sobre las \(K\) clases, como máximo \textbf{2} (el
caso binario es el conocido \(\tfrac1n\sum_i(\hat p_i-y_i)^2\in[0,1]\)).
Cuanto \textbf{menor}, mejor; sirve para \textbf{comparar modelos} sobre
las mismas observaciones.

\begin{tcolorbox}[enhanced jigsaw, coltitle=black, rightrule=.15mm, arc=.35mm, leftrule=.75mm, colback=white, colbacktitle=quarto-callout-tip-color!10!white, left=2mm, breakable, toptitle=1mm, toprule=.15mm, bottomtitle=1mm, titlerule=0mm, title=\textcolor{quarto-callout-tip-color}{\faLightbulb}\hspace{0.5em}{🔧 En R. Brier score multiclase}, bottomrule=.15mm, opacityback=0, opacitybacktitle=0.6, colframe=quarto-callout-tip-color-frame]

A partir de la matriz de probabilidades predichas (\texttt{fitted}) y la
clase observada en \emph{one-hot}:

\begin{Shaded}
\begin{Highlighting}[]
\NormalTok{P   }\OtherTok{\textless{}{-}} \FunctionTok{fitted}\NormalTok{(m\_nom)                         }\CommentTok{\# n x K; columnas = niveles de la respuesta}
\NormalTok{obs }\OtherTok{\textless{}{-}} \FunctionTok{factor}\NormalTok{(glow}\SpecialCharTok{$}\NormalTok{raterisk, }\AttributeTok{levels =} \FunctionTok{colnames}\NormalTok{(P))}
\NormalTok{Y   }\OtherTok{\textless{}{-}} \FunctionTok{model.matrix}\NormalTok{(}\SpecialCharTok{\textasciitilde{}} \DecValTok{0} \SpecialCharTok{+}\NormalTok{ obs)               }\CommentTok{\# one{-}hot}
\FunctionTok{colnames}\NormalTok{(Y) }\OtherTok{\textless{}{-}} \FunctionTok{levels}\NormalTok{(obs)                   }\CommentTok{\# alinea columnas con las de P}
\NormalTok{brier }\OtherTok{\textless{}{-}} \FunctionTok{mean}\NormalTok{(}\FunctionTok{rowSums}\NormalTok{((P }\SpecialCharTok{{-}}\NormalTok{ Y)}\SpecialCharTok{\^{}}\DecValTok{2}\NormalTok{))            }\CommentTok{\# 0 = perfecto; máximo 2}
\NormalTok{brier}
\end{Highlighting}
\end{Shaded}

Con \texttt{yardstick} se calcula directamente ---la respuesta y las
columnas de probabilidad deben llevar el nombre de cada clase---:

\begin{Shaded}
\begin{Highlighting}[]
\FunctionTok{library}\NormalTok{(yardstick); }\FunctionTok{library}\NormalTok{(tibble)}
\NormalTok{prob }\OtherTok{\textless{}{-}} \FunctionTok{as\_tibble}\NormalTok{(}\FunctionTok{fitted}\NormalTok{(m\_nom))             }\CommentTok{\# columnas = niveles de la respuesta}
\NormalTok{prob}\SpecialCharTok{$}\NormalTok{obs }\OtherTok{\textless{}{-}} \FunctionTok{factor}\NormalTok{(glow}\SpecialCharTok{$}\NormalTok{raterisk, }\AttributeTok{levels =} \FunctionTok{names}\NormalTok{(prob))}
\FunctionTok{brier\_class}\NormalTok{(prob, }\AttributeTok{truth =}\NormalTok{ obs, Same, Less, Greater)}
\end{Highlighting}
\end{Shaded}

\end{tcolorbox}

\hypertarget{sec-u13-ordinal}{%
\subsection{3.4 Categorías ordenadas: el modelo de odds
proporcionales}\label{sec-u13-ordinal}}

Consideramos ahora la respuesta individualizada \texttt{raterisk}, que
\textbf{tiene orden}
(\texttt{Less\ \textless{}\ Same\ \textless{}\ Greater}), y el modelo
politómico anterior lo ignoraba. Cuando el orden existe, conviene
aprovecharlo con el modelo de \textbf{odds proporcionales}
(\emph{cumulative logit}), que modela los \textbf{logits acumulados} con
umbrales ordenados y \textbf{un único} vector de efectos: \[
\operatorname{logit}\Pr(Y\le k)= \operatorname{log}\frac{Pr(Y\le k)}{Pr(Y> k)}=\theta_k-\bigl(\beta \cdot\text{age}+\gamma \cdot \mathbf1(\text{priorfrac}=\text{Yes})\bigr),
\] con
\(k\in\{\text{Less},\,\text{Same}\},\quad \theta_{\text{Less}}<\theta_{\text{Same}}\).

Frente al nominal de la \protect\hyperlink{sec-u13-nominal}{Sección 3.3}
---que estimaba \textbf{un vector de efectos por categoría}--- aquí hay
dos piezas bien diferenciadas:

\begin{itemize}
\item
  los \textbf{umbrales} (interceptos ordenados), uno por cada corte
  entre categorías contiguas. Con tres niveles hay \textbf{dos},

  \begin{itemize}
  \tightlist
  \item
    corte \texttt{Less\textbar{}Same}, es decir
    \(\Pr(Y\le\text{Less})\): \(\theta_{\text{Less}}\) y
  \item
    corte \texttt{Same\textbar{}Greater}, \(\Pr(Y\le\text{Same})\):
    \(\theta_{\text{Same}}\),
  \end{itemize}

  ordenados \(\theta_{\text{Less}}<\theta_{\text{Same}}\) para
  garantizar que las probabilidades acumuladas crezcan y, con ellas, que
  las de cada categoría sean no negativas---;
\item
  un \textbf{único vector de efectos} \((\beta,\gamma)\)
  \textbf{compartido por los dos cortes} ---de ahí «odds
  proporcionales»---: las covariables \textbf{desplazan en bloque} toda
  la distribución sobre la escala ordenada, sin alterar la separación
  entre umbrales.
\end{itemize}

\textbf{El signo y su lectura.}

El predictor lineal entra \textbf{restado}, \(\theta_k-\eta\), con
\(\eta=\beta \cdot\text{age}+\gamma \cdot \mathbf1(\text{priorfrac}=\text{Yes})\)
-que es la convención de \texttt{MASS::polr}-. No es un capricho: hace
que un coeficiente \textbf{positivo} empuje hacia las categorías
\textbf{altas}. Si \(\beta>0\), sube \(\eta\), baja \(\theta_k-\eta\) y
baja \(\Pr(Y\le k)\) en \textbf{todos} los cortes ---es decir, sube
\(\Pr(Y>k)\)---.

Conviene, pues, leer el efecto sobre las \textbf{odds de superar un
corte}: - una unidad más de \(x\) multiplica
\(\operatorname{odds}(Y>k)\) por \(e^{\beta}\) ---el \textbf{mismo}
factor en los dos cortes--- y, equivalentemente,
\(\operatorname{odds}(Y\le k)\) por \(e^{-\beta}\) - el hecho de que una
paciente haya tenido fractura previa frente a que no la haya tenido,
resta \(\gamma\) a \(\operatorname{logit}\Pr(Y\le k)\), de modo que
multiplica por \(e^{-\gamma}\) las odds acumuladas
\(\operatorname{odds}(Y\le k)\) y por \(e^{+\gamma}\) las odds de
superar el corte \(\operatorname{odds}(Y>k)\), y es el mismo factor en
los dos cortes.

\textbf{En GLOW}, los umbrales estimados son
\(\theta_{\text{Less}}\approx-2{,}06\) y
\(\theta_{\text{Same}}\approx-0{,}44\), y el vector de efectos
\((\beta_{\text{age}},\gamma_{\text{priorfrac}})\approx(-0{,}023,\ 0{,}90)\).
Se lee así: la \textbf{fractura previa} (\(\gamma>0\),
\(e^{\gamma}\approx 2{,}46\)) multiplica por \(\approx 2{,}5\) las odds
de percibirse en una categoría \textbf{más alta} (más \texttt{Greater},
menos \texttt{Less}), por igual en ambos cortes; la \textbf{edad}
(\(\beta<0\), \(e^{\beta}\approx 0{,}98\) por año) empuja levemente
hacia categorías \textbf{más bajas} (\texttt{Less}). Ambos efectos
coinciden en dirección con los del nominal, señal de que la
proporcionalidad de odds es razonable aquí.

Se ajusta con \texttt{MASS::polr}, que \textbf{exige} una respuesta de
tipo \textbf{factor ordenado}:

\begin{Shaded}
\begin{Highlighting}[]
\NormalTok{glow\_ord }\OtherTok{\textless{}{-}}\NormalTok{ glow }\SpecialCharTok{|\textgreater{}}
  \FunctionTok{mutate}\NormalTok{(}\AttributeTok{raterisk =} \FunctionTok{ordered}\NormalTok{(raterisk, }\AttributeTok{levels =} \FunctionTok{c}\NormalTok{(}\StringTok{"Less"}\NormalTok{, }\StringTok{"Same"}\NormalTok{, }\StringTok{"Greater"}\NormalTok{)))}
\NormalTok{m\_ord }\OtherTok{\textless{}{-}}\NormalTok{ MASS}\SpecialCharTok{::}\FunctionTok{polr}\NormalTok{(raterisk }\SpecialCharTok{\textasciitilde{}}\NormalTok{ age }\SpecialCharTok{+}\NormalTok{ priorfrac, }\AttributeTok{data =}\NormalTok{ glow\_ord, }\AttributeTok{Hess =} \ConstantTok{TRUE}\NormalTok{)}
\FunctionTok{tidy}\NormalTok{(m\_ord)        }\CommentTok{\# dos umbrales (theta) y los efectos (beta), comunes a ambos cortes}
\end{Highlighting}
\end{Shaded}

\begin{verbatim}
# A tibble: 4 x 5
  term         estimate std.error statistic coef.type  
  <chr>           <dbl>     <dbl>     <dbl> <chr>      
1 age           -0.0228   0.00974    -2.34  coefficient
2 priorfracYes   0.900    0.204       4.41  coefficient
3 Less|Same     -2.06     0.665      -3.10  scale      
4 Same|Greater  -0.439    0.658      -0.667 scale      
\end{verbatim}

\begin{Shaded}
\begin{Highlighting}[]
\FunctionTok{exp}\NormalTok{(}\FunctionTok{coef}\NormalTok{(m\_ord))      }\CommentTok{\# odds ratios acumulados}
\end{Highlighting}
\end{Shaded}

\begin{verbatim}
         age priorfracYes 
   0.9774424    2.4588903 
\end{verbatim}

Frente al nominal de 3.3, el ordinal \textbf{gana parsimonia}: un solo
\(\beta\) por covariable en vez de uno por categoría, a cambio del
\textbf{supuesto de proporcionalidad de odds} (el mismo efecto en todos
los cortes), que conviene contrastar (un test de líneas paralelas).

\begin{tcolorbox}[enhanced jigsaw, coltitle=black, rightrule=.15mm, arc=.35mm, leftrule=.75mm, colback=white, colbacktitle=quarto-callout-tip-color!10!white, left=2mm, breakable, toptitle=1mm, toprule=.15mm, bottomtitle=1mm, titlerule=0mm, title=\textcolor{quarto-callout-tip-color}{\faLightbulb}\hspace{0.5em}{🔧 En R. Ajustar e inspeccionar el ordinal (odds proporcionales)}, bottomrule=.15mm, opacityback=0, opacitybacktitle=0.6, colframe=quarto-callout-tip-color-frame]

\textbf{Respuesta.} Debe ser un \textbf{factor ordenado}
(\texttt{ordered()} con los niveles en su orden natural):

\begin{Shaded}
\begin{Highlighting}[]
\NormalTok{glow}\SpecialCharTok{$}\NormalTok{raterisk }\OtherTok{\textless{}{-}} \FunctionTok{ordered}\NormalTok{(glow}\SpecialCharTok{$}\NormalTok{raterisk, }\AttributeTok{levels =} \FunctionTok{c}\NormalTok{(}\StringTok{"Less"}\NormalTok{, }\StringTok{"Same"}\NormalTok{, }\StringTok{"Greater"}\NormalTok{))}
\end{Highlighting}
\end{Shaded}

\textbf{Ajuste (motores equivalentes).}

\begin{Shaded}
\begin{Highlighting}[]
\NormalTok{m\_ord  }\OtherTok{\textless{}{-}}\NormalTok{ MASS}\SpecialCharTok{::}\FunctionTok{polr}\NormalTok{(raterisk }\SpecialCharTok{\textasciitilde{}}\NormalTok{ age }\SpecialCharTok{+}\NormalTok{ priorfrac, }\AttributeTok{data =}\NormalTok{ glow, }\AttributeTok{Hess =} \ConstantTok{TRUE}\NormalTok{)}
\NormalTok{m\_clm  }\OtherTok{\textless{}{-}}\NormalTok{ ordinal}\SpecialCharTok{::}\FunctionTok{clm}\NormalTok{(raterisk }\SpecialCharTok{\textasciitilde{}}\NormalTok{ age }\SpecialCharTok{+}\NormalTok{ priorfrac, }\AttributeTok{data =}\NormalTok{ glow)}
\NormalTok{m\_vglm }\OtherTok{\textless{}{-}}\NormalTok{ VGAM}\SpecialCharTok{::}\FunctionTok{vglm}\NormalTok{(raterisk }\SpecialCharTok{\textasciitilde{}}\NormalTok{ age }\SpecialCharTok{+}\NormalTok{ priorfrac,}
                     \AttributeTok{family =} \FunctionTok{cumulative}\NormalTok{(}\AttributeTok{parallel =} \ConstantTok{TRUE}\NormalTok{), }\AttributeTok{data =}\NormalTok{ glow)  }\CommentTok{\# parallel = odds proporcionales}
\end{Highlighting}
\end{Shaded}

\textbf{Resúmenes.}

\begin{itemize}
\tightlist
\item
  \texttt{summary(m\_ord)} da \textbf{un único bloque de pendientes}
  (común a los cortes) y los \textbf{umbrales} (\texttt{Intercepts} /
  \texttt{zeta});
\item
  \texttt{broom::tidy(m\_ord)} los ordena (marca \texttt{coef.type}
  =\texttt{coefficient} / \texttt{scale})\\
\item
  \texttt{exp(coef(m\_ord))} da los \textbf{odds ratios acumulados}.
\end{itemize}

\texttt{polr} no imprime \(p\)-valores; se sacan del valor \(t\) o con
\texttt{ordinal::clm} (que sí los da).

\textbf{Bondad de ajuste y comparación.}

\begin{itemize}
\tightlist
\item
  Verosimilitud y criterios: \texttt{logLik()}, \texttt{AIC()},
  \texttt{BIC()}, \texttt{deviance()}.
\item
  \textbf{Modelos anidados} con el LRT (\texttt{anova},
  \texttt{lmtest::lrtest}).
\item
  Y clave aquí, el \textbf{supuesto de proporcionalidad de odds} (líneas
  paralelas) se contrasta con \texttt{brant::brant(m\_ord)} o
  \texttt{ordinal::nominal\_test(m\_clm)}.
\item
  Comparar con el \textbf{nominal} por AIC/BIC indica si compensa la
  parsimonia.
\end{itemize}

\textbf{Ajustes y predicciones.}

\begin{itemize}
\tightlist
\item
  \texttt{predict(m\_ord,\ type\ =\ "probs")} da la \textbf{matriz de
  probabilidades},
\item
  \texttt{predict(m\_ord,\ type\ =\ "class")} la categoría modal, y
\item
  \texttt{ggeffects::ggpredict(m\_ord,\ terms\ =\ "age{[}all{]}")} las
  dibuja.
\end{itemize}

\end{tcolorbox}

\begin{Shaded}
\begin{Highlighting}[]
\FunctionTok{library}\NormalTok{(ggeffects)}
\FunctionTok{plot}\NormalTok{(}\FunctionTok{ggpredict}\NormalTok{(m\_ord, }\AttributeTok{terms =} \StringTok{"age [all]"}\NormalTok{))}
\end{Highlighting}
\end{Shaded}

\begin{figure}[H]

{\centering \includegraphics{caso1_binaria_files/figure-pdf/fig-u13-pred-ordinal-1.pdf}

}

\caption{\label{fig-u13-pred-ordinal}Probabilidad predicha de cada nivel
de raterisk según la edad (odds proporcionales, priorfrac = No). Las
curvas están ligadas por un único efecto de la edad común a todos los
cortes.}

\end{figure}

La Figura~\ref{fig-u13-pred-ordinal} lo deja ver: frente a las
\textbf{curvas libres} del nominal (Figura~\ref{fig-u13-pred-nominal}),
aquí están \textbf{ligadas} por un único efecto de la edad ---la
categoría central (\texttt{Same}) forma una joroba y las extremas son
monótonas---. Si ambas imágenes se parecen, el supuesto de
proporcionalidad es razonable y ganamos parsimonia; si difieren mucho,
el ordinal estaría forzando.

Con ambos modelos ajustados, la comparación numérica resume las
\textbf{discrepancias} entre nominal y ordinal:

\begin{Shaded}
\begin{Highlighting}[]
\CommentTok{\# Ajuste y parsimonia (el ordinal usa menos parámetros)}
\NormalTok{ajuste }\OtherTok{\textless{}{-}}\NormalTok{ tibble}\SpecialCharTok{::}\FunctionTok{tibble}\NormalTok{(}
  \AttributeTok{Modelo     =} \FunctionTok{c}\NormalTok{(}\StringTok{"Nominal (multinom)"}\NormalTok{, }\StringTok{"Ordinal (polr)"}\NormalTok{),}
  \AttributeTok{Parametros =} \FunctionTok{c}\NormalTok{(}\FunctionTok{length}\NormalTok{(}\FunctionTok{coef}\NormalTok{(m\_nom)), }\FunctionTok{length}\NormalTok{(}\FunctionTok{coef}\NormalTok{(m\_ord)) }\SpecialCharTok{+} \FunctionTok{length}\NormalTok{(m\_ord}\SpecialCharTok{$}\NormalTok{zeta)),}
  \AttributeTok{logLik     =} \FunctionTok{c}\NormalTok{(}\FunctionTok{as.numeric}\NormalTok{(}\FunctionTok{logLik}\NormalTok{(m\_nom)), }\FunctionTok{as.numeric}\NormalTok{(}\FunctionTok{logLik}\NormalTok{(m\_ord))),}
  \AttributeTok{AIC        =} \FunctionTok{c}\NormalTok{(}\FunctionTok{AIC}\NormalTok{(m\_nom), }\FunctionTok{AIC}\NormalTok{(m\_ord)),}
  \AttributeTok{BIC        =} \FunctionTok{c}\NormalTok{(}\FunctionTok{BIC}\NormalTok{(m\_nom), }\FunctionTok{BIC}\NormalTok{(m\_ord))}
\NormalTok{)}

\CommentTok{\# Odds ratios: el nominal da uno por categoría; el ordinal, uno (acumulado) común}
\NormalTok{or\_nom }\OtherTok{\textless{}{-}} \FunctionTok{exp}\NormalTok{(}\FunctionTok{coef}\NormalTok{(m\_nom))      }\CommentTok{\# filas Less/Greater x (Intercept, age, priorfracYes)}
\NormalTok{or\_ord }\OtherTok{\textless{}{-}} \FunctionTok{exp}\NormalTok{(}\FunctionTok{coef}\NormalTok{(m\_ord))      }\CommentTok{\# age, priorfracYes (acumulado)}
\NormalTok{discrepancias }\OtherTok{\textless{}{-}}\NormalTok{ tibble}\SpecialCharTok{::}\FunctionTok{tibble}\NormalTok{(}
  \AttributeTok{Covariable              =} \FunctionTok{c}\NormalTok{(}\StringTok{"age"}\NormalTok{, }\StringTok{"priorfracYes"}\NormalTok{),}
  \StringTok{\textasciigrave{}}\AttributeTok{Nominal Less|Same}\StringTok{\textasciigrave{}}     \OtherTok{=}\NormalTok{ or\_nom[}\StringTok{"Less"}\NormalTok{,    }\FunctionTok{c}\NormalTok{(}\StringTok{"age"}\NormalTok{, }\StringTok{"priorfracYes"}\NormalTok{)],}
  \StringTok{\textasciigrave{}}\AttributeTok{Nominal Greater|Same}\StringTok{\textasciigrave{}}  \OtherTok{=}\NormalTok{ or\_nom[}\StringTok{"Greater"}\NormalTok{, }\FunctionTok{c}\NormalTok{(}\StringTok{"age"}\NormalTok{, }\StringTok{"priorfracYes"}\NormalTok{)],}
  \StringTok{\textasciigrave{}}\AttributeTok{Ordinal (acumulado)}\StringTok{\textasciigrave{}}   \OtherTok{=}\NormalTok{ or\_ord[}\FunctionTok{c}\NormalTok{(}\StringTok{"age"}\NormalTok{, }\StringTok{"priorfracYes"}\NormalTok{)]}
\NormalTok{)}

\NormalTok{ajuste}
\end{Highlighting}
\end{Shaded}

\begin{verbatim}
# A tibble: 2 x 5
  Modelo             Parametros logLik   AIC   BIC
  <chr>                   <int>  <dbl> <dbl> <dbl>
1 Nominal (multinom)          6  -536. 1084. 1109.
2 Ordinal (polr)              4  -536. 1081. 1098.
\end{verbatim}

\begin{Shaded}
\begin{Highlighting}[]
\NormalTok{discrepancias}
\end{Highlighting}
\end{Shaded}

\begin{verbatim}
# A tibble: 2 x 4
  Covariable   `Nominal Less|Same` `Nominal Greater|Same` `Ordinal (acumulado)`
  <chr>                      <dbl>                  <dbl>                 <dbl>
1 age                        1.01                   0.978                 0.977
2 priorfracYes               0.654                  2.23                  2.46 
\end{verbatim}

La tabla resume el contraste. En \textbf{ajuste}, el ordinal emplea
\textbf{4 parámetros} (dos umbrales + dos pendientes) frente a los
\textbf{6} del nominal (dos bloques de tres); si la proporcionalidad de
odds se sostiene, su \texttt{AIC}/\texttt{BIC} salen
\textbf{competitivos o mejores} pese a algo menos de verosimilitud. En
\textbf{efectos}, la diferencia estructural es nítida: el nominal estima
\textbf{un OR por categoría} y el ordinal los \textbf{colapsa en uno
solo, acumulado}. Para \texttt{priorfrac}, el nominal da \textbf{0,65}
sobre \texttt{Less} y \textbf{2,23} sobre \texttt{Greater} (frente a
\texttt{Same}): leídos en el orden, ambos empujan \textbf{hacia arriba}
(más \texttt{Greater}, menos \texttt{Less}), justo lo que resume el OR
acumulado \textbf{2,46}. Para \texttt{age} ocurre lo simétrico
(\textbf{1,01} y \textbf{0,98} → empuje \textbf{hacia abajo}, hacia
\texttt{Less}; acumulado \textbf{0,98}). Como las dos lecturas nominales
son \textbf{coherentes con el orden} en ambas covariables, el resumen
ordinal es fiel y la Figura~\ref{fig-u13-pred-ordinal} se parece a la
Figura~\ref{fig-u13-pred-nominal}; si un OR claramente \(>1\) conviviera
con otro \(<1\) \textbf{incoherentes} con el orden, el ordinal estaría
forzando ---y el \textbf{test de líneas paralelas} lo delataría---.

\begin{tcolorbox}[enhanced jigsaw, coltitle=black, rightrule=.15mm, arc=.35mm, leftrule=.75mm, colback=white, colbacktitle=quarto-callout-note-color!10!white, left=2mm, breakable, toptitle=1mm, toprule=.15mm, bottomtitle=1mm, titlerule=0mm, title=\textcolor{quarto-callout-note-color}{\faInfo}\hspace{0.5em}{Puente a la Unidad 4}, bottomrule=.15mm, opacityback=0, opacitybacktitle=0.6, colframe=quarto-callout-note-color-frame]

La \protect\hyperlink{sec-u14}{Unidad 4} introduce los \textbf{efectos
aleatorios}, sobre la respuesta \textbf{binaria}, que es donde los
desarrollamos. Añadirlos a un \textbf{ordinal} como este es una
extensión directa (un \emph{ordinal mixto}, con \texttt{ordinal::clmm}):
cada grupo desvía los umbrales y el resto del modelo no cambia. El hilo
ordinal, además, reaparece en la \protect\hyperlink{sec-u15}{Unidad 5}:
modelar el \textbf{tiempo discreto a un evento} es, de hecho, un ordinal
\emph{continuation-ratio}, allí con \textbf{fragilidad} de centro.

\end{tcolorbox}

\begin{tcolorbox}[enhanced jigsaw, coltitle=black, rightrule=.15mm, arc=.35mm, leftrule=.75mm, colback=white, colbacktitle=quarto-callout-important-color!10!white, left=2mm, breakable, toptitle=1mm, toprule=.15mm, bottomtitle=1mm, titlerule=0mm, title=\textcolor{quarto-callout-important-color}{\faExclamation}\hspace{0.5em}{Tarea del estudiante (1.3) --- resuelve las Preguntas 3 y 4 del caso}, bottomrule=.15mm, opacityback=0, opacitybacktitle=0.6, colframe=quarto-callout-important-color-frame]

Sobre la cohorte de tu equipo, responde dos preguntas del cliente
(\protect\hyperlink{sec-u16}{Unidad 6}).

\textbf{Pregunta 3 --- tipo de fractura (\texttt{clase\_nom}).}

\begin{enumerate}
\def\labelenumi{\arabic{enumi}.}
\tightlist
\item
  \textbf{Identifica y ajusta (RA3).} Reconoce que el tipo es una
  respuesta \textbf{nominal} (sin orden) y ajusta el modelo adecuado
  (\texttt{nnet::multinom}) con \texttt{x1} y \texttt{x2}.
\item
  \textbf{Interpreta y predice.} Lee los efectos como \textbf{odds
  ratios relativos} frente a la categoría base (¿hacia qué tipo empuja
  cada factor?), predice el tipo más probable para un par de perfiles y
  evalúa (p.~ej. \textbf{Brier multiclase}).
\end{enumerate}

\textbf{Pregunta 4 --- gravedad (\texttt{sever\_ord}).}

\begin{enumerate}
\def\labelenumi{\arabic{enumi}.}
\setcounter{enumi}{2}
\tightlist
\item
  \textbf{Identifica y ajusta (RA3).} Reconoce que la gravedad es una
  respuesta \textbf{ordinal} y ajusta el modelo de \textbf{odds
  proporcionales}; comprueba el supuesto de proporcionalidad.
\item
  \textbf{Interpreta y responde.} Lee el \textbf{odds ratio acumulado}
  común (¿el tratamiento hace las fracturas menos graves?) y responde al
  cliente.
\end{enumerate}

\emph{Apoyo --- control de calidad (RA2, RA4).} Como técnica de bondad
de ajuste, agrupa la binaria de la Pregunta 1 en su versión
\textbf{binomial} y aplica el \textbf{contraste χ²} (deviance y
Pearson), interpretándolo y discutiendo el efecto del tamaño de celda.

\end{tcolorbox}

\hypertarget{sec-u14}{%
\section{4 · Efectos Aleatorios y Modelos Mixtos}\label{sec-u14}}

Hasta aquí, cada coeficiente ha sido un \textbf{efecto fijo}: un único
número válido para toda la población. Pero los datos rara vez son una
masa homogénea ---los pacientes se agrupan en centros, los alumnos en
aulas, las medidas repetidas en cada sujeto---, y los grupos
\textbf{difieren entre sí}, tanto en su nivel de partida como en cómo
responden a las covariables. Ignorar ese agrupamiento (o gastarse un
parámetro por grupo) lleva a problemas conocidos; los \textbf{efectos
aleatorios} son la vía intermedia: modelar la \textbf{variación entre
grupos} como una distribución de la que cada grupo es una muestra. Es el
primero de los dos hilos en espiral del curso ---los efectos mixtos---,
que abrimos aquí sobre el GLMM logístico y reactivaremos en cada caso.

En nuestros datos ese agrupamiento tiene nombre: la variable
\textbf{\texttt{centro}}. Tanto en \textbf{GLOW} como en la
\textbf{cohorte} simulada, cada paciente pertenece a un centro
hospitalario, y los centros \textbf{no son intercambiables}: difieren en
\textbf{criterios diagnósticos} ---cómo se define y registra una
fractura, qué umbrales se aplican--- y en \textbf{factores de entorno}
ligados a su \textbf{ubicación} ---perfil de las pacientes, exposición,
hábitos, recursos del área---. Dos mujeres con la misma edad y los
mismos antecedentes pueden, por tanto, tener un riesgo algo distinto
según \textbf{dónde} se las atienda.

En la Figura~\ref{fig-u14-prop-centro} se muestran las diferencias entre
centros respecto a las tasas de fractura, en las bases de datos
\texttt{cohorte} y GLOW. De partida, se aprecia cierta variabilidad de
base entre los centros, mucho más notoria en \texttt{cohorte} que en
GLOW.

\begin{Shaded}
\begin{Highlighting}[]
\FunctionTok{library}\NormalTok{(patchwork)}
\NormalTok{lims     }\OtherTok{\textless{}{-}} \FunctionTok{c}\NormalTok{(}\DecValTok{0}\NormalTok{, }\FloatTok{0.62}\NormalTok{)}
\NormalTok{col\_coh  }\OtherTok{\textless{}{-}} \StringTok{"\#2c7fb8"}   \CommentTok{\# azul  → cohorte}
\NormalTok{col\_glow }\OtherTok{\textless{}{-}} \StringTok{"\#e6810a"}   \CommentTok{\# naranja → GLOW}

\NormalTok{prop\_coh }\OtherTok{\textless{}{-}}\NormalTok{ cohorte }\SpecialCharTok{|\textgreater{}}
  \FunctionTok{group\_by}\NormalTok{(centro) }\SpecialCharTok{|\textgreater{}}
  \FunctionTok{summarise}\NormalTok{(}\AttributeTok{prop =} \FunctionTok{mean}\NormalTok{(ever), }\AttributeTok{n =} \FunctionTok{n}\NormalTok{(), }\AttributeTok{.groups =} \StringTok{"drop"}\NormalTok{)}
\NormalTok{p\_coh }\OtherTok{\textless{}{-}} \FunctionTok{ggplot}\NormalTok{(prop\_coh, }\FunctionTok{aes}\NormalTok{(}\FunctionTok{reorder}\NormalTok{(}\FunctionTok{factor}\NormalTok{(centro), prop), prop)) }\SpecialCharTok{+}
  \FunctionTok{geom\_hline}\NormalTok{(}\AttributeTok{yintercept =} \FunctionTok{mean}\NormalTok{(cohorte}\SpecialCharTok{$}\NormalTok{ever), }\AttributeTok{linetype =} \StringTok{"dashed"}\NormalTok{, }\AttributeTok{color =} \StringTok{"grey50"}\NormalTok{) }\SpecialCharTok{+}
  \FunctionTok{geom\_col}\NormalTok{(}\AttributeTok{fill =}\NormalTok{ col\_coh) }\SpecialCharTok{+}
  \FunctionTok{coord\_flip}\NormalTok{() }\SpecialCharTok{+} \FunctionTok{scale\_y\_continuous}\NormalTok{(}\AttributeTok{limits =}\NormalTok{ lims) }\SpecialCharTok{+}
  \FunctionTok{labs}\NormalTok{(}\AttributeTok{x =} \StringTok{"Centro"}\NormalTok{, }\AttributeTok{y =} \StringTok{"Proporción de fractura (ever)"}\NormalTok{, }\AttributeTok{title =} \StringTok{"Cohorte (24 centros)"}\NormalTok{)}

\NormalTok{prop\_glow }\OtherTok{\textless{}{-}}\NormalTok{ glow }\SpecialCharTok{|\textgreater{}}
  \FunctionTok{mutate}\NormalTok{(}\AttributeTok{frac01 =} \FunctionTok{as.integer}\NormalTok{(fracture }\SpecialCharTok{==} \StringTok{"Yes"}\NormalTok{)) }\SpecialCharTok{|\textgreater{}}
  \FunctionTok{group\_by}\NormalTok{(site\_id) }\SpecialCharTok{|\textgreater{}}
  \FunctionTok{summarise}\NormalTok{(}\AttributeTok{prop =} \FunctionTok{mean}\NormalTok{(frac01), }\AttributeTok{n =} \FunctionTok{n}\NormalTok{(), }\AttributeTok{.groups =} \StringTok{"drop"}\NormalTok{)}
\NormalTok{p\_glow }\OtherTok{\textless{}{-}} \FunctionTok{ggplot}\NormalTok{(prop\_glow, }\FunctionTok{aes}\NormalTok{(}\FunctionTok{reorder}\NormalTok{(}\FunctionTok{factor}\NormalTok{(site\_id), prop), prop)) }\SpecialCharTok{+}
  \FunctionTok{geom\_hline}\NormalTok{(}\AttributeTok{yintercept =} \FunctionTok{mean}\NormalTok{(glow}\SpecialCharTok{$}\NormalTok{fracture }\SpecialCharTok{==} \StringTok{"Yes"}\NormalTok{), }\AttributeTok{linetype =} \StringTok{"dashed"}\NormalTok{, }\AttributeTok{color =} \StringTok{"grey50"}\NormalTok{) }\SpecialCharTok{+}
  \FunctionTok{geom\_col}\NormalTok{(}\AttributeTok{fill =}\NormalTok{ col\_glow) }\SpecialCharTok{+}
  \FunctionTok{coord\_flip}\NormalTok{() }\SpecialCharTok{+} \FunctionTok{scale\_y\_continuous}\NormalTok{(}\AttributeTok{limits =} \FunctionTok{c}\NormalTok{(}\DecValTok{0}\NormalTok{,}\FloatTok{0.35}\NormalTok{)) }\SpecialCharTok{+}
  \FunctionTok{labs}\NormalTok{(}\AttributeTok{x =} \StringTok{"Centro (site\_id)"}\NormalTok{, }\AttributeTok{y =} \StringTok{"Proporción de fractura"}\NormalTok{, }\AttributeTok{title =} \StringTok{"GLOW (6 centros)"}\NormalTok{)}

\NormalTok{p\_coh }\SpecialCharTok{+}\NormalTok{ p\_glow}
\end{Highlighting}
\end{Shaded}

\begin{figure}[H]

{\centering \includegraphics{caso1_binaria_files/figure-pdf/fig-u14-prop-centro-1.pdf}

}

\caption{\label{fig-u14-prop-centro}Proporción observada de fractura por
centro. Izquierda (azul): cohorte simulada (24 centros). Derecha
(naranja): GLOW (6 centros, datos reales). Línea discontinua: proporción
global. Ambos paneles comparten el eje de proporción: la cohorte se
despliega (\textasciitilde0,12 a \textasciitilde0,60; los centros
estructuran el riesgo), mientras que en GLOW los seis se apiñan (el
centro apenas aporta, ICC ≈ 0,7\%). El color distingue cada conjunto, no
codifica ninguna variable.}

\end{figure}

Ahora bien, el centro \textbf{no nos interesa como predictor}: no
queremos ---ni tendría sentido clínico--- predecir si una paciente se
fracturará \textbf{según a qué hospital pertenezca}, ni en los seis
centros de GLOW, ni en los veinticuatro de la cohorte, pues no son un
fin en sí mismos; son una \textbf{muestra} de los muchos centros
posibles. Lo que \textbf{sí} debemos reconocer es que el centro induce
\textbf{agrupamiento}: las pacientes de un mismo centro se
\textbf{parecen más entre sí} de lo que explicaría el azar ---no son del
todo independientes--- y, a la vez, hay \textbf{diferencias sistemáticas
entre centros} ---unos con un riesgo basal más alto que otros---. Gastar
un coeficiente \textbf{fijo} por centro malgastaría parámetros y no
generalizaría a centros nuevos; ignorarlo falsearía la precisión.

La vía adecuada es modelar esa variación \textbf{entre centros} como
\textbf{efectos aleatorios} ---el objeto de esta unidad---: estimamos
\textbf{cuánto varían} los centros (una varianza), no el valor concreto
de cada uno. Trabajamos sobre la \textbf{cohorte} simulada, cuya verdad
conocida permite comprobar si lo recuperamos bien.

\begin{tcolorbox}[enhanced jigsaw, coltitle=black, rightrule=.15mm, arc=.35mm, leftrule=.75mm, colback=white, colbacktitle=quarto-callout-note-color!10!white, left=2mm, breakable, toptitle=1mm, toprule=.15mm, bottomtitle=1mm, titlerule=0mm, title=\textcolor{quarto-callout-note-color}{\faInfo}\hspace{0.5em}{Sobre esta unidad}, bottomrule=.15mm, opacityback=0, opacitybacktitle=0.6, colframe=quarto-callout-note-color-frame]

RA implicados:

\begin{itemize}
\tightlist
\item
  \textbf{RA6 · Efectos aleatorios.} Ajustar e interpretar GLMM
  (intercepto y pendiente): componentes de varianza, ICC,
  \emph{shrinkage} y condicional vs marginal.
\item
  \textbf{RA4 · Diagnóstico y bondad de ajuste} (del mixto). Residuos
  simulados (DHARMa), \(R^2\) marginal/condicional y solidez del ajuste.
\end{itemize}

(Agresti, 2015; Bates et~al., 2015)

\end{tcolorbox}

\hypertarget{sec-u14-que-son}{%
\subsection{4.1 Qué son y para qué sirven}\label{sec-u14-que-son}}

La idea central es un cambio de pregunta. Un efecto \textbf{fijo}
responde \emph{«¿cuál es el efecto?»} con \textbf{un número para toda la
población}. Un efecto \textbf{aleatorio} no estima un número por grupo,
sino la \textbf{distribución} de la que esos números proceden: en lugar
de \emph{«¿cuál es el efecto del grupo 7?»}, pregunta \emph{«¿cuánto
varían los efectos entre grupos?»} ---su \textbf{varianza}---. Se modela
la \textbf{dispersión} entre grupos (en nuestro caso, \texttt{centros}),
no cada grupo por separado.

\textbf{Dos cosas pueden variar de un grupo a otro:}

\begin{itemize}
\tightlist
\item
  El \textbf{nivel de base} → un \textbf{intercepto aleatorio}
  \(u_{0j}\): cada centro parte de un riesgo basal distinto (la curva se
  \textbf{desplaza} arriba o abajo).
\item
  La \textbf{sensibilidad a una covariable} → una \textbf{pendiente
  aleatoria} \(u_{1j}\): en cada centro, el efecto de cierta covariable
  continua es distinto (la curva se \textbf{inclina} de forma
  diferente).
\end{itemize}

\textbf{¿Para qué sirven?} Cuatro cosas concretas:

\begin{enumerate}
\def\labelenumi{\arabic{enumi}.}
\tightlist
\item
  \textbf{Respetar la no independencia.} Las observaciones de un mismo
  grupo se parecen entre sí; ignorarlo da errores estándar
  \textbf{demasiado pequeños} (falsa precisión). El efecto aleatorio
  modela esa correlación intragrupo correctamente.
\item
  \textbf{Partial pooling (tomar prestado).} Las estimaciones de cada
  grupo se \textbf{encogen} hacia la media global, tanto más cuanto
  menos datos aporta el grupo: estabiliza los grupos pequeños sin
  imponer a todos un único valor rígido.
\item
  \textbf{Cuantificar la heterogeneidad.} Las \textbf{varianzas}
  \(\sigma^2\) son, ellas mismas, la respuesta a una pregunta
  sustantiva: ¿cuánto difieren los grupos en su nivel, o en el efecto de
  una covariable?
\item
  \textbf{Generalizar.} Al tratar los grupos como una \textbf{muestra}
  de una población de grupos, las conclusiones valen para grupos
  \textbf{no observados}, no solo para los de la muestra.
\end{enumerate}

En una frase: el efecto fijo dice \emph{cuánto}; el efecto aleatorio
dice \emph{cuánto varía ese cuánto entre grupos} ---y lo resume en una
varianza---.

Los efectos aleatorios se pueden incorporar en un modelo lineal
(generalizado) a través de la interceptación, de la pendiente o de
ambas. Veamos sus implicaciones y las modelizaciones asociadas.

\hypertarget{sec-u14-intercepto}{%
\subsection{4.2 Intercepto aleatorio}\label{sec-u14-intercepto}}

Cuando hablamos de una afectación en la interceptación, nos estamos
refiriendo a la modelización del \textbf{nivel de base}. Para tratar
este nivel de base, considerando un posible efecto de agrupación
relativo a \texttt{centro} en nuestra base de datos, hay tres
estrategias:

\begin{itemize}
\tightlist
\item
  \textbf{Ignorarlo} (\emph{complete pooling}): un solo intercepto para
  todos. Subestima la incertidumbre y puede sesgar los efectos.
\item
  \textbf{Un parámetro por centro} (\emph{no pooling}): tantos
  interceptos como centros. Gasta muchos grados de libertad y no
  generaliza a centros nuevos.
\item
  \textbf{Intercepto aleatorio} (\emph{partial pooling}): los
  interceptos de centro son \textbf{extracciones de una distribución},
  \(u_j\sim\mathcal{N}(0,\sigma_u^2)\); no estimamos cada \(u_j\), sino
  su \textbf{varianza} \(\sigma_u^2\), asumiendo que todos están
  centrados en torno a una interceptación media global (y por lo tanto
  estas desviaciones de la interceptación global tienen media cero).
\end{itemize}

El modelo, para la binaria de la cohorte, es un \textbf{GLMM logístico}:
\[
\operatorname{logit}\Pr(\text{ever}_{ij}=1)=\beta_0+\beta_1 x_{1ij}+\beta_2 x_{2ij}+u_j,
\qquad u_j\sim\mathcal{N}(0,\sigma_u^2).
\]

Ese término \(u_j\) no es una abstracción: ya teníamos su \textbf{huella
visual} en la cohorte. La Figura~\ref{fig-u14-prop-centro} mostraba que
la proporción de fractura \textbf{varía de un centro a otro} ---de
\(\approx 0{,}12\) a \(\approx 0{,}60\) (en torno a una media global
explícita con la vertical punteada)---, justo la dispersión de
\textbf{niveles basales} que \(u_j\) recoge (y que en GLOW era mínima).
Hay que tener en cuenta que en general, una proporción por centro
mezclaría el nivel basal con la \textbf{composición} de covariables del
centro; si estas estuvieran correladas con \texttt{centro}, el gráfico
estaría confundiendo efectos.

En \textbf{forma matricial} ---la notación habitual de los modelos
mixtos, que conviene reconocer desde ya---, el predictor lineal de todas
las observaciones se escribe \[
\boldsymbol\eta=\mathbf{X}\boldsymbol\beta+\mathbf{Z}\mathbf{u},\qquad
\mathbf{u}\sim\mathcal{N}(\mathbf{0},\mathbf{G}),\qquad g(\boldsymbol\mu)=\boldsymbol\eta,
\] donde \(\mathbf{X}\) (\(n\times p\)) y \(\boldsymbol\beta\) son el
diseño y los coeficientes \textbf{fijos}; \(\mathbf{Z}\) (\(n\times q\))
y \(\mathbf{u}\) son el diseño y los efectos \textbf{aleatorios}, con
\(\mathbf{G}\) su matriz de covarianzas; y \(g\) es el enlace (aquí,
logit).

Para el \textbf{intercepto aleatorio},
\(\mathbf{u}=(u_1,\dots,u_J)^\top\) son los \(J\) interceptos de centro,
\(\mathbf{Z}\) es la matriz \textbf{indicadora} que asigna cada
observación a su centro, y \(\mathbf{G}=\sigma_u^2\,\mathbf{I}_J\). La
parte \(\mathbf{X}\boldsymbol\beta\) es lo de siempre; lo nuevo es el
término \(\mathbf{Z}\mathbf{u}\).

\hypertarget{intercepto-aleatorio-en-cohorte}{%
\subsubsection{\texorpdfstring{Intercepto aleatorio en
\texttt{cohorte}}{Intercepto aleatorio en cohorte}}\label{intercepto-aleatorio-en-cohorte}}

Para interpretar los resultados que genera un modelo mixto, nada mejor
que ajustarlo en unos datos y ver qué produce. Utilizamos
\texttt{cohorte} para intentar predecir la respuesta binaria relativa a
si se dió o no una fractura, \texttt{evento}, en función de la
fragilidad ósea (\texttt{x1}) y el tratamiento preventivo (\texttt{x2})
recibido por cada paciente, teniendo en cuenta el \texttt{centro}
clínico de referencia.

La primera tarea pasa pues, por inspeccionar los datos. Así, acompañamos
al gráfico inicial en el que representamos las probabilidades de
fractura por centros en Figura~\ref{fig-u14-prop-centro}, con la
Figura~\ref{fig-cohorte-eda-box}, donde mostramos la relación de la
respuesta con las predictivas \texttt{x1} y \texttt{x2}.

\begin{Shaded}
\begin{Highlighting}[]
\NormalTok{cohorte }\SpecialCharTok{|\textgreater{}}
  \FunctionTok{mutate}\NormalTok{(}\AttributeTok{Desenlace   =} \FunctionTok{factor}\NormalTok{(evento, }\AttributeTok{levels =} \FunctionTok{c}\NormalTok{(}\DecValTok{0}\NormalTok{, }\DecValTok{1}\NormalTok{), }\AttributeTok{labels =} \FunctionTok{c}\NormalTok{(}\StringTok{"Sin fractura"}\NormalTok{, }\StringTok{"Fractura"}\NormalTok{)),}
         \AttributeTok{Tratamiento =} \FunctionTok{factor}\NormalTok{(x2,     }\AttributeTok{levels =} \FunctionTok{c}\NormalTok{(}\DecValTok{0}\NormalTok{, }\DecValTok{1}\NormalTok{), }\AttributeTok{labels =} \FunctionTok{c}\NormalTok{(}\StringTok{"No"}\NormalTok{, }\StringTok{"Sí"}\NormalTok{))) }\SpecialCharTok{|\textgreater{}}
  \FunctionTok{ggplot}\NormalTok{(}\FunctionTok{aes}\NormalTok{(}\AttributeTok{x =}\NormalTok{ x1, }\AttributeTok{y =}\NormalTok{ Desenlace, }\AttributeTok{fill =}\NormalTok{ Tratamiento)) }\SpecialCharTok{+}
  \FunctionTok{geom\_boxplot}\NormalTok{(}\AttributeTok{alpha =} \FloatTok{0.75}\NormalTok{, }\AttributeTok{outlier.alpha =} \FloatTok{0.25}\NormalTok{) }\SpecialCharTok{+}
  \FunctionTok{labs}\NormalTok{(}\AttributeTok{x =} \StringTok{"Fragilidad ósea (x1, en z)"}\NormalTok{, }\AttributeTok{y =} \StringTok{"Desenlace"}\NormalTok{, }\AttributeTok{fill =} \StringTok{"Tratamiento"}\NormalTok{)}
\end{Highlighting}
\end{Shaded}

\begin{figure}[H]

{\centering \includegraphics[width=1\textwidth,height=\textheight]{caso1_binaria_files/figure-pdf/fig-cohorte-eda-box-1.pdf}

}

\caption{\label{fig-cohorte-eda-box}Fragilidad ósea (x1) según el
desenlace, por brazo de tratamiento (x2). Quienes fracturan presentan
mayor fragilidad; entre las tratadas, las que llegan a fracturar son aún
más frágiles (el desplazamiento de umbral propio de un factor
protector).}

\end{figure}

El patrón en la Figura~\ref{fig-cohorte-eda-box} es \textbf{nítido}: las
cajas de \textbf{Fractura} están claramente \textbf{desplazadas a la
derecha} ---mayor fragilidad--- respecto a las de \textbf{Sin fractura}
(la mediana de \texttt{x1} pasa de \(\approx-0{,}19\) a
\(\approx0{,}69\)), que es el efecto de la fragilidad sobre el riesgo. Y
dentro de cada desenlace, las \textbf{tratadas} (\texttt{x2} = Sí)
quedan algo \textbf{más a la derecha} que las no tratadas: entre las que
\textbf{fracturan}, las tratadas son \textbf{aún más frágiles} (mediana
\(\approx0{,}88\) frente a \(\approx0{,}64\)). Tiene sentido ---como el
tratamiento \textbf{protege}, solo «caen» las de fragilidad más alta---:
es el \textbf{desplazamiento de umbral} propio de un factor protector.

Ajustemos pues, el modelo de interceptación aleatoria, pero también el
modelo \emph{pooled}, en el que ignoramos el centro, pues de este modo
podremos valorar el efecto de incorporar el \texttt{centro}.

\begin{tcolorbox}[enhanced jigsaw, coltitle=black, rightrule=.15mm, arc=.35mm, leftrule=.75mm, colback=white, colbacktitle=quarto-callout-tip-color!10!white, left=2mm, breakable, toptitle=1mm, toprule=.15mm, bottomtitle=1mm, titlerule=0mm, title=\textcolor{quarto-callout-tip-color}{\faLightbulb}\hspace{0.5em}{🔧 En R. Ajustar un GLMM de intercepto aleatorio}, bottomrule=.15mm, opacityback=0, opacitybacktitle=0.6, colframe=quarto-callout-tip-color-frame]

Mismo \texttt{glmer} que un \texttt{glm}, declarando el efecto aleatorio
con \texttt{(\ ...\ \textbar{}\ grupo\ )}. Para un \textbf{intercepto
aleatorio} por \texttt{centro}:

\begin{Shaded}
\begin{Highlighting}[]
\FunctionTok{library}\NormalTok{(lme4)}
\NormalTok{m\_int }\OtherTok{\textless{}{-}} \FunctionTok{glmer}\NormalTok{(ever }\SpecialCharTok{\textasciitilde{}}\NormalTok{ x1 }\SpecialCharTok{+}\NormalTok{ x2 }\SpecialCharTok{+}\NormalTok{ (}\DecValTok{1} \SpecialCharTok{|}\NormalTok{ centro), }\AttributeTok{family =}\NormalTok{ binomial, }\AttributeTok{data =}\NormalTok{ cohorte)}
\FunctionTok{summary}\NormalTok{(m\_int)        }\CommentTok{\# añade el bloque \textquotesingle{}Random effects\textquotesingle{} con la varianza de centro}
\end{Highlighting}
\end{Shaded}

La sintaxis completa ---pendientes, cruzados, anidados--- en la
\protect\hyperlink{sec-u14-extensiones}{Sección 4.5}.

\end{tcolorbox}

\textbf{Ajuste de los modelos pooled e intercepto aleatorio}

\begin{Shaded}
\begin{Highlighting}[]
\FunctionTok{library}\NormalTok{(lme4)}
\NormalTok{m\_pool }\OtherTok{\textless{}{-}} \FunctionTok{glm}\NormalTok{  (ever }\SpecialCharTok{\textasciitilde{}}\NormalTok{ x1 }\SpecialCharTok{+}\NormalTok{ x2,                }\AttributeTok{family =}\NormalTok{ binomial, }\AttributeTok{data =}\NormalTok{ cohorte)  }\CommentTok{\# complete pooling}
\NormalTok{m\_int  }\OtherTok{\textless{}{-}} \FunctionTok{glmer}\NormalTok{(ever }\SpecialCharTok{\textasciitilde{}}\NormalTok{ x1 }\SpecialCharTok{+}\NormalTok{ x2 }\SpecialCharTok{+}\NormalTok{ (}\DecValTok{1} \SpecialCharTok{|}\NormalTok{ centro),  }\AttributeTok{family =}\NormalTok{ binomial, }\AttributeTok{data =}\NormalTok{ cohorte)  }\CommentTok{\# intercepto aleatorio}

\FunctionTok{summary}\NormalTok{(m\_pool)}
\end{Highlighting}
\end{Shaded}

\begin{verbatim}

Call:
glm(formula = ever ~ x1 + x2, family = binomial, data = cohorte)

Coefficients:
            Estimate Std. Error z value Pr(>|z|)    
(Intercept) -0.84845    0.10582  -8.018 1.08e-15 ***
x1           1.01615    0.09679  10.499  < 2e-16 ***
x2          -0.63341    0.17022  -3.721 0.000198 ***
---
Signif. codes:  0 '***' 0.001 '**' 0.01 '*' 0.05 '.' 0.1 ' ' 1

(Dispersion parameter for binomial family taken to be 1)

    Null deviance: 1055.58  on 876  degrees of freedom
Residual deviance:  903.66  on 874  degrees of freedom
AIC: 909.66

Number of Fisher Scoring iterations: 4
\end{verbatim}

\textbf{Cómo se lee la salida de \texttt{glmer}}

El resumen del ajuste \texttt{glmer} es:

\begin{Shaded}
\begin{Highlighting}[]
\FunctionTok{summary}\NormalTok{(m\_int)}
\end{Highlighting}
\end{Shaded}

\begin{verbatim}
Generalized linear mixed model fit by maximum likelihood (Laplace
  Approximation) [glmerMod]
 Family: binomial  ( logit )
Formula: ever ~ x1 + x2 + (1 | centro)
   Data: cohorte

      AIC       BIC    logLik -2*log(L)  df.resid 
    899.0     918.1    -445.5     891.0       873 

Scaled residuals: 
    Min      1Q  Median      3Q     Max 
-2.0791 -0.5968 -0.3734  0.6784 10.9504 

Random effects:
 Groups Name        Variance Std.Dev.
 centro (Intercept) 0.1959   0.4426  
Number of obs: 877, groups:  centro, 24

Fixed effects:
            Estimate Std. Error z value Pr(>|z|)    
(Intercept)  -0.8955     0.1422  -6.297 3.04e-10 ***
x1            1.0803     0.1017  10.620  < 2e-16 ***
x2           -0.6605     0.1741  -3.793 0.000149 ***
---
Signif. codes:  0 '***' 0.001 '**' 0.01 '*' 0.05 '.' 0.1 ' ' 1

Correlation of Fixed Effects:
   (Intr) x1    
x1 -0.189       
x2 -0.422 -0.140
\end{verbatim}

El bloque de \textbf{efectos fijos} es idéntico al de \texttt{glm}
---coeficientes con su error estándar, \(z\) y \(p\), que se interpretan
igual (log-odds, \(e^{\beta}\) = OR)---.

La novedad es el bloque \textbf{\texttt{Random\ effects}}, encima: ahí
están las \textbf{varianzas} y desviaciones típicas de los efectos
aleatorios ---en \texttt{m\_int}, una sola fila,
\texttt{centro\ (Intercept)}--- y el número de grupos. \texttt{glm} no
tiene ese bloque: estima un único intercepto y nada más.

Con el ajuste delante, recorremos las \textbf{tres lecturas} de un
intercepto aleatorio.

\hypertarget{componentes-de-varianza}{%
\subsubsection{1. Componentes de
varianza}\label{componentes-de-varianza}}

La varianza \(\sigma_u^2\) resume cuánto varían los grupos entre sí.
Surge un nuevo índice, común en modelos mixtos: el ICC. En la escala
latente logística, el \textbf{coeficiente de correlación intraclase
(ICC)} mide qué fracción de la variación total ocurre \textbf{entre}
centros.

El \textbf{coeficiente de correlación intraclase (ICC)} mide qué
fracción de la variación total ocurre \textbf{entre} grupos (y no dentro
de ellos); equivalentemente, es la \textbf{correlación} entre dos
observaciones del \textbf{mismo} grupo. En un modelo con intercepto
aleatorio, descompone la varianza en una parte \textbf{entre grupos}
(\(\sigma_u^2\)) y una parte \textbf{dentro} (residual,
\(\sigma_\varepsilon^2\)):

\[
\text{ICC}=\frac{\sigma_u^2}{\sigma_u^2+\sigma_\varepsilon^2}\in[0,1].
\] \(\text{ICC}\to 0\) indica grupos prácticamente indistinguibles (el
agrupamiento apenas importa); \(\text{ICC}\to 1\), grupos muy
heterogéneos y observaciones fuertemente correladas dentro de cada uno.

En un \textbf{GLMM logístico} no hay una varianza residual libre: en la
\textbf{escala latente}, el término de error sigue una logística
estándar, cuya varianza es \textbf{fija}, \(\pi^2/3\approx 3{,}29\). Por
eso el \textbf{ICC latente} se calcula como \[
\text{ICC}=\frac{\sigma_u^2}{\sigma_u^2+\pi^2/3}.
\] Con enlace \textbf{probit} la varianza latente es \(1\), así que
\(\text{ICC}=\sigma_u^2/(\sigma_u^2+1)\).

\begin{tcolorbox}[enhanced jigsaw, coltitle=black, rightrule=.15mm, arc=.35mm, leftrule=.75mm, colback=white, colbacktitle=quarto-callout-tip-color!10!white, left=2mm, breakable, toptitle=1mm, toprule=.15mm, bottomtitle=1mm, titlerule=0mm, title=\textcolor{quarto-callout-tip-color}{\faLightbulb}\hspace{0.5em}{🔧 En R. Calcular el ICC de un GLMM}, bottomrule=.15mm, opacityback=0, opacitybacktitle=0.6, colframe=quarto-callout-tip-color-frame]

\textbf{A mano}, desde la varianza de \texttt{VarCorr}. En la escala
latente de un GLMM \textbf{logístico} la varianza residual es fija,
\(\pi^2/3\) (con \textbf{probit}, \(1\)):

\begin{Shaded}
\begin{Highlighting}[]
\NormalTok{vc   }\OtherTok{\textless{}{-}} \FunctionTok{as.data.frame}\NormalTok{(}\FunctionTok{VarCorr}\NormalTok{(m\_int))}
\NormalTok{s2\_u }\OtherTok{\textless{}{-}}\NormalTok{ vc}\SpecialCharTok{$}\NormalTok{vcov[}\DecValTok{1}\NormalTok{]                  }\CommentTok{\# sigma\_u\^{}2 (varianza del intercepto de centro)}
\NormalTok{s2\_u }\SpecialCharTok{/}\NormalTok{ (s2\_u }\SpecialCharTok{+}\NormalTok{ pi}\SpecialCharTok{\^{}}\DecValTok{2} \SpecialCharTok{/} \DecValTok{3}\NormalTok{)            }\CommentTok{\# ICC latente (logit)}
\end{Highlighting}
\end{Shaded}

\textbf{Con \texttt{performance}}, directo ---y devuelve dos
versiones---:

\begin{Shaded}
\begin{Highlighting}[]
\FunctionTok{library}\NormalTok{(performance)}
\FunctionTok{icc}\NormalTok{(m\_int)}
\end{Highlighting}
\end{Shaded}

\begin{itemize}
\tightlist
\item
  \textbf{Adjusted ICC}:
  \(\sigma_u^2/(\sigma_u^2+\sigma^2_{\text{resid}})\), solo efectos
  aleatorios. Es el que \textbf{coincide} con la fórmula manual y el que
  interpretamos como «cuánto pesa el centro frente a las diferencias
  entre pacientes».
\item
  \textbf{Conditional ICC}: mete además la varianza de los
  \textbf{efectos fijos} en el denominador, así que sale \textbf{menor}.
\end{itemize}

Vale limpio para un \textbf{intercepto aleatorio}; con \textbf{pendiente
aleatoria} el ICC depende de la covariable y conviene usar
\texttt{performance::icc()}.

\end{tcolorbox}

\begin{center}\rule{0.5\linewidth}{0.5pt}\end{center}

En nuestro ejemplo, el ICC responde a una pregunta concreta: de toda la
variación en el riesgo de fractura, ¿cuánta se debe a en qué clínica te
atienden, y cuánta a las diferencias entre pacientes (su fragilidad, su
tratamiento, su azar individual)?

Comparemos pues, la varianza del efecto aleatorio y el ICC, y puesto que
los datos los hemos simulado con unos valores dados, podemos contrastar
las estimaciones con la verdad del DGP (\emph{Data Generating Process}).

\begin{Shaded}
\begin{Highlighting}[]
\CommentTok{\# matriz de varianzas{-}covarianzas de los efectos aleatorios}
\NormalTok{vc }\OtherTok{\textless{}{-}} \FunctionTok{as.data.frame}\NormalTok{(}\FunctionTok{VarCorr}\NormalTok{(m\_int))}
\NormalTok{vc}\SpecialCharTok{$}\NormalTok{vcov[}\DecValTok{1}\NormalTok{]    }\CommentTok{\# varianza  sigma\_u\^{}2}
\end{Highlighting}
\end{Shaded}

\begin{verbatim}
[1] 0.1959143
\end{verbatim}

\begin{Shaded}
\begin{Highlighting}[]
\NormalTok{vc}\SpecialCharTok{$}\NormalTok{sdcor[}\DecValTok{1}\NormalTok{]   }\CommentTok{\# desviación sigma\_u}
\end{Highlighting}
\end{Shaded}

\begin{verbatim}
[1] 0.4426221
\end{verbatim}

\begin{Shaded}
\begin{Highlighting}[]
\CommentTok{\# y de ahí el ICC latente:}
\NormalTok{vc}\SpecialCharTok{$}\NormalTok{vcov[}\DecValTok{1}\NormalTok{] }\SpecialCharTok{/}\NormalTok{ (vc}\SpecialCharTok{$}\NormalTok{vcov[}\DecValTok{1}\NormalTok{] }\SpecialCharTok{+}\NormalTok{ pi}\SpecialCharTok{\^{}}\DecValTok{2}\SpecialCharTok{/}\DecValTok{3}\NormalTok{)}
\end{Highlighting}
\end{Shaded}

\begin{verbatim}
[1] 0.05620383
\end{verbatim}

\begin{Shaded}
\begin{Highlighting}[]
\CommentTok{\# verdad simulada}
\FunctionTok{attr}\NormalTok{(cohorte, }\StringTok{"verdad"}\NormalTok{)}\SpecialCharTok{$}\NormalTok{binaria[}\FunctionTok{c}\NormalTok{(}\StringTok{"sigma\_u"}\NormalTok{, }\StringTok{"icc\_latente"}\NormalTok{)]}
\end{Highlighting}
\end{Shaded}

\begin{verbatim}
$sigma_u
[1] 0.7

$icc_latente
[1] 0.1296342
\end{verbatim}

Aquí el modelo estima \(\hat\sigma_u\approx0{,}44\) (ICC
\(\approx0{,}06\)), \textbf{por debajo} de la verdad del DGP
(\(\sigma_u=0{,}70\), ICC \(=0{,}13\)): la \textbf{subestimación} típica
de las varianzas con respuesta binaria y pocos grupos (24).

Es pues, un resultado que merece comentarse: el ICC cuantifica cuánta de
la variación del riesgo se debe al centro (≈ 6--13 \% aquí) frente a las
diferencias entre pacientes (el resto). El centro pesa poco, pero no
nada: conviene modelarlo ---para no falsear la precisión--- aunque el
grueso del riesgo se juegue a nivel de paciente.

\hypertarget{condicional-vs-marginal}{%
\subsubsection{2. Condicional vs
marginal}\label{condicional-vs-marginal}}

Con un enlace \textbf{no lineal} como el logit, el efecto de una
covariable tiene \textbf{dos versiones que no coinciden} (en un modelo
lineal sí lo harían). Geométricamente se entiende bien con la
Figura~\ref{fig-u14-cond-marg}: cada centro tiene su propia curva
logística con la \textbf{misma pendiente} ---el efecto
\textbf{condicional}, «para un centro dado», que es el que estima
\texttt{glmer}--- pero \textbf{desplazada} por su intercepto \(u_j\). Al
\textbf{promediar} esas curvas sobre la población de centros, la
sigmoide sale \textbf{más plana}, y su pendiente ---el efecto
\textbf{marginal}, poblacional--- es \textbf{menor en magnitud}.

Comparemos pues los coeficientes, de nuevo entre los dos modelos y con
la verdad del DGP.

\begin{Shaded}
\begin{Highlighting}[]
\CommentTok{\# verdad}
\FunctionTok{attr}\NormalTok{(cohorte,}\StringTok{"verdad"}\NormalTok{)}\SpecialCharTok{$}\NormalTok{binaria[}\FunctionTok{c}\NormalTok{(}\StringTok{"beta"}\NormalTok{)]}
\end{Highlighting}
\end{Shaded}

\begin{verbatim}
$beta
   x1    x2 
 0.85 -0.65 
\end{verbatim}

\begin{Shaded}
\begin{Highlighting}[]
\CommentTok{\# coeficientes estimados}
\NormalTok{sp }\OtherTok{\textless{}{-}} \FunctionTok{coef}\NormalTok{(}\FunctionTok{summary}\NormalTok{(m\_pool))   }\CommentTok{\# tabla de efectos fijos del pooled}
\NormalTok{si }\OtherTok{\textless{}{-}} \FunctionTok{coef}\NormalTok{(}\FunctionTok{summary}\NormalTok{(m\_int))    }\CommentTok{\# tabla de efectos fijos del mixto}

\NormalTok{coefs }\OtherTok{\textless{}{-}} \FunctionTok{rbind}\NormalTok{(}
  \FunctionTok{data.frame}\NormalTok{(}\AttributeTok{modelo =} \StringTok{"pooled (glm)"}\NormalTok{,  }\AttributeTok{term =} \FunctionTok{rownames}\NormalTok{(sp),}
             \AttributeTok{estimate =}\NormalTok{ sp[, }\StringTok{"Estimate"}\NormalTok{], }\AttributeTok{se =}\NormalTok{ sp[, }\StringTok{"Std. Error"}\NormalTok{]),}
  \FunctionTok{data.frame}\NormalTok{(}\AttributeTok{modelo =} \StringTok{"mixto (glmer)"}\NormalTok{, }\AttributeTok{term =} \FunctionTok{rownames}\NormalTok{(si),}
             \AttributeTok{estimate =}\NormalTok{ si[, }\StringTok{"Estimate"}\NormalTok{], }\AttributeTok{se =}\NormalTok{ si[, }\StringTok{"Std. Error"}\NormalTok{])}
\NormalTok{)}

\FunctionTok{library}\NormalTok{(ggplot2)}
\FunctionTok{ggplot}\NormalTok{(coefs, }\FunctionTok{aes}\NormalTok{(estimate, term, }\AttributeTok{color =}\NormalTok{ modelo)) }\SpecialCharTok{+}
  \FunctionTok{geom\_vline}\NormalTok{(}\AttributeTok{xintercept =} \DecValTok{0}\NormalTok{, }\AttributeTok{linetype =} \StringTok{"dashed"}\NormalTok{, }\AttributeTok{color =} \StringTok{"grey60"}\NormalTok{) }\SpecialCharTok{+}
  \FunctionTok{geom\_pointrange}\NormalTok{(}\FunctionTok{aes}\NormalTok{(}\AttributeTok{xmin =}\NormalTok{ estimate }\SpecialCharTok{{-}} \FloatTok{1.96} \SpecialCharTok{*}\NormalTok{ se, }\AttributeTok{xmax =}\NormalTok{ estimate }\SpecialCharTok{+} \FloatTok{1.96} \SpecialCharTok{*}\NormalTok{ se),}
                  \AttributeTok{position =} \FunctionTok{position\_dodge}\NormalTok{(}\AttributeTok{width =} \FloatTok{0.5}\NormalTok{)) }\SpecialCharTok{+}
  \FunctionTok{labs}\NormalTok{(}\AttributeTok{x =} \StringTok{"Coeficiente (log{-}odds), ± 1,96·EE"}\NormalTok{, }\AttributeTok{y =} \ConstantTok{NULL}\NormalTok{, }\AttributeTok{color =} \ConstantTok{NULL}\NormalTok{)}
\end{Highlighting}
\end{Shaded}

\begin{figure}[H]

{\centering \includegraphics{caso1_binaria_files/figure-pdf/fig-u14-glmm-coef-1.pdf}

}

\caption{\label{fig-u14-glmm-coef}Efectos fijos del modelo \emph{pooled}
(\texttt{glm}) frente al mixto de intercepto aleatorio (\texttt{glmer}):
punto = coeficiente (log-odds) y bigotes = ±1,96·EE (intervalo de Wald
al 95 \%). Dos lecturas a la vez: los coeficientes del mixto son algo
\textbf{mayores en magnitud} (la des-atenuación marginal→condicional) y
sus bigotes, \textbf{más anchos} ---errores estándar honestos que tienen
en cuenta el agrupamiento por centro, que el \emph{pooled} ignoraba---.}

\end{figure}

Se ve en los números: \texttt{m\_int} da \(x_1=1{,}08\),
\(x_2=-0{,}66\), algo \textbf{mayores} que el \emph{pooled}
(\(x_1=1{,}02\), \(x_2=-0{,}63\)) ---la atenuación marginal→condicional,
leve aquí porque \(\sigma_u\) es pequeño,
\(\beta_{\text{marg}}\approx\beta_{\text{cond}}/\sqrt{1+0{,}346\,\sigma_u^2}\)---.
Y, sobre todo, los \textbf{errores estándar crecen} en el mixto (el del
intercepto pasa de \(0{,}106\) a \(0{,}142\)): \texttt{m\_int}
\textbf{tiene en cuenta} el agrupamiento que el \emph{pooled} ignoraba,
que daba una precisión \textbf{falsamente alta}.

\begin{Shaded}
\begin{Highlighting}[]
\NormalTok{b0 }\OtherTok{\textless{}{-}} \DecValTok{0}\NormalTok{; b1 }\OtherTok{\textless{}{-}} \FloatTok{1.5}\NormalTok{; sigma\_u }\OtherTok{\textless{}{-}} \DecValTok{2}          \CommentTok{\# sigma\_u exagerado para que el aplanamiento se vea}
\NormalTok{x  }\OtherTok{\textless{}{-}} \FunctionTok{seq}\NormalTok{(}\SpecialCharTok{{-}}\DecValTok{6}\NormalTok{, }\DecValTok{6}\NormalTok{, }\AttributeTok{length.out =} \DecValTok{200}\NormalTok{)}

\CommentTok{\# Familia de curvas condicionales: varios centros (distinto u\_j), MISMA pendiente}
\NormalTok{cond\_fam }\OtherTok{\textless{}{-}}\NormalTok{ tidyr}\SpecialCharTok{::}\FunctionTok{expand\_grid}\NormalTok{(}\AttributeTok{x =}\NormalTok{ x, }\AttributeTok{u =} \FunctionTok{c}\NormalTok{(}\SpecialCharTok{{-}}\DecValTok{4}\NormalTok{, }\SpecialCharTok{{-}}\DecValTok{2}\NormalTok{, }\DecValTok{0}\NormalTok{, }\DecValTok{2}\NormalTok{, }\DecValTok{4}\NormalTok{)) }\SpecialCharTok{|\textgreater{}}
\NormalTok{  dplyr}\SpecialCharTok{::}\FunctionTok{mutate}\NormalTok{(}\AttributeTok{p =} \FunctionTok{plogis}\NormalTok{(b0 }\SpecialCharTok{+}\NormalTok{ b1 }\SpecialCharTok{*}\NormalTok{ x }\SpecialCharTok{+}\NormalTok{ u))}

\CommentTok{\# Curva marginal: promedio de plogis(b0 + b1*x + u) sobre u \textasciitilde{} N(0, sigma\_u\^{}2)}
\NormalTok{ug }\OtherTok{\textless{}{-}} \FunctionTok{seq}\NormalTok{(}\SpecialCharTok{{-}}\DecValTok{4} \SpecialCharTok{*}\NormalTok{ sigma\_u, }\DecValTok{4} \SpecialCharTok{*}\NormalTok{ sigma\_u, }\AttributeTok{length.out =} \DecValTok{401}\NormalTok{)}
\NormalTok{w  }\OtherTok{\textless{}{-}} \FunctionTok{dnorm}\NormalTok{(ug, }\DecValTok{0}\NormalTok{, sigma\_u); w }\OtherTok{\textless{}{-}}\NormalTok{ w }\SpecialCharTok{/} \FunctionTok{sum}\NormalTok{(w)}
\NormalTok{destacadas }\OtherTok{\textless{}{-}}\NormalTok{ dplyr}\SpecialCharTok{::}\FunctionTok{bind\_rows}\NormalTok{(}
\NormalTok{  tibble}\SpecialCharTok{::}\FunctionTok{tibble}\NormalTok{(}\AttributeTok{x =}\NormalTok{ x, }\AttributeTok{p =} \FunctionTok{plogis}\NormalTok{(b0 }\SpecialCharTok{+}\NormalTok{ b1 }\SpecialCharTok{*}\NormalTok{ x),}
                 \AttributeTok{tipo =} \StringTok{"condicional (un centro)"}\NormalTok{),}
\NormalTok{  tibble}\SpecialCharTok{::}\FunctionTok{tibble}\NormalTok{(}\AttributeTok{x =}\NormalTok{ x, }\AttributeTok{p =} \FunctionTok{sapply}\NormalTok{(x, \textbackslash{}(xi) }\FunctionTok{sum}\NormalTok{(w }\SpecialCharTok{*} \FunctionTok{plogis}\NormalTok{(b0 }\SpecialCharTok{+}\NormalTok{ b1 }\SpecialCharTok{*}\NormalTok{ xi }\SpecialCharTok{+}\NormalTok{ ug))),}
                 \AttributeTok{tipo =} \StringTok{"marginal (promedio)"}\NormalTok{)}
\NormalTok{)}

\FunctionTok{ggplot}\NormalTok{() }\SpecialCharTok{+}
  \FunctionTok{geom\_line}\NormalTok{(}\AttributeTok{data =}\NormalTok{ cond\_fam,   }\FunctionTok{aes}\NormalTok{(x, p, }\AttributeTok{group =}\NormalTok{ u), }\AttributeTok{color =} \StringTok{"grey75"}\NormalTok{, }\AttributeTok{linewidth =} \FloatTok{0.5}\NormalTok{) }\SpecialCharTok{+}
  \FunctionTok{geom\_line}\NormalTok{(}\AttributeTok{data =}\NormalTok{ destacadas, }\FunctionTok{aes}\NormalTok{(x, p, }\AttributeTok{color =}\NormalTok{ tipo), }\AttributeTok{linewidth =} \FloatTok{1.2}\NormalTok{) }\SpecialCharTok{+}
  \FunctionTok{scale\_color\_manual}\NormalTok{(}\AttributeTok{values =} \FunctionTok{c}\NormalTok{(}\StringTok{"condicional (un centro)"} \OtherTok{=} \StringTok{"\#2c7fb8"}\NormalTok{,}
                                \StringTok{"marginal (promedio)"}     \OtherTok{=} \StringTok{"\#e6810a"}\NormalTok{)) }\SpecialCharTok{+}
  \FunctionTok{labs}\NormalTok{(}\AttributeTok{x =} \StringTok{"covariable"}\NormalTok{, }\AttributeTok{y =} \StringTok{"Probabilidad de fractura"}\NormalTok{, }\AttributeTok{color =} \ConstantTok{NULL}\NormalTok{) }\SpecialCharTok{+}
  \FunctionTok{theme}\NormalTok{(}\AttributeTok{legend.position =} \StringTok{"top"}\NormalTok{)}
\end{Highlighting}
\end{Shaded}

\begin{figure}[H]

{\centering \includegraphics{caso1_binaria_files/figure-pdf/fig-u14-cond-marg-1.pdf}

}

\caption{\label{fig-u14-cond-marg}Por qué \emph{condicional} ≠
\emph{marginal} con enlace no lineal. En gris, varias curvas logísticas
con la misma pendiente desplazadas por el intercepto de centro; en azul,
una de ellas (un centro típico: el efecto condicional, lo que estima
glmer). Su promedio sobre los centros (naranja) es la curva marginal,
más plana: la pendiente poblacional es menor que la de cada centro. El
efecto se ha exagerado (sigma\_u grande) para hacerlo visible; con
sigma\_u pequeño, como en la cohorte, ambas curvas casi coinciden.}

\end{figure}

\hypertarget{encogimiento-o-shrinkage-partial-pooling.}{%
\subsubsection{\texorpdfstring{3. Encogimiento o \emph{shrinkage}
(\emph{partial
pooling}).}{3. Encogimiento o shrinkage (partial pooling).}}\label{encogimiento-o-shrinkage-partial-pooling.}}

Los interceptos de centro \(\hat{u}_j\) \textbf{no} son las desviaciones
crudas de cada centro (las del \emph{no pooling}), sino un
\textbf{compromiso} entre los \textbf{datos del propio centro} y la
\textbf{media global} ---el \(0\) hacia el que
\(u_j\sim\mathcal{N}(0,\sigma_u^2)\) los empuja---, ponderado por
\textbf{cuánta información} aporta el centro: con muchos datos apenas se
encoge; con pocos, su estimación ruidosa se \textbf{arrima a 0}, tomando
fuerza prestada del conjunto.

\begin{tcolorbox}[enhanced jigsaw, coltitle=black, rightrule=.15mm, arc=.35mm, leftrule=.75mm, colback=white, colbacktitle=quarto-callout-tip-color!10!white, left=2mm, breakable, toptitle=1mm, toprule=.15mm, bottomtitle=1mm, titlerule=0mm, title=\textcolor{quarto-callout-tip-color}{\faLightbulb}\hspace{0.5em}{🔧 En R. Obtener los efectos aleatorios estimados (\texttt{ranef})}, bottomrule=.15mm, opacityback=0, opacitybacktitle=0.6, colframe=quarto-callout-tip-color-frame]

\texttt{ranef} devuelve los efectos aleatorios \textbf{predichos} (las
desviaciones \(\hat u_j\) de cada grupo, ya \textbf{encogidas} hacia 0):

\begin{Shaded}
\begin{Highlighting}[]
\FunctionTok{ranef}\NormalTok{(m\_int)}\SpecialCharTok{$}\NormalTok{centro            }\CommentTok{\# data.frame: una fila por centro, columna (Intercept) = u\_j}
\end{Highlighting}
\end{Shaded}

Para acompañarlos de su \textbf{incertidumbre} (y poder dibujar
intervalos), pide las varianzas condicionales y pásalo a
\texttt{data.frame}:

\begin{Shaded}
\begin{Highlighting}[]
\NormalTok{re }\OtherTok{\textless{}{-}} \FunctionTok{as.data.frame}\NormalTok{(}\FunctionTok{ranef}\NormalTok{(m\_int, }\AttributeTok{condVar =} \ConstantTok{TRUE}\NormalTok{))}
\CommentTok{\# columnas: grpvar (factor de agrupamiento), term, grp (centro), condval (u\_j), condsd (su EE)}
\end{Highlighting}
\end{Shaded}

Con eso, el \emph{caterpillar} es inmediato (puntos = \texttt{condval},
bigotes = ±1,96·\texttt{condsd}):

\begin{Shaded}
\begin{Highlighting}[]
\FunctionTok{library}\NormalTok{(ggplot2)}
\FunctionTok{ggplot}\NormalTok{(re, }\FunctionTok{aes}\NormalTok{(}\FunctionTok{reorder}\NormalTok{(grp, condval), condval)) }\SpecialCharTok{+}
  \FunctionTok{geom\_hline}\NormalTok{(}\AttributeTok{yintercept =} \DecValTok{0}\NormalTok{, }\AttributeTok{linetype =} \StringTok{"dashed"}\NormalTok{) }\SpecialCharTok{+}
  \FunctionTok{geom\_pointrange}\NormalTok{(}\FunctionTok{aes}\NormalTok{(}\AttributeTok{ymin =}\NormalTok{ condval }\SpecialCharTok{{-}} \FloatTok{1.96} \SpecialCharTok{*}\NormalTok{ condsd, }\AttributeTok{ymax =}\NormalTok{ condval }\SpecialCharTok{+} \FloatTok{1.96} \SpecialCharTok{*}\NormalTok{ condsd)) }\SpecialCharTok{+}
  \FunctionTok{coord\_flip}\NormalTok{() }\SpecialCharTok{+}
  \FunctionTok{labs}\NormalTok{(}\AttributeTok{x =} \StringTok{"Centro"}\NormalTok{, }\AttributeTok{y =} \StringTok{"Intercepto aleatorio (log{-}odds)"}\NormalTok{)}
\end{Highlighting}
\end{Shaded}

Dos matices: \texttt{ranef} da \textbf{desviaciones} respecto a la media
(centradas en 0); si quieres el \textbf{coeficiente por grupo} completo
(fijo + aleatorio), usa \texttt{coef(m\_int)\$centro}.

Y en un modelo con \textbf{pendiente aleatoria}
(\texttt{(1\ +\ x1\ \textbar{}\ centro)}),
\texttt{ranef(m\_int)\$centro} trae \textbf{dos} columnas
---\texttt{(Intercept)} y \texttt{x1}---, una desviación por efecto
aleatorio.

\end{tcolorbox}

El \emph{caterpillar} o gráfico de oruga en la
Figura~\ref{fig-u14-cohorte-caterpillar} lo hace visible:

\begin{Shaded}
\begin{Highlighting}[]
\NormalTok{re\_coh }\OtherTok{\textless{}{-}} \FunctionTok{as.data.frame}\NormalTok{(}\FunctionTok{ranef}\NormalTok{(m\_int, }\AttributeTok{condVar =} \ConstantTok{TRUE}\NormalTok{))   }\CommentTok{\# columnas: grp, condval, condsd}
\FunctionTok{ggplot}\NormalTok{(re\_coh, }\FunctionTok{aes}\NormalTok{(}\FunctionTok{reorder}\NormalTok{(grp, condval), condval)) }\SpecialCharTok{+}
  \FunctionTok{geom\_hline}\NormalTok{(}\AttributeTok{yintercept =} \DecValTok{0}\NormalTok{, }\AttributeTok{linetype =} \StringTok{"dashed"}\NormalTok{) }\SpecialCharTok{+}
  \FunctionTok{geom\_pointrange}\NormalTok{(}\FunctionTok{aes}\NormalTok{(}\AttributeTok{ymin =}\NormalTok{ condval }\SpecialCharTok{{-}} \FloatTok{1.96} \SpecialCharTok{*}\NormalTok{ condsd,}
                      \AttributeTok{ymax =}\NormalTok{ condval }\SpecialCharTok{+} \FloatTok{1.96} \SpecialCharTok{*}\NormalTok{ condsd), }\AttributeTok{linewidth =} \FloatTok{0.3}\NormalTok{) }\SpecialCharTok{+}
  \FunctionTok{coord\_flip}\NormalTok{() }\SpecialCharTok{+}
  \FunctionTok{labs}\NormalTok{(}\AttributeTok{x =} \StringTok{"Centro"}\NormalTok{, }\AttributeTok{y =} \StringTok{"Intercepto aleatorio (log{-}odds)"}\NormalTok{)}
\end{Highlighting}
\end{Shaded}

\begin{figure}[H]

{\centering \includegraphics[width=0.85\textwidth,height=\textheight]{caso1_binaria_files/figure-pdf/fig-u14-cohorte-caterpillar-1.pdf}

}

\caption{\label{fig-u14-cohorte-caterpillar}Efecto de centro modelizado
en la cohorte: interceptos aleatorios de los 24 centros (desviación en
log-odds respecto al riesgo medio) con IC 95\%, ordenados. La dispersión
del conjunto es lo que el mixto resume en sigma\_u; los centros cuyo
intervalo no cruza el 0 se separan del promedio.}

\end{figure}

En la Figura~\ref{fig-u14-cohorte-caterpillar} cada punto es la
desviación estimada de un centro respecto al riesgo medio; el
\textbf{shrinkage} acerca a \(0\) a los centros con \textbf{intervalos
más anchos} (menos datos), y la \textbf{amplitud} del conjunto es la
traducción visual de \(\hat\sigma_u\). Un centro con dos pacientes no
dispara un intercepto extremo: se estabiliza.

\begin{center}\rule{0.5\linewidth}{0.5pt}\end{center}

Para comparar finalmente los dos modelos, con y sin efecto aleatorio
para centro, nos basamos en las medidas residuales que obtenemos:

\begin{Shaded}
\begin{Highlighting}[]
\CommentTok{\# 1) Varianza entre centros: sigma\_u\^{}2 (y sigma\_u)}
\NormalTok{vc   }\OtherTok{\textless{}{-}} \FunctionTok{as.data.frame}\NormalTok{(}\FunctionTok{VarCorr}\NormalTok{(m\_int))}
\NormalTok{s2\_u }\OtherTok{\textless{}{-}}\NormalTok{ vc}\SpecialCharTok{$}\NormalTok{vcov[}\DecValTok{1}\NormalTok{]          }\CommentTok{\# sigma\_u\^{}2}
\NormalTok{s\_u  }\OtherTok{\textless{}{-}}\NormalTok{ vc}\SpecialCharTok{$}\NormalTok{sdcor[}\DecValTok{1}\NormalTok{]         }\CommentTok{\# sigma\_u}
\FunctionTok{c}\NormalTok{(}\AttributeTok{sigma2\_u =}\NormalTok{ s2\_u, }\AttributeTok{sigma\_u =}\NormalTok{ s\_u)}
\end{Highlighting}
\end{Shaded}

\begin{verbatim}
 sigma2_u   sigma_u 
0.1959143 0.4426221 
\end{verbatim}

\begin{Shaded}
\begin{Highlighting}[]
\CommentTok{\# 2) Varianza residual (latente, fija) e ICC}
\NormalTok{resid\_var }\OtherTok{\textless{}{-}}\NormalTok{ pi}\SpecialCharTok{\^{}}\DecValTok{2} \SpecialCharTok{/} \DecValTok{3}       \CommentTok{\# \textasciitilde{}3.29  (en logit)}
\NormalTok{icc       }\OtherTok{\textless{}{-}}\NormalTok{ s2\_u }\SpecialCharTok{/}\NormalTok{ (s2\_u }\SpecialCharTok{+}\NormalTok{ resid\_var)}
\FunctionTok{c}\NormalTok{(}\AttributeTok{var\_residual =}\NormalTok{ resid\_var, }\AttributeTok{ICC =}\NormalTok{ icc)}
\end{Highlighting}
\end{Shaded}

\begin{verbatim}
var_residual          ICC 
  3.28986813   0.05620383 
\end{verbatim}

\begin{Shaded}
\begin{Highlighting}[]
\CommentTok{\# (equivalente con paquete:  performance::icc(m\_int)  {-}\textgreater{} \textquotesingle{}adjusted ICC\textquotesingle{})}

\CommentTok{\# 3) AIC de los dos modelos (menor = mejor)}
\FunctionTok{AIC}\NormalTok{(m\_pool, m\_int)          }\CommentTok{\# data.frame con df y AIC de cada uno}
\end{Highlighting}
\end{Shaded}

\begin{verbatim}
       df      AIC
m_pool  3 909.6627
m_int   4 898.9709
\end{verbatim}

\textbf{¿Por qué preferimos el intercepto aleatorio?} Lo dicen sus
propios resultados, sin recurrir a la verdad del proceso generador. El
modelo estima una \textbf{varianza entre centros no nula},
\(\hat\sigma_u^2\approx0{,}19\) (\(\hat\sigma_u\approx0{,}44\)): hay
\textbf{heterogeneidad real} entre clínicas, justo lo que el
\emph{pooled} da por inexistente (\(\sigma_u^2=0\)). En la escala
latente, esa varianza convive con una \textbf{residual fija}
\(\pi^2/3\approx3{,}29\), de modo que el \textbf{ICC} \(\approx0{,}06\):
alrededor de un \textbf{6 \%} de la variación del riesgo es
\textbf{entre} centros y el \textbf{94 \%} restante, \textbf{entre
pacientes}. Es una fracción \textbf{modesta pero suficiente} para
invalidar el supuesto de independencia del \emph{pooled} (de ahí que sus
errores estándar fueran demasiado pequeños). Y la comparación formal
cierra el argumento: el \textbf{AIC del mixto es menor} que el del
\emph{pooled} (898.97 frente a 909.66), así que el único parámetro extra
---la varianza \(\sigma_u^2\)--- \textbf{compensa con creces}. En suma:
aunque el grueso del riesgo se juega a nivel de paciente, existe una
capa de variación \textbf{entre clínicas que es real}; modelarla (mixto)
da una descripción más fiel y una inferencia honesta, mientras que
ignorarla (\emph{pooled}) infla la precisión.

Queda abierta una pregunta: ¿pesa la \textbf{fragilidad ósea} lo mismo
en todas las clínicas, o su efecto sobre el riesgo de fractura es
\textbf{más fuerte en unas que en otras}? Lo vemos a continuación la
\protect\hyperlink{sec-u14-pendiente}{Sección 4.3}.

\hypertarget{sec-u14-pendiente}{%
\subsection{4.3 Pendiente aleatoria: cuando el efecto varía por
grupo}\label{sec-u14-pendiente}}

El intercepto aleatorio deja que el \textbf{nivel} cambie por centro,
pero asume que el \textbf{efecto} de cada covariable es el mismo en
todos. A menudo no lo es: la edad, por ejemplo, puede pesar más en unos
centros que en otros. Eso se modela con una \textbf{pendiente
aleatoria}, añadiendo una desviación de grupo al coeficiente:

\[
\operatorname{logit}\Pr(\text{ever}_{ij}=1)=\beta_0+\beta_1 x_{1ij}+\beta_2 x_{2ij}+u_{0j}+u_{1j}\,x_{1ij},
\qquad
\begin{pmatrix}u_{0j}\\ u_{1j}\end{pmatrix}\sim\mathcal{N}(\mathbf{0},\boldsymbol\Sigma).
\]

Cómo se lee:

\begin{itemize}
\tightlist
\item
  \(\beta_1\) es el efecto \textbf{medio} de \(x_1\) en la población de
  centros; el efecto \textbf{en el centro \(j\)} es \(\beta_1+u_{1j}\).
\item
  \(\sigma_{u_1}^2\) (la varianza de las pendientes) dice \textbf{cuánto
  varía} ese efecto entre centros: si es \(\approx 0\), \(x_1\) actúa
  casi igual en todos; si es grande, en unos centros pesa mucho y en
  otros poco.
\item
  La \textbf{correlación} \(\rho\) entre \(u_{0j}\) y \(u_{1j}\) liga
  nivel y sensibilidad: \(\rho>0\) indica que los centros con mayor
  riesgo basal tienen, además, un efecto de \(x_1\) más fuerte.
\end{itemize}

Las pendientes aleatorias son \textbf{más exigentes} con los datos
(necesitan suficientes grupos y variación de \(x_1\) dentro de cada uno)
y es frecuente que el ajuste no converja o resulte \textbf{singular};
por eso se empieza por el intercepto y se añade la pendiente solo si
está justificada.

En forma matricial, una pendiente aleatoria solo \textbf{añade columnas
a \(\mathbf{Z}\)} (una por efecto aleatorio y grupo) y agranda
\(\mathbf{G}\), que pasa a tener un bloque \(\boldsymbol\Sigma\) de
\(2\times2\) por grupo.

La pendiente aleatoria admite \textbf{dos formas}, y conviene
distinguirlas:

\begin{itemize}
\tightlist
\item
  \textbf{Solo pendiente}, \texttt{(0\ +\ x1\ \textbar{}\ centro)}: la
  \textbf{sensibilidad} a \texttt{x1} varía por centro, pero el
  \textbf{nivel basal} es común a todos. Hipótesis \textbf{poco
  habitual} ---rara vez varía el efecto sin variar el nivel--- y a
  menudo da ajustes \textbf{singulares}.
\item
  \textbf{Intercepto + pendiente},
  \texttt{(1\ +\ x1\ \textbar{}\ centro)}: \textbf{ambos} varían y
  pueden estar \textbf{correlados} (la \(\rho\) entre \(u_{0j}\) y
  \(u_{1j}\)). Es la forma \textbf{estándar}: si dejamos variar la
  pendiente, lo natural es permitir también que el nivel cambie, y
  además, si tenemos ya indicios de interceptación aleatoria de un
  ajuste previo, lo natural es incorporar también la pendiente para
  comprobar si el ajuste mejora.
\end{itemize}

\textbf{¿Tiene sentido en la cohorte?} Ya previamente concluimos con el
sentido de incluir intercepto aleatorio. Veamos si ganamos algo más al
incorporar la pendiente. Para verlo, agrupamos \texttt{x1} en tramos y,
en cada centro, calculamos la proporción de fracturas por tramo: si el
\textbf{efecto} de \texttt{x1} variara por centro, las
\textbf{pendientes} de esas líneas diferirían de un panel a otro.

\begin{Shaded}
\begin{Highlighting}[]
\FunctionTok{library}\NormalTok{(tidyverse)}
\NormalTok{n\_bins }\OtherTok{\textless{}{-}} \DecValTok{4}
\NormalTok{br  }\OtherTok{\textless{}{-}} \FunctionTok{seq}\NormalTok{(}\FunctionTok{min}\NormalTok{(cohorte}\SpecialCharTok{$}\NormalTok{x1), }\FunctionTok{max}\NormalTok{(cohorte}\SpecialCharTok{$}\NormalTok{x1), }\AttributeTok{length.out =}\NormalTok{ n\_bins }\SpecialCharTok{+} \DecValTok{1}\NormalTok{)}
\NormalTok{mid }\OtherTok{\textless{}{-}}\NormalTok{ (}\FunctionTok{head}\NormalTok{(br, }\SpecialCharTok{{-}}\DecValTok{1}\NormalTok{) }\SpecialCharTok{+} \FunctionTok{tail}\NormalTok{(br, }\SpecialCharTok{{-}}\DecValTok{1}\NormalTok{)) }\SpecialCharTok{/} \DecValTok{2}              \CommentTok{\# punto medio de cada bin (posición en X)}
\NormalTok{cohorte }\SpecialCharTok{|\textgreater{}}
  \FunctionTok{mutate}\NormalTok{(}\AttributeTok{bin    =} \FunctionTok{cut}\NormalTok{(x1, }\AttributeTok{breaks =}\NormalTok{ br, }\AttributeTok{include.lowest =} \ConstantTok{TRUE}\NormalTok{),  }\CommentTok{\# bins de igual longitud}
         \AttributeTok{x1\_bin =}\NormalTok{ mid[}\FunctionTok{as.integer}\NormalTok{(bin)]) }\SpecialCharTok{|\textgreater{}}
  \FunctionTok{group\_by}\NormalTok{(centro, x1\_bin) }\SpecialCharTok{|\textgreater{}}
  \FunctionTok{summarise}\NormalTok{(}\AttributeTok{prop =} \FunctionTok{mean}\NormalTok{(evento), }\AttributeTok{.groups =} \StringTok{"drop"}\NormalTok{) }\SpecialCharTok{|\textgreater{}}           \CommentTok{\# proporción de fracturas por centro y bin}
  \FunctionTok{ggplot}\NormalTok{(}\FunctionTok{aes}\NormalTok{(x1\_bin, prop, }\AttributeTok{color =} \FunctionTok{factor}\NormalTok{(centro), }\AttributeTok{group =}\NormalTok{ centro)) }\SpecialCharTok{+}
  \FunctionTok{geom\_line}\NormalTok{(}\AttributeTok{alpha =} \FloatTok{0.5}\NormalTok{) }\SpecialCharTok{+}
  \FunctionTok{geom\_point}\NormalTok{(}\AttributeTok{size =} \FloatTok{1.6}\NormalTok{, }\AttributeTok{alpha =} \FloatTok{0.8}\NormalTok{) }\SpecialCharTok{+}
  \FunctionTok{facet\_wrap}\NormalTok{(}\SpecialCharTok{\textasciitilde{}}\NormalTok{ centro, }\AttributeTok{ncol =} \DecValTok{6}\NormalTok{) }\SpecialCharTok{+}                              \CommentTok{\# 6 columnas × 4 filas}
  \FunctionTok{labs}\NormalTok{(}\AttributeTok{x =} \StringTok{"x1 (fragilidad ósea)"}\NormalTok{, }\AttributeTok{y =} \StringTok{"Proporción de fracturas"}\NormalTok{) }\SpecialCharTok{+}
  \FunctionTok{guides}\NormalTok{(}\AttributeTok{color =} \StringTok{"none"}\NormalTok{)                                        }\CommentTok{\# sin leyenda de centros}
\end{Highlighting}
\end{Shaded}

\begin{figure}[H]

{\centering \includegraphics{caso1_binaria_files/figure-pdf/fig-u14-x1-centro-1.pdf}

}

\caption{\label{fig-u14-x1-centro}Proporción de fracturas frente a la
fragilidad ósea (x1), en bins de igual longitud, con una faceta por
centro (rejilla 6×4). La tendencia creciente común refleja el efecto
(fijo) de x1; las diferencias de nivel entre paneles, el intercepto
aleatorio. Con pocas pacientes por centro y bin las curvas son ruidosas,
así que las diferencias de pendiente no son interpretables a ojo.}

\end{figure}

En la Figura~\ref{fig-u14-x1-centro} se aprecia con claridad una
\textbf{tendencia creciente común} ---el efecto de \texttt{x1},
compartido por todos los centros--- y \textbf{diferencias de nivel}
entre paneles (unos sistemáticamente más altos): la huella del
\textbf{intercepto aleatorio}. Las diferencias de \textbf{pendiente}
(algún panel parece plano o decreciente) son, en cambio,
\textbf{difíciles de separar del ruido de muestreo}: con tan pocas
pacientes por centro y tramo, cada punto se apoya en pocos casos y
oscila con facilidad. A ojo no podemos decidir si hay pendientes
aleatorias reales; lo \textbf{contrastaremos formalmente} a
continuación, ajustando los modelos.

\begin{tcolorbox}[enhanced jigsaw, coltitle=black, rightrule=.15mm, arc=.35mm, leftrule=.75mm, colback=white, colbacktitle=quarto-callout-tip-color!10!white, left=2mm, breakable, toptitle=1mm, toprule=.15mm, bottomtitle=1mm, titlerule=0mm, title=\textcolor{quarto-callout-tip-color}{\faLightbulb}\hspace{0.5em}{🔧 En R. Añadir una pendiente aleatoria}, bottomrule=.15mm, opacityback=0, opacitybacktitle=0.6, colframe=quarto-callout-tip-color-frame]

\begin{Shaded}
\begin{Highlighting}[]
\NormalTok{m\_slope0 }\OtherTok{\textless{}{-}} \FunctionTok{glmer}\NormalTok{(ever }\SpecialCharTok{\textasciitilde{}}\NormalTok{ x1 }\SpecialCharTok{+}\NormalTok{ x2 }\SpecialCharTok{+}\NormalTok{ (}\DecValTok{0} \SpecialCharTok{+}\NormalTok{ x1 }\SpecialCharTok{|}\NormalTok{ centro), }\AttributeTok{family =}\NormalTok{ binomial, }\AttributeTok{data =}\NormalTok{ cohorte)  }\CommentTok{\# solo pendiente}
\NormalTok{m\_slope  }\OtherTok{\textless{}{-}} \FunctionTok{glmer}\NormalTok{(ever }\SpecialCharTok{\textasciitilde{}}\NormalTok{ x1 }\SpecialCharTok{+}\NormalTok{ x2 }\SpecialCharTok{+}\NormalTok{ (}\DecValTok{1} \SpecialCharTok{+}\NormalTok{ x1 }\SpecialCharTok{|}\NormalTok{ centro), }\AttributeTok{family =}\NormalTok{ binomial, }\AttributeTok{data =}\NormalTok{ cohorte)  }\CommentTok{\# intercepto + pendiente}
\FunctionTok{VarCorr}\NormalTok{(m\_slope)   }\CommentTok{\# sigma\_u0, sigma\_u1 y su correlacion rho}
\end{Highlighting}
\end{Shaded}

\end{tcolorbox}

Añadimos pues pendiente aleatoria al modelo.

\begin{Shaded}
\begin{Highlighting}[]
\NormalTok{m\_slope  }\OtherTok{\textless{}{-}} \FunctionTok{glmer}\NormalTok{(ever }\SpecialCharTok{\textasciitilde{}}\NormalTok{ x1 }\SpecialCharTok{+}\NormalTok{ x2 }\SpecialCharTok{+}\NormalTok{ (}\DecValTok{1} \SpecialCharTok{+}\NormalTok{ x1 }\SpecialCharTok{|}\NormalTok{ centro), }\AttributeTok{family =}\NormalTok{ binomial, }\AttributeTok{data =}\NormalTok{ cohorte)  }\CommentTok{\# intercepto + pendiente}

\FunctionTok{VarCorr}\NormalTok{(m\_slope)        }\CommentTok{\# sigma\_u0, sigma\_u1 y su correlacion}
\end{Highlighting}
\end{Shaded}

\begin{verbatim}
 Groups Name        Std.Dev. Corr   
 centro (Intercept) 0.53421         
        x1          0.39714  -0.797 
\end{verbatim}

En el modelo \textbf{intercepto + pendiente}, \texttt{VarCorr} da una
desviación de pendiente \(\hat\sigma_{u_1}\approx0{,}40\) con una
\textbf{correlación casi extrema} entre los efectos aleatorios
(intercepto y pendiente) \(\rho\approx-0{,}80\) ---señales de un ajuste
\textbf{inestable}, frecuente cuando la pendiente no aporta señal
real---. La pregunta, entonces, no es \emph{¿se puede?} sino
\emph{¿merece la pena?}; y eso se decide \textbf{comparando}.

\textbf{La estrategia de selección, sistematizada.}

Con efectos aleatorios se procede \textbf{de menos a más} y se contrasta
cada salto:

\begin{enumerate}
\def\labelenumi{\arabic{enumi}.}
\tightlist
\item
  \textbf{Empezar por el intercepto} (\texttt{(1\ \textbar{}\ g)}): casi
  siempre justificado cuando hay agrupamiento.
\item
  \textbf{Añadir la pendiente} (\texttt{(1\ +\ x\ \textbar{}\ g)})
  \textbf{solo si} mejora el ajuste de forma apreciable.
\item
  \textbf{Comparar} por criterios de información (\textbf{AIC/BIC},
  menor es mejor) y por \textbf{LRT} (\texttt{anova}) entre modelos
  anidados; vigilar avisos de \textbf{ajuste singular} o de no
  convergencia.
\item
  Ante la duda, \textbf{parsimonia}: el modelo más simple que describe
  bien los datos.
\end{enumerate}

\begin{tcolorbox}[enhanced jigsaw, coltitle=black, rightrule=.15mm, arc=.35mm, leftrule=.75mm, colback=white, colbacktitle=quarto-callout-tip-color!10!white, left=2mm, breakable, toptitle=1mm, toprule=.15mm, bottomtitle=1mm, titlerule=0mm, title=\textcolor{quarto-callout-tip-color}{\faLightbulb}\hspace{0.5em}{🔧 En R. Comparar y elegir}, bottomrule=.15mm, opacityback=0, opacitybacktitle=0.6, colframe=quarto-callout-tip-color-frame]

\begin{Shaded}
\begin{Highlighting}[]
\FunctionTok{AIC}\NormalTok{(m\_pool, m\_int, m\_slope)   }\CommentTok{\# escalera de complejidad (menor AIC, mejor)}
\FunctionTok{anova}\NormalTok{(m\_int, m\_slope)          }\CommentTok{\# LRT: ¿aporta la pendiente sobre el intercepto?}
\end{Highlighting}
\end{Shaded}

\end{tcolorbox}

\begin{Shaded}
\begin{Highlighting}[]
\FunctionTok{AIC}\NormalTok{(m\_pool, m\_int, m\_slope)   }\CommentTok{\# pooled  \textless{}  intercepto  \textless{}  intercepto + pendiente}
\end{Highlighting}
\end{Shaded}

\begin{verbatim}
        df      AIC
m_pool   3 909.6627
m_int    4 898.9709
m_slope  6 898.0549
\end{verbatim}

\begin{Shaded}
\begin{Highlighting}[]
\FunctionTok{anova}\NormalTok{(m\_int, m\_slope)          }\CommentTok{\# LRT de la pendiente}
\end{Highlighting}
\end{Shaded}

\begin{verbatim}
Data: cohorte
Models:
m_int: ever ~ x1 + x2 + (1 | centro)
m_slope: ever ~ x1 + x2 + (1 + x1 | centro)
        npar    AIC    BIC  logLik -2*log(L) Chisq Df Pr(>Chisq)  
m_int      4 898.97 918.08 -445.49    890.97                      
m_slope    6 898.05 926.71 -443.03    886.05 4.916  2    0.08561 .
---
Signif. codes:  0 '***' 0.001 '**' 0.01 '*' 0.05 '.' 0.1 ' ' 1
\end{verbatim}

Los números \textbf{deciden}: el \texttt{AIC} apenas cambia al añadir la
pendiente (\(899{,}0\to898{,}1\), despreciable) y el \textbf{LRT} frente
a \texttt{m\_int} no es significativo (\(p\approx0{,}08\), 2 g.l.). Como
el DGP \textbf{no} incluye pendiente aleatoria, ese
\(\hat\sigma_{u_1}=0{,}40\) era \textbf{ruido sobreajustado}.

\textbf{Conclusión:} nos quedamos con el \textbf{intercepto aleatorio}
(\texttt{m\_int}) ---mejora claramente al \emph{pooled}
(\protect\hyperlink{sec-u14-intercepto}{Sección 4.2}) y la pendiente no
se justifica---.

\hypertarget{sec-u14-diagnostico}{%
\subsection{4.4 Bondad de ajuste y diagnóstico del modelo
mixto}\label{sec-u14-diagnostico}}

Ajustado el modelo, ¿\textbf{ajusta bien}? Un GLMM hereda las
dificultades del GLM binario ---el \(\chi^2\) de la deviance no es un
test de bondad de ajuste válido con respuesta \textbf{individual}
(\protect\hyperlink{sec-u13-gof-binomial}{Sección 3.2})--- y añade la
del efecto aleatorio: los residuos ordinarios (Pearson, deviance) se
\textbf{enredan} con los \(u_j\) y dejan de ser informativos. El
instrumental, por tanto, es el que afinamos en la \textbf{Unidad 2},
ahora \textbf{sobre el modelo mixto}: lo que se \textbf{extrae} del
ajuste, los \textbf{residuos simulados}, una medida de \textbf{varianza
explicada} adaptada a los mixtos y unos \textbf{chequeos del propio
ajuste}.

\textbf{Lo que se extrae del ajuste.}

Casi todo el diagnóstico se apoya en cantidades que se sacan del modelo
ya ajustado ---coeficientes, varianzas, valores ajustados, predicciones
y residuos---:

\begin{tcolorbox}[enhanced jigsaw, coltitle=black, rightrule=.15mm, arc=.35mm, leftrule=.75mm, colback=white, colbacktitle=quarto-callout-tip-color!10!white, left=2mm, breakable, toptitle=1mm, toprule=.15mm, bottomtitle=1mm, titlerule=0mm, title=\textcolor{quarto-callout-tip-color}{\faLightbulb}\hspace{0.5em}{🔧 En R. Extraer de un GLMM ajustado \texttt{m}}, bottomrule=.15mm, opacityback=0, opacitybacktitle=0.6, colframe=quarto-callout-tip-color-frame]

\emph{Coeficientes y varianzas:}

\begin{itemize}
\tightlist
\item
  \texttt{fixef(m)} --- efectos \textbf{fijos} (\(\beta\));
  \texttt{vcov(m)}, su matriz de covarianzas (de donde salen los EE).
\item
  \texttt{ranef(m)} --- efectos \textbf{aleatorios} predichos (\(u_j\),
  encogidos); \texttt{coef(m)} --- el coeficiente \textbf{por grupo}
  (fijo + aleatorio).
\item
  \texttt{VarCorr(m)} --- \textbf{varianzas y correlaciones} de los
  efectos aleatorios (\(\sigma_u\), \(\rho\)).
\item
  \texttt{confint(m,\ method\ =\ "Wald")} (o \texttt{"profile"}, más
  fiable) --- IC de los \(\beta\) y de \(\sigma_u\).
\end{itemize}

\emph{Ajustes y predicciones:}

\begin{itemize}
\tightlist
\item
  \texttt{fitted(m)} --- probabilidades \textbf{ajustadas} (incluyen el
  efecto aleatorio de cada grupo).
\item
  \texttt{predict(m,\ type\ =\ "response")} --- predicción en
  probabilidad; el argumento \textbf{\texttt{re.form}} es la clave:

  \begin{itemize}
  \tightlist
  \item
    \texttt{re.form\ =\ NULL} (por defecto) → \textbf{condicional} al
    grupo (usa su \(u_j\)).
  \item
    \texttt{re.form\ =\ NA} (o \texttt{\textasciitilde{}0}) →
    \textbf{marginal/poblacional} (ignora los efectos aleatorios: «un
    centro medio»).
  \end{itemize}
\item
  \texttt{predict(m,\ newdata\ =\ nd,\ type\ =\ "response",\ re.form\ =\ \textasciitilde{}0)}
  --- para \textbf{perfiles nuevos} o centros no observados.
\item
  \texttt{residuals(m,\ type\ =\ "pearson")} (o \texttt{"deviance"}).
\end{itemize}

\emph{Ajuste global:}

\begin{itemize}
\tightlist
\item
  \texttt{summary(m)} --- todo junto (efectos fijos + bloque
  \texttt{Random\ effects}); \texttt{AIC(m)}, \texttt{BIC(m)},
  \texttt{logLik(m)}, \texttt{deviance(m)}.
\item
  Salida ordenada con \texttt{broom.mixed::tidy(m)}, \texttt{glance(m)},
  \texttt{augment(m)}.
\end{itemize}

\end{tcolorbox}

\textbf{Residuos simulados (DHARMa).}

Con respuesta binaria y efectos aleatorios, los residuos de Pearson o
deviance crudos \textbf{no} se leen (el caso individual de la
\protect\hyperlink{sec-u12-diagnostico}{Sección 2.6}). La solución, ya
conocida, son los \textbf{residuos simulados} de \texttt{DHARMa}
(Hartig, 2024): el paquete \textbf{simula} muchas respuestas desde el
modelo ajustado y, comparando cada observación con su distribución
simulada, produce \textbf{residuos cuantil escalados} que, bajo un
modelo correcto, son \textbf{uniformes} \(U(0,1)\) ---y eso vale
\textbf{igual para un GLMM}---. Estos residuos nos permiten testar
uniformidad de los residuos (con Kolmogorov-Smirnov),
sobre/infra-dispersión (se testa que el parámetro de dispersión es 1),
exceso de valores atípicos y heterocedasticidad.

\begin{tcolorbox}[enhanced jigsaw, coltitle=black, rightrule=.15mm, arc=.35mm, leftrule=.75mm, colback=white, colbacktitle=quarto-callout-tip-color!10!white, left=2mm, breakable, toptitle=1mm, toprule=.15mm, bottomtitle=1mm, titlerule=0mm, title=\textcolor{quarto-callout-tip-color}{\faLightbulb}\hspace{0.5em}{🔧 En R. Residuos simulados (DHARMa) en modelos mixtos}, bottomrule=.15mm, opacityback=0, opacitybacktitle=0.6, colframe=quarto-callout-tip-color-frame]

\texttt{DHARMa} genera \textbf{residuos cuantil escalados} ---uniformes
\(U(0,1)\) bajo un modelo correcto--- por \textbf{simulación}, y
funciona igual para \texttt{glmer}/\texttt{glmmTMB} que para
\texttt{glm}. El flujo es: simular una vez y reutilizar el objeto.

\begin{Shaded}
\begin{Highlighting}[]
\FunctionTok{library}\NormalTok{(DHARMa)}
\NormalTok{sim }\OtherTok{\textless{}{-}} \FunctionTok{simulateResiduals}\NormalTok{(m, }\AttributeTok{n =} \DecValTok{1000}\NormalTok{)   }\CommentTok{\# m = modelo glmer/glmmTMB}
\FunctionTok{plot}\NormalTok{(sim)                                }\CommentTok{\# Q{-}Q (uniformidad, KS) + residuos vs predicho}
\end{Highlighting}
\end{Shaded}

\textbf{Tests formales} sobre \texttt{sim}:

\begin{Shaded}
\begin{Highlighting}[]
\FunctionTok{testUniformity}\NormalTok{(sim)     }\CommentTok{\# KS: ¿residuos uniformes? (ajuste global)}
\FunctionTok{testDispersion}\NormalTok{(sim)     }\CommentTok{\# sobre/infradispersión (cociente \textasciitilde{}1 = ok)}
\FunctionTok{testOutliers}\NormalTok{(sim)       }\CommentTok{\# exceso de valores atípicos}
\FunctionTok{testQuantiles}\NormalTok{(sim)      }\CommentTok{\# patrón / heterocedasticidad (cuantiles vs predicho)}
\end{Highlighting}
\end{Shaded}

\textbf{Lo propio de un mixto:}

\begin{itemize}
\tightlist
\item
  \textbf{Residuos frente a una covariable o al grupo}, para detectar
  estructura no modelada (no linealidad, un centro desviado):
\end{itemize}

\begin{Shaded}
\begin{Highlighting}[]
  \FunctionTok{plotResiduals}\NormalTok{(sim, }\AttributeTok{form =}\NormalTok{ datos}\SpecialCharTok{$}\NormalTok{x1)       }\CommentTok{\# vs un predictor {-}\textgreater{} ¿forma funcional?}
  \FunctionTok{plotResiduals}\NormalTok{(sim, }\AttributeTok{form =}\NormalTok{ datos}\SpecialCharTok{$}\NormalTok{centro)   }\CommentTok{\# vs el factor de agrupamiento}
\end{Highlighting}
\end{Shaded}

\begin{itemize}
\tightlist
\item
  \textbf{Agregando por grupo} ---un residuo por centro, para mirar el
  nivel del efecto aleatorio---:
\end{itemize}

\begin{Shaded}
\begin{Highlighting}[]
\NormalTok{  sim\_g }\OtherTok{\textless{}{-}} \FunctionTok{recalculateResiduals}\NormalTok{(sim, }\AttributeTok{group =}\NormalTok{ datos}\SpecialCharTok{$}\NormalTok{centro)}
  \FunctionTok{plot}\NormalTok{(sim\_g)}
\end{Highlighting}
\end{Shaded}

\begin{itemize}
\tightlist
\item
  DHARMa calcula los residuos \textbf{condicionando} a los efectos
  aleatorios estimados (lo habitual); para diagnosticar la
  \textbf{estructura aleatoria} en sí, la agregación por grupo de arriba
  es la vía práctica.
\end{itemize}

Notas: con respuesta \textbf{binaria} o de \textbf{conteo} (muchos
empates) DHARMa lo gestiona con \emph{randomización}
(\texttt{integerResponse\ =\ TRUE}, por defecto lo detecta); y
\texttt{refit\ =\ TRUE} calcula los residuos por \textbf{bootstrap
paramétrico} ---más exacto con \textbf{pocos grupos}, pero lento---.

\end{tcolorbox}

Se revisan pues el \textbf{Q--Q} (uniformidad, test KS), los
\textbf{residuos frente a lo predicho} (patrones, dispersión,
\emph{outliers}) y, lo propio de un mixto, \textbf{agrupando por centro}
para mirar el nivel del efecto aleatorio:

\begin{Shaded}
\begin{Highlighting}[]
\FunctionTok{library}\NormalTok{(DHARMa)}
\NormalTok{sim }\OtherTok{\textless{}{-}} \FunctionTok{simulateResiduals}\NormalTok{(m\_int)}
\FunctionTok{plot}\NormalTok{(sim)}
\end{Highlighting}
\end{Shaded}

\begin{figure}[H]

{\centering \includegraphics{caso1_binaria_files/figure-pdf/fig-u14-dharma-1.pdf}

}

\caption{\label{fig-u14-dharma}Residuos simulados (DHARMa) del GLMM de
intercepto aleatorio. Izquierda: Q--Q de los residuos escalados, que
bajo un buen ajuste siguen la diagonal uniforme. Derecha: residuos
frente a lo predicho, que deberían repartirse sin patrón. Curvatura en
el Q--Q o un abanico a la derecha señalarían forma funcional inadecuada
o dispersión mal modelada.}

\end{figure}

\begin{Shaded}
\begin{Highlighting}[]
\FunctionTok{testDispersion}\NormalTok{(sim)                          }\CommentTok{\# sobre/infradispersión}
\end{Highlighting}
\end{Shaded}

\begin{figure}[H]

{\centering \includegraphics{caso1_binaria_files/figure-pdf/u14-dharma-tests-1.pdf}

}

\end{figure}

\begin{verbatim}

    DHARMa nonparametric dispersion test via sd of residuals fitted vs.
    simulated

data:  simulationOutput
dispersion = 0.99685, p-value = 0.976
alternative hypothesis: two.sided
\end{verbatim}

\begin{Shaded}
\begin{Highlighting}[]
\FunctionTok{plotResiduals}\NormalTok{(sim, }\AttributeTok{form =}\NormalTok{ cohorte}\SpecialCharTok{$}\NormalTok{centro)    }\CommentTok{\# residuos agregados por centro}
\end{Highlighting}
\end{Shaded}

\begin{figure}[H]

{\centering \includegraphics{caso1_binaria_files/figure-pdf/u14-dharma-tests-2.pdf}

}

\end{figure}

\textbf{Conclusión: no hay problemas.} El Q--Q se pega a la diagonal y
los tres tests (Kolmogorov-Smirnov, Dispersion y outliers) salen
\textbf{no significativos} (KS \(p=0{,}71\), dispersión \(p=0{,}98\),
\emph{outliers} \(p=0{,}18\); DHARMa rotula \emph{«no significant
problems detected»}) ---el intercepto aleatorio reproduce bien los
datos, también al \textbf{agrupar por centro}---. Si fallara, lo
veríamos como una \textbf{curvatura} sistemática en el Q--Q,
\textbf{sobredispersión} o un grupo de centros desviado, que apuntarían
a una \textbf{covariable omitida}, una \textbf{forma funcional}
inadecuada o una \textbf{estructura aleatoria} insuficiente.

\textbf{¿Cuánto explica el modelo? \(R^2\) marginal y condicional.}

En un modelo mixto la varianza explicada se parte en dos (Nakagawa): el
\(R^2\) \textbf{marginal} ---solo los \textbf{efectos fijos}--- y el
\(R^2\) \textbf{condicional} ---fijos \textbf{más} aleatorios---. La
\textbf{diferencia}
\(R^2_{\text{condicional}}-R^2_{\text{marginal}}=\dfrac{\sigma^2_u}{\sigma^2_f+\sigma^2_u+\sigma^2_\varepsilon}\)
\textbf{es el ICC no ajustado}: la fracción de la varianza total que
aporta el \textbf{agrupamiento} (los centros) más allá de los efectos
fijos. Es justo el \emph{Unadjusted/Conditional ICC} que devuelve
\texttt{performance::icc(m)}.

\begin{Shaded}
\begin{Highlighting}[]
\NormalTok{performance}\SpecialCharTok{::}\FunctionTok{r2}\NormalTok{(m\_int)     }\CommentTok{\# R2 marginal (fijos) y condicional (fijos + aleatorios)}
\end{Highlighting}
\end{Shaded}

\begin{verbatim}
# R2 for Mixed Models

  Conditional R2: 0.311
     Marginal R2: 0.270
\end{verbatim}

\begin{Shaded}
\begin{Highlighting}[]
\NormalTok{performance}\SpecialCharTok{::}\FunctionTok{icc}\NormalTok{(m\_int)    }\CommentTok{\# icc no ajustado = R2(cond){-}R2(marg)}
\end{Highlighting}
\end{Shaded}

\begin{verbatim}
# Intraclass Correlation Coefficient

    Adjusted ICC: 0.056
  Unadjusted ICC: 0.041
\end{verbatim}

En la cohorte, \(R^2_{\text{marginal}}=0{,}27\) y
\(R^2_{\text{condicional}}=0{,}31\): los efectos fijos (\texttt{x1},
\texttt{x2}) explican el \textbf{grueso} (≈ 27 \%) y la estructura de
centro \textbf{añade unos 4 puntos} (hasta el 31 \%) ---una contribución
\textbf{modesta}, coherente con el ICC pequeño---.

\textbf{El ajuste, ¿es sólido? Singularidad y convergencia.}

Un GLMM puede \textbf{converger mal} o quedar \textbf{singular}
---típico al pedir más estructura aleatoria de la que los datos
sostienen, como la pendiente casi singular de la
\protect\hyperlink{sec-u14-pendiente}{Sección 4.3}---. Conviene
comprobarlo siempre:

\begin{Shaded}
\begin{Highlighting}[]
\FunctionTok{isSingular}\NormalTok{(m\_int)                    }\CommentTok{\# TRUE = ajuste singular (estructura aleatoria no sostenida)}
\end{Highlighting}
\end{Shaded}

\begin{verbatim}
[1] FALSE
\end{verbatim}

\begin{Shaded}
\begin{Highlighting}[]
\NormalTok{performance}\SpecialCharTok{::}\FunctionTok{check\_singularity}\NormalTok{(m\_int)}
\end{Highlighting}
\end{Shaded}

\begin{verbatim}
[1] FALSE
\end{verbatim}

Aquí \texttt{isSingular(m\_int)} devuelve \textbf{\texttt{FALSE}}: el
intercepto aleatorio está \textbf{bien sostenido} por los datos (a
diferencia de lo que ocurría al forzar la pendiente en la
\protect\hyperlink{sec-u14-pendiente}{Sección 4.3}). Un ajuste
\textbf{singular} (o con avisos de convergencia) \textbf{no sería de
fiar} y obligaría a \textbf{simplificar} la parte aleatoria. Para una
panorámica de golpe, \texttt{performance::check\_model(m\_int)} reúne
varios diagnósticos ---incluida la \textbf{normalidad de los efectos
aleatorios} (\(u_j\sim\mathcal N\))--- en un solo gráfico.

\begin{tcolorbox}[enhanced jigsaw, coltitle=black, rightrule=.15mm, arc=.35mm, leftrule=.75mm, colback=white, colbacktitle=quarto-callout-tip-color!10!white, left=2mm, breakable, toptitle=1mm, toprule=.15mm, bottomtitle=1mm, titlerule=0mm, title=\textcolor{quarto-callout-tip-color}{\faLightbulb}\hspace{0.5em}{🔧 En R. Bondad de ajuste y solidez de un GLMM}, bottomrule=.15mm, opacityback=0, opacitybacktitle=0.6, colframe=quarto-callout-tip-color-frame]

Para los \textbf{residuos}, el tip de DHARMa de arriba. El resto del
diagnóstico:

\begin{Shaded}
\begin{Highlighting}[]
\CommentTok{\# Varianza explicada (Nakagawa)}
\NormalTok{performance}\SpecialCharTok{::}\FunctionTok{r2}\NormalTok{(m\_int)        }\CommentTok{\# R2 marginal (fijos) y condicional (fijos + aleatorios)}
\NormalTok{performance}\SpecialCharTok{::}\FunctionTok{icc}\NormalTok{(m\_int)       }\CommentTok{\# ICC ajustado y no ajustado (= R2\_cond {-} R2\_marg)}

\CommentTok{\# Solidez del ajuste}
\NormalTok{performance}\SpecialCharTok{::}\FunctionTok{check\_model}\NormalTok{(m\_int)   }\CommentTok{\# panel de diagnósticos del mixto}
\FunctionTok{isSingular}\NormalTok{(m\_int)                 }\CommentTok{\# ¿ajuste singular? (estructura aleatoria no sostenida)}
\end{Highlighting}
\end{Shaded}

\end{tcolorbox}

\hypertarget{sec-u14-extensiones}{%
\subsection{4.5 Extensiones: estructuras aleatorias más
ricas}\label{sec-u14-extensiones}}

La parte aleatoria puede crecer más allá de un intercepto. Conviene
reconocer las estructuras habituales ---aunque en la cohorte, con
pacientes simplemente agrupadas en centros, no todas tengan cabida---
porque aparecen en cuanto el diseño se complica.

\textbf{Correlados vs no correlados.} Al dejar variar intercepto y
pendiente, \texttt{(1\ +\ x\ \textbar{}\ g)} estima además su
\textbf{correlación} \(\rho\). Que \(\rho>0\) implicaría que los centros
de mayor riesgo basal responden, además, más a \texttt{x}. Si esa
correlación no interesa o \textbf{desestabiliza} el ajuste
(correlaciones de \(\pm1\), ajustes singulares), la especificación
\texttt{(1\ +\ x\ \textbar{}\textbar{}\ g)} los deja
\textbf{independientes} (sin \(\rho\)), que es un modelo más parco, pero
recomendable cuando no tenemos bastante información (datos) para
estimarla.

\textbf{Cruzados vs anidados.} Dos agrupamientos pueden combinarse de
dos maneras, y la diferencia es \textbf{de diseño}, no de sintaxis:

\begin{itemize}
\tightlist
\item
  \textbf{Anidado:} cada nivel de un factor pertenece a \textbf{uno
  solo} del otro. Pacientes ⊂ \textbf{centros} ⊂ \textbf{regiones}: un
  centro está en una única región. Se escribe
  \texttt{(1\ \textbar{}\ region/centro)}. Y \textbf{un anidamiento lo
  es por diseño}: si los centros se numeran \(1,2,\dots\)
  \textbf{dentro} de cada región (de modo que el «centro 1» de una
  región no es el de otra), hay que declararlo anidado ---o usar
  identificadores \textbf{únicos}--- o \texttt{lme4} confundiría centros
  distintos que comparten etiqueta.
\item
  \textbf{Cruzado:} los factores se combinan \textbf{libremente} y
  ninguno anida en el otro ---una misma paciente valorada por
  \textbf{varios médicos} (médicos y pacientes se cruzan), o la rejilla
  \textbf{centro × periodo}---. Se escribe
  \texttt{(1\ \textbar{}\ g1)\ +\ (1\ \textbar{}\ g2)}.
\end{itemize}

Regla práctica:

\begin{itemize}
\tightlist
\item
  \emph{¿cada nivel de A vive dentro de un único nivel de B?} →
  \textbf{anidado};
\item
  \emph{¿se combinan sin restricción?} → \textbf{cruzado}.
\end{itemize}

\begin{tcolorbox}[enhanced jigsaw, coltitle=black, rightrule=.15mm, arc=.35mm, leftrule=.75mm, colback=white, colbacktitle=quarto-callout-tip-color!10!white, left=2mm, breakable, toptitle=1mm, toprule=.15mm, bottomtitle=1mm, titlerule=0mm, title=\textcolor{quarto-callout-tip-color}{\faLightbulb}\hspace{0.5em}{🔧 En R. La sintaxis \texttt{(\ \textbar{}\ )} para cada estructura}, bottomrule=.15mm, opacityback=0, opacitybacktitle=0.6, colframe=quarto-callout-tip-color-frame]

Un GLMM binomial se ajusta con \textbf{\texttt{lme4::glmer}} (o
\texttt{glmmTMB::glmmTMB}, más rápido y flexible), con
\texttt{family\ =\ binomial}, \textbf{igual para datos individuales y
agrupados} ---solo cambia la respuesta, como en \texttt{glm}---:

\begin{itemize}
\tightlist
\item
  \textbf{Individual} (0/1 por sujeto):
  \texttt{glmer(y\ \textasciitilde{}\ x1\ +\ x2\ +\ (1\ \textbar{}\ g),\ family\ =\ binomial,\ data\ =\ d)}.
\item
  \textbf{Agrupado} (éxitos/ensayos por fila):
  \texttt{glmer(cbind(exitos,\ fracasos)\ \textasciitilde{}\ x1\ +\ x2\ +\ (1\ \textbar{}\ g),\ family\ =\ binomial,\ data\ =\ d)}
  (o la proporción con \texttt{weights\ =\ m}).
\end{itemize}

La parte \texttt{(\ ...\ \textbar{}\ grupo\ )} declara \textbf{qué}
varía por grupo:

\begin{itemize}
\tightlist
\item
  \texttt{(1\ \textbar{}\ g)} --- \textbf{intercepto} aleatorio por
  \texttt{g}.
\item
  \texttt{(1\ +\ x\ \textbar{}\ g)} (o \texttt{(x\ \textbar{}\ g)}) ---
  intercepto \textbf{y} pendiente de \texttt{x}, \textbf{correlados}.
\item
  \texttt{(0\ +\ x\ \textbar{}\ g)} --- \textbf{solo} pendiente de
  \texttt{x} (sin intercepto aleatorio).
\item
  \texttt{(1\ +\ x\ \textbar{}\textbar{}\ g)} --- intercepto y
  pendiente, \textbf{sin} correlación.
\item
  \texttt{(1\ \textbar{}\ g1)\ +\ (1\ \textbar{}\ g2)} --- dos
  agrupamientos \textbf{cruzados}; \texttt{(1\ \textbar{}\ g1/g2)} para
  \textbf{anidados}.
\end{itemize}

\end{tcolorbox}

En la cohorte la estructura es \textbf{simple} ---pacientes en centros,
un solo nivel---, así que estas extensiones quedan como repertorio para
diseños más ricos (las reactivaremos en el Caso 2). El \textbf{protocolo
de selección} de la \protect\hyperlink{sec-u14-pendiente}{Sección 4.3}
sigue valiendo: empezar simple, añadir estructura solo si AIC/LRT la
respaldan, y vigilar los ajustes singulares.

\hypertarget{sec-u14-glmm-glow}{%
\subsection{4.6 El modelo mixto sobre datos reales: GLOW por
centro}\label{sec-u14-glmm-glow}}

La cohorte tiene la ventaja de la \textbf{verdad conocida}, pero 24
centros son muchos para \emph{verlos}. GLOW ofrece el contrapunto
\textbf{real e intuitivo}: 500 mujeres seguidas en \textbf{6 centros}
(\texttt{site\_id}). Las pacientes de un mismo centro comparten contexto
asistencial, así que cabe el \textbf{mismo GLMM logístico} con
intercepto aleatorio por centro, ahora sobre la respuesta real
\texttt{fracture}:

\begin{Shaded}
\begin{Highlighting}[]
\FunctionTok{library}\NormalTok{(lme4)}
\NormalTok{glow\_m }\OtherTok{\textless{}{-}}\NormalTok{ glow }\SpecialCharTok{|\textgreater{}} \FunctionTok{mutate}\NormalTok{(}\AttributeTok{site\_id =} \FunctionTok{factor}\NormalTok{(site\_id))}

\NormalTok{m\_glow\_pool }\OtherTok{\textless{}{-}} \FunctionTok{glm}\NormalTok{  (fracture }\SpecialCharTok{\textasciitilde{}}\NormalTok{ age }\SpecialCharTok{+}\NormalTok{ priorfrac,                 }\AttributeTok{family =}\NormalTok{ binomial, }\AttributeTok{data =}\NormalTok{ glow\_m)}
\NormalTok{m\_glow\_mix  }\OtherTok{\textless{}{-}} \FunctionTok{glmer}\NormalTok{(fracture }\SpecialCharTok{\textasciitilde{}}\NormalTok{ age }\SpecialCharTok{+}\NormalTok{ priorfrac }\SpecialCharTok{+}\NormalTok{ (}\DecValTok{1} \SpecialCharTok{|}\NormalTok{ site\_id), }\AttributeTok{family =}\NormalTok{ binomial, }\AttributeTok{data =}\NormalTok{ glow\_m)}

\FunctionTok{summary}\NormalTok{(m\_glow\_mix)}
\end{Highlighting}
\end{Shaded}

\begin{verbatim}
Generalized linear mixed model fit by maximum likelihood (Laplace
  Approximation) [glmerMod]
 Family: binomial  ( logit )
Formula: fracture ~ age + priorfrac + (1 | site_id)
   Data: glow_m

      AIC       BIC    logLik -2*log(L)  df.resid 
    536.2     553.1    -264.1     528.2       496 

Scaled residuals: 
     Min       1Q   Median       3Q      Max 
-1.21215 -0.55258 -0.43669 -0.06166  2.62011 

Random effects:
 Groups  Name        Variance Std.Dev.
 site_id (Intercept) 0.02237  0.1496  
Number of obs: 500, groups:  site_id, 6

Fixed effects:
             Estimate Std. Error z value Pr(>|z|)    
(Intercept)  -4.20529    0.85331  -4.928  8.3e-07 ***
age           0.04085    0.01223   3.341 0.000835 ***
priorfracYes  0.84935    0.23558   3.605 0.000312 ***
---
Signif. codes:  0 '***' 0.001 '**' 0.01 '*' 0.05 '.' 0.1 ' ' 1

Correlation of Fixed Effects:
            (Intr) age   
age         -0.985       
priorfracYs  0.165 -0.258
\end{verbatim}

\begin{Shaded}
\begin{Highlighting}[]
\CommentTok{\# Varianza entre centros e ICC latente}
\NormalTok{vc }\OtherTok{\textless{}{-}} \FunctionTok{as.data.frame}\NormalTok{(}\FunctionTok{VarCorr}\NormalTok{(m\_glow\_mix))}
\FunctionTok{c}\NormalTok{(}\AttributeTok{sigma\_u =} \FunctionTok{sqrt}\NormalTok{(vc}\SpecialCharTok{$}\NormalTok{vcov[}\DecValTok{1}\NormalTok{]), }\AttributeTok{ICC =}\NormalTok{ vc}\SpecialCharTok{$}\NormalTok{vcov[}\DecValTok{1}\NormalTok{] }\SpecialCharTok{/}\NormalTok{ (vc}\SpecialCharTok{$}\NormalTok{vcov[}\DecValTok{1}\NormalTok{] }\SpecialCharTok{+}\NormalTok{ pi}\SpecialCharTok{\^{}}\DecValTok{2} \SpecialCharTok{/} \DecValTok{3}\NormalTok{))}
\end{Highlighting}
\end{Shaded}

\begin{verbatim}
    sigma_u         ICC 
0.149575578 0.006754597 
\end{verbatim}

\begin{Shaded}
\begin{Highlighting}[]
\CommentTok{\# Efectos fijos: agrupado (glm) vs mixto (glmer)}
\FunctionTok{cbind}\NormalTok{(}\AttributeTok{pooled =} \FunctionTok{coef}\NormalTok{(m\_glow\_pool), }\AttributeTok{mixto =} \FunctionTok{fixef}\NormalTok{(m\_glow\_mix))}
\end{Highlighting}
\end{Shaded}

\begin{verbatim}
                  pooled       mixto
(Intercept)  -4.21429456 -4.20528747
age           0.04119282  0.04085258
priorfracYes  0.83883503  0.84934528
\end{verbatim}

\begin{Shaded}
\begin{Highlighting}[]
\NormalTok{re\_glow }\OtherTok{\textless{}{-}} \FunctionTok{as.data.frame}\NormalTok{(}\FunctionTok{ranef}\NormalTok{(m\_glow\_mix, }\AttributeTok{condVar =} \ConstantTok{TRUE}\NormalTok{))   }\CommentTok{\# grp, condval, condsd}
\FunctionTok{ggplot}\NormalTok{(re\_glow, }\FunctionTok{aes}\NormalTok{(}\FunctionTok{reorder}\NormalTok{(grp, condval), condval)) }\SpecialCharTok{+}
  \FunctionTok{geom\_hline}\NormalTok{(}\AttributeTok{yintercept =} \DecValTok{0}\NormalTok{, }\AttributeTok{linetype =} \StringTok{"dashed"}\NormalTok{) }\SpecialCharTok{+}
  \FunctionTok{geom\_pointrange}\NormalTok{(}\FunctionTok{aes}\NormalTok{(}\AttributeTok{ymin =}\NormalTok{ condval }\SpecialCharTok{{-}} \FloatTok{1.96} \SpecialCharTok{*}\NormalTok{ condsd,}
                      \AttributeTok{ymax =}\NormalTok{ condval }\SpecialCharTok{+} \FloatTok{1.96} \SpecialCharTok{*}\NormalTok{ condsd)) }\SpecialCharTok{+}
  \FunctionTok{coord\_flip}\NormalTok{() }\SpecialCharTok{+}
  \FunctionTok{labs}\NormalTok{(}\AttributeTok{x =} \StringTok{"Centro (site\_id)"}\NormalTok{, }\AttributeTok{y =} \StringTok{"Intercepto aleatorio (log{-}odds)"}\NormalTok{)}
\end{Highlighting}
\end{Shaded}

\begin{figure}[H]

{\centering \includegraphics[width=0.75\textwidth,height=\textheight]{caso1_binaria_files/figure-pdf/fig-u14-glow-caterpillar-1.pdf}

}

\caption{\label{fig-u14-glow-caterpillar}Efecto de centro modelizado en
GLOW: intercepto aleatorio de cada uno de los 6 centros (desviación en
log-odds respecto al riesgo medio) con IC 95\%. Los seis se apiñan en
torno a 0 y todos los intervalos lo cruzan: la variación entre centros
es mínima (ICC latente ≈ 0,7\%).}

\end{figure}

Y el resultado es \textbf{claro}: la variación entre centros es
\textbf{mínima}. El modelo estima \(\hat\sigma_u=0{,}15\) (log-odds) y
un \textbf{ICC latente de apenas \(0{,}7\,\%\)} ---casi todo el riesgo
está \textbf{dentro} de los centros, no entre ellos---. El
\emph{caterpillar} lo confirma: los \textbf{seis} interceptos se apiñan
en torno a \(0\) (de \(-0{,}13\) a \(+0{,}09\)) y \textbf{todos} los
intervalos lo cruzan, así que ninguno se separa del riesgo medio. Y como
el efecto aleatorio es tan pequeño, los efectos fijos del mixto son
\textbf{prácticamente idénticos} a los del \emph{pooled} (edad
\(\approx+0{,}04\) por año; fractura previa
\(\widehat{\text{OR}}\approx e^{0{,}85}\!\approx\!2{,}3\)). En GLOW,
\textbf{agrupar por centro apenas aporta}: un hallazgo legítimo ---no
toda jerarquía esconde heterogeneidad--- y el contraste perfecto con la
cohorte, donde el centro \textbf{sí} estructuraba el riesgo
(\(\hat\sigma_u=0{,}44\), ICC \(\approx 6\,\%\)).

Esta base hace explícitos dos \textbf{límites} que la cohorte
disimulaba. Con \textbf{solo 6 grupos}, la varianza \(\sigma_u\) se
estima con \textbf{mucha incertidumbre} ---por eso la cohorte, con 24
centros y verdad conocida, sigue siendo el banco de pruebas riguroso---;
y un \(\hat\sigma_u\) tan pequeño (y, en otros casos, \textbf{pegado al
borde}) es el tipo de situación de \textbf{frontera} que el enfoque
bayesiano maneja con más suavidad (callout al cierre). GLOW gana en
\textbf{intuición} (pocos centros, reales, dibujables); la cohorte, en
\textbf{rigor}.

\begin{tcolorbox}[enhanced jigsaw, coltitle=black, rightrule=.15mm, arc=.35mm, leftrule=.75mm, colback=white, colbacktitle=quarto-callout-note-color!10!white, left=2mm, breakable, toptitle=1mm, toprule=.15mm, bottomtitle=1mm, titlerule=0mm, title=\textcolor{quarto-callout-note-color}{\faInfo}\hspace{0.5em}{Puente a la Unidad 5}, bottomrule=.15mm, opacityback=0, opacitybacktitle=0.6, colframe=quarto-callout-note-color-frame]

Esta misma cohorte guarda una pregunta que aquí hemos dejado fuera: no
solo \textbf{si} se fractura, sino \textbf{cuándo}. Las variables
\texttt{tiempo} (periodo de la revisión) y \texttt{evento} (1 = fractura
observada, 0 = \textbf{censura}) abren el \textbf{análisis de
supervivencia}, que la \protect\hyperlink{sec-u15}{Unidad 5} aborda
\textbf{sin salir del marco GLM}: el tiempo discreto se reescribe como
una regresión binomial con enlace \texttt{cloglog} sobre datos
\textbf{persona-periodo}. Y el hilo mixto \textbf{vuelve} allí: el mismo
intercepto aleatorio por centro reaparece como \textbf{fragilidad}
---\texttt{glmer(...,\ family\ =\ binomial("cloglog"))}---, de modo que
los efectos aleatorios que hemos introducido aquí se reactivan sobre el
\textbf{reloj de la supervivencia}
(\protect\hyperlink{sec-u15-glm}{Sección 5.3}).

\end{tcolorbox}

\begin{tcolorbox}[enhanced jigsaw, coltitle=black, rightrule=.15mm, arc=.35mm, leftrule=.75mm, colback=white, colbacktitle=quarto-callout-note-color!10!white, left=2mm, breakable, toptitle=1mm, toprule=.15mm, bottomtitle=1mm, titlerule=0mm, title=\textcolor{quarto-callout-note-color}{\faInfo}\hspace{0.5em}{Los efectos aleatorios encajan mejor en el enfoque bayesiano}, bottomrule=.15mm, opacityback=0, opacitybacktitle=0.6, colframe=quarto-callout-note-color-frame]

La distinción \textbf{fijo vs aleatorio} es, en parte, una incomodidad
\textbf{frecuentista}: los aleatorios se tratan como variables latentes
que hay que \textbf{integrar} de la verosimilitud (con aproximaciones
---Laplace, cuadratura---), y su varianza puede quedar \textbf{pegada al
borde} ---los ajustes \emph{singulares} o el \(\hat\sigma_u\)
subestimado que vimos---. En el marco \textbf{bayesiano} no hay tal
dicotomía: \textbf{todos} los parámetros tienen distribución, y un
efecto aleatorio es solo un parámetro con un \textbf{prior jerárquico}
(\(u_j\sim\mathcal{N}(0,\sigma_u^2)\), con su propio prior para
\(\sigma_u^2\)). El \emph{partial pooling} \textbf{emerge solo} de esa
jerarquía, la incertidumbre de la varianza se propaga entera, los
problemas de frontera se manejan con suavidad y estructuras complejas
(como el nominal mixto) encajan sin forzar. Herramientas: \texttt{brms},
\texttt{rstanarm}, \texttt{MCMCglmm}.

\end{tcolorbox}

\begin{tcolorbox}[enhanced jigsaw, coltitle=black, rightrule=.15mm, arc=.35mm, leftrule=.75mm, colback=white, colbacktitle=quarto-callout-important-color!10!white, left=2mm, breakable, toptitle=1mm, toprule=.15mm, bottomtitle=1mm, titlerule=0mm, title=\textcolor{quarto-callout-important-color}{\faExclamation}\hspace{0.5em}{Tarea del estudiante (1.4) --- resuelve la Pregunta 5 del caso}, bottomrule=.15mm, opacityback=0, opacitybacktitle=0.6, colframe=quarto-callout-important-color-frame]

Responde la \textbf{Pregunta 5} del cliente
(\protect\hyperlink{sec-u16}{Unidad 6}) sobre la cohorte de tu equipo:
\emph{¿cuánto del riesgo depende de la clínica, hay centros mejores o
peores, y qué esperar de uno nuevo?}

Tareas:

\begin{enumerate}
\def\labelenumi{\arabic{enumi}.}
\tightlist
\item
  \textbf{Reconoce el agrupamiento (RA6).} Ajusta el riesgo de fractura
  con \textbf{intercepto aleatorio por centro} (\texttt{glmer},
  \texttt{(1\ \textbar{}\ centro)}) y compáralo con la versión
  \emph{pooled} (\texttt{glm}): interpreta el \textbf{ICC} (cuánto pesa
  la clínica) y explica por qué el \emph{pooled} \textbf{atenúa} los
  efectos (condicional vs marginal), contrastándolos con la
  \textbf{verdad del DGP}.
\item
  \textbf{Centros mejores y peores.} Interpreta el \emph{caterpillar} de
  los interceptos por centro (\texttt{ranef}) y el
  \textbf{\emph{shrinkage}}: qué clínicas se desvían y por qué las
  extremas se encogen.
\item
  \textbf{¿Pendiente aleatoria?} Prueba
  \texttt{(1\ +\ x1\ \textbar{}\ centro)} y razona ---varianza estimada,
  correlación con el intercepto, convergencia--- si los datos la
  justifican.
\item
  \textbf{Diagnóstico (RA4).} Residuos simulados
  (\texttt{DHARMa::simulateResiduals}, Q--Q y dispersión), \(R^2\)
  \textbf{marginal} y \textbf{condicional} (\texttt{performance::r2}) y
  \texttt{isSingular} para confirmar que el ajuste es sólido.
\item
  \textbf{Responde al cliente,} incluida la expectativa para \textbf{una
  clínica nueva} (predicción marginal).
\end{enumerate}

\emph{Opcional (dato real):} repite el punto 1 sobre GLOW con
\texttt{(1\ \textbar{}\ site\_id)} y comenta por qué allí, con tan pocos
centros, el agrupamiento apenas aporta.

\end{tcolorbox}

\hypertarget{sec-u15}{%
\section{5 · Supervivencia: del hazard al GLM en tiempo
discreto}\label{sec-u15}}

Las unidades anteriores respondían a \emph{«¿ocurre el evento?»}: una
fractura, sí o no, modelada con la logística. Pero a esa pregunta le
falta otra ---\emph{«¿y \textbf{cuándo}?»}---. Pasar de la ocurrencia al
\textbf{tiempo hasta el evento} abre el \textbf{análisis de
supervivencia}, el segundo hilo en espiral del curso. Sin salir del
marco GLM, el tiempo-a-evento en \textbf{tiempo discreto} se reescribe
como una regresión binomial con enlace \texttt{cloglog}, y nos permitirá
integrar efectos de agrupación con modelos mixtos como los que ya vimos
en la \protect\hyperlink{sec-u14}{Unidad 4}---.

\textbf{Los datos: un seguimiento por periodos.}

GLOW se queda corto aquí: registra si hubo fractura en el primer año,
pero no \textbf{cuándo} ni a quién se dejó de observar. Por eso
trabajamos con una \protect\hyperlink{sec-datos-cohorte}{\textbf{cohorte
simulada}} ---inspirada en el mismo contexto clínico (fracturas en
mujeres seguidas en varias clínicas) y descrita al detalle, con todas
sus variables y su proceso generador, en la Presentación del caso--- que
añade justo lo que falta: \textbf{el tiempo}.

Y lo añade de una forma muy concreta que conviene fijar desde el
principio, porque determina toda la teoría que sigue: a cada mujer no la
observamos de forma continua, sino en una serie de revisiones
\(t=1,2,\dots,6\) (una al año). En cada revisión a la que llega sin
haberse fracturado ---es decir, mientras sigue \emph{en riesgo}---
anotamos una sola cosa: ¿fractura en este periodo, sí o no? La secuencia
continúa hasta que ocurre una de dos cosas: o registramos la fractura (y
ahí termina el seguimiento), o la mujer sale del estudio sin haberla
tenido (\textbf{censura}). Dicho de otro modo, el dato de cada paciente
es una secuencia de respuestas binarias, periodo a periodo: eso es,
literalmente, un \textbf{proceso estocástico en tiempo discreto}, y es
la materia prima con la que construiremos el modelo.

Resumido en variables, de cada paciente tenemos:

\begin{itemize}
\tightlist
\item
  el \textbf{tiempo hasta la fractura o salida del estudio ---por
  finalización sin rotura---} (\texttt{tiempo}, la revisión en que se
  detecta),
\item
  si llegamos a observar el evento o no (\texttt{evento}: 1 = fractura
  observada; 0 = \textbf{censura}, p.~ej. fin administrativo del
  seguimiento)
\item
  y sus \textbf{covariables de base} ---constantes a lo largo del
  seguimiento---, con lectura clínica: \texttt{x1} es un marcador de
  \textbf{fragilidad ósea} (baja densidad mineral, estandarizada en
  puntuación z), un factor de riesgo continuo ---a mayor valor, hueso
  más frágil y más riesgo a fractura (\(\beta_1>0\))---; \texttt{x2} es
  un \textbf{factor binario protector} ---recibir o no un tratamiento
  preventivo, \(\beta_2<0\)---; y \texttt{centro} es la \textbf{clínica}
  a la que pertenece cada paciente ---el agrupamiento que trataremos
  como \textbf{fragilidad inicial}, con una interceptación aleatoria---.
  Al ser simulada, conocemos la verdad del proceso generador y podremos
  comprobar que el ajuste la recupera.
\end{itemize}

\textbf{Qué cambia respecto a antes.}

Hasta ahora estimábamos la \textbf{probabilidad} de un desenlace
binario. Ahora el objeto es \textbf{dinámico}, e interesa:

\begin{itemize}
\tightlist
\item
  el \textbf{riesgo (\emph{hazard}) a lo largo del tiempo}: ¿qué riesgo
  de fracturar hay en cada periodo, entre quienes llegan sin fractura?,
\item
  la \textbf{curva de supervivencia} \(S(t)\): ¿qué fracción de la
  muestra/población sigue sin fractura pasado \(t\)?
\item
  y cómo las covariables \textbf{aceleran o retrasan} el evento (los
  \emph{hazard ratios}).
\end{itemize}

Y además, aparece un ingrediente nuevo, la \textbf{censura}: hay que
aprovechar la información de las pacientes a las que \textbf{no} vimos
fracturar en la ventana de observación, sin sesgarla. Todo ello, como
veremos, cabe dentro del GLM.

Comenzaremos definiendo y describiendo los conceptos básicos en el
análisis de supervivencia, y seguiremos implementando este análisis de
un modo no paramétrico y después a través de un modelo lineal
generalizado binario.

\begin{tcolorbox}[enhanced jigsaw, coltitle=black, rightrule=.15mm, arc=.35mm, leftrule=.75mm, colback=white, colbacktitle=quarto-callout-note-color!10!white, left=2mm, breakable, toptitle=1mm, toprule=.15mm, bottomtitle=1mm, titlerule=0mm, title=\textcolor{quarto-callout-note-color}{\faInfo}\hspace{0.5em}{Sobre esta unidad}, bottomrule=.15mm, opacityback=0, opacitybacktitle=0.6, colframe=quarto-callout-note-color-frame]

RA implicados:

\begin{itemize}
\tightlist
\item
  \textbf{RA7 · Tiempo-a-evento y censura.} Reformular la supervivencia
  en tiempo discreto como un GLM \texttt{cloglog} persona-periodo,
  tratar la censura y predecir hazard y supervivencia.
\item
  \textbf{RA3 · Interpretación en la escala correcta} (vía hazard
  ratios).
\end{itemize}

(Dobson \& Barnett, 2018, cap. de supervivencia; Aitkin \& Clayton,
1980; Allison, 1982; Cox, 1972; Prentice \& Gloeckler, 1978)

\end{tcolorbox}

\hypertarget{sec-u15-hazard}{%
\subsection{5.1 La función de supervivencia y la función de riesgo en
tiempo discreto}\label{sec-u15-hazard}}

Dado que nuestros datos son un \textbf{seguimiento por periodos},
describimos el tiempo-a-evento con las herramientas que le corresponden:
las de \textbf{tiempo discreto}. No necesitaremos ninguna otra.

Sea \(T\) el \textbf{tiempo hasta la fractura}, medido en \textbf{número
de revisión}. Como no observamos el instante exacto, sino solo el
periodo en que se detecta el evento, \(T\) es una \textbf{variable
aleatoria discreta} con valores en \(\{1,2,\dots\}\). Toda su
información se resume en dos funciones ---conocer una determina la
otra---:

\begin{itemize}
\item
  \textbf{Función de riesgo por periodo (hazard discreto).} De los
  sujetos que llegan al periodo \(t\) sin evento (el \textbf{conjunto de
  riesgo} en \(t\)), \(h_t\) es la \textbf{fracción que fractura durante
  ese periodo}: \[
  h_t=\Pr(T=t\mid T\ge t)\in[0,1].
  \] Es exactamente la respuesta a la pregunta que hacemos en cada
  revisión ---entre quienes llegan sin fractura, ¿qué proporción
  fractura \emph{ahora}?---. Nótese lo importante: \(h_t\) es una
  \textbf{probabilidad} (acotada por 1). En lenguaje de tablas de vida,
  la \textbf{probabilidad condicional de fallo} \(q_t\).
\item
  \textbf{Función de supervivencia discreta.} \(S_t=\Pr(T>t)\), la
  probabilidad de ``seguir sin fractura'' pasado el periodo \(t\). Se
  reconstruye periodo a periodo, encadenando la probabilidad de
  \emph{no} fallar en cada uno: \[
  S_t=\prod_{k=1}^{t}(1-h_k).
  \] Se interpreta directamente (\(S_5=0.8\): el 80 \% llega sin
  fractura más allá de la quinta revisión), cumple \(S_0=1\) y es
  \textbf{no creciente}.
\end{itemize}

Ambas se enlazan por \(\;h_t=1-\dfrac{S_t}{S_{t-1}}\;\): el hazard es la
\textbf{caída relativa} de la supervivencia de un periodo al siguiente.
Y la probabilidad de que la fractura caiga \textbf{exactamente} en el
periodo \(t\) combina las dos ideas ---sobrevivir a los anteriores y
fallar en \(t\)---: \[
\Pr(T=t)=h_t\prod_{k<t}(1-h_k).
\]

\begin{tcolorbox}[enhanced jigsaw, coltitle=black, rightrule=.15mm, arc=.35mm, leftrule=.75mm, colback=white, colbacktitle=quarto-callout-note-color!10!white, left=2mm, breakable, toptitle=1mm, toprule=.15mm, bottomtitle=1mm, titlerule=0mm, title=\textcolor{quarto-callout-note-color}{\faInfo}\hspace{0.5em}{Un apunte}, bottomrule=.15mm, opacityback=0, opacitybacktitle=0.6, colframe=quarto-callout-note-color-frame]

El hazard por periodo \(h_t\) es una \textbf{probabilidad}, no una tasa:
esto es lo propio del tiempo discreto y es justo lo que nos permitirá
tratarlo con un GLM binomial. (El análisis en tiempo \emph{continuo},
donde el riesgo es una tasa instantánea, es el terreno del modelo de
Cox; lo veremos en detalle en el Caso 3.)

\end{tcolorbox}

\textbf{La conexión con Bernoulli (y con el GLM).}

En nuestros datos, cada periodo en riesgo es un \textbf{ensayo de
Bernoulli} con probabilidad \(h_t\): el sujeto ``falla'' (evento) o
``sobrevive'' ese periodo. Una paciente \textbf{censurada} es aquella
que sale del seguimiento sin haber tenido una fractura (se acaba el
periodo de observación); la censura no está vinculada al riesgo de
fractura.

Disponer los datos con \textbf{una fila por sujeto y periodo en riesgo}
(persona-periodo) convierte el problema en una \textbf{secuencia de
ensayos de Bernoulli} con probabilidad \(h_t\). La verosimilitud de
supervivencia se \textbf{factoriza} exactamente como un producto de
Bernoulli sobre los periodos,

\[ P(T=t)=\begin{cases}
h_t \prod_{s<t} (1-h_s) & \text{ si evento en } t\\
\prod_{s \leq c} (1-h_s) & \text{ si censura}
\end{cases}\]

así que un \textbf{GLM binomial ordinario} (con enlace cloglog,
\protect\hyperlink{sec-u15-glm}{Sección 5.3}) maximiza la verosimilitud
correcta, sin maquinaria extra.

\begin{tcolorbox}[enhanced jigsaw, coltitle=black, rightrule=.15mm, arc=.35mm, leftrule=.75mm, colback=white, colbacktitle=quarto-callout-note-color!10!white, left=2mm, breakable, toptitle=1mm, toprule=.15mm, bottomtitle=1mm, titlerule=0mm, title=\textcolor{quarto-callout-note-color}{\faInfo}\hspace{0.5em}{¿Hace falta un efecto aleatorio por paciente?}, bottomrule=.15mm, opacityback=0, opacitybacktitle=0.6, colframe=quarto-callout-note-color-frame]

Aunque cada sujeto aporte \textbf{varias filas} en el formato
persona-periodo, \textbf{no}. Las filas de un mismo sujeto son
\textbf{condicionalmente independientes} dadas las covariables y los
interceptos por periodo ---es la factorización de la verosimilitud de
arriba, no una aproximación---, de modo que un GLM ordinario ya maximiza
la verosimilitud correcta. Los efectos aleatorios entran en \textbf{otro
nivel}: el de la \textbf{clínica} (\texttt{centro}), como
\emph{fragilidad compartida} (la ajustamos en la
\protect\hyperlink{sec-u15-glm}{Sección 5.3}). Una \emph{fragilidad
individual} por paciente sería una extensión posible ---no un
requisito--- y, además, difícil de identificar con tan pocos periodos
por sujeto.

\end{tcolorbox}

\hypertarget{sec-u15-ph}{%
\subsection{5.2 Precedentes en clave de supervivencia: Kaplan--Meier y
Cox}\label{sec-u15-ph}}

El análisis de supervivencia tiene un instrumental \textbf{propio},
anterior al GLM, que conviene conocer antes de construir el nuestro:
aporta el contexto y el vocabulario. Lo recorremos de lo más descriptivo
a lo más estructurado.

Primero, \textbf{Kaplan--Meier}: la forma no paramétrica de \emph{dejar
hablar a los datos}, que dibuja la curva de supervivencia y compara
grupos sin suponer nada sobre el riesgo. Después, los \textbf{modelos de
riesgos proporcionales}, el marco común que permite pasar de
\emph{describir} a \emph{explicar} ---cómo las covariables modulan el
riesgo--- y del que cuelga nuestro GLM. Su versión clásica en tiempo
continuo es el \textbf{modelo de Cox}, que aquí solo presentamos como
referencia: su desarrollo queda para el Caso 3.

Con estos precedentes tendremos el vocabulario para construir, en la
\protect\hyperlink{sec-u15-glm}{Sección 5.3}, un GLM que ---al ser la
cohorte simulada--- podremos \textbf{validar contra la verdad del
proceso generador}. Y ese GLM no se queda en describir: predice el
riesgo base y la supervivencia periodo a periodo, incorpora la
fragilidad entre clínicas como efecto aleatorio y puede relajar la
proporcionalidad cuando haga falta, todo dentro del único marco que ya
manejamos.

\hypertarget{sec-u15-km}{%
\subsubsection{Kaplan--Meier: dejar hablar a los
datos}\label{sec-u15-km}}

En supervivencia, lo primero es \textbf{mirar los datos}, sin
comprometerse todavía con ningún modelo. Antes de imponer una fórmula al
riesgo o un efecto a las covariables, interesa ver la forma de la curva
de supervivencia ---cómo y cuándo cae--- y si dos grupos
(tratamiento/control) se separan. Ese es el terreno de los métodos
\textbf{no paramétricos}: en lugar de estimar los parámetros de un
modelo, leen la curva directamente de los datos, con el mínimo de
supuestos. Su herramienta de referencia ---y casi siempre el primer paso
de cualquier análisis de supervivencia--- es el estimador de
\textbf{Kaplan--Meier}.

Antes de nada, conviene \textbf{ver la forma de los datos}. La
\texttt{cohorte} tiene \textbf{una fila por mujer}: sus covariables de
base (\texttt{x1}, \texttt{x2}, \texttt{centro}), el \textbf{periodo} en
que se observa el desenlace (\texttt{tiempo}) y si este fue
\textbf{fractura o censura} (\texttt{evento}: 1 = fractura, 0 =
censura).

\begin{Shaded}
\begin{Highlighting}[]
\FunctionTok{head}\NormalTok{(cohorte[, }\FunctionTok{c}\NormalTok{(}\StringTok{"id"}\NormalTok{, }\StringTok{"centro"}\NormalTok{, }\StringTok{"x1"}\NormalTok{, }\StringTok{"x2"}\NormalTok{, }\StringTok{"tiempo"}\NormalTok{, }\StringTok{"evento"}\NormalTok{)])}
\end{Highlighting}
\end{Shaded}

\begin{verbatim}
  id centro         x1 x2 tiempo evento
1  1      1 -1.6736304  0      4      1
2  2      1  1.0929380  1      3      1
3  3      1 -0.3526607  1      6      0
4  4      1  0.2976267  0      4      0
5  5      1 -0.3236454  1      6      0
6  6      1 -1.6879930  1      4      0
\end{verbatim}

El estimador de \textbf{Kaplan--Meier} (Kaplan \& Meier, 1958)
reconstruye la curva de supervivencia \textbf{sin imponer ninguna forma
funcional}: en cada instante con eventos multiplica la fracción que
``sobrevive'', \[
\widehat S(t)=\prod_{t_k\le t}\Bigl(1-\frac{d_k}{n_k}\Bigr),
\] con \(d_k\) eventos y \(n_k\) sujetos en riesgo en \(t_k\). Fíjate en
que es exactamente el producto \(\prod(1-h_k)\) de la
\protect\hyperlink{sec-u15-hazard}{Sección 5.1}, con el hazard estimado
de forma empírica, \(\widehat h_k=d_k/n_k\): Kaplan--Meier es el
\textbf{hermano no paramétrico} de la supervivencia discreta
\(S_t=\prod(1-h_k)\), y es especialmente útil para comparar curvas de
supervivencia entre grupos/perfiles. Su visualización nos permite
responder preguntas descriptivas del tipo : ¿qué grupo sobrevive más?,
¿en qué momento empieza a agravarse el riesgo en un perfil?, \ldots{}

\begin{Shaded}
\begin{Highlighting}[]
\FunctionTok{library}\NormalTok{(survival)}
\FunctionTok{library}\NormalTok{(survminer)}

\NormalTok{km }\OtherTok{\textless{}{-}} \FunctionTok{survfit}\NormalTok{(}\FunctionTok{Surv}\NormalTok{(tiempo, evento) }\SpecialCharTok{\textasciitilde{}}\NormalTok{ x2, }\AttributeTok{data =}\NormalTok{ cohorte)}
\FunctionTok{ggsurvplot}\NormalTok{(km, }\AttributeTok{data =}\NormalTok{ cohorte, }\AttributeTok{conf.int =} \ConstantTok{TRUE}\NormalTok{, }\AttributeTok{pval =} \ConstantTok{TRUE}\NormalTok{,}
           \AttributeTok{legend.title =} \StringTok{"Tratamiento (x2)"}\NormalTok{,}
           \AttributeTok{legend.labs  =} \FunctionTok{c}\NormalTok{(}\StringTok{"No"}\NormalTok{, }\StringTok{"Sí"}\NormalTok{),}
           \AttributeTok{xlab =} \StringTok{"Periodo de revisión"}\NormalTok{, }\AttributeTok{ylab =} \StringTok{"Supervivencia S(t)"}\NormalTok{)}
\end{Highlighting}
\end{Shaded}

\begin{figure}[H]

{\centering \includegraphics{caso1_binaria_files/figure-pdf/fig-u15-km-1.pdf}

}

\caption{\label{fig-u15-km}Kaplan--Meier de la cohorte por brazo de
tratamiento (x2): describe la caída de S(t) sin modelo, y el log-rank
contrasta los dos grupos. KM solo admite grupos discretos ---no la
fragilidad continua x1--- y no entrega hazard ratios; para eso hace
falta un modelo.}

\end{figure}

En la cohorte el patrón es el esperado: las tratadas (\texttt{x2} = Sí)
se mantienen por encima de las no tratadas y la separación crece con los
periodos ---arrancan casi juntas y, al sexto, la supervivencia ronda
\(0.73\) frente a \(\approx0.56\)---.

\textbf{¿Es esa separación real o podría ser ruido muestral?}

Para contrastarlo se usa la prueba \textbf{log-rank} (Mantel, 1966), el
contraste no paramétrico que acompaña a Kaplan--Meier. Su idea: en cada
periodo con eventos cuenta cuántas fracturas se observan en un grupo y
cuántas cabría esperar si las dos curvas fueran iguales (bajo \(H_0\):
misma supervivencia en ambos grupos), y acumula esas diferencias
``observado − esperado'' a lo largo del seguimiento. Si se apilan de
forma sistemática en un grupo, la separación difícilmente es casual y el
\(p\)-valor sale pequeño. No supone ninguna forma funcional ---por eso
es el compañero natural de KM---, pero, como la propia curva, no mide el
tamaño del efecto ni ajusta por covariables: solo responde \emph{sí/no}
a ``¿difieren los grupos?''.

En R, el \(p\)-valor que \texttt{ggsurvplot(...,\ pval\ =\ TRUE)}
imprime sobre la curva es justo el del log-rank; el contraste explícito
se obtiene con \texttt{survdiff}:

\begin{Shaded}
\begin{Highlighting}[]
\FunctionTok{survdiff}\NormalTok{(}\FunctionTok{Surv}\NormalTok{(tiempo, evento) }\SpecialCharTok{\textasciitilde{}}\NormalTok{ x2, }\AttributeTok{data =}\NormalTok{ cohorte)}
\end{Highlighting}
\end{Shaded}

Sobre la cohorte da \(p=0.0027\): la ventaja de las tratadas \textbf{no}
es ruido.

Con todo, KM \textbf{describe} y \textbf{compara}, pero se queda como
una primera herramienta descriptiva: \textbf{no ajusta por covariables},
no admite predictores \textbf{continuos} (la fragilidad \texttt{x1} se
queda fuera) y no entrega un \textbf{tamaño de efecto} interpretable.
Para eso hace falta una \textbf{regresión}.

\hypertarget{sec-u15-prop}{%
\subsubsection{Riesgos proporcionales, en tiempo
discreto}\label{sec-u15-prop}}

Esa regresión exige \textbf{imponer una estructura}: decir cómo entran
las covariables en el riesgo. La estructura que adoptamos ---y que nos
llevará directamente al GLM binario--- es la de \textbf{riesgos
proporcionales}.

La idea es que \textbf{una covariable multiplica el riesgo de cada
periodo por un factor constante} ---el mismo en la primera revisión que
en la sexta---, de modo que su efecto se resume en un solo número. Falta
precisar \emph{sobre qué cantidad} actúa ese factor. La respuesta sale
de la propia \(S_t\). Como \(S_t=\prod_{k\le t}(1-h_k)\), tomando
logaritmos cada periodo aporta ``\(-\log(1-h_k)\ge 0\)'' a la caída de
la supervivencia acumulada; llamemos a eso el \textbf{aporte de riesgo}
del periodo. La hipótesis de riesgos proporcionales dice que la
covariable multiplica ese aporte por un factor común
\(\exp(\mathbf{x}'\boldsymbol\beta)\): \[
-\log\bigl(1-h_t(\mathbf{x})\bigr)=\exp(\mathbf{x}'\boldsymbol\beta)\,\bigl[-\log(1-h_{0t})\bigr],
\] donde \(h_{0t}\) es el riesgo del periodo para la paciente de
\textbf{referencia} o basal (en la que no hay tratamiento,
\(\mathbf{x}=\mathbf 0\)). Y ahora el paso clave: tomando \(\log\) a
ambos lados, aparece el enlace \textbf{clog-log}, \[
\underbrace{\log\bigl(-\log(1-h_t)\bigr)}_{\operatorname{cloglog}(h_t)}
=\underbrace{\log\bigl(-\log(1-h_{0t})\bigr)}_{\alpha_t}+\mathbf{x}'\boldsymbol\beta.
\] Es decir: la hipótesis de riesgos proporcionales, escrita sobre el
hazard por periodo, \emph{es} un GLM binomial con enlace cloglog ---un
intercepto \(\alpha_t\) por revisión (el riesgo base) más un \(\beta\)
por covariable---. No hemos salido del tiempo discreto ni del GLM
binario: el enlace no se \emph{elige}, se \emph{deduce} del supuesto
(Allison, 1982; Prentice \& Gloeckler, 1978). Y en el modelo que
ajustaremos, estimar el hazard base será, literalmente, estimar esos
\(\alpha_t\) ---uno por revisión---; en la
\protect\hyperlink{sec-u15-glm}{Sección 5.3} veremos que crecen,
recuperando el riesgo base creciente del simulador.

Dos lecturas inmediatas de ese factor constante:

\begin{itemize}
\tightlist
\item
  \textbf{Sobre la supervivencia.} Si el aporte de cada periodo se
  multiplica por \(\exp(\mathbf{x}'\boldsymbol\beta)\), la supervivencia
  entera se eleva a esa potencia:
  \(S_t(\mathbf{x})=S_{0t}^{\exp(\mathbf{x}'\boldsymbol\beta)}\). En la
  escala \(\operatorname{cloglog}\) de la supervivencia, las curvas de
  dos grupos van \textbf{paralelas} (separadas por una distancia fija) y
  nunca se cruzan ---esa constancia es el supuesto---.
\item
  \textbf{El hazard ratio.} Cada \(\exp(\beta_k)\) es el factor por el
  que se multiplica el riesgo ante un aumento unitario de \(x_k\):
  \(>1\) aumenta el riesgo, \(<1\) protege (reduce riesgo), \(=1\)
  indiferente. Un solo número describe el efecto en los seis periodos.
\end{itemize}

La Figura~\ref{fig-u15-ph-ilustra} lo hace visible: a la izquierda, en
la escala \(\operatorname{cloglog}\), esa misma diferencia es una
distancia vertical fija ---las curvas paralelas de los riesgos
proporcionales---; a la derecha, en la escala del hazard, la brecha
entre grupos se ensancha (el efecto es multiplicativo).

\begin{Shaded}
\begin{Highlighting}[]
\NormalTok{beta }\OtherTok{\textless{}{-}} \FunctionTok{log}\NormalTok{(}\DecValTok{2}\NormalTok{)                                  }\CommentTok{\# efecto de la covariable: HR = exp(beta) = 2}
\NormalTok{base }\OtherTok{\textless{}{-}} \FunctionTok{tibble}\NormalTok{(}\AttributeTok{periodo =} \DecValTok{1}\SpecialCharTok{:}\DecValTok{6}\NormalTok{,}
               \AttributeTok{h0 =} \FunctionTok{c}\NormalTok{(}\FloatTok{0.05}\NormalTok{, }\FloatTok{0.07}\NormalTok{, }\FloatTok{0.09}\NormalTok{, }\FloatTok{0.11}\NormalTok{, }\FloatTok{0.13}\NormalTok{, }\FloatTok{0.15}\NormalTok{))  }\CommentTok{\# hazard base creciente (ilustrativo)}

\NormalTok{base }\SpecialCharTok{|\textgreater{}}
  \FunctionTok{mutate}\NormalTok{(}\StringTok{\textasciigrave{}}\AttributeTok{Referencia (x = 0)}\StringTok{\textasciigrave{}} \OtherTok{=}\NormalTok{ h0,}
         \StringTok{\textasciigrave{}}\AttributeTok{Grupo (x = 1)}\StringTok{\textasciigrave{}}      \OtherTok{=} \DecValTok{1} \SpecialCharTok{{-}}\NormalTok{ (}\DecValTok{1} \SpecialCharTok{{-}}\NormalTok{ h0)}\SpecialCharTok{\^{}}\FunctionTok{exp}\NormalTok{(beta)) }\SpecialCharTok{|\textgreater{}}   \CommentTok{\# PH discreto exacto}
  \FunctionTok{pivot\_longer}\NormalTok{(}\FunctionTok{c}\NormalTok{(}\StringTok{\textasciigrave{}}\AttributeTok{Referencia (x = 0)}\StringTok{\textasciigrave{}}\NormalTok{, }\StringTok{\textasciigrave{}}\AttributeTok{Grupo (x = 1)}\StringTok{\textasciigrave{}}\NormalTok{),}
               \AttributeTok{names\_to =} \StringTok{"grupo"}\NormalTok{, }\AttributeTok{values\_to =} \StringTok{"hazard"}\NormalTok{) }\SpecialCharTok{|\textgreater{}}
  \FunctionTok{mutate}\NormalTok{(}\StringTok{\textasciigrave{}}\AttributeTok{Escala hazard: h\_t}\StringTok{\textasciigrave{}}                 \OtherTok{=}\NormalTok{ hazard,}
         \StringTok{\textasciigrave{}}\AttributeTok{Escala cloglog: log({-}log(1 {-} h\_t))}\StringTok{\textasciigrave{}} \OtherTok{=} \FunctionTok{log}\NormalTok{(}\SpecialCharTok{{-}}\FunctionTok{log}\NormalTok{(}\DecValTok{1} \SpecialCharTok{{-}}\NormalTok{ hazard))) }\SpecialCharTok{|\textgreater{}}
  \FunctionTok{pivot\_longer}\NormalTok{(}\FunctionTok{c}\NormalTok{(}\StringTok{\textasciigrave{}}\AttributeTok{Escala hazard: h\_t}\StringTok{\textasciigrave{}}\NormalTok{, }\StringTok{\textasciigrave{}}\AttributeTok{Escala cloglog: log({-}log(1 {-} h\_t))}\StringTok{\textasciigrave{}}\NormalTok{),}
               \AttributeTok{names\_to =} \StringTok{"escala"}\NormalTok{, }\AttributeTok{values\_to =} \StringTok{"valor"}\NormalTok{) }\SpecialCharTok{|\textgreater{}}
  \FunctionTok{ggplot}\NormalTok{(}\FunctionTok{aes}\NormalTok{(periodo, valor, }\AttributeTok{color =}\NormalTok{ grupo)) }\SpecialCharTok{+}
  \FunctionTok{geom\_line}\NormalTok{(}\AttributeTok{linewidth =} \DecValTok{1}\NormalTok{) }\SpecialCharTok{+}
  \FunctionTok{geom\_point}\NormalTok{(}\AttributeTok{size =} \DecValTok{2}\NormalTok{) }\SpecialCharTok{+}
  \FunctionTok{facet\_wrap}\NormalTok{(}\SpecialCharTok{\textasciitilde{}}\NormalTok{escala, }\AttributeTok{scales =} \StringTok{"free\_y"}\NormalTok{) }\SpecialCharTok{+}
  \FunctionTok{scale\_color\_manual}\NormalTok{(}\AttributeTok{values =} \FunctionTok{c}\NormalTok{(}\StringTok{"Referencia (x = 0)"} \OtherTok{=} \StringTok{"grey55"}\NormalTok{,}
                                \StringTok{"Grupo (x = 1)"}      \OtherTok{=} \StringTok{"\#2c7fb8"}\NormalTok{)) }\SpecialCharTok{+}
  \FunctionTok{labs}\NormalTok{(}\AttributeTok{x =} \StringTok{"Periodo de revisión"}\NormalTok{, }\AttributeTok{y =} \ConstantTok{NULL}\NormalTok{, }\AttributeTok{color =} \ConstantTok{NULL}\NormalTok{) }\SpecialCharTok{+}
  \FunctionTok{theme\_minimal}\NormalTok{(}\AttributeTok{base\_size =} \DecValTok{12}\NormalTok{) }\SpecialCharTok{+}
  \FunctionTok{theme}\NormalTok{(}\AttributeTok{legend.position =} \StringTok{"top"}\NormalTok{)}
\end{Highlighting}
\end{Shaded}

\begin{figure}[H]

{\centering \includegraphics{caso1_binaria_files/figure-pdf/fig-u15-ph-ilustra-1.pdf}

}

\caption{\label{fig-u15-ph-ilustra}Qué significa `riesgos
proporcionales' en tiempo discreto (riesgo base creciente de juguete, HR
= 2). Izquierda (escala cloglog): ese efecto es una distancia vertical
fija (β); las curvas van paralelas y nunca se cruzan. Si esa distancia
cambiara con los periodos, la proporcionalidad fallaría. Derecha (hazard
h\_t): el grupo con la covariable alta es el de referencia con su aporte
de riesgo escalado por exp(β) cada periodo ---la brecha se ensancha
porque el efecto es multiplicativo---.}

\end{figure}

\begin{tcolorbox}[enhanced jigsaw, coltitle=black, rightrule=.15mm, arc=.35mm, leftrule=.75mm, colback=white, colbacktitle=quarto-callout-note-color!10!white, left=2mm, breakable, toptitle=1mm, toprule=.15mm, bottomtitle=1mm, titlerule=0mm, title=\textcolor{quarto-callout-note-color}{\faInfo}\hspace{0.5em}{Cuándo falla la proporcionalidad}, bottomrule=.15mm, opacityback=0, opacitybacktitle=0.6, colframe=quarto-callout-note-color-frame]

El supuesto de proporcionalidad se puede romper de dos formas.

\begin{itemize}
\tightlist
\item
  En un efecto que se desvanece, el hazard ratio se mueve hacia 1 sin
  llegar a cruzarlo: la ventaja se diluye y las curvas convergen, pero
  no se cortan. Es lo que ocurre con una vacuna cuya protección decae
  con el tiempo, o por depleción de susceptibles ---los pacientes más
  frágiles tienen el evento pronto y queda una población seleccionada
  más robusta, de modo que el efecto del factor de riesgo se atenúa---.
\item
  En curvas que se cruzan, el hazard ratio pasa de \(>1\) a \(<1\) y la
  ventaja se invierte. El caso de manual es cirugía frente a tratamiento
  conservador: riesgo alto al principio (mortalidad perioperatoria) y
  menor después, así que las curvas de supervivencia se cortan.
\end{itemize}

La \textbf{señal práctica}: si las curvas de Kaplan--Meier convergen sin
tocarse, es un efecto que se desvanece; si se cortan, es un cruce. En
cualquiera de los dos casos ---o si el efecto de una covariable cambia
con el tiempo, o si el tratamiento tiene un efecto inicial fuerte que
luego desaparece--- un único hazard ratio deja de describir bien el
efecto, y el modelo de Cox básico (o nuestro GLM con un solo HR) puede
ser inadecuado: hay que recurrir a sus extensiones ---coeficientes que
varían en el tiempo, estratificación---.

\end{tcolorbox}

\textbf{¿Es razonable asumir proporcionalidad aquí?}

La proporcionalidad es un supuesto, pero \textbf{se puede comprobar}, y
con una covariable a mano el chequeo es directo. Para la binaria
\texttt{x2}, la propia curva de Kaplan--Meier ya da una primera
respuesta: bajo proporcionalidad las curvas de los dos grupos \textbf{no
se cruzan} y se separan de forma ordenada ---justo lo que veíamos, las
tratadas siempre por encima con la brecha creciendo---, así que no hay
señal de alarma. El chequeo gráfico más fino es el \textbf{diagrama
log--log}: representar `\(\log(-\log\widehat S(t))\)' frente a
`\(\log t\)' ---es decir, el \texttt{cloglog} de la supervivencia de
Kaplan--Meier--- debería dar curvas aproximadamente paralelas si el
supuesto se cumple. Que el diagnóstico natural sea precisamente el
\texttt{cloglog} no es casualidad: es el enlace del modelo que
construiremos.

\begin{Shaded}
\begin{Highlighting}[]
\FunctionTok{ggsurvplot}\NormalTok{(km, }\AttributeTok{data =}\NormalTok{ cohorte, }\AttributeTok{fun =} \StringTok{"cloglog"}\NormalTok{,}
           \AttributeTok{legend.title =} \StringTok{"Tratamiento"}\NormalTok{, }\AttributeTok{legend.labs =} \FunctionTok{c}\NormalTok{(}\StringTok{"No"}\NormalTok{, }\StringTok{"Sí"}\NormalTok{),}
           \AttributeTok{xlab =} \StringTok{"log(periodo)"}\NormalTok{, }\AttributeTok{ylab =} \StringTok{"log({-}log S(t))"}\NormalTok{)}
\end{Highlighting}
\end{Shaded}

\begin{figure}[H]

{\centering \includegraphics{caso1_binaria_files/figure-pdf/fig-u15-km-loglog-1.pdf}

}

\caption{\label{fig-u15-km-loglog}Diagrama log-log de Kaplan-Meier por
tratamiento: log(-log S(t)) frente a log(tiempo). Bajo riesgos
proporcionales las dos curvas son aproximadamente paralelas. Con solo
seis periodos el chequeo es grueso, pero no se aprecia cruce ni
convergencia, así que la proporcionalidad para x2 es asumible.}

\end{figure}

Para la covariable \textbf{continua} \texttt{x1}, Kaplan--Meier no sirve
directamente (habría que \textbf{categorizarla} en grupos); ahí el
chequeo general son los \textbf{residuos de Schoenfeld} del modelo
ajustado ---una tendencia temporal en ellos delata la
no-proporcionalidad---, que dejamos para cuando tengamos el ajuste.

\begin{tcolorbox}[enhanced jigsaw, coltitle=black, rightrule=.15mm, arc=.35mm, leftrule=.75mm, colback=white, colbacktitle=quarto-callout-tip-color!10!white, left=2mm, breakable, toptitle=1mm, toprule=.15mm, bottomtitle=1mm, titlerule=0mm, title=\textcolor{quarto-callout-tip-color}{\faLightbulb}\hspace{0.5em}{🔧 En R. Estimar y dibujar la curva de Kaplan--Meier}, bottomrule=.15mm, opacityback=0, opacitybacktitle=0.6, colframe=quarto-callout-tip-color-frame]

Hace falta el paquete \textbf{\texttt{survival}} (y
\textbf{\texttt{survminer}} para dibujarla). Los datos van \textbf{a
nivel de paciente} ---una fila por mujer---, con el \textbf{tiempo}
hasta el evento y el \textbf{indicador} de evento/censura (1 = evento, 0
= censura). La respuesta se construye con \texttt{Surv(tiempo,\ evento)}
y la curva se estima con \texttt{survfit}:

\begin{Shaded}
\begin{Highlighting}[]
\FunctionTok{library}\NormalTok{(survival); }\FunctionTok{library}\NormalTok{(survminer)}
\NormalTok{km }\OtherTok{\textless{}{-}} \FunctionTok{survfit}\NormalTok{(}\FunctionTok{Surv}\NormalTok{(tiempo, evento) }\SpecialCharTok{\textasciitilde{}}\NormalTok{ x2, }\AttributeTok{data =}\NormalTok{ cohorte)  }\CommentTok{\# \textasciitilde{} 1 para la curva global}
\FunctionTok{ggsurvplot}\NormalTok{(km, }\AttributeTok{data =}\NormalTok{ cohorte, }\AttributeTok{pval =} \ConstantTok{TRUE}\NormalTok{)               }\CommentTok{\# añade la prueba log{-}rank}
\end{Highlighting}
\end{Shaded}

A la derecha de \texttt{\textasciitilde{}} va la variable de grupo (aquí
\texttt{x2}), y \texttt{evento} debe codificarse 1 = evento, 0 =
censura; con \texttt{fun\ =\ "cloglog"} se obtiene el diagrama log--log
de arriba.

\end{tcolorbox}

\hypertarget{sec-u15-cox}{%
\subsubsection{Cox: el mismo supuesto en tiempo
continuo}\label{sec-u15-cox}}

El supuesto de riesgos proporcionales que acabamos de escribir en tiempo
discreto tiene una versión célebre en \textbf{tiempo continuo}: el
\textbf{modelo de Cox} (Cox, 1972), el estándar del análisis de
supervivencia. Aquí \textbf{no} entramos en su maquinaria ---la
verosimilitud parcial y su carácter semiparamétrico se ven en detalle en
el Caso 3---; aquí nos basta una idea: nuestro GLM binario es su
contrapartida en tiempo discreto y comparte con él la lectura de los
efectos como \emph{hazard ratios}. Reencontraremos a Cox, ya con su
maquinaria, en el Caso 3.

Con Kaplan--Meier para describir y el marco de riesgos proporcionales
para explicar, ya tenemos todo lo que necesitamos. En la
\protect\hyperlink{sec-u15-glm}{Sección 5.3} construimos nuestro modelo
dentro del GLM, lo ajustamos, comprobamos que recupera la verdad del
simulador y derivamos las conclusiones clínicas.

\hypertarget{sec-u15-glm}{%
\subsection{5.3 La modelización GLM del problema de
supervivencia}\label{sec-u15-glm}}

Las técnicas de la \protect\hyperlink{sec-u15-ph}{Sección 5.2} eran el
precedente; ahora construimos nuestro modelo, y lo hacemos como en la
práctica: a partir de los datos, sin suponer que conocemos el proceso
que los generó. De las 877 mujeres de la cohorte tenemos tres piezas
---un \textbf{evento} binario (fractura sí/no), un \textbf{tiempo}
discreto (la revisión en que se observa) y una \textbf{censura}
(seguimientos que terminan sin fractura)--- y el objetivo es
\textbf{integrarlas} en un modelo que estime el riesgo, respete la
censura y dé efectos interpretables, todo \textbf{dentro del marco GLM}.
La sección lo construye paso a paso: qué es la variable de interés y qué
supuestos hacemos, qué enlace usamos y por qué, cómo se ajusta e
interpreta sobre la cohorte, qué aporta como modelo de supervivencia y
qué conclusiones clínicas se derivan.

\hypertarget{sec-u15-anatomia}{%
\subsubsection{\texorpdfstring{De los datos al modelo: qué es \(T\) y
qué
asumimos}{De los datos al modelo: qué es T y qué asumimos}}\label{sec-u15-anatomia}}

En la práctica, a partir de los datos observados hay que construir un
modelo razonable. ¿Qué tenemos, exactamente, de cada paciente de la
cohorte?

\begin{itemize}
\tightlist
\item
  un \textbf{indicador de fractura} (ocurrió / no ocurrió),
  \texttt{ever},
\item
  el \textbf{periodo de revisión} en que se observó el desenlace,
  \texttt{tiempo}, y
\item
  una marca de \textbf{censura} (si el seguimiento terminó sin
  fractura), \texttt{evento}.
\end{itemize}

\textbf{La variable de interés es un tiempo discreto.}

Llamemos \(T\) al \textbf{tiempo hasta la fractura}. Como no vemos el
instante exacto, sino solo la revisión en que se detecta, \(T\) es una
variable aleatoria discreta que toma valores en los periodos
\(\{1,2,\dots,6\}\). De cada paciente observamos su último periodo y si
terminó en \textbf{evento} (\(\delta=1\)) o en \textbf{censura}
(\(\delta=0\), y entonces solo sabemos que \(T\) es mayor que ese
periodo).

\textbf{No modelamos \(T\) directamente, sino su hazard.}

La distribución de \(T\) con censura es incómoda de manejar; en cambio,
el \textbf{hazard discreto} ---la probabilidad de fracturar en la
revisión \(t\) entre quienes llegan sin fractura a ella
(\protect\hyperlink{sec-u15-hazard}{Sección 5.1})--- se deja estimar de
forma natural: \[
h_t=\Pr(T=t\mid T\ge t)
\] Esto es así porque descompone la historia de cada paciente en una
\textbf{secuencia de preguntas binarias}: en cada revisión a la que
llega, ¿fractura, sí o no? Y una secuencia de respuestas `sí/no' es
justo lo que un GLM binomial sabe modelar.

Esa descomposición nos pide \textbf{reorganizar los datos}. En la
\texttt{cohorte} tenemos una fila por mujer ---su \texttt{tiempo} y su
\texttt{evento}---; el hazard por periodo, en cambio, vive en el
\textbf{formato persona-periodo} que anticipamos en la
\protect\hyperlink{sec-u15-hazard}{Sección 5.1}: \textbf{una fila por
paciente y revisión en riesgo}, con una respuesta binaria
\texttt{y\ =\ 1} solo en la revisión donde se detecta la fractura. La
\textbf{censura encaja sin esfuerzo}: una paciente censurada aporta las
filas de los periodos que sí se observaron (en todas, \texttt{y\ =\ 0})
y deja de aportar después ---no hay que inventarle un valor de \(T\)---.

La transformación la hace \texttt{expandir\_persona\_periodo()}
(definida en la Presentación del caso): recorre cada mujer y genera sus
filas, periodo a periodo, hasta el último que se observó.

\begin{Shaded}
\begin{Highlighting}[]
\NormalTok{pp }\OtherTok{\textless{}{-}} \FunctionTok{expandir\_persona\_periodo}\NormalTok{(cohorte)   }\CommentTok{\# de 1 fila por mujer a 1 fila por mujer y periodo en riesgo}
\FunctionTok{head}\NormalTok{(pp, }\DecValTok{8}\NormalTok{)}
\end{Highlighting}
\end{Shaded}

\begin{verbatim}
  id centro         x1 x2 periodo y
1  1      1 -1.6736304  0       1 0
2  1      1 -1.6736304  0       2 0
3  1      1 -1.6736304  0       3 0
4  1      1 -1.6736304  0       4 1
5  2      1  1.0929380  1       1 0
6  2      1  1.0929380  1       2 0
7  2      1  1.0929380  1       3 1
8  3      1 -0.3526607  1       1 0
\end{verbatim}

Compárese con la \texttt{cohorte} de antes: una mujer que fractura en su
tercera revisión pasa a ser \textbf{tres filas}
(\texttt{y\ =\ 0,\ 0,\ 1}); una censurada en el cuarto periodo sin
fractura, \textbf{cuatro filas} todas con \texttt{y\ =\ 0}. Cada fila es
el \textbf{ensayo de Bernoulli} de un periodo, y sobre este banco
ajustamos el GLM.

\textbf{Qué asumimos}:

\begin{itemize}
\tightlist
\item
  \textbf{Independencia condicional.} Las distintas filas de una misma
  paciente cuentan como ensayos independientes dadas las covariables y
  el periodo ---no porque sean sujetos distintos, sino por la
  factorización de la verosimilitud de supervivencia (el «punto fino» de
  la \protect\hyperlink{sec-u15-hazard}{Sección 5.1})---. Por eso basta
  un GLM ordinario, sin efecto aleatorio de paciente.
\item
  \textbf{Censura no informativa.} El motivo de abandonar el seguimiento
  \textbf{no} está ligado al riesgo de fracturar; bajo ese supuesto, la
  censura no introduce sesgo.
\item
  \textbf{Riesgos proporcionales.} Asumimos que las covariables
  \textbf{multiplican} el hazard por un factor \textbf{constante en el
  tiempo} (el supuesto de la \protect\hyperlink{sec-u15-ph}{Sección
  5.2}). Esto permite resumir cada efecto en un único hazard ratio y,
  como veremos, \textbf{fija el enlace}. Si no se cumpliera, se
  relajaría con interacciones covariable × periodo.
\item
  \emph{Forma del riesgo base libre.} No imponemos ninguna fórmula a
  cómo cambia el riesgo con las revisiones: dejamos \textbf{un parámetro
  por periodo}, \(\alpha_t\) (uno por revisión). Es lo más flexible
  posible ---el análogo discreto de ``no especificar \(h_0(t)\)'' que
  hacía Cox---.
\end{itemize}

Con esas piezas, el modelo casi se escribe solo: un \textbf{GLM
binomial} para el hazard de cada periodo, con predictor lineal \[
\eta_t(\mathbf{x})=\alpha_t+\mathbf{x}'\boldsymbol\beta
\] ---un intercepto por revisión (\(\alpha_t\), el riesgo base) más los
efectos de las covariables (\(\boldsymbol\beta\))---. Queda \textbf{una
sola decisión}: el \textbf{enlace} que conecta el hazard (una
probabilidad en \([0,1]\)) con ese predictor. Y no es libre: el supuesto
de \textbf{riesgos proporcionales} lo determina ---y ya lo dedujimos en
la \protect\hyperlink{sec-u15-prop}{Sección 5.2}---.

\hypertarget{sec-u15-cloglog}{%
\subsubsection{El enlace complementary log-log y por qué es el
natural}\label{sec-u15-cloglog}}

El hazard es una \textbf{probabilidad}; el predictor lineal
\(\alpha_t+\mathbf{x}'\boldsymbol\beta\) vive en toda la recta real. Ya
sabemos, por la \protect\hyperlink{sec-u15-prop}{Sección 5.2}, qué
\textbf{enlace} los conecta bajo riesgos proporcionales: el
\textbf{complementary log-log}, \[
\operatorname{cloglog}(p)=\log(-\log(1-p))=\eta,\qquad p=1-\exp(-\exp(\eta)).
\] A diferencia del logit y el probit, es \textbf{asimétrico} (se
aproxima a 1 más despacio que a 0).

\begin{Shaded}
\begin{Highlighting}[]
\FunctionTok{tibble}\NormalTok{(}\AttributeTok{eta =} \FunctionTok{seq}\NormalTok{(}\SpecialCharTok{{-}}\DecValTok{4}\NormalTok{, }\DecValTok{4}\NormalTok{, }\AttributeTok{by =} \FloatTok{0.02}\NormalTok{)) }\SpecialCharTok{|\textgreater{}}
  \FunctionTok{mutate}\NormalTok{(}\AttributeTok{logit   =} \FunctionTok{plogis}\NormalTok{(eta),}
         \AttributeTok{cloglog =} \DecValTok{1} \SpecialCharTok{{-}} \FunctionTok{exp}\NormalTok{(}\SpecialCharTok{{-}}\FunctionTok{exp}\NormalTok{(eta))) }\SpecialCharTok{|\textgreater{}}
  \FunctionTok{pivot\_longer}\NormalTok{(}\FunctionTok{c}\NormalTok{(logit, cloglog), }\AttributeTok{names\_to =} \StringTok{"enlace"}\NormalTok{, }\AttributeTok{values\_to =} \StringTok{"p"}\NormalTok{) }\SpecialCharTok{|\textgreater{}}
  \FunctionTok{ggplot}\NormalTok{(}\FunctionTok{aes}\NormalTok{(eta, p, }\AttributeTok{color =}\NormalTok{ enlace)) }\SpecialCharTok{+}
  \FunctionTok{geom\_line}\NormalTok{(}\AttributeTok{linewidth =} \DecValTok{1}\NormalTok{) }\SpecialCharTok{+}
  \FunctionTok{labs}\NormalTok{(}\AttributeTok{x =} \FunctionTok{expression}\NormalTok{(eta), }\AttributeTok{y =} \FunctionTok{expression}\NormalTok{(p }\SpecialCharTok{==}\NormalTok{ g}\SpecialCharTok{\^{}}\NormalTok{\{}\SpecialCharTok{{-}}\DecValTok{1}\NormalTok{\}(eta)), }\AttributeTok{color =} \StringTok{"Enlace"}\NormalTok{)}
\end{Highlighting}
\end{Shaded}

\begin{figure}[H]

{\centering \includegraphics{caso1_binaria_files/figure-pdf/fig-u15-enlaces-1.pdf}

}

\caption{\label{fig-u15-enlaces}De los enlaces comparados en la Unidad
2, aquí retenemos solo logit y cloglog ---los relevantes en
supervivencia discreta---. La asimetría del cloglog (se aproxima a 1 más
despacio que a 0) es lo que lo hace el enlace natural de los riesgos
proporcionales.}

\end{figure}

Recordemos por qué este enlace y no el logit: como dedujimos en la
\protect\hyperlink{sec-u15-prop}{Sección 5.2}, escribir la hipótesis de
riesgos proporcionales sobre el hazard por periodo da
\textbf{exactamente}
\(\operatorname{cloglog}(h_t)=\alpha_t+\mathbf{x}'\boldsymbol\beta\),
con un intercepto \(\alpha_t\) que absorbe el riesgo base (Allison,
1982; Prentice \& Gloeckler, 1978). Con \texttt{logit} se obtiene la
aproximación ``logística discreta'', razonable cuando el riesgo por
periodo es pequeño, pero solo el \texttt{cloglog} \textbf{preserva
exactamente} los riesgos proporcionales subyacentes ---y por eso
reproduce los \textbf{mismos hazard ratios que Cox}---.

\hypertarget{sec-u15-lab}{%
\subsubsection{Ajuste en tiempo discreto e
interpretación}\label{sec-u15-lab}}

Con los datos ya en formato \textbf{persona-periodo} (\texttt{pp}),
ajustar el hazard es, literalmente, un \textbf{GLM binomial con enlace
cloglog}. Construido el modelo, ahora aprovechamos que la cohorte es
\textbf{simulada} para hacer algo que con datos reales no podríamos:
\textbf{contrastar las estimaciones con la verdad} del proceso generador
---la validación que anticipábamos---.

\begin{Shaded}
\begin{Highlighting}[]
\NormalTok{m\_pp }\OtherTok{\textless{}{-}} \FunctionTok{glm}\NormalTok{(y }\SpecialCharTok{\textasciitilde{}}\NormalTok{ periodo }\SpecialCharTok{+}\NormalTok{ x1 }\SpecialCharTok{+}\NormalTok{ x2, }\AttributeTok{family =} \FunctionTok{binomial}\NormalTok{(}\StringTok{"cloglog"}\NormalTok{), }\AttributeTok{data =}\NormalTok{ pp)}
\FunctionTok{summary}\NormalTok{(m\_pp)}
\end{Highlighting}
\end{Shaded}

\begin{verbatim}

Call:
glm(formula = y ~ periodo + x1 + x2, family = binomial("cloglog"), 
    data = pp)

Coefficients:
            Estimate Std. Error z value Pr(>|z|)    
(Intercept) -2.97695    0.15603 -19.079  < 2e-16 ***
periodo2     0.38666    0.18723   2.065 0.038909 *  
periodo3     0.49718    0.18982   2.619 0.008814 ** 
periodo4     0.44472    0.21705   2.049 0.040472 *  
periodo5     0.58974    0.24694   2.388 0.016931 *  
periodo6     0.77511    0.30201   2.567 0.010272 *  
x1           0.84153    0.06866  12.256  < 2e-16 ***
x2          -0.45364    0.13269  -3.419 0.000629 ***
---
Signif. codes:  0 '***' 0.001 '**' 0.01 '*' 0.05 '.' 0.1 ' ' 1

(Dispersion parameter for binomial family taken to be 1)

    Null deviance: 1810.7  on 3427  degrees of freedom
Residual deviance: 1640.8  on 3420  degrees of freedom
AIC: 1656.8

Number of Fisher Scoring iterations: 6
\end{verbatim}

\begin{Shaded}
\begin{Highlighting}[]
\FunctionTok{exp}\NormalTok{(}\FunctionTok{coef}\NormalTok{(m\_pp)[}\FunctionTok{c}\NormalTok{(}\StringTok{"x1"}\NormalTok{, }\StringTok{"x2"}\NormalTok{)])               }\CommentTok{\# hazard ratios estimados}
\end{Highlighting}
\end{Shaded}

\begin{verbatim}
       x1        x2 
2.3199063 0.6353109 
\end{verbatim}

\begin{Shaded}
\begin{Highlighting}[]
\FunctionTok{exp}\NormalTok{(}\FunctionTok{attr}\NormalTok{(cohorte, }\StringTok{"verdad"}\NormalTok{)}\SpecialCharTok{$}\NormalTok{binaria}\SpecialCharTok{$}\NormalTok{beta)    }\CommentTok{\# HR verdaderos: e\^{}0.85, e\^{}{-}0.65}
\end{Highlighting}
\end{Shaded}

\begin{verbatim}
       x1        x2 
2.3396469 0.5220458 
\end{verbatim}

Los coeficientes de \texttt{periodo} son los \textbf{interceptos por
periodo} \(\alpha_t\), y aquí \textbf{crecen}
(\(0.39,\ 0.50,\ \dots,\ 0.78\) sobre el periodo 1): el modelo recupera
el \textbf{riesgo base creciente} del DGP. Los de \texttt{x1} y
\texttt{x2}, exponenciados, son \textbf{hazard ratios}: \texttt{x1} da
\(\widehat{\text{HR}}\approx2.32\), que \textbf{clava} la verdad
(\(e^{0.85}\approx2.34\)), y \texttt{x2} da \(\approx0.64\), algo
\textbf{atenuado} frente a su verdad (\(e^{-0.65}\approx0.52\)) ---el
efecto protector se queda corto, en parte por omitir la fragilidad---.
La \textbf{censura} no exige nada extra: un sujeto censurado simplemente
\textbf{deja de aportar filas} tras su último periodo en riesgo, sin un
\texttt{y\ =\ 1}; aporta lo que se observó (que no falló mientras
estuvo) y nada más.

El riesgo base por periodo y la supervivencia que implica:

\begin{Shaded}
\begin{Highlighting}[]
\NormalTok{base }\OtherTok{\textless{}{-}} \FunctionTok{tibble}\NormalTok{(}\AttributeTok{periodo =} \FunctionTok{factor}\NormalTok{(}\FunctionTok{levels}\NormalTok{(pp}\SpecialCharTok{$}\NormalTok{periodo), }\AttributeTok{levels =} \FunctionTok{levels}\NormalTok{(pp}\SpecialCharTok{$}\NormalTok{periodo)),}
               \AttributeTok{x1 =} \DecValTok{0}\NormalTok{, }\AttributeTok{x2 =} \DecValTok{0}\NormalTok{)}
\NormalTok{base}\SpecialCharTok{$}\NormalTok{h0 }\OtherTok{\textless{}{-}} \FunctionTok{predict}\NormalTok{(m\_pp, }\AttributeTok{newdata =}\NormalTok{ base, }\AttributeTok{type =} \StringTok{"response"}\NormalTok{)   }\CommentTok{\# hazard base por periodo}
\NormalTok{base}\SpecialCharTok{$}\NormalTok{t  }\OtherTok{\textless{}{-}} \FunctionTok{seq\_len}\NormalTok{(}\FunctionTok{nrow}\NormalTok{(base))}
\NormalTok{base}\SpecialCharTok{$}\NormalTok{S  }\OtherTok{\textless{}{-}} \FunctionTok{cumprod}\NormalTok{(}\DecValTok{1} \SpecialCharTok{{-}}\NormalTok{ base}\SpecialCharTok{$}\NormalTok{h0)                               }\CommentTok{\# supervivencia base}

\NormalTok{pH }\OtherTok{\textless{}{-}} \FunctionTok{ggplot}\NormalTok{(base, }\FunctionTok{aes}\NormalTok{(t, h0)) }\SpecialCharTok{+}
  \FunctionTok{geom\_col}\NormalTok{(}\AttributeTok{fill =} \StringTok{"grey70"}\NormalTok{) }\SpecialCharTok{+}
  \FunctionTok{labs}\NormalTok{(}\AttributeTok{x =} \StringTok{"Periodo"}\NormalTok{, }\AttributeTok{y =} \FunctionTok{expression}\NormalTok{(hazard}\SpecialCharTok{\textasciitilde{}}\NormalTok{base}\SpecialCharTok{\textasciitilde{}}\ErrorTok{\textasciitilde{}}\NormalTok{h[}\DecValTok{0}\NormalTok{][t]), }\AttributeTok{title =} \StringTok{"Riesgo base por periodo"}\NormalTok{)}
\NormalTok{pS }\OtherTok{\textless{}{-}} \FunctionTok{ggplot}\NormalTok{(base, }\FunctionTok{aes}\NormalTok{(t, S)) }\SpecialCharTok{+}
  \FunctionTok{geom\_line}\NormalTok{(}\AttributeTok{linewidth =} \DecValTok{1}\NormalTok{) }\SpecialCharTok{+} \FunctionTok{geom\_point}\NormalTok{(}\AttributeTok{size =} \DecValTok{2}\NormalTok{) }\SpecialCharTok{+}
  \FunctionTok{scale\_y\_continuous}\NormalTok{(}\AttributeTok{limits =} \FunctionTok{c}\NormalTok{(}\DecValTok{0}\NormalTok{, }\DecValTok{1}\NormalTok{)) }\SpecialCharTok{+}
  \FunctionTok{labs}\NormalTok{(}\AttributeTok{x =} \StringTok{"Periodo"}\NormalTok{, }\AttributeTok{y =} \FunctionTok{expression}\NormalTok{(supervivencia}\SpecialCharTok{\textasciitilde{}}\ErrorTok{\textasciitilde{}}\NormalTok{S[t]), }\AttributeTok{title =} \StringTok{"Supervivencia base"}\NormalTok{)}
\NormalTok{pH }\SpecialCharTok{+}\NormalTok{ pS}
\end{Highlighting}
\end{Shaded}

\begin{figure}[H]

{\centering \includegraphics{caso1_binaria_files/figure-pdf/fig-u15-hazard-base-1.pdf}

}

\caption{\label{fig-u15-hazard-base}Riesgo base por periodo (h0t,
evaluado en x1 = x2 = 0) y la supervivencia base St = prod(1 - h0k) que
implica. En este DGP el riesgo base es creciente, de modo que la
supervivencia cae cada vez más deprisa.}

\end{figure}

\textbf{Fragilidad: el centro como efecto aleatorio.}

Las pacientes se agrupan en clínicas; añadir un \textbf{intercepto
aleatorio por centro} ---la \emph{fragilidad}, en el lenguaje de la
supervivencia--- es, otra vez, un \textbf{GLMM cloglog}, y reactiva la
\protect\hyperlink{sec-u14}{Unidad 4} sobre el reloj de supervivencia.
Y, como pedíamos, se ajusta \textbf{sin bayesiana}, con \texttt{lme4}:

\begin{Shaded}
\begin{Highlighting}[]
\FunctionTok{library}\NormalTok{(lme4)}
\NormalTok{m\_frail }\OtherTok{\textless{}{-}} \FunctionTok{glmer}\NormalTok{(y }\SpecialCharTok{\textasciitilde{}}\NormalTok{ periodo }\SpecialCharTok{+}\NormalTok{ x1 }\SpecialCharTok{+}\NormalTok{ x2 }\SpecialCharTok{+}\NormalTok{ (}\DecValTok{1} \SpecialCharTok{|}\NormalTok{ centro),}
                 \AttributeTok{family =} \FunctionTok{binomial}\NormalTok{(}\StringTok{"cloglog"}\NormalTok{), }\AttributeTok{data =}\NormalTok{ pp)}
\FunctionTok{VarCorr}\NormalTok{(m\_frail)               }\CommentTok{\# sigma de la fragilidad (verdad del DGP: 0.70)}
\end{Highlighting}
\end{Shaded}

\begin{verbatim}
 Groups Name        Std.Dev.
 centro (Intercept) 0.367   
\end{verbatim}

\begin{Shaded}
\begin{Highlighting}[]
\FunctionTok{fixef}\NormalTok{(m\_frail)[}\FunctionTok{c}\NormalTok{(}\StringTok{"x1"}\NormalTok{, }\StringTok{"x2"}\NormalTok{)]  }\CommentTok{\# efectos fijos, ahora condicionales al centro}
\end{Highlighting}
\end{Shaded}

\begin{verbatim}
        x1         x2 
 0.8879172 -0.4681168 
\end{verbatim}

\texttt{VarCorr} estima la \textbf{varianza de la fragilidad}:
\(\hat\sigma_u\approx0.37\), \textbf{por debajo} de la verdad (\(0.70\))
---la misma subestimación de las varianzas que con datos binarios y
pocos grupos (24 centros)---; mide cuánto varía el riesgo base de unas
clínicas a otras. Al \textbf{condicionar al centro}, los efectos crecen
algo en magnitud (\(x_1: 0.84\!\to\!0.89\); \(x_2: -0.45\!\to\!-0.47\)):
es el salto marginal→condicional de 1.4, ahora sobre el hazard, aunque
leve porque \(\hat\sigma_u\) es pequeño.

\begin{tcolorbox}[enhanced jigsaw, coltitle=black, rightrule=.15mm, arc=.35mm, leftrule=.75mm, colback=white, colbacktitle=quarto-callout-tip-color!10!white, left=2mm, breakable, toptitle=1mm, toprule=.15mm, bottomtitle=1mm, titlerule=0mm, title=\textcolor{quarto-callout-tip-color}{\faLightbulb}\hspace{0.5em}{🔧 En R. El GLM de supervivencia en persona-periodo}, bottomrule=.15mm, opacityback=0, opacitybacktitle=0.6, colframe=quarto-callout-tip-color-frame]

\begin{Shaded}
\begin{Highlighting}[]
\NormalTok{pp }\OtherTok{\textless{}{-}} \FunctionTok{expandir\_persona\_periodo}\NormalTok{(cohorte)                                   }\CommentTok{\# 1 fila por paciente y periodo}
\NormalTok{m  }\OtherTok{\textless{}{-}} \FunctionTok{glm}\NormalTok{(y }\SpecialCharTok{\textasciitilde{}}\NormalTok{ periodo }\SpecialCharTok{+}\NormalTok{ x1 }\SpecialCharTok{+}\NormalTok{ x2, }\AttributeTok{family =} \FunctionTok{binomial}\NormalTok{(}\StringTok{"cloglog"}\NormalTok{), }\AttributeTok{data =}\NormalTok{ pp) }\CommentTok{\# hazard discreto}
\end{Highlighting}
\end{Shaded}

\begin{itemize}
\tightlist
\item
  \textbf{Formato persona-periodo} --- \texttt{y\ =\ 1} solo en el
  periodo del evento; la \textbf{censura se trata sola} (la paciente
  aporta filas \texttt{y\ =\ 0} hasta su último periodo observado, sin
  fila de evento).
\item
  \textbf{El factor \texttt{periodo}} estima el \textbf{hazard base}
  \(\alpha_t\) (un intercepto por revisión); las covariables entran como
  en cualquier GLM y su \(e^{\hat\beta}\) es un \textbf{\emph{hazard
  ratio}} con \texttt{cloglog} (un \textbf{odds ratio} si se usa
  \texttt{binomial("logit")}, la logística de tiempo discreto).
\item
  \textbf{Del hazard a la supervivencia} ---
  \texttt{predict(m,\ newdata,\ type\ =\ "response")} da \(\hat h_t\)
  por periodo; la supervivencia se acumula con \texttt{cumprod(1\ -\ h)}
  (y el riesgo acumulado es \texttt{1\ -\ S}).
\item
  \textbf{Fragilidad} --- un intercepto aleatorio por grupo,
  \texttt{+\ (1\ \textbar{}\ centro)}, lo convierte en un \textbf{GLMM
  cloglog}: \texttt{lme4::glmer(...,\ family\ =\ binomial("cloglog"))},
  con \texttt{VarCorr(m)} para \(\hat\sigma_u\) y \texttt{fixef(m)} para
  los efectos condicionales.
\end{itemize}

\end{tcolorbox}

\hypertarget{sec-u15-logit}{%
\subsubsection{Cloglog frente a logit: ¿cuánto importa el
enlace?}\label{sec-u15-logit}}

Elegimos el \texttt{cloglog} por una razón estructural: preserva los
riesgos proporcionales. Pero los datos persona-periodo son binarios, así
que nada impide ajustarlos con el enlace por defecto, el \textbf{logit}
---la \emph{regresión logística de tiempo discreto}---. La diferencia es
de \textbf{interpretación}: con \texttt{cloglog} los \(\exp(\beta)\) son
\textbf{hazard ratios} (riesgos proporcionales); con \texttt{logit} son
\textbf{odds ratios} (un supuesto de \emph{odds proporcionales}).
Ajustemos los dos y comparémoslos:

\begin{Shaded}
\begin{Highlighting}[]
\NormalTok{m\_pp\_logit }\OtherTok{\textless{}{-}} \FunctionTok{glm}\NormalTok{(y }\SpecialCharTok{\textasciitilde{}}\NormalTok{ periodo }\SpecialCharTok{+}\NormalTok{ x1 }\SpecialCharTok{+}\NormalTok{ x2, }\AttributeTok{family =} \FunctionTok{binomial}\NormalTok{(}\StringTok{"logit"}\NormalTok{), }\AttributeTok{data =}\NormalTok{ pp)}

\CommentTok{\# cloglog {-}\textgreater{} hazard ratios ; logit {-}\textgreater{} odds ratios}
\FunctionTok{cbind}\NormalTok{(}\AttributeTok{cloglog\_HR =} \FunctionTok{exp}\NormalTok{(}\FunctionTok{coef}\NormalTok{(m\_pp)[}\FunctionTok{c}\NormalTok{(}\StringTok{"x1"}\NormalTok{, }\StringTok{"x2"}\NormalTok{)]),}
      \AttributeTok{logit\_OR   =} \FunctionTok{exp}\NormalTok{(}\FunctionTok{coef}\NormalTok{(m\_pp\_logit)[}\FunctionTok{c}\NormalTok{(}\StringTok{"x1"}\NormalTok{, }\StringTok{"x2"}\NormalTok{)]))}
\end{Highlighting}
\end{Shaded}

\begin{verbatim}
   cloglog_HR  logit_OR
x1  2.3199063 2.4420551
x2  0.6353109 0.6116684
\end{verbatim}

\begin{Shaded}
\begin{Highlighting}[]
\FunctionTok{c}\NormalTok{(}\AttributeTok{AIC\_cloglog =} \FunctionTok{AIC}\NormalTok{(m\_pp), }\AttributeTok{AIC\_logit =} \FunctionTok{AIC}\NormalTok{(m\_pp\_logit))}
\end{Highlighting}
\end{Shaded}

\begin{verbatim}
AIC_cloglog   AIC_logit 
   1656.826    1657.275 
\end{verbatim}

Como vemos en las tabla, la diferencia en los coeficientes estimados
para ambos modelos es mínima: para la fragilidad, HR \(=2.32\) (cloglog)
frente a OR \(=2.44\) (logit); para el tratamiento, \(0.64\) frente a
\(0.61\). El patrón es el esperado ---el odds ratio se aleja un poco más
de 1 que el hazard ratio (algo más alto cuando el efecto es de riesgo,
algo más bajo cuando es protector), porque con hazards por periodo que
son moderados, odds y riesgo se parecen pero no son idénticos---. Y los
AIC casi empatan (\(1656.8\) con cloglog frente a \(1657.3\) con logit),
con una ventaja mínima para el cloglog. En estos datos, el enlace apenas
mueve los resultados.

\begin{Shaded}
\begin{Highlighting}[]
\NormalTok{base2 }\OtherTok{\textless{}{-}} \FunctionTok{tibble}\NormalTok{(}\AttributeTok{periodo =} \FunctionTok{factor}\NormalTok{(}\FunctionTok{levels}\NormalTok{(pp}\SpecialCharTok{$}\NormalTok{periodo), }\FunctionTok{levels}\NormalTok{(pp}\SpecialCharTok{$}\NormalTok{periodo)), }\AttributeTok{x1 =} \DecValTok{0}\NormalTok{, }\AttributeTok{x2 =} \DecValTok{0}\NormalTok{)}
\NormalTok{pred2 }\OtherTok{\textless{}{-}}\NormalTok{ base2 }\SpecialCharTok{|\textgreater{}}
  \FunctionTok{mutate}\NormalTok{(}\AttributeTok{cloglog =} \FunctionTok{predict}\NormalTok{(m\_pp,       }\AttributeTok{newdata =}\NormalTok{ base2, }\AttributeTok{type =} \StringTok{"response"}\NormalTok{),}
         \AttributeTok{logit   =} \FunctionTok{predict}\NormalTok{(m\_pp\_logit, }\AttributeTok{newdata =}\NormalTok{ base2, }\AttributeTok{type =} \StringTok{"response"}\NormalTok{)) }\SpecialCharTok{|\textgreater{}}
  \FunctionTok{pivot\_longer}\NormalTok{(}\FunctionTok{c}\NormalTok{(cloglog, logit), }\AttributeTok{names\_to =} \StringTok{"enlace"}\NormalTok{, }\AttributeTok{values\_to =} \StringTok{"h"}\NormalTok{) }\SpecialCharTok{|\textgreater{}}
  \FunctionTok{mutate}\NormalTok{(}\AttributeTok{t =} \FunctionTok{as.integer}\NormalTok{(periodo)) }\SpecialCharTok{|\textgreater{}}
  \FunctionTok{arrange}\NormalTok{(enlace, t) }\SpecialCharTok{|\textgreater{}}
  \FunctionTok{group\_by}\NormalTok{(enlace) }\SpecialCharTok{|\textgreater{}}
  \FunctionTok{mutate}\NormalTok{(}\AttributeTok{S =} \FunctionTok{cumprod}\NormalTok{(}\DecValTok{1} \SpecialCharTok{{-}}\NormalTok{ h)) }\SpecialCharTok{|\textgreater{}}
  \FunctionTok{ungroup}\NormalTok{()}

\FunctionTok{ggplot}\NormalTok{(pred2, }\FunctionTok{aes}\NormalTok{(t, S, }\AttributeTok{color =}\NormalTok{ enlace)) }\SpecialCharTok{+}
  \FunctionTok{geom\_line}\NormalTok{(}\AttributeTok{linewidth =} \DecValTok{1}\NormalTok{) }\SpecialCharTok{+} \FunctionTok{geom\_point}\NormalTok{(}\AttributeTok{size =} \DecValTok{2}\NormalTok{) }\SpecialCharTok{+}
  \FunctionTok{scale\_y\_continuous}\NormalTok{(}\AttributeTok{limits =} \FunctionTok{c}\NormalTok{(}\DecValTok{0}\NormalTok{, }\DecValTok{1}\NormalTok{)) }\SpecialCharTok{+}
  \FunctionTok{labs}\NormalTok{(}\AttributeTok{x =} \StringTok{"Periodo"}\NormalTok{, }\AttributeTok{y =} \StringTok{"Supervivencia base S(t)"}\NormalTok{, }\AttributeTok{color =} \StringTok{"Enlace"}\NormalTok{)}
\end{Highlighting}
\end{Shaded}

\begin{figure}[H]

{\centering \includegraphics{caso1_binaria_files/figure-pdf/fig-u15-logit-cmp-1.pdf}

}

\caption{\label{fig-u15-logit-cmp}Supervivencia base (x1 = x2 = 0)
reconstruida desde el GLM con enlace cloglog y con logit. Las dos curvas
son casi indistinguibles: con hazards por periodo moderados ambos
enlaces dan prácticamente lo mismo. La preferencia por el cloglog es
estructural (preserva los riesgos proporcionales y reproduce a Cox), no
numérica.}

\end{figure}

El gráfico lo confirma: la supervivencia base reconstruida con uno y
otro enlace es casi la misma curva. La conclusión es doble.
Numéricamente, aquí da casi igual; pero conceptualmente la elección no
es indiferente: nos quedamos con el \texttt{cloglog} por lo que
significa ---es el único que preserva los riesgos proporcionales y, por
tanto, el que reproduce a Cox y permite leer los efectos como hazard
ratios---. El \texttt{logit} sería la opción si quisiéramos razonar en
odds.

\hypertarget{sec-u15-comparativa}{%
\subsubsection{Qué aporta el GLM en
supervivencia}\label{sec-u15-comparativa}}

Merece la pena recoger lo que hemos ganado al plantear la supervivencia
como un GLM binomial con enlace \texttt{cloglog}, sin salir del marco
que ya manejábamos:

\begin{itemize}
\tightlist
\item
  Modela el riesgo base: los interceptos \(\alpha_t\) son el hazard de
  cada periodo, así que el modelo no solo estima los efectos, sino que
  predice el hazard y la supervivencia periodo a periodo (lo que hicimos
  en los gráficos anteriores).
\item
  Trata la censura sin nada extra: una paciente censurada aporta sus
  filas observadas y deja de aportar después; la verosimilitud es la
  correcta por construcción.
\item
  Maneja el tiempo discreto de forma exacta: cuando muchas mujeres
  fracturan en el mismo periodo, el formato persona-periodo lo recoge
  tal cual, sin aproximaciones.
\item
  Se extiende a la fragilidad: añadir \texttt{(1\ \textbar{}\ centro)}
  es un GLMM \texttt{cloglog} que se ajusta con \texttt{glmer},
  reactivando los efectos mixtos de la
  \protect\hyperlink{sec-u14}{Unidad 4} ---como vimos arriba---.
\end{itemize}

Todo ello leyendo los efectos como \emph{hazard ratios}, la escala del
estándar clínico. Ese estándar es el modelo de Cox, la versión en tiempo
continuo del mismo supuesto de riesgos proporcionales, que
desarrollaremos ---y compararemos formalmente con este GLM--- en el Caso
3. Aquí la validación la da el propio simulador: los \emph{hazard
ratios} estimados recuperan los valores verdaderos del DGP
(\(\approx 2{,}3\) y \(0{,}64\) frente a \(e^{0{,}85}\) y
\(e^{-0{,}65}\)), que es cuanto necesitamos para confiar en el modelo.

\hypertarget{sec-u15-clinica}{%
\subsubsection{Conclusiones para la práctica
clínica}\label{sec-u15-clinica}}

Traducidas a quien trata pacientes, las estimaciones del modelo dicen:

\begin{itemize}
\tightlist
\item
  \textbf{La fragilidad ósea manda.} Cada desviación típica de baja
  densidad mineral \textbf{multiplica por \textasciitilde2,3} el riesgo
  de fractura en cada periodo: es el factor dominante. Identificar y
  vigilar de cerca a las pacientes de alta fragilidad es la palanca de
  mayor impacto.
\item
  \textbf{El tratamiento ayuda, pero no basta por sí solo.} Reduce el
  riesgo alrededor de un tercio (HR \(0.64\)), un efecto real y estable
  en el tiempo, pero parcial: en pacientes muy frágiles el riesgo sigue
  siendo alto pese al tratamiento. Conviene combinarlo con el manejo de
  la fragilidad.
\item
  \textbf{El riesgo no se relaja con el seguimiento.} El hazard base es
  creciente: la probabilidad de fractura por periodo aumenta con el
  tiempo, así que la vigilancia no debe disminuir en las revisiones
  tardías.
\item
  \textbf{Importa el centro.} La \emph{fragilidad} entre clínicas indica
  que el riesgo basal varía de un centro a otro: parte del riesgo
  individual depende del contexto asistencial, no solo de la paciente.
\item
  \textbf{Un solo número resume el efecto.} Al cumplirse los riesgos
  proporcionales, el hazard ratio es estable en los seis periodos: una
  única cifra (×2,3 por fragilidad; ×0,64 por tratamiento) describe el
  efecto durante todo el seguimiento, lo que facilita comunicárselo a la
  paciente.
\end{itemize}

El siguiente gráfico hace tangible la decisión clínica: el
\textbf{riesgo acumulado de fractura} (\(1-S_t\)) a lo largo de los
periodos para tres perfiles de paciente.

\begin{Shaded}
\begin{Highlighting}[]
\NormalTok{arquetipos }\OtherTok{\textless{}{-}} \FunctionTok{tibble}\NormalTok{(}
  \AttributeTok{arquetipo =} \FunctionTok{factor}\NormalTok{(}\FunctionTok{c}\NormalTok{(}\StringTok{"Alta fragilidad, sin tratamiento"}\NormalTok{,}
                       \StringTok{"Alta fragilidad, con tratamiento"}\NormalTok{,}
                       \StringTok{"Fragilidad media, con tratamiento"}\NormalTok{),}
                     \AttributeTok{levels =} \FunctionTok{c}\NormalTok{(}\StringTok{"Alta fragilidad, sin tratamiento"}\NormalTok{,}
                                \StringTok{"Alta fragilidad, con tratamiento"}\NormalTok{,}
                                \StringTok{"Fragilidad media, con tratamiento"}\NormalTok{)),}
  \AttributeTok{x1 =} \FunctionTok{c}\NormalTok{(}\DecValTok{1}\NormalTok{, }\DecValTok{1}\NormalTok{, }\DecValTok{0}\NormalTok{), }\AttributeTok{x2 =} \FunctionTok{c}\NormalTok{(}\DecValTok{0}\NormalTok{, }\DecValTok{1}\NormalTok{, }\DecValTok{1}\NormalTok{))}

\NormalTok{grid }\OtherTok{\textless{}{-}}\NormalTok{ tidyr}\SpecialCharTok{::}\FunctionTok{expand\_grid}\NormalTok{(}\AttributeTok{periodo =} \FunctionTok{factor}\NormalTok{(}\FunctionTok{levels}\NormalTok{(pp}\SpecialCharTok{$}\NormalTok{periodo), }\FunctionTok{levels}\NormalTok{(pp}\SpecialCharTok{$}\NormalTok{periodo)),}
                           \AttributeTok{arquetipo =}\NormalTok{ arquetipos}\SpecialCharTok{$}\NormalTok{arquetipo) }\SpecialCharTok{|\textgreater{}}
  \FunctionTok{left\_join}\NormalTok{(arquetipos, }\AttributeTok{by =} \StringTok{"arquetipo"}\NormalTok{)}
\NormalTok{grid}\SpecialCharTok{$}\NormalTok{h }\OtherTok{\textless{}{-}} \FunctionTok{predict}\NormalTok{(m\_pp, }\AttributeTok{newdata =}\NormalTok{ grid, }\AttributeTok{type =} \StringTok{"response"}\NormalTok{)}
\NormalTok{grid }\OtherTok{\textless{}{-}}\NormalTok{ grid }\SpecialCharTok{|\textgreater{}}
  \FunctionTok{mutate}\NormalTok{(}\AttributeTok{t =} \FunctionTok{as.integer}\NormalTok{(periodo)) }\SpecialCharTok{|\textgreater{}}
  \FunctionTok{arrange}\NormalTok{(arquetipo, t) }\SpecialCharTok{|\textgreater{}}
  \FunctionTok{group\_by}\NormalTok{(arquetipo) }\SpecialCharTok{|\textgreater{}}
  \FunctionTok{mutate}\NormalTok{(}\AttributeTok{S =} \FunctionTok{cumprod}\NormalTok{(}\DecValTok{1} \SpecialCharTok{{-}}\NormalTok{ h), }\AttributeTok{riesgo\_acum =} \DecValTok{1} \SpecialCharTok{{-}}\NormalTok{ S) }\SpecialCharTok{|\textgreater{}}
  \FunctionTok{ungroup}\NormalTok{()}

\FunctionTok{ggplot}\NormalTok{(grid, }\FunctionTok{aes}\NormalTok{(t, riesgo\_acum, }\AttributeTok{color =}\NormalTok{ arquetipo)) }\SpecialCharTok{+}
  \FunctionTok{geom\_line}\NormalTok{(}\AttributeTok{linewidth =} \DecValTok{1}\NormalTok{) }\SpecialCharTok{+} \FunctionTok{geom\_point}\NormalTok{(}\AttributeTok{size =} \DecValTok{2}\NormalTok{) }\SpecialCharTok{+}
  \FunctionTok{scale\_y\_continuous}\NormalTok{(}\AttributeTok{labels =}\NormalTok{ scales}\SpecialCharTok{::}\NormalTok{percent) }\SpecialCharTok{+}
  \FunctionTok{labs}\NormalTok{(}\AttributeTok{x =} \StringTok{"Periodo de revisión"}\NormalTok{, }\AttributeTok{y =} \StringTok{"Riesgo acumulado de fractura (1 {-} S)"}\NormalTok{,}
       \AttributeTok{color =} \ConstantTok{NULL}\NormalTok{) }\SpecialCharTok{+}
  \FunctionTok{theme}\NormalTok{(}\AttributeTok{legend.position =} \StringTok{"top"}\NormalTok{)}
\end{Highlighting}
\end{Shaded}

\begin{figure}[H]

{\centering \includegraphics{caso1_binaria_files/figure-pdf/fig-u15-clinica-1.pdf}

}

\caption{\label{fig-u15-clinica}Riesgo acumulado de fractura (1 - St)
predicho por el GLM cloglog para tres perfiles. La distancia vertical
entre `alta fragilidad sin tratar' y `con tratamiento' es la reducción
de riesgo absoluta atribuible al tratamiento; crece con el seguimiento.
Los valores exactos se leen sobre el eje al renderizar.}

\end{figure}

Los números ponen cifras a la decisión: tras los seis periodos, una
paciente de alta fragilidad sin tratar acumula un riesgo de fractura de
≈68 \%; con tratamiento baja a ≈52 \% (unos 16 puntos menos); y si
además su fragilidad fuera media, el riesgo caería a ≈27 \%. La
fragilidad pesa más que el fármaco ---reducirla, cuando se puede, mueve
más la aguja que el tratamiento por sí solo---.

Estas conclusiones son didácticas ---descansan en una cohorte
simulada---, pero el flujo es el real: de los hazard ratios y la curva
de supervivencia a una recomendación accionable. Con datos reales, el
mismo modelo daría las cifras sobre las que decidir.

\begin{tcolorbox}[enhanced jigsaw, coltitle=black, rightrule=.15mm, arc=.35mm, leftrule=.75mm, colback=white, colbacktitle=quarto-callout-important-color!10!white, left=2mm, breakable, toptitle=1mm, toprule=.15mm, bottomtitle=1mm, titlerule=0mm, title=\textcolor{quarto-callout-important-color}{\faExclamation}\hspace{0.5em}{Tarea del estudiante (1.5) --- resuelve la Pregunta 2 del caso}, bottomrule=.15mm, opacityback=0, opacitybacktitle=0.6, colframe=quarto-callout-important-color-frame]

Responde la \textbf{Pregunta 2} del cliente
(\protect\hyperlink{sec-u16}{Unidad 6}) sobre la cohorte de tu equipo:
\emph{¿cuándo se fracturan, cada cuánto conviene revisar, y retrasa el
tratamiento la fractura?}

Tareas:

\begin{enumerate}
\def\labelenumi{\arabic{enumi}.}
\tightlist
\item
  \textbf{Prepara y ajusta (RA7).} Reconoce que es un problema de
  \textbf{tiempo-a-evento con censura}, expande a
  \textbf{persona-periodo} (\texttt{pp}) y ajusta el GLM con enlace
  \textbf{cloglog}; reporta los \textbf{hazard ratios} de \texttt{x1} y
  \texttt{x2} (con su IC).
\item
  \textbf{Cuándo, y cada cuánto revisar.} Representa el \textbf{hazard
  base por periodo} (\(\hat h_{0t}\)) y la \textbf{supervivencia}
  \(S_t\); a partir de su forma (¿riesgo creciente?, ¿se acelera la
  caída?), propón al cliente un \textbf{calendario de revisiones}.
\item
  \textbf{Censura.} Explica por qué, siendo \textbf{no informativa}, no
  sesga la estimación: qué aporta ---y qué no--- una paciente censurada.
\item
  \textbf{Fragilidad de centro.} Añade \texttt{(1\ \textbar{}\ centro)}
  y comenta qué cambia (o no) en los efectos fijos y sus errores
  estándar.
\item
  \textbf{Predicción y respuesta.} Traza la \textbf{supervivencia} (o el
  riesgo acumulado \(1-S_t\)) para un par de perfiles (fragilidad
  alta/media, con y sin tratamiento), di si el tratamiento
  \textbf{retrasa} la fractura y responde al cliente.
\end{enumerate}

\emph{Control de calidad:} contrasta tus estimaciones con la
\textbf{verdad del DGP} (\texttt{attr(cohorte,\ "verdad")}) ---ajusta a
ciegas y valida al final---.

\end{tcolorbox}

\hypertarget{sec-u16}{%
\section{6 · Estudio de caso, exposición y evaluación}\label{sec-u16}}

\begin{tcolorbox}[enhanced jigsaw, coltitle=black, rightrule=.15mm, arc=.35mm, leftrule=.75mm, colback=white, colbacktitle=quarto-callout-note-color!10!white, left=2mm, breakable, toptitle=1mm, toprule=.15mm, bottomtitle=1mm, titlerule=0mm, title=\textcolor{quarto-callout-note-color}{\faInfo}\hspace{0.5em}{Sobre esta unidad}, bottomrule=.15mm, opacityback=0, opacitybacktitle=0.6, colframe=quarto-callout-note-color-frame]

Cierre del Caso 1: trabajo acompañado del estudio de caso y
\textbf{exposiciones} con Q\&A individual. No hay teoría nueva; se
integra y se evalúa todo lo anterior. Esta unidad reúne la
\textbf{maquinaria de evaluación} (actividades, instrucciones,
rúbricas). \textbf{RA implicados:} RA8 (desarrollar, presentar y
defender estudios de caso reproducibles) y, de forma acumulativa,
RA1--RA7.

\end{tcolorbox}

\hypertarget{sec-u16-caso}{%
\subsection{6.1 El estudio de caso (entregable)}\label{sec-u16-caso}}

\hypertarget{el-encargo}{%
\subsubsection{El encargo}\label{el-encargo}}

Te contrata la dirección médica de \textbf{OsteoRed}, una red de 24
clínicas de salud ósea que, dentro de un programa de prevención, ha
seguido durante seis años a una cohorte de mujeres. No son estadísticos:
te traen sus datos y sus preguntas, y esperan \textbf{respuestas y
recomendaciones}, no una clase de modelos.

De cada mujer registraron su \textbf{fragilidad ósea} (un marcador de
densidad mineral, estandarizado; \texttt{x1}), si seguía el
\textbf{tratamiento preventivo} (\texttt{x2}), la \textbf{clínica} a la
que acude (\texttt{centro}) y, a lo largo de hasta \textbf{seis
revisiones anuales}, si sufrió una \textbf{fractura} y \textbf{cuándo}
---o si terminó el seguimiento sin fracturarse--- (\texttt{ever},
\texttt{tiempo}, \texttt{evento}); y, cuando la hubo, su \textbf{tipo}
(\texttt{clase\_nom}: A/B/C) y su \textbf{gravedad}
(\texttt{sever\_ord}: leve/moderada/grave).

\begin{quote}
\textbf{Lo que la dirección de OsteoRed quiere saber}

\begin{enumerate}
\def\labelenumi{\arabic{enumi}.}
\tightlist
\item
  «¿Podemos anticipar qué mujeres acabarán fracturándose? ¿Influyen de
  verdad la fragilidad ósea y nuestro tratamiento, y cuánto? ¿Podríamos
  estimar el riesgo de una paciente nueva? Y el tratamiento, ¿protege
  por igual a todas o depende de lo frágil que esté cada una?»
\item
  «No nos basta con saber \emph{si} se fracturan: necesitamos saber
  \textbf{cuándo}. ¿En qué momento del seguimiento se dispara el riesgo?
  ¿Cada cuánto deberíamos citarlas a revisión? ¿El tratamiento
  \textbf{retrasa} la fractura? Y cuidado: muchas terminan el
  seguimiento sin fracturarse, no queremos que eso distorsione las
  cuentas.»
\item
  «Cuando hay fractura, la clasificamos en tres \textbf{tipos} (A, B,
  C). ¿Los mismos factores empujan hacia unos tipos más que otros?
  ¿Podríamos anticipar el tipo más probable?»
\item
  «Cada fractura tiene una \textbf{gravedad} (leve, moderada, grave).
  ¿Qué la agrava? ¿Nuestro tratamiento, además de evitar fracturas, hace
  que las que ocurren sean \textbf{menos graves}?»
\item
  «Trabajamos en \textbf{24 clínicas} y sospechamos que no todas rinden
  igual. ¿El riesgo depende de la clínica o es todo cosa de la paciente?
  ¿Hay centros que lo hacen claramente mejor o peor? Si abrimos uno
  nuevo, ¿qué deberíamos esperar?»
\end{enumerate}
\end{quote}

\hypertarget{tu-trabajo}{%
\subsubsection{Tu trabajo}\label{tu-trabajo}}

Entregas un \textbf{informe de consultoría reproducible} que convierta
esas preguntas en análisis. Nadie te dice qué modelo usar:
\textbf{elegirlo y defenderlo es parte de la nota}. Para \textbf{cada}
pregunta del cliente:

\begin{enumerate}
\def\labelenumi{\arabic{enumi}.}
\tightlist
\item
  \textbf{Identifica el objetivo} y qué tipo de respuesta tienes entre
  manos.
\item
  \textbf{Propón un análisis adecuado y justifícalo:} por qué ese modelo
  y no otro, qué enlace, si procede una interacción, si hay que
  reconocer el agrupamiento por clínica o el tiempo con censura.
\item
  \textbf{Resuélvelo} con rigor, con las decisiones de selección y el
  diagnóstico que lo respalden.
\item
  \textbf{Interprétalo} en la escala que el cliente entienda:
  probabilidades, veces más riesgo, plazos\ldots{}
\item
  \textbf{Responde al cliente} con una recomendación accionable.
\end{enumerate}

De forma transversal, el informe incluye \textbf{predicción} donde el
cliente la pide (el riesgo de una paciente nueva, el calendario de
revisiones, qué esperar de una clínica nueva), el \textbf{control de
calidad} de cada análisis y la \textbf{reproducibilidad}: un
\texttt{.qmd} que corre de principio a fin (RA8).

\hypertarget{el-formato-un-informe-tuxe9cnico}{%
\subsubsection{El formato: un informe
técnico}\label{el-formato-un-informe-tuxe9cnico}}

OsteoRed pregunta en términos \textbf{ejecutivos} (qué hacer, a quién
vigilar, cada cuánto), pero tu entregable es un \textbf{informe
técnico}, no un resumen ejecutivo. La diferencia está en el lector: un
resumen ejecutivo va directo a conclusiones y recomendaciones, con lo
técnico al mínimo; un informe técnico se escribe para alguien que podría
\textbf{auditar o reproducir} tu trabajo, así que documenta el
razonamiento completo ---objetivo, método elegido y su justificación,
supuestos, ajuste, diagnóstico e interpretación--- antes de llegar a la
recomendación.

En la práctica: traduce cada pregunta ejecutiva del cliente a uno o
varios \textbf{objetivos analíticos precisos}; organiza el informe
siguiendo esas preguntas y, dentro de cada una, recorre el ciclo
(objetivo → método justificado → resultados → diagnóstico →
interpretación → respuesta y recomendación); y \textbf{devuelve cada
conclusión al lenguaje del cliente}. El rigor técnico y la lectura
accionable conviven: lo técnico sostiene la recomendación, no la
sustituye.

\begin{tcolorbox}[enhanced jigsaw, coltitle=black, rightrule=.15mm, arc=.35mm, leftrule=.75mm, colback=white, colbacktitle=quarto-callout-important-color!10!white, left=2mm, breakable, toptitle=1mm, toprule=.15mm, bottomtitle=1mm, titlerule=0mm, title=\textcolor{quarto-callout-important-color}{\faExclamation}\hspace{0.5em}{Apartado de Reflexión sobre el uso de la IA (obligatorio)}, bottomrule=.15mm, opacityback=0, opacitybacktitle=0.6, colframe=quarto-callout-important-color-frame]

Cierra el informe con un apartado de \textbf{Reflexión ética sobre la
IA} que recoja una reflexión sobre el uso que se ha hecho de la IA para
la realización del trabajo: qué herramientas usasteis y para qué, los
\textbf{prompts relevantes}, y ---sobre todo--- cómo
\textbf{verificasteis y validasteis} lo que produjeron (qué
comprobasteis, qué corregisteis, qué descartasteis).

Es la evidencia de vuestro proceso y del uso responsable de la IA.

\end{tcolorbox}

\begin{tcolorbox}[enhanced jigsaw, coltitle=black, rightrule=.15mm, arc=.35mm, leftrule=.75mm, colback=white, colbacktitle=quarto-callout-important-color!10!white, left=2mm, breakable, toptitle=1mm, toprule=.15mm, bottomtitle=1mm, titlerule=0mm, title=\textcolor{quarto-callout-important-color}{\faExclamation}\hspace{0.5em}{Anexo de reproducibilidad (obligatorio)}, bottomrule=.15mm, opacityback=0, opacitybacktitle=0.6, colframe=quarto-callout-important-color-frame]

Añade al final un \textbf{anexo de reproducibilidad} que permita a otra
persona regenerar el informe: el \textbf{entorno}
(\texttt{sessionInfo()} o \texttt{renv}, con las versiones de los
paquetes), la \textbf{semilla} usada, las \textbf{instrucciones para
ejecutar} el \texttt{.qmd} de principio a fin y la \textbf{procedencia y
estructura de los datos} (\texttt{cohorte} y su expansión
persona-periodo).

\end{tcolorbox}

\begin{tcolorbox}[enhanced jigsaw, coltitle=black, rightrule=.15mm, arc=.35mm, leftrule=.75mm, colback=white, colbacktitle=quarto-callout-important-color!10!white, left=2mm, breakable, toptitle=1mm, toprule=.15mm, bottomtitle=1mm, titlerule=0mm, title=\textcolor{quarto-callout-important-color}{\faExclamation}\hspace{0.5em}{Importante}, bottomrule=.15mm, opacityback=0, opacitybacktitle=0.6, colframe=quarto-callout-important-color-frame]

Además, el líder del equipo ha de completar y presentar, en un documento
adicional: - La \textbf{planificación y reparto} de tareas asumidas en
el equipo (roles, tareas, cómo os organizasteis en el tiempo, qué
revisiones o hitos hubo y cómo resolvisteis los bloqueos. - La
evaluación de los compañeros, anotando el grado de implicación (del 0 al
100\%) en el trabajo. - El checklist de tareas realizadas (según la
propuesta especificada abajo) cumplimentado.

\end{tcolorbox}

\begin{tcolorbox}[enhanced jigsaw, coltitle=black, rightrule=.15mm, arc=.35mm, leftrule=.75mm, colback=white, colbacktitle=quarto-callout-tip-color!10!white, left=2mm, breakable, toptitle=1mm, toprule=.15mm, bottomtitle=1mm, titlerule=0mm, title=\textcolor{quarto-callout-tip-color}{\faLightbulb}\hspace{0.5em}{Control de calidad interno: ajusta a ciegas y valida contra el DGP}, bottomrule=.15mm, opacityback=0, opacitybacktitle=0.6, colframe=quarto-callout-tip-color-frame]

Esto es tuyo, no del cliente. Como la cohorte es simulada, los
parámetros que la generaron están en \texttt{attr(cohorte,\ "verdad")}:
los efectos \(\beta\) de \texttt{x1} y \texttt{x2}, la varianza entre
centros \(\sigma_u\), el riesgo base por periodo \(\alpha_t\) y los
parámetros del nominal y del ordinal. La disciplina que se te pide:
\textbf{ajusta y decide sin mirar esa ``verdad''} y, solo al final,
compárala con tus estimaciones, mostrando el contraste en el informe.

Ten presente que esta comprobación es un lujo \textbf{didáctico},
posible solo porque los datos son simulados. Con datos reales no existe
tal ``verdad'': allí la confianza en el modelo descansa en el
diagnóstico, la validación fuera de muestra y el conocimiento del
problema. Recógelo así en tu Reflexión.

\end{tcolorbox}

\begin{tcolorbox}[enhanced jigsaw, coltitle=black, rightrule=.15mm, arc=.35mm, leftrule=.75mm, colback=white, colbacktitle=quarto-callout-tip-color!10!white, left=2mm, breakable, toptitle=1mm, toprule=.15mm, bottomtitle=1mm, titlerule=0mm, title=\textcolor{quarto-callout-tip-color}{\faLightbulb}\hspace{0.5em}{Cada equipo, una base de datos distinta}, bottomrule=.15mm, opacityback=0, opacitybacktitle=0.6, colframe=quarto-callout-tip-color-frame]

Cada equipo trabaja sobre \textbf{su propia cohorte}, generada con una
configuración diferente ---pero coherente--- de los parámetros del DGP.
Las diez configuraciones y los datos por equipo los produce el script
\texttt{R/generar\_datos\_equipos.R}, que deja en
\texttt{datos\_equipos/equipo\_XX/} la cohorte (\texttt{cohorte.rds} /
\texttt{.csv}) y su expansión persona-periodo (\texttt{pp.rds} /
\texttt{.csv}). El enunciado y las preguntas del cliente son los mismos
para todos; lo que cambia de un equipo a otro son las
\textbf{conclusiones numéricas} (efectos, riesgos, plazos). La clave con
los parámetros verdaderos (\texttt{configuraciones\_equipos.csv}) la
guarda el/la docente.

\end{tcolorbox}

\hypertarget{sec-u16-checklist}{%
\subsection{6.2 Checklist de tareas}\label{sec-u16-checklist}}

Antes de entregar, comprueba dos cosas: que has \textbf{respondido cada
pregunta del cliente} y que no te has saltado ninguna de las
\textbf{tareas que un buen analista} hace en un proyecto. Verifica cada
casilla y marca las que hayas comprobado e incluido en la entrega.

\hypertarget{respondida-cada-pregunta-del-cliente}{%
\subsubsection{¿Respondida cada pregunta del
cliente?}\label{respondida-cada-pregunta-del-cliente}}

\begin{itemize}
\tightlist
\item[$\square$]
  \textbf{P1 · Riesgo de fractura:} quién se fractura, papel de la
  fragilidad y del tratamiento, riesgo de una paciente nueva, y si el
  efecto del tratamiento depende de la fragilidad.
\item[$\square$]
  \textbf{P2 · Cuándo ocurre:} evolución del riesgo a lo largo de las
  revisiones, calendario de seguimiento y efecto del tratamiento sobre
  los plazos, con la censura bien tratada.
\item[$\square$]
  \textbf{P3 · Tipo de fractura:} qué factores empujan hacia cada tipo y
  anticipación del más probable.
\item[$\square$]
  \textbf{P4 · Gravedad:} qué agrava la fractura y si el tratamiento la
  reduce.
\item[$\square$]
  \textbf{P5 · Peso de la clínica:} cuánto del riesgo depende del
  centro, centros mejores o peores y qué esperar de uno nuevo.
\end{itemize}

\hypertarget{tareas-analuxedticas-del-proyecto}{%
\subsubsection{Tareas analíticas del
proyecto}\label{tareas-analuxedticas-del-proyecto}}

\emph{Comprensión y datos}

\begin{itemize}
\tightlist
\item[$\square$]
  Traducir cada pregunta del cliente a objetivos analíticos precisos.
\item[$\square$]
  Explorar y depurar los datos (EDA): tipos de variable, escalas,
  atípicos y ausentes.
\end{itemize}

\emph{Modelización}

\begin{itemize}
\tightlist
\item[$\square$]
  Identificar el tipo de respuesta y elegir un modelo adecuado a cada
  objetivo, justificándolo.
\item[$\square$]
  Especificar bien: enlace, interacciones pertinentes, agrupamiento
  (efectos aleatorios), tiempo con censura, proporcionalidad.
\item[$\square$]
  Comprobar la modificación de un efecto por otra variable cuando el
  cliente lo pregunta (interacción).
\item[$\square$]
  Ajustar y comprobar la convergencia y la estabilidad del ajuste.
\end{itemize}

\emph{Evaluación}

\begin{itemize}
\tightlist
\item[$\square$]
  Diagnosticar: residuos, calibración, separación, bondad de ajuste
  (χ²), sobredispersión y los supuestos propios del modelo
  (proporcionalidad de odds en el ordinal, riesgos proporcionales en
  supervivencia, convergencia y no-singularidad en los mixtos).
\item[$\square$]
  Evaluar la capacidad predictiva con la métrica adecuada: AUC y
  calibración (binaria), Brier multiclase (nominal),
  hazard/supervivencia por periodo.
\item[$\square$]
  Seleccionar y comparar modelos con criterios (AIC/BIC) y contrastes
  (LRT), justificando la elección.
\end{itemize}

\emph{Resultados}

\begin{itemize}
\tightlist
\item[$\square$]
  Interpretar los efectos en la escala propia de cada respuesta ---odds
  ratios y efectos marginales (binaria), OR relativos (nominal), OR
  acumulado (ordinal), hazard ratios (supervivencia)---, en términos
  relativos y absolutos, con sus intervalos de confianza.
\item[$\square$]
  Predecir donde el cliente lo pide (perfiles de paciente, paciente
  nueva, clínica nueva, calendario).
\item[$\square$]
  Control de calidad: contrastar las estimaciones con la verdad del DGP.
\end{itemize}

\emph{Cierre}

\begin{itemize}
\tightlist
\item[$\square$]
  Responder a cada pregunta con una recomendación accionable, en el
  lenguaje del cliente.
\item[$\square$]
  Comunicar: informe técnico claro y ordenado, figuras legibles y
  referencias (APA 7).
\item[$\square$]
  Reproducibilidad: el \texttt{.qmd} corre de principio a fin; semilla,
  entorno y datos documentados en el anexo.
\item[$\square$]
  Incluir la reflexión sobre el uso de la IA.
\item[$\square$]
  Adjuntar el documento de planificación y reparto de tareas,
  evaluación, valoración y checklist de tareas realizadas.
\end{itemize}

\hypertarget{sec-repro}{%
\section{Reproducibilidad}\label{sec-repro}}

\begin{Shaded}
\begin{Highlighting}[]
\FunctionTok{sessionInfo}\NormalTok{()}
\end{Highlighting}
\end{Shaded}

\begin{verbatim}
R version 4.6.0 (2026-04-24)
Platform: x86_64-apple-darwin20
Running under: macOS Sequoia 15.7.7

Matrix products: default
BLAS:   /Library/Frameworks/R.framework/Versions/4.6-x86_64/Resources/lib/libRblas.0.dylib 
LAPACK: /Library/Frameworks/R.framework/Versions/4.6-x86_64/Resources/lib/libRlapack.dylib;  LAPACK version 3.12.1

locale:
[1] en_US.UTF-8/en_US.UTF-8/en_US.UTF-8/C/en_US.UTF-8/en_US.UTF-8

time zone: Europe/Madrid
tzcode source: internal

attached base packages:
[1] stats     graphics  grDevices utils     datasets  methods   base     

other attached packages:
 [1] survminer_0.5.2        ggpubr_0.6.3           survival_3.8-6        
 [4] nnet_7.3-20            yardstick_1.4.0        ggeffects_2.3.2       
 [7] emmeans_2.0.3          pROC_1.19.0.1          car_3.1-5             
[10] carData_3.0-6          marginaleffects_0.32.0 GGally_2.4.0          
[13] MuMIn_1.48.19          lubridate_1.9.5        forcats_1.0.1         
[16] stringr_1.6.0          dplyr_1.2.1            purrr_1.2.2           
[19] readr_2.2.0            tidyr_1.3.2            tibble_3.3.1          
[22] ggplot2_4.0.3          tidyverse_2.0.0        performance_0.17.0    
[25] arm_1.15-3             lme4_2.0-1             Matrix_1.7-5          
[28] MASS_7.3-65            DHARMa_0.5.0           see_0.14.0            
[31] patchwork_1.3.2        aplore3_0.9            broom_1.0.13          

loaded via a namespace (and not attached):
 [1] RColorBrewer_1.1-3     rstudioapi_0.18.0      jsonlite_2.0.0        
 [4] shape_1.4.6.1          datawizard_1.3.1       magrittr_2.0.5        
 [7] TH.data_1.1-5          estimability_1.5.1     jomo_2.7-6            
[10] farver_2.1.2           logistf_1.26.1         nloptr_2.2.1          
[13] rmarkdown_2.31         vctrs_0.7.3            minqa_1.2.8           
[16] rstatix_0.7.3          tinytex_0.59           htmltools_0.5.9       
[19] haven_2.5.5            Formula_1.2-5          mitml_0.4-5           
[22] plyr_1.8.9             sandwich_3.1-2         zoo_1.8-15            
[25] mime_0.13              lifecycle_1.0.5        iterators_1.0.14      
[28] pkgconfig_2.0.3        gap_1.15.2             R6_2.6.1              
[31] fastmap_1.2.0          rbibutils_2.4.1        shiny_1.13.0          
[34] digest_0.6.39          qgam_2.0.0             labeling_0.4.3        
[37] timechange_0.4.0       abind_1.4-8            mgcv_1.9-4            
[40] compiler_4.6.0         withr_3.0.2            doParallel_1.0.17     
[43] S7_0.2.2               backports_1.5.1        ggstats_0.13.0        
[46] ggsignif_0.6.4         pan_1.9                tools_4.6.0           
[49] otel_0.2.0             httpuv_1.6.17          glue_1.8.1            
[52] nlme_3.1-169           promises_1.5.0         grid_4.6.0            
[55] checkmate_2.3.4        generics_0.1.4         operator.tools_1.6.3.1
[58] gtable_0.3.6           tzdb_0.5.0             formula.tools_1.7.1   
[61] data.table_1.18.4      hms_1.1.4              utf8_1.2.6            
[64] foreach_1.5.2          pillar_1.11.1          later_1.4.8           
[67] splines_4.6.0          lattice_0.22-9         tidyselect_1.2.1      
[70] knitr_1.51             reformulas_0.4.4       gridExtra_2.3         
[73] stats4_4.6.0           xfun_0.58              stringi_1.8.7         
[76] yaml_2.3.12            boot_1.3-32            evaluate_1.0.5        
[79] codetools_0.2-20       cli_3.6.6              rpart_4.1.27          
[82] xtable_1.8-8           Rdpack_2.6.6           Rcpp_1.1.1-1.1        
[85] coda_0.19-4.1          parallel_4.6.0         gap.datasets_0.0.6    
[88] glmnet_5.0             mvtnorm_1.4-0          scales_1.4.0          
[91] insight_1.5.1          rlang_1.2.0            multcomp_1.4-31       
[94] mice_3.19.0           
\end{verbatim}

\begin{Shaded}
\begin{Highlighting}[]
\CommentTok{\# renv::snapshot()   \# fijar el estado del entorno al cerrar el caso}
\end{Highlighting}
\end{Shaded}




\end{document}
